% !TeX program = pdflatex
% Expansions about the short-range end of the long-range Ising model in d = 1, 2, 3.
% Public edition (SRI/note). Build: latexmk -pdf -synctex=1 sri_short_range.tex (pdfLaTeX, standard packages only,
% embedded bibliography, no external input files). Every quoted number is reproduced as described in Appendix
% "Reproduction" (SRI/code/<package>, documentation in SRI/docs/<package>).
\documentclass[11pt,a4paper]{article}
\usepackage[T1]{fontenc}
\usepackage{lmodern}
\usepackage[margin=23mm,headheight=14pt]{geometry}
\usepackage{amsmath,amssymb,mathtools,booktabs,longtable,array}
\usepackage{microtype,xcolor,fancyhdr,enumitem}
\usepackage[font=small,labelfont=bf]{caption}
\usepackage[colorlinks=true,linkcolor=teal!60!black,citecolor=teal!60!black,urlcolor=blue!60!black]{hyperref}
\usepackage{xurl}
\usepackage{needspace,etoolbox}
\pretocmd{\section}{\needspace{10\baselineskip}}{}{}
\pretocmd{\subsection}{\needspace{6\baselineskip}}{}{}
\hypersetup{pdftitle={Expansions about the short-range end of the long-range Ising model in one, two and three dimensions},
 pdfauthor={LRI project note}}
% The MD5 of this source is written into the PDF information dictionary; tools/check_note.py compares it with the
% source to decide whether the PDF is current.
\edef\SRINoteSourceMDfive{\pdfmdfivesum file {sri_short_range.tex}}
\pdfinfo{/SRINoteSourceMD5 (\SRINoteSourceMDfive)}
\pagestyle{fancy}\fancyhf{}
\fancyhead[L]{\small LRI: expansions about the short-range end}
\fancyhead[R]{\small 17 October 2026}\fancyfoot[C]{\thepage}
\setlength{\parskip}{3.5pt}
\setlength{\emergencystretch}{2.5em}
\setlist{itemsep=2pt,topsep=4pt}
\numberwithin{equation}{section}
\allowdisplaybreaks[2]
% == open items: \TODO marks what is still to be filled in before the export (tools/check_note.py lists them)
\newcommand{\TODO}[1]{\textcolor{red!75!black}{\textbf{[TODO:}~#1\textbf{]}}}
% == notation
\newcommand{\dd}{\,\mathrm d}
\newcommand{\De}{\Delta_{\mathcal E}}
\newcommand{\Ds}{\Delta_\sigma}
\newcommand{\CSR}{C^{\rm SR}}
\newcommand{\db}{\bar\delta}
\newcommand{\AC}{\operatorname{AC}}
\newcommand{\FP}{\operatorname{FP}}
\newcommand{\Li}{\operatorname{Li}}
\newcommand{\Res}{\operatorname*{Res}_{\delta=0}}
\newcommand{\cI}{\mathcal I}
\newcommand{\Og}{\mathcal O_g}
\newcommand{\Oh}{\mathcal O_h}
\newcommand{\hMS}{h_{\rm MS}}
\newcommand{\ket}[1]{\lvert#1\rangle}
\newcommand{\file}[1]{\path{#1}}
\newcommand{\sci}[2]{#1\times10^{#2}}
\newcommand{\pFq}[3]{{}_4F_3\!\left(\begin{matrix}#1\\[1pt]#2\end{matrix};#3\right)}

\title{\textbf{Expansions about the short-range end\\ of the long-range Ising model\\
in one, two and three dimensions}\\[3mm]
\large The dual of the chain through eighth order, the crossover coefficients $c_2$ ($d=2$) and $c_1$ ($d=3$),
their checks and their status}
\author{LRI project note}
\date{17 October 2026}

\begin{document}
\maketitle

\begin{abstract}
\noindent
We collect the expansions of the thermal scaling dimension $\De$ of the critical long-range Ising model about the
point where it joins the short-range Ising universality class, in $d=1,2,3$. In $d=1$ the strong--weak dual of the
chain of Benedetti, Lauria, Maz\'a\v c and van Vliet gives $\De=1+\sum_n\CSR_n\Ds^{n/2}$; its published coefficients
are $\CSR_1=-2$ and $\CSR_2=\frac12$. In the matched two-coupling prescription defined in Sec.~\ref{sec:dual}, the
archived integral expressions for $\CSR_3$ to $\CSR_7$ reduce exactly to rational combinations of zeta values,
$\CSR_5$ to $\CSR_7$ by exact word-algebra certificates. The step from the kink integrals to those expressions cannot
be rerun, because 21 stage programs of the archive are missing; it is backed by an independent re-derivation of the
inputs by Mellin--Hepp sector integration, which reproduces $\CSR_3$ to $\CSR_7$ within statistical errors of
$\sci{3}{-14}$ to $\sci{3}{-5}$ (largest pull $0.9$ over the 25 inputs of $\CSR_5$ to $\CSR_7$).
$\CSR_8=247.381204805\ldots$ is a numerical record printed to about 85 digits, of which a second evaluation of one
input cross-checks about 57 decimal places; its rational part, 88, is exact, and a weight-seven zeta-value formula
matches all printed digits but is not proved. The independent re-derivation gives $\CSR_8=242.8\pm2.4$ (one
standard error), 1.9 to 4.0 standard errors below the record depending on the error estimator; the difference sits
in one input integral. A third calculation, in a scheme defined by the magnetization at the integrable point $s=1$,
where the Bethe-ansatz solution gives the whole $s=1$ part, shares no input with the other two and gives
$\CSR_8=247.38118\pm0.13$, an empirical numerical uncertainty. It agrees with the record to $\sci{2.5}{-5}$, reproduces
$\CSR_3$ to $\CSR_7$ to $\sci{7}{-26}$, and so indicates that the low independent value is the one in error; the
re-derivation's own input comparison locates the difference in that integral. In $d=2$, $c_1=0$, and $c_2=-4.0662701481325\ldots$ is a closed form in Gamma values and three
${}_4F_3(1)$ constants whose identification with the physical coefficient rests on one unproved uniform-remainder
step. A Monte Carlo integration of the regulated integral gives $-4.075\pm0.019$ (one standard error), and
fourth-order perturbation theory with the truncated cylinder Hamiltonian gives $-4.0670(24)$, where the error is a
linear sum of systematic diagnostics; its relative difference from the closed form is $\sci{1.9}{-4}$. In $d=3$,
$c_1=0.2732$, with the estimated range $[0.2690,0.2760]$, conditional on assumed gaps in the omitted spectrum and on
assumed ranges of the input data. Every
result carries its status word and caveats. The appendices give in full the derivations, proofs, methods and checks
of the supplied packages behind these results (the independent re-derivation, the Monte Carlo and the cylinder
continuum analysis are described in the main text), including the formal arguments that two couplings close under renormalization of the dual and
that its prescription extracts the local beta functions. Appendix~\ref{app:repro} maps each number to the package and
command that reproduce it, and a last appendix gives the magnetization-scheme calculation of $\CSR_8$ that the paper
quotes.
\end{abstract}

\clearpage
\setcounter{tocdepth}{2}
\begingroup\small\tableofcontents\endgroup
\clearpage

%====================================================================
\section{Introduction and conventions}\label{sec:intro}
%====================================================================
\subsection{The short-range end}
The critical long-range Ising (LRI) model in $d$ dimensions has couplings $J(r)\propto r^{-(d+s)}$. Its spin field has
the exact dimension $\Ds=(d-s)/2$ throughout the long-range regime, and its fixed point joins the short-range Ising
(SRI) universality class at
\begin{equation}
 s_*=d-2\Ds^{\rm SR}.
\end{equation}
Near $s_*$ the natural small parameter is
\begin{equation}
 \delta=\frac{s_*-s}{2}=\Ds-\Ds^{\rm SR}.
 \label{eq:delta}
\end{equation}
This note collects what is known about the thermal dimension $\De$ near that end in $d=1,2,3$. These are the results
quoted in Sec.~3.2, ``Near the short-range end'', of the accompanying paper~\cite{Paper} (its table of short-range
coefficients, the closed forms of the dual of the chain, the $c_2$ row of its table of perturbative Hamiltonian
truncation, and the closed-form appendix for $c_2$), with the derivations, checks and status of each. The paper's
resummation of the expansion about the generalized-free-field end, $s=d/2$, also imposes these values at the other
end of the long-range window, $s=s_*$.

Two different expansions apply. In $d=1$ the short-range end is $s_*=1$, where $\Ds^{\rm SR}=0$ (for $s>1$ the chain
has no finite-temperature transition); the expansion is that of a strong--weak dual and proceeds in $\sqrt{\Ds}$
(Sec.~\ref{sec:d1}). In $d=2,3$ the fixed point near $s_*$ is the SRI conformal field theory (CFT) coupled to a
generalized free field (GFF), and the expansion proceeds in integer powers of $\delta$
(Secs.~\ref{sec:d2} and~\ref{sec:d3}). Table~\ref{tab:conv} lists the conventions.

\begin{table}[htbp]\centering\small
\renewcommand{\arraystretch}{1.25}\setlength{\tabcolsep}{4pt}
\begin{tabular}{@{}cllll@{}}
\toprule
$d$ & $s_*$ & $\Ds^{\rm SR}$, $\De^{\rm SR}$ & variable & expansion\\
\midrule
1 & $1$ & $0$, $1$ & $x=\Ds=\frac{1-s}{2}$ & $\De=1+\sum_{n\ge1}\CSR_n\,x^{n/2}$\\
2 & $\frac74$ & $\frac18$, $1$ & $\delta=\frac{7/4-s}{2}=\Ds-\frac18$ & $\De=1+c_1\delta+c_2\delta^2+O(\delta^3)$\\
3 & $1.963702388(48)$ & $0.518148806(24)$, $1.41262528(29)$ & $\delta=\frac{s_*-s}{2}$ & $\De=\De^{\rm SR}+c_1\delta+O(\delta^2)$\\
\bottomrule
\end{tabular}
\caption{Conventions. The $d=3$ data are those of Ref.~\cite{Chang:2024whx}. In $d=1$ the variable $\delta$ of
Eq.~\eqref{eq:delta} is $x$; the dual of Sec.~\ref{sec:d1} uses $\db=1-s=2x$, which its computational sources write
$\delta$ (as in the input name $N_\delta$). A coefficient of $(s_*-s)^n$ is $c_n/2^n$: the coefficient of $(7/4-s)^2$ in
$d=2$ is $c_2/4$, and that of $s_*-s$ in $d=3$ is $c_1/2$. The paper writes $C^{\rm SR}_n(d)$ for all three:
$C^{\rm SR}_n(1)=\CSR_n$ multiplies $\Ds^{n/2}$, while $C^{\rm SR}_n(2)=c_n$ and $C^{\rm SR}_n(3)=c_n$ multiply
$\delta^n$ (so $c_2=C^{\rm SR}_2(2)$ and $c_1=C^{\rm SR}_1(3)$).}
\label{tab:conv}
\end{table}

The symbol $b_0$ follows the sources and has two unrelated meanings: in $d=1$ it is the charge of the compact field of
the dual (Sec.~\ref{sec:dual}); in $d=2$ and $d=3$ it is the one-loop coefficient of the beta function of the
crossover coupling (Secs.~\ref{sec:rg2} and~\ref{sec:d3}). A few other letters are reused, always with a local
definition: $h$ is the current coupling of the dual in Sec.~\ref{sec:d1} and the $g^4$ coefficient of
$\gamma_{\mathcal E}$ or of the cylinder observable in Sec.~\ref{sec:d2} ($h_0$, $h(\lambda)$); $W$ is an input of
$\CSR_8$, the Gaussian source factor of Sec.~\ref{sec:d2} and, as $W_k$, the cylinder counterterms; $\kappa$ is the
reference inverse length of the dual (Appendix~\ref{app:dualA}), the finite redefinition parameter of
Sec.~\ref{sec:rg2} and, as $\kappa_r$, a counterterm coefficient of Appendix~\ref{app:cyl}; and $\chi$ is the
generalized free field in $d=2,3$, while the bare current of the dual is $\chi_B$ and its renormalized current primary
(Appendices~\ref{app:dualA} and~\ref{app:dualC}) is written $\chi$ there.

\subsection{Sources and who produced them}\label{sec:sources}
The note rests on seven packages of the release. Four were written by another AI system, OpenAI GPT-6 astra (via
ChatGPT), and supplied by the author: a research compendium for the $d=1$ dual (its derivations, code, exact
expressions, certificates and numerical records of $\CSR_3$ to $\CSR_8$; ``the compendium''; its original
archive, of which the release ships a pruned copy, is ``the archive''), a closed-form study of $c_2$ in $d=2$ (``the
study'', also called the closed-form or final-method study), a base study of $c_2$ with the truncated cylinder Hamiltonian, and a package for $c_1$
in $d=3$ (``the $c_1$ package''). Claude Code (Anthropic) sessions reviewed each of them (``the review'' of each),
re-derived every input of $\CSR_5$ to $\CSR_8$ independently, ran the Monte Carlo of $c_2$, and wrote the cluster
layer and the continuum analysis of the cylinder calculation, and computed $\CSR_8$ in the magnetization scheme
(Sec.~\ref{sec:integ}); OpenAI Codex reviewed the plans and the code of those two calculations. Two formal arguments about the dual, that two couplings close under its renormalization and that its
prescription extracts the local beta functions, are also by GPT-6 astra.

The main text states each result with its basis and caveats; the appendices give the derivations behind them, edited
to the notation of this note: the compendium in Appendices~\ref{app:dualA} (the model, the regulator and the
finite-response renormalization), \ref{app:dualB} (the third to sixth orders) and \ref{app:dualC} (the exact
certificates and the seventh and eighth orders); the two arguments about the dual in Appendix~\ref{app:clos}; the
study in Appendix~\ref{app:cform}; the method of the cylinder calculation in Appendix~\ref{app:cyl}; and the $c_1$
package in Appendix~\ref{app:cone}. Section~\ref{sec:prov} gives the authorship, and Appendix~\ref{app:repro} the
folder of each package. A third calculation of $\CSR_8$, in a scheme defined by the magnetization at the integrable
point $s=1$, was made for the project after the other packages: Sec.~\ref{sec:integ} summarizes it,
Appendix~\ref{app:integ} gives it in full, and its code and records are the package \file{integ_d1}.

\subsection{Status words}\label{sec:words}
Each quoted result carries one of the following status words.
\begin{description}[leftmargin=1.5em]
\item[literature] a published value, cited where it is used and quoted as published; where this note re-derives it
 as a check, Table~\ref{tab:main} says so;
\item[proved closed form within the prescription] ($d=1$) an exact expression into which the archived integral
 expressions of the matched two-coupling prescription (Sec.~\ref{sec:dual}) reduce exactly, with exact certificates
 where the reduction is computer-assisted; it covers neither the step from the kink integrals to those expressions nor
 the prescription itself;
\item[closed form conditional on one unproved step] ($c_2$) an exact expression whose evaluation is proved, but whose
 identification with the physical coefficient needs a step that is not proved;
\item[numerical] a number obtained by numerical integration, extrapolation or sampling, with the kind of its error
 stated: a statistical standard error, a linear sum of systematic diagnostics that is not a standard deviation, or,
 for a record without an error estimate, the precision to which it is cross-checked;
\item[conjectured] an exact expression supported numerically but without a proof;
\item[conditional] a statement or range that holds only under explicitly stated, unproved assumptions.
\end{description}
The file \file{results/sr_values.json} uses the short forms literature, closed form, numerical, conjectured and
conditional, with the qualification in a note or caveat field; it lists $c_1$ in $d=3$ as a numerical retained value
and an estimated range with the status conditional. Errors are always labelled with their kind. A pull is (our value $-$ claimed
value)/(quoted standard error).

\subsection{Results and their status}
Table~\ref{tab:main} gives every result of this note that the paper quotes, and three that it does not print (the
largest pull of the recomputed inputs, the independent value of $\CSR_8$ and the agreement of the two independent
codes for $c_1$), with its status word, its basis and its caveats. The summaries and checks follow in
Secs.~\ref{sec:d1} to~\ref{sec:d3}, the derivations in Appendices~\ref{app:dualA} to~\ref{app:cone}
and~\ref{app:integ}, and the authorship in Sec.~\ref{sec:prov}.

{\small
\renewcommand{\arraystretch}{1.2}
\begin{longtable}{@{}>{\raggedright\arraybackslash}p{.12\textwidth}>{\raggedright\arraybackslash}p{.215\textwidth}>{\raggedright\arraybackslash}p{.165\textwidth}>{\raggedright\arraybackslash}p{.41\textwidth}@{}}
\caption{The results near the short-range end, with the status words of Sec.~\ref{sec:words}. In $d=1$ ``the
prescription'' is the matched two-coupling prescription of Sec.~\ref{sec:dual}. The paper prints six significant
digits; the longer strings are evaluations of the closed forms and of the $\CSR_8$ record, rounded to 15 significant
digits (all digits are in \texttt{results/sr\_values.json}).}\label{tab:main}\\
\toprule
result & value & status & basis and caveats\\
\midrule\endfirsthead
\toprule
result & value & status & basis and caveats\\
\midrule\endhead
\bottomrule\endfoot
\multicolumn{4}{@{}l}{$d=1$}\\
$\CSR_1$, $\CSR_2$ & $-2$, $\frac12$ & literature & Refs.~\cite{Benedetti:2024wgx,Benedetti:2025nzp}; re-derived from the leading flow (Sec.~\ref{sec:invariants}).\\
$\CSR_3$ to $\CSR_7$ & Eq.~\eqref{eq:closed}: & proved closed form within the prescription & The archived integral expressions of the beta inputs reduce exactly to zeta values: $\CSR_3$, $\CSR_4$ by analytic period evaluations (with four archived mixed response coefficients), $\CSR_5$ to $\CSR_7$ by exact word-algebra certificates (Sec.~\ref{sec:closed}). The step from the kink integrals to these expressions cannot be rerun (Appendix~\ref{app:cannot}); the independent re-derivation reproduces $\CSR_3$ to $\CSR_7$ numerically (Sec.~\ref{sec:indep}). Prescription caveat: Secs.~\ref{sec:dual} and~\ref{sec:d1caveats}, Appendix~\ref{app:clos}.\\
\quad$\CSR_3$ & $5.72736813369645\ldots$ & & \\
\quad$\CSR_4$ & $-12.1961914926697\ldots$ & & \\
\quad$\CSR_5$ & $29.2378115809900\ldots$ & & \\
\quad$\CSR_6$ & $-58.6656974742253\ldots$ & & \\
\quad$\CSR_7$ & $51.0689886129281\ldots$ & & \\
$\pi_0\CSR_8$ (rational part) & $88$ & proved closed form within the prescription & Exact rational projection, from two separate exact calculations (Sec.~\ref{sec:C8}); the same two qualifications as $\CSR_3$ to $\CSR_7$.\\
$\CSR_8$ & $247.381204804994\ldots$ & numerical & Record printed to about 85 digits (printing floor $\sci{2}{-85}$); a second evaluation of its input $P_6$ cross-checks about 57 decimal places; six inputs are single numerical records; no error bound (Sec.~\ref{sec:C8}).\\
Eq.~\eqref{eq:C8cand} & weight-seven formula for $\CSR_8$ & conjectured & Matches the record to $\sci{3.1}{-86}$; no certificate; the record is its only high-precision support (the magnetization-scheme value agrees to $\sci{2.5}{-5}$).\\
inputs of $\CSR_5$ to $\CSR_7$, independent & largest pull $0.9$ & numerical & One standard error (the largest of three statistical estimators); 25 inputs, largest pull $0.915$, rms $0.34$; the coefficients within $0.78$ ($\CSR_6$); the paper says only that the recomputed inputs match.\\
$\CSR_8$, independent & $242.8\pm2.4$ & numerical & One standard error (largest of three estimators); 1.9 to 4.0 standard errors below the record, depending on the estimator; the difference sits in one input integral, $A_9[\varepsilon^0\ell^1]$ (Sec.~\ref{sec:indep}); the magnetization-scheme value (next row) indicates that this value is the one in error; not printed by the paper.\\
$\CSR_8$, magnetization scheme & $247.38118\pm0.13$ & numerical & Empirical numerical uncertainty, a linear sum of the differences between three runs and a validated roundoff estimate, not a standard deviation or a bound; no input shared with the two hard-core calculations, no nine-kink integral; reproduces $\CSR_3$ to $\CSR_7$ to $\sci{7}{-26}$; $\sci{2.5}{-5}$ from the record; two check criteria amended after their data were seen; the paper quotes it as $247.38\pm0.13$ (its Sec.~3.2); package \file{integ_d1} (Sec.~\ref{sec:integ}, Appendix~\ref{app:integ}).\\
\midrule
\multicolumn{4}{@{}l}{$d=2$}\\
$c_1$ & $0$ & literature & BRRZ~\cite{Behan:2017emf}; re-derived ($I_2(0)=0$, Sec.~\ref{sec:d2}).\\
$c_2$ & $-4.06627014813252\ldots$ & closed form conditional on one unproved step & The evaluation of $\cI$ is proved; the identification with the physical coefficient needs the uniform-remainder step (Sec.~\ref{sec:unproved}); a formal result in critically tuned analytic minimal subtraction.\\
$c_2$, Monte Carlo & $-4.075\pm0.019$ & numerical & One standard error (statistical); $0.5$ standard errors from the closed form (Sec.~\ref{sec:mc}).\\
$c_2$, cylinder & $-4.0670(24)$ & numerical & Linear sum of systematic diagnostics (64\% regulator model), not a standard deviation; $\sci{1.9}{-4}$ relative from the closed form, which was known but is not an input; error rule revised after the first result; 11 of the 25 points are archived records of the base study (Sec.~\ref{sec:cyl}).\\
\midrule
\multicolumn{4}{@{}l}{$d=3$}\\
$c_1$ & $0.2732$; estimated range $[0.2690,\,0.2760]$ & conditional & Retained 60-primary value $0.27315242\ldots$ (Sec.~\ref{sec:retained}). The range holds if no omitted primary lies below the gaps 8 ($\mathbb Z_2$-even) and 7 ($\mathbb Z_2$-odd), the inputs lie within twice their quoted errors, a published constraint on $\lambda_{\epsilon\epsilon\epsilon'}$ applies and the chosen margins and allowances suffice; it is neither a bound nor a confidence interval, and mild relaxations extend it (Secs.~\ref{sec:range} and~\ref{sec:howcond}).\\
$c_1$, two independent codes & agree to $10^{-12}$ & numerical & Absolute difference in $c_1$ of the review's two re-implementations from the package (Sec.~\ref{sec:retained}); not printed by the paper.\\
$c_1$, BRRZ & $0.27(4)$ & literature & Ref.~\cite{Behan:2017emf}.\\
\end{longtable}
}

%====================================================================
\section{One dimension: the dual of the chain}\label{sec:d1}
%====================================================================
\subsection{The dual, its two couplings and the prescription}\label{sec:dual}
Near $s=1$ the domain walls of the critical chain form a dilute gas of alternating kinks and antikinks with a
logarithmic interaction~\cite{Anderson:1970kyb,Kosterlitz:1976zz}. Benedetti, Lauria, Maz\'a\v c and van Vliet recast
this gas as a two-state system (a qubit) coupled to a compact Gaussian field $\phi\sim\phi+2\pi/b_0$ of negative
dimension~\cite{Benedetti:2024wgx,Benedetti:2025nzp}. With $\db=1-s$ and $f_{\db}(r)=(r^{\db}-1)/\db$, the field has
covariance $-2f_{\db}(|X-Y|)$ on neutral test functions, and the path-ordered evolution
\begin{equation}
 \mathrm P\exp\int\dd X\,\bigl[g\,\Og+h\,\Oh\bigr],\qquad
 \Og=\hat\sigma_+{:}e^{\mathrm i b_0\phi}{:}+\hat\sigma_-{:}e^{-\mathrm i b_0\phi}{:},\qquad
 \Oh=\hat\sigma_3\chi_B,\qquad \chi_B=\tfrac{\mathrm i}{\sqrt2}\partial\phi,
 \label{eq:dual}
\end{equation}
reproduces the alternating gas because $\hat\sigma_\pm^2=0$ (Appendices~\ref{app:dualA:model}
and~\ref{app:dualA:gas} derive the gas and this representation). The bare current $\chi_B$ is a different field from the generalized free
field $\chi$ of Secs.~\ref{sec:d2} and~\ref{sec:d3} (and from the renormalized current primary that
Appendix~\ref{app:dualA:dictionary} writes $\chi$); the symbol table of the dual is Table~\ref{tab:dualA:notation}. The critical chain is the sector invariant under a $U(1)$
that shifts $\phi$ and rotates the qubit. The fugacity $g$ and the current coupling $h$ flow; the current can be
resummed into the effective charge $b=b_0-\sqrt2h$.

\paragraph{The matched two-coupling prescription.} The coefficients $\CSR_3$ to $\CSR_8$ are computed in one fixed
prescription, taken from the compendium and derived in Appendices~\ref{app:dualA:regulator}
to~\ref{app:dualA:recurrence}:
\begin{enumerate}[label=(\roman*),leftmargin=2em]
\item \emph{Regulator.} Kinks at ordered positions $0<X_1<\cdots<X_n<1$ on an interval of unit length, with a hard core
 $X_{i+1}-X_i\ge\varepsilon$ on consecutive kinks only (none between a kink and an endpoint), and $\ell=\log\varepsilon$
 (Eq.~\eqref{eq:dualA:domain}).
\item \emph{Exact current resummation.} Before the hard core is imposed, a field-dependent change of frame removes the
 bulk current interaction. It shifts the wall charge to $b=b_0-\sqrt2h$ and leaves charges $\pm h/\sqrt2$ at the two
 endpoints of the interval, whose contraction with the walls is an \emph{endpoint source} $\eta=\sqrt2\,h\,b$
 (Appendix~\ref{app:dualA:regulator}).
\item \emph{Matching at the cutoff.} Normal ordering at the cutoff defines the charge $B=b\,\varepsilon^{\db/2}$ and
 the fugacity $g$, keeping the $\db$-dependent normal-ordering factor of the wall interaction exactly (finite-tail
 matching); the essential couplings are
 \begin{equation}
  u=g^2,\qquad v=B^2=1+z,\qquad \beta_u=2uF(u,z,\db),\qquad \beta_z=(1+z)H(u,z,\db),
 \end{equation}
 in units with $b_0=1$, so that $z=0$ is the marginal point (Appendix~\ref{app:dualA:matching}).
\item \emph{Vacuum normalization.} The ordered moments of the regulated gas (Sec.~\ref{sec:inputs1}) are divided by the
 vacuum moments, as complete series in $\varepsilon$ and $\ell$, before any coefficient is extracted
 (Appendix~\ref{app:dualA:moments}).
\item \emph{Two projected responses.} The beta functions of $u$ and $z$ are solved, order by order, from two responses
 only: the odd (charged) response and the current response, each projected on its coefficient
 $[\varepsilon^0\ell^1]$, through a finite-response Jacobian that carries their finite parts $[\varepsilon^0\ell^0]$
 (Appendix~\ref{app:dualA:recurrence}).
\end{enumerate}
``Matched'' refers to (iii), and ``two-coupling'' to the assumption that $(u,z)$ close under renormalization: no
further coupling enters their flow at the orders computed. At leading order
(Appendix~\ref{app:dualA:dictionary})
\begin{equation}
 F=z+u+\cdots,\qquad H=-\db+4u+\cdots,\qquad u_*=\frac{\db}{4}+O(\db^2),\quad z_*=-\frac{\db}{4}+O(\db^2),
\end{equation}
so the fugacity at the fixed point is of order $\sqrt{\db}$, which is why the series runs in $\sqrt{\Ds}$.

\paragraph{What ``proved within the prescription'' covers.} In this prescription the beta inputs are fixed
regulated integrals. A \emph{proved closed form within the prescription} means that the archived exact expressions of
these integrals reduce exactly to rational combinations of zeta values (Sec.~\ref{sec:closed}). Two steps lie outside
it. First, the step from the kink integrals to the archived expressions cannot be rerun, because 21 stage programs of
the archive are missing (Appendix~\ref{app:cannot}); it is backed numerically by the independent re-derivation of
Sec.~\ref{sec:indep}. Second, whether the prescription computes the coefficients of the dual (that the dual
describes the chain is the published duality proposal), which rests on the two-coupling closure and on the prescription extracting the local beta functions; two unrefereed arguments support
both formally (Appendix~\ref{app:clos}: the closure statement in Appendix~\ref{app:clos:closure}, the prescription
statement in Appendix~\ref{app:clos:prescription}). Their status is stated once, in the caveat ``Two-coupling
closure'' of Sec.~\ref{sec:d1caveats}, together with the evidence that bears on it.

\subsection{Stability invariants and the thermal series}\label{sec:invariants}
At the fixed point $F=H=0$, so the trace and determinant of the stability matrix in $(u,z)$ are
\begin{equation}
 T=2u_*F_u+(1+z_*)H_z,\qquad D=2u_*(1+z_*)\bigl(F_uH_z-F_zH_u\bigr),
\end{equation}
with the derivatives taken before the fixed point is substituted. They are invariant under regular changes of
couplings, although individual beta coefficients are not (Appendix~\ref{app:dualA:invariants}). With
\begin{equation}
 T(\db)=\frac{\db}{2}+\sum_{m\ge2}T_m\db^m,\qquad D(\db)=-2\db+\sum_{m\ge2}D_m\db^m,
\end{equation}
the two dimensions are
\begin{equation}
 \Delta_\pm=1+\frac T2\pm\sqrt{\frac{T^2}{4}-D},\qquad \De=\Delta_-,
 \label{eq:eigen}
\end{equation}
assuming, as in Ref.~\cite{Paper}, that the relevant operator is not crossed by another level. The square root is
$\sqrt{2\db}$ times a power series in $\db$, so the odd powers of $\sqrt{\Ds}$ come from the eigenvalue splitting and
the even powers from $T/2$. At leading order $\Delta_\pm=1\pm\sqrt{2\db}+\db/4+O(\db^{3/2})$~\cite{Benedetti:2025nzp},
which is
\begin{equation}
 \CSR_1=-2,\qquad \CSR_2=\frac12.
\end{equation}
Expanding Eq.~\eqref{eq:eigen} gives the exact conversions
\begin{equation}
\begin{aligned}
 \CSR_3&=D_2-\tfrac1{16}, & \CSR_4&=2T_2, & \CSR_5&=2D_3-\tfrac{T_2}{2}+\tfrac14\bigl(D_2-\tfrac1{16}\bigr)^2,\\
 \CSR_6&=4T_3, & \CSR_7&=4D_4-T_3-T_2^2+\tfrac12\CSR_3\CSR_5, & \CSR_8&=8T_4,
\end{aligned}
\label{eq:conv}
\end{equation}
and in general $\CSR_{2m}=2^{m-1}T_m$. These are algebraic identities; they do not determine an uncomputed invariant
($\CSR_9$ needs $D_5$).

\subsection{Beta inputs and finite-response renormalization}\label{sec:inputs1}
The invariants are polynomials in Taylor coefficients of $F$ and $H$ at the marginal point $v=1$, $\db=0$, which
the computation calls beta inputs. They come from ordered moments of the regulated gas, odd moments $A_{2m+1}$,
vacuum moments $Z_{2m}$ and current moments $Q_{2m}$, with up to nine kinks and their derivatives in the charge $v$,
the tail parameter $\db$ and the endpoint source $\eta$. After the vacuum normalization (iv), the coefficients
$[\varepsilon^0\ell^0]$ and $[\varepsilon^0\ell^1]$ are taken, and the finite-response Jacobian (v) turns them into
beta coefficients (Appendices~\ref{app:dualA:moments} to~\ref{app:dualA:recurrence}; the inputs of each order,
evaluated or given by their archived exact expressions, are in Appendices~\ref{app:dualB} and~\ref{app:dualC}). Five inputs ($a,p,b,c,f$) fix $\CSR_3$ and $\CSR_4$, 25 fix $\CSR_7$, and ten more
($W,S,U,J_9,L_3,M_{21},M_{12},N_2,N_\delta,K_v$) enter $\CSR_8$. For example (the input $b=b_5$ is not the charge $b$)
\begin{align}
 a&=-4\zeta(2),\qquad c=-16,\qquad p=8-2\zeta(2)-7\zeta(3)-8\log2,\notag\\
 b&=-18-6\zeta(2)+\tfrac{28}{3}\zeta(3)+32\log2-12\log3-4\log^22-8\Li_2(-\tfrac12),\\
 f&=32-\tfrac{136}{3}\zeta(3)-64\log2+48\log3+16\log^22+32\Li_2(-\tfrac12),\notag
\end{align}
and $D_2=\frac12-\frac a2-\frac c8$, $T_2=-\frac a8+\frac b4-\frac c{32}+\frac f{16}+\frac p2$. The regulator
constants $\log2$, $\log3$ and $\Li_2(-\frac12)$ are genuine parts of the individual inputs and cancel from the
invariants. The eighth order is affine in the ten new inputs,
\begin{equation}
 \CSR_8=B_{\rm old}+4W+2S+\tfrac34U+\tfrac14J_9+\tfrac3{32}L_3-\tfrac14M_{21}+\tfrac12M_{12}-\tfrac1{16}N_2
 +\tfrac18N_\delta+\tfrac1{32}K_v,
 \label{eq:C8assembly}
\end{equation}
where the polynomial $B_{\rm old}$ in the 25 earlier inputs (printed in Appendix~\ref{app:dualC:Bold}) has the value
$-861.66356723547977920\ldots$.

\subsection{Closed forms through the seventh order}\label{sec:closed}
The invariants reduce to
\begin{align}
 D_2&=\tfrac52+2\zeta(2),\qquad T_2=2-2\zeta(2)-4\zeta(3),\qquad
 D_3=\tfrac{21}{4}-13\zeta(2)+11\zeta(3)+11\zeta(4),\notag\\
 T_3&=\tfrac{13}{2}-\tfrac{19}{2}\zeta(2)+26\zeta(3)-20\zeta(4)-68\zeta(5)+28\zeta(2)\zeta(3),\\
 D_4&=13-\tfrac{179}{8}\zeta(2)+\tfrac{47}{2}\zeta(3)-\tfrac{493}{8}\zeta(4)+289\zeta(5)-130\zeta(2)\zeta(3)
 +\tfrac{499}{4}\zeta(6)-76\zeta(3)^2,\notag
\end{align}
and Eq.~\eqref{eq:conv} gives
\begin{equation}
\begin{aligned}
 \CSR_3&=\frac{39}{16}+2\zeta(2),\qquad \CSR_4=4-4\zeta(2)-8\zeta(3),\\
 \CSR_5&=\frac{11249}{1024}-\frac{361}{16}\zeta(2)+24\zeta(3)+\frac{49}{2}\zeta(4),\\
 \CSR_6&=26-38\zeta(2)+104\zeta(3)-80\zeta(4)-272\zeta(5)+112\zeta(2)\zeta(3),\\
 \CSR_7&=\frac{1798583}{32768}-\frac{90637}{1024}\zeta(2)+\frac{453}{4}\zeta(3)-\frac{16835}{64}\zeta(4)+1224\zeta(5)\\
 &\quad-540\zeta(2)\zeta(3)+\frac{4335}{8}\zeta(6)-320\zeta(3)^2.
\end{aligned}
\label{eq:closed}
\end{equation}
Their values are in Table~\ref{tab:main}. Each $\CSR_n$ with $n\le7$ lies in the rational span of multiple zeta
values of weight at most $n-1$.

\paragraph{How the closed forms are established.} $\CSR_3$ and $\CSR_4$ follow from analytic evaluations of the
four- and five-insertion periods, closed by a finite-part trilogarithm identity proved with explicit primitives and a
real-branch Spence--Kummer identity; the four mixed response coefficients of Eq.~\eqref{eq:dualB:mixed} that attach the
periods to the inputs are archived values, not derived here (Appendix~\ref{app:dualB}, which also assembles $\CSR_5$ and $\CSR_6$ from
their inputs and evaluates the analytic inputs $A=\tfrac12b_{3,vv}$, $E=b_{3,v\db}$ (Appendix~\ref{app:dualB:conv}) and the two
endpoint constants of $\CSR_5$; the other inputs of
these orders are archived exact expressions). From $\CSR_5$ on, the exact archived expressions of the beta inputs, rational
functions of the cutoff multiplying iterated integrals, are lifted to a free rational vector space of words of
regularized hyperlogarithm constants. The difference between the archived expression and the proposed closed form is
then written as an explicit rational combination of shuffle products of proved zero relations (shuffle, stuffle,
duality, a convergent confluence lemma and a beta-integral identity), and a standalone verifier checks it word by
word in exact rational arithmetic, with no numerical period and no assumption of numerical independence
(Appendix~\ref{app:dualC:certificates}):
\begin{center}\small
\begin{tabular}{@{}lrrr@{}}
\toprule
target & weights & nonzero words & certificate rows\\
\midrule
$D_3$ (hence $\CSR_5$) & 2--4 & 170 & 476\\
$\CSR_6$ & 2--5 & 1\,945 & 56\,381\\
$D_4$, $\CSR_7$ (with all product terms) & 2--6 & 25\,212 & 588\,878\\
\bottomrule
\end{tabular}
\end{center}
A separate standard-library checker parses the archived input expressions, lifts them into the word algebra and
verifies the trace assembly and the target coefficients, so the certificates are tied to the inputs. The certificates
prove the reduction of the archived regulated-integral expressions to zeta values. They do not reintegrate the
original kink kernels, and they do not test the prescription itself (Sec.~\ref{sec:dual}).

\subsection{The eighth order}\label{sec:C8}
\paragraph{Rational part.} The formal weight-zero projection $\pi_0$, which removes positive-weight constant periods
after integration, is computed exactly (Appendix~\ref{app:dualC:C8rational}; the trace reduction of the eighth order
is in Appendix~\ref{app:dualC:C8setup}). A decomposition into four raw residues (of the six-kink and eight-kink current
moments, the seven-kink tail and the nine-kink fugacity) gives, after the first two are combined with the known terms,
\begin{equation}
 \pi_0\CSR_8=-\frac{111389}{24}-\frac{910891}{8}+\frac{1423087}{12}=88,
\end{equation}
and a separate literal full-Jacobian calculation gives $\pi_0T(\db)=\frac{\db}2+2\db^2+\frac{13}2\db^3+11\db^4+\cdots$,
hence $\pi_0\CSR_8=8\times11=88$. None of the raw integrals was inferred from the answer or from an earlier
conjectured value, 160, which this result excludes.

\paragraph{Full value.} Four of the ten new inputs were computed in four separate stages: the six-kink mixed logarithm
$P_6=1261.944897843140491372852392951531767542124984174\ldots$ (an exact period expression with 66\,404 terms,
evaluated at 70 and 100 working digits), from which $N_\delta=-4P_6+11218.914021881803142\ldots$; the seven-kink
tail coefficient $U=-1296.687960291526491656\ldots$; the eight-kink current derivative
$K_v=-1486.962109278823406128\ldots$; and the ninth-order fugacity $J_9=-273.724515999324646514\ldots$, a
cancellation between pieces of order $6\times10^5$. The other six ($W,S,L_3,M_{21},M_{12},N_2$) enter as
\emph{inherited numerical records}: decimal values from one recorded 100-digit evaluation, with no surviving exact
expression and no second evaluation (Appendix~\ref{app:dualC:C8packages} describes all ten). The assembly reduces to $\CSR_8=C_{\rm fixed}-\frac12P_6$ with
$C_{\rm fixed}=878.353653726563794093781271\ldots$, and gives
\begin{equation}
 \CSR_8=247.3812048049935484073550747660999986521591512209\ldots .
 \label{eq:C8num}
\end{equation}

\paragraph{How many digits.} Four facts bear on the digits of Eq.~\eqref{eq:C8num}; none is an interval bound.
\begin{itemize}
\item \emph{Cross-check.} The two evaluations of $P_6$ give values of $\CSR_8$ that differ by $\sci{2.2}{-57}$, the
 precision of the 70-digit run. This cross-checks about 57 decimal places, and only through $P_6$.
\item \emph{Printing floor.} The record is assembled exactly from 36 numbers (the 35 inputs, with $N_\delta$ given as
 $P_6$ and its remainder), all decimal strings of 85 to 96 decimals except the exact $c=-16$. Changing each by half a
 unit in its last printed digit changes $\CSR_8$ by about $\sci{2}{-85}$ in all, about 87 significant digits, quoted
 as about 85. This is the precision to which the inputs are written, not a demonstrated accuracy: the six inherited
 inputs have neither exact expressions nor second evaluations.
\item \emph{The formula.} The 100-digit value agrees with the conjectured Eq.~\eqref{eq:C8cand} below to
 $\sci{3.1}{-86}$. This is consistent with both, but the formula has no high-precision support other than this record,
 so the agreement is not independent evidence for the record's digits.
\item \emph{A third value.} The magnetization-scheme value of Sec.~\ref{sec:integ}, $247.38118\pm0.13$, agrees with the
 record to $\sci{2.5}{-5}$; it bears on the first few digits only.
\end{itemize}

\paragraph{The weight-seven candidate.} An integer-relation search (Appendix~\ref{app:dualC:C8assembly}) on the
70-digit value, in the basis
$\{\zeta(2),\zeta(3),\zeta(4),\zeta(5),\zeta(2)\zeta(3),\zeta(6),\zeta(3)^2,\zeta(7),\zeta(2)\zeta(5),\zeta(3)\zeta(4)\}$
with the rational part 88 subtracted, gives
\begin{equation}
\begin{aligned}
 \CSR_8\overset{?}{=}{}&88-178\zeta(2)+110\zeta(3)+113\zeta(4)+4892\zeta(5)-2376\zeta(2)\zeta(3)-\frac{8630}{3}\zeta(6)
 +1872\zeta(3)^2\\
 &-12168\zeta(7)+6048\zeta(2)\zeta(5)+1536\zeta(3)\zeta(4).
\end{aligned}
\label{eq:C8cand}
\end{equation}
The 100-digit evaluation differs from Eq.~\eqref{eq:C8cand} by $\sci{3.1}{-86}$, and a search with the rational
part left free recovers 88. In the review of the compendium (Sec.~\ref{sec:review1}) the relation stayed the only
one when 26 further constants were added to the basis. A looser earlier search had produced a different relation
that failed at higher precision. Equation~\eqref{eq:C8cand} is therefore conjectured: the residual is strong
numerical evidence, not an interval enclosure or an exact certificate. A proof would assemble every input of
$\CSR_8$ as an exact expression in one regularized-period convention (which requires recovering or regenerating the
six inherited inputs), subtract Eq.~\eqref{eq:C8cand} and certify the difference at every weight through seven,
products included, as was done for $\CSR_7$.

\subsection{Independent re-derivation of the inputs}\label{sec:indep}
The compendium's raw integral expressions had been replayed only with its own code. Every input of $\CSR_5$ to
$\CSR_8$ was therefore recomputed from scratch.
\begin{itemize}
\item \emph{Definitions.} The 35 inputs were written as Taylor coefficients of the auxiliary beta functions from the
 compendium's text alone (the definitions now in Appendices~\ref{app:dualA:recurrence} and~\ref{app:dualC:C8setup}).
\item \emph{Renormalization map.} The response-to-beta map, the vacuum normalization, the fixed point and the
 invariants were derived again from first principles (symbolically); the map agrees with the compendium's (Eq.~\eqref{eq:dualA:recurrence}). All
 required raw coefficients must be supplied; a missing one is an error, never a zero.
\item \emph{Engine.} A new numerical engine computes the $\varepsilon$ expansion of the hard-core moments: a Mellin
 transform in the cutoff, Hepp sectors (total orderings of the gaps) with the divergent variables subtracted
 analytically, and rank-one lattice quasi-Monte Carlo with six independent random shifts per job.
\item \emph{Direct check.} The raw moments were also integrated directly at finite cutoffs (nested Gauss--Legendre,
 relative accuracy about $10^{-14}$) and compared with the engine's series.
\item \emph{Independence.} No integration code or raw data of the compendium was used. Its numbers were read by a
 script only in a comparison step that refused to run before the independent results had been written and hashed.
 Some compendium numbers were visible earlier in task descriptions and reports; the result records disclose each
 case, and none entered a computation.
\end{itemize}
Three error estimators were computed: linear propagation with the replicate covariance, the standard error of six
replicate assemblies (one per random shift), and a sector-pooled estimate from the scatter of the shifts within each
integration sector, summed over sectors. The protocol, fixed before the comparison, quotes the largest of the three;
all three are statistical estimates, not bounds.

\paragraph{Through the seventh order} (production about $\sci{3.3}{3}$ core-hours on the symmetry cluster). All 25
inputs agree with the compendium: the rms pull is $0.34$, and the largest, $0.915$, is that of the input the compendium
calls $\Omega$ (the coefficient of $u^3z$ in $F$, Eq.~\eqref{eq:dualC:C7jets}). The coefficients agree within $0.78$ standard errors:
\begin{center}\small
\begin{tabular}{@{}lllr@{}}
\toprule
& independent & closed form & pull\\
\midrule
$\CSR_5$ & $29.237811563\pm\sci{1.35}{-7}$ & $29.2378115809900\ldots$ & $-0.13$\\
$\CSR_6$ & $-58.66569861\pm\sci{1.45}{-6}$ & $-58.6656974742253\ldots$ & $-0.78$\\
$D_4$ & $-2.53463048\pm\sci{7.85}{-6}$ & $-2.53462979953825\ldots$ & $-0.09$\\
$\CSR_7$ & $51.0689860\pm\sci{3.14}{-5}$ & $51.0689886129281\ldots$ & $-0.08$\\
\bottomrule
\end{tabular}
\end{center}
The errors are one standard error. Five checks use no compendium number. They are the proved low-order closed forms
(17 quantities, among them $\CSR_3$ and $\CSR_4$ with errors $\sci{3.3}{-14}$ and $\sci{9.8}{-10}$; largest pull
$0.52$), the cancellation of 702 negative powers of $\varepsilon$ in the normalized responses (largest pull $2.37$),
39 higher-logarithm predictions (largest pull $0.96$), a renormalization-group second-logarithm test derived for this
work (six responses, largest pull $1.09$), and the direct finite-cutoff comparison. In the last, a moment passes when
its two leading coefficients agree with the direct fit, that is, no pull reaches 3 and their joint
two-degree-of-freedom $\chi^2$ stays below $11.83$; 44 of the 47 moments pass. The other three are flagged: two depend
on the fit window, and one comes from a single production coefficient that is 4 replicate standard errors low and
changes $\CSR_7$ by $\sci{7}{-10}$. $\CSR_7$ and $D_4$ are effectively one test, dominated by the eight-kink current
moment.

\paragraph{The eighth order} (about $\sci{1.9}{4}$ core-hours, on the symmetry cluster and a personal desktop). All
ten new inputs and all 201 raw scalars are the project's own. Table~\ref{tab:C8inputs} compares them with the
compendium.

\begin{table}[htbp]\centering\small
\renewcommand{\arraystretch}{1.15}
\begin{tabular}{@{}lrllr@{}}
\toprule
input & coefficient & compendium & independent & pull\\
\midrule
$W$ & $4$ & $-37.70811151495$ & $-37.70811151495\pm\sci{1.1}{-10}$ & $-0.02$\\
$S$ & $2$ & $878.1837337477$ & $878.1837337543\pm\sci{3.1}{-9}$ & $+2.13$\\
$U$ & $\frac34$ & $-1296.687960$ & $-1296.687551\pm0.0015$ & $+0.27$\\
$J_9$ & $\frac14$ & $-273.72452$ & $-291.94\pm9.48$ & $-1.92$\\
$L_3$ & $\frac3{32}$ & $-0.6386928106$ & $-0.6386928106\pm\sci{9}{-11}$ & $-0.02$\\
$M_{21}$ & $-\frac14$ & $-171.2355765656$ & $-171.2355765656\pm\sci{1.2}{-9}$ & $-0.01$\\
$M_{12}$ & $\frac12$ & $-422.7150579699$ & $-422.7150579699\pm\sci{6.1}{-9}$ & $0.00$\\
$N_2$ & $-\frac1{16}$ & $189.742012494$ & $189.742009475\pm\sci{8.5}{-6}$ & $-0.35$\\
$N_\delta$ & $\frac18$ & $6171.134430509$ & $6171.134427764\pm\sci{9.1}{-6}$ & $-0.30$\\
$K_v$ & $\frac1{32}$ & $-1486.962109$ & $-1487.004626\pm0.158$ & $-0.27$\\
\midrule
$\CSR_8-\frac14J_9$ & & $315.812334$ & $315.811301\pm0.0051$ & $-0.20$\\
$\CSR_8$ & & $247.381204805$ & $242.825\pm2.371$ & $-1.92$\\
\bottomrule
\end{tabular}
\caption{The new inputs of $\CSR_8$, with their coefficients in Eq.~\eqref{eq:C8assembly}: compendium values and the
independent re-derivation (one standard error, the largest of three estimators). The $+2.13$ of $S$ is
$\sci{7.5}{-12}$ relative and moves $\CSR_8$ by $\sci{1.3}{-8}$.}
\label{tab:C8inputs}
\end{table}

The result is
\begin{equation}
 \CSR_8=242.83\pm2.37\qquad(\text{independent; one standard error}),
\end{equation}
$1.92$ standard errors below the record with the quoted error (rounded, $242.8\pm2.4$; for the other
estimators see the next paragraph). Everything
except the $J_9$ term is tested sharply: $\CSR_8-\frac14J_9$ agrees at $-0.20\sigma$, which checks $B_{\rm old}$, the
seven inputs of the re-derivation's first tier (the six inherited ones and $N_\delta$), $U$ and $K_v$ together to
$\sci{2}{-5}$ of $\CSR_8$. The whole difference sits in $J_9$, and within it in one raw coefficient,
$A_9[\varepsilon^0\ell^1]$ of the nine-kink odd moment: independent minus compendium is $-20.4\pm9.3$
($-2.19\sigma$, $\sci{3.3}{-5}$ relative), while the other four pieces of $J_9$ agree.

\paragraph{How large the tension is.} The size of the pull depends on the error estimator, because $\sigma(\CSR_8)$
comes almost entirely (93\%) from $A_9[\varepsilon^0\ell^1]$. With the same central value $242.825$:
\begin{center}\small
\begin{tabular}{@{}llr@{}}
\toprule
estimator & $\sigma(\CSR_8)$ & pull\\
\midrule
sector-pooled (quoted, the largest of the three) & $2.37$ & $-1.92$\\
per job the larger of the replicate and pooled projections & $2.42$ & $-1.89$\\
chunk-pooled (the post-assembly statistics check) & & $-2.09$\\
linear propagation & $1.37$ & $-3.32$\\
replicate spread (five degrees of freedom) & $1.15$ & $-3.98$\\
\bottomrule
\end{tabular}
\end{center}
The independent value therefore lies 1.9 to 4.0 standard errors below the record, depending on the estimator. Through
the seventh order the same rule, the largest of the three, gives an rms pull of $0.34$ over the 25 inputs, so there it
is conservative. For $A_9$ the replicate error is $0.52$ of the sector-pooled one, and an archive-free comparison of
production with an independent probe on a subsample of its sectors differs by $+2.16$ standard errors, so
$\sigma(A_9)$, and with it $\sigma(\CSR_8)$, may be underestimated. Which value is right is indicated by the third
calculation of Sec.~\ref{sec:integ}, which agrees with the record. A second independent evaluation of
$A_9[\varepsilon^0\ell^1]$ with fresh random shifts (about $\sci{9}{3}$ core-hours) would show whether the low value
is a fluctuation or a systematic error. A 1\%
measurement of $\CSR_8$ cannot test the eleven rational coefficients of Eq.~\eqref{eq:C8cand} in any case; that would
need $A_9$ to about $10^{-9}$ relative, beyond the present method.

\subsection{The eighth order from the integrable point}\label{sec:integ}
A third calculation of $\CSR_8$ uses neither the matched two-coupling prescription nor any input of the two
calculations above. At $s=1$ the wall gas is the Coulomb gas of the anisotropic Kondo model~\cite{Anderson:1969xx,Anderson:1970kyb}, equivalently of the ohmic spin--boson model~\cite{Leggett:1987zz}, which is solvable by Bethe
ansatz~\cite{Andrei:1983,Tsvelick:1983}. The calculation (Appendix~\ref{app:integ}) has three parts.
\begin{itemize}
\item \emph{Scheme.} The gas, without a hard core, is put in a uniform magnetic field, and its couplings $\hat u$ and
 $\hat z$ are defined from the magnetization (the magnetization scheme). The flow of $\hat z$ is then exact,
 $\hat H=4\hat u-\db$, and only one function $\hat F$ has to be computed; $T$ and $D$ follow as in
 Sec.~\ref{sec:invariants}, with a factor $1-\db/2$ that converts the field scale to a length scale.
\item \emph{The exact part at $s=1$.} The whole $\db=0$ part of $\hat F$, to all orders in the fugacity, follows from
 the magnetization at $s=1$, derived from the Bethe-ansatz equations by a Wiener--Hopf solution as a Mellin--Barnes
 integral (Eq.~\eqref{eq:integ:MB}). It contains the counterpart of $J_9$ exactly, so no nine-kink integral is needed.
\item \emph{Kink clusters.} The $\db$-dependent coefficients come from connected clusters of two to eight kinks with
 logarithm insertions, continued analytically in the charge $v$ to $v=1$ by sector decomposition. Their poles at $v=1$
 must cancel in $\hat F$. These cancellations, the exact clusters at $s=1$ and the closed forms of $\CSR_3$ to $\CSR_7$
 are the checks.
\end{itemize}
The same assembly reproduces the closed forms of $\CSR_3$ to $\CSR_7$ to at most $\sci{7}{-26}$ relative. $\CSR_8$ needs
two new coefficients, a six-kink one, $f_{202}=172415.83018563\pm\sci{1.7}{-7}$, and an eight-kink one,
$f_{301}=-514275.75\pm0.17$, which enter as $2f_{202}+\frac34f_{301}$. The result is
\begin{equation}
 \CSR_8=247.38118\pm0.13\qquad(\text{magnetization scheme; empirical numerical uncertainty}).
 \label{eq:C8integ}
\end{equation}
The central value is given to five decimals because it is compared with the record below. The uncertainty is $\frac34$
of that of $f_{301}$, a linear sum of the differences between three complete runs (two
quadrature grids and two contour radii) and of a roundoff estimate validated by recomputing 56 sectors in double-double
arithmetic. It is not a standard deviation or a bound, and it cannot exclude an error common to all three runs. The
central value was frozen before any comparison; two criteria of the checks were amended after their data had been seen
(Appendix~\ref{app:integ:disclosures}). The code and records of this calculation are the package
\file{code/integ_d1/}, whose replay rebuilds every record from the shipped engine outputs (Appendix~\ref{app:repro}).

\paragraph{Comparison.} Equation~\eqref{eq:C8integ} agrees with the record~\eqref{eq:C8num} to $\sci{2.5}{-5}$
($10^{-7}$ relative, $0.0002$ of its uncertainty) and lies $4.56$ above the independent value, 1.9 of that value's
standard error. It shares no input with either hard-core calculation and never uses $A_9$. Since $\CSR_8-\frac14J_9$ of
the independent re-derivation agrees with the compendium (Sec.~\ref{sec:indep}), the low independent value lies in its
$J_9$, and within $J_9$ in $A_9[\varepsilon^0\ell^1]$, the one piece that deviates. Whether that deviation is a
fluctuation with an underestimated error or a systematic error is not known. At $\sci{5}{-4}$ relative,
Eq.~\eqref{eq:C8integ} supports the leading digits of the record but can neither identify nor certify the rational
coefficients of Eq.~\eqref{eq:C8cand}.

\paragraph{What the integrable point tests directly.} Of the $s=1$ coefficients of the flow up to the order that
$\CSR_8$ needs, only two combinations are invariant under changes of coupling, and both agree with the compendium; one
of them, $b+\frac f4-\frac a2-\frac c8-A$, vanishes identically with the closed forms of Sec.~\ref{sec:inputs1}
(Appendix~\ref{app:integ:gauge}). All the others, $J_9$ among them, are coordinate choices of the hard-core scheme, so
the $s=1$ solution alone cannot test $A_9$.

\subsection{Review of the compendium}\label{sec:review1}
The compendium was reviewed on 2 October 2026 by a multi-agent workflow (ten review dimensions and adversarial
verification). It found no error in any coefficient or proof claim: the conversions and invariants of $\CSR_3$ to
$\CSR_7$ were recomputed from their stated formulas, the $\CSR_5$ to $\CSR_7$ certificates were found sound (mutation
tests: altered rows are rejected), and every task of the archive's reproduction script passed on Linux, including a
fresh evaluation of the exact $P_6$ expression that reproduces it byte for byte. These tasks replay the archived
expressions, certificates and records; none regenerates the raw kink integrals, which is impossible because 21
ordered-integration stage programs are missing from the archive (Appendix~\ref{app:cannot}). The other defects concern
the archive as well: two files whose names differ only by case (they collide on case-insensitive file systems), a
reconstruction step that does not recreate two historical paths, and bibliographic errors. None affects a quoted
number.

\subsection{Caveats}\label{sec:d1caveats}
\begin{itemize}
\item \emph{Two-coupling closure.} The expansion assumes that the renormalization of the dual involves only the
 couplings $g$ and $b$, and that the prescription of Sec.~\ref{sec:dual} extracts their $\beta$ functions. Two
 unrefereed arguments, given in full in Appendix~\ref{app:clos}, support both formally, coefficientwise in $g$, $h$
 and $\db$ and so at every order. The first (Appendix~\ref{app:clos:closure}) argues that, for a symmetry-preserving
 local subtraction, every bulk ultraviolet counterterm lies in $\mathrm{span}\{\mathbf 1,\Og,\Oh\}$, so no third bulk coupling is needed ($h$ and $b$ are two coordinates on one
 direction). A literal two-operator truncation would nevertheless be wrong: finite irrelevant terms, such as
 ${:}\chi_B^2{:}$ from $\Og\Og$, feed the marginal $\beta$ functions, and the full multi-kink integrals of the archive
 include them (Appendices~\ref{app:clos:counterexample} and~\ref{app:clos:notmean}). The second
 (Appendix~\ref{app:clos:prescription}) argues that the prescription extracts the local $\beta$ functions. A
 divergence at a fixed endpoint of the interval can only multiply the identity, so it cancels when the responses are
 divided by the vacuum moments (Appendix~\ref{app:clos:endpoint}), while the finite endpoint-source terms remain
 and are those of the archive's recurrence (Eq.~\eqref{eq:dualA:recurrence}, Appendix~\ref{app:clos:uniqueness}); the two-response system then has the local $\beta$ functions as its solution, up to a finite change of
 $(g,b)$, which leaves $T$ and $D$, and so the series, unchanged. Neither argument checks the historical high-order
 integrations, which are tested only numerically, by the independent re-derivation and by the magnetization scheme of
 Sec.~\ref{sec:integ} (the certificates concern the later reduction to zeta values), nor the identification of the dual
 with the chain (Appendix~\ref{app:clos:scope}). The magnetization scheme, with a different regulator and different
 responses but the same two-coupling description, reproduces $\CSR_3$ to $\CSR_7$ to $\sci{7}{-26}$ and the record of
 $\CSR_8$ within its uncertainty: numerical support for the prescription, not a test of the closure. Independently, the NRG (next item) confirms $\CSR_3$ to $\CSR_5$ without using the dual.
\item \emph{The NRG test in numbers.} The numerical renormalization group (NRG) of the accompanying
 paper~\cite{Paper,NRGnote} does not use the dual. At $\Ds=0.01$, the residual of its $\De$ after each added order
 of the series is $\sci{4.7}{-3}$, $\sci{9.8}{-4}$, $\sci{2.4}{-4}$ and $\sci{5.3}{-5}$ through $\CSR_2$ to $\CSR_5$;
 the NRG error there is $\sci{4.7}{-5}$ (a linear sum of chain-depth, truncation and discretization parts, not a
 standard deviation), so the last residual is 1.1 times that error. The test is of the partial sums, and with them of
 $\CSR_3$ to $\CSR_5$; the term of $\CSR_6$ ($\sci{5.9}{-5}$) is of the size of the NRG error, and the higher terms
 are below it.
\item \emph{Level crossing.} $\De$ is identified with $\Delta_-$, assuming that no other level crosses it.
\item \emph{$\CSR_8$ is not proved.} Its record is printed to about 85 digits and cross-checked to about 57 decimal
 places (Sec.~\ref{sec:C8}), and the zeta-value form is conjectured. The magnetization-scheme value agrees with the
 record to $\sci{2.5}{-5}$, within its empirical uncertainty of $0.13$ (Sec.~\ref{sec:integ}); the independent value
 lies 1.9 to 4.0 standard errors below, with the difference traced to one input integral.
\item \emph{Numerical errors are not bounds.} The compendium's high-precision evaluations are checked by repeated runs
 with different settings (Appendix~\ref{app:dualC:numerics}), not by interval arithmetic; the independent errors are statistical, and the
 magnetization-scheme uncertainty is an empirical numerical estimate.
\item \emph{Side results.} Appendices~\ref{app:dualC:conjectures} to~\ref{app:dualC:logs} also record general-formula
 tests and results derived from published data of the dual (the shadow relation, the logarithmic endpoint). They hold
 under the regularity assumptions stated there, are not quoted by the paper and are not results of Table~\ref{tab:main}.
\end{itemize}

%====================================================================
\section{Two dimensions: \texorpdfstring{$c_2$}{c2}}\label{sec:d2}
%====================================================================
\subsection{The crossover theory}
In $d=2$ the long-range fixed point meets the short-range Ising CFT at $s_*=2-2\Ds^{\rm SR}=7/4$. Near $s_*$,
Behan, Rastelli, Rychkov and Zan (BRRZ) describe it as the Ising CFT coupled to a generalized free scalar $\chi$ of
dimension $\Delta_\chi=15/8-\delta$ through $g\int\sigma\chi$~\cite{Behan:2017dwr,Behan:2017emf}; the relevant
thermal coupling is tuned to the critical surface. With unit-normalized fields, $\Delta_\sigma^{\rm SR}=1/8$,
$\De^{\rm SR}=1$, $c=1/2$, $C_{\sigma\sigma\epsilon}=1/2$ and $C_{\epsilon\epsilon\epsilon}=0$. The exact Ising
correlators $\langle\sigma_1\cdots\sigma_{2n}\rangle$ and $\langle\epsilon\epsilon\sigma_1\cdots\sigma_{2n}\rangle$
needed here follow from a neutral-charge (bosonization) formula~\cite{Ardonne:2010zz}. Appendix~\ref{app:cform} derives
every step of this section in full: the conventions in Appendix~\ref{app:cform:conv} and the correlators in
Appendix~\ref{app:cform:correlators}.

\emph{First order.} With $\AC$ the analytic continuation in the exponents (defined, with the circular finite part
$\FP$ used below, in Appendix~\ref{app:cform:conv}), the first-order energy shift is proportional to
\begin{equation}
\begin{aligned}
 I_2(\delta)&=\AC\!\int_{\mathbb C}|1-z|^{-4+2\delta}\Bigl(\frac{|1+z|^2}{4|z|}-1\Bigr)\dd^2z\\
 &=\frac{\pi \delta^2}{4(1-\delta)^2}\frac{\Gamma(\delta)\Gamma(\frac12-\delta)}{\Gamma(1-\delta)\Gamma(\frac12+\delta)}
 =\frac\pi4\bigl[\delta+(2+4\log2)\delta^2+\cdots\bigr],
\end{aligned}
\end{equation}
which vanishes at $\delta=0$ (Appendix~\ref{app:cform:lower}). This reproduces the known cancellation $c_1=0$~\cite{Behan:2017emf}, and the expansion
starts at $c_2\delta^2$.

\subsection{Renormalization-group reduction}\label{sec:rg2}
In pole-only analytic minimal subtraction (Appendix~\ref{app:cform:rg}), $\beta(g)=-\delta g+b_0g^3+O(g^5)$ and
$\gamma_{\mathcal E}=\hMS g^4+O(g^6)$, so $g_*^2=\delta/b_0$ and
\begin{equation}
 c_2=\frac{\hMS}{b_0^2}=-\frac{4z_{4,1}}{b_0^2},
\end{equation}
where $z_{4,1}$ is the coefficient of $g^4/\delta$ in $\log Z_{\mathcal E}$. The one-loop coefficient is exact:
\begin{equation}
 b_0=\frac{9\pi^2}{8}\Bigl[\frac{\Gamma(3/4)}{\Gamma(1/4)}\Bigr]^2=1.26840430893914254379\ldots,
\end{equation}
obtained from the four-spin correlator by a quadratic change of variables and the complex beta integral
(Appendix~\ref{app:cform:lower}); it agrees
with the BRRZ value $1.268404$. The split between a $g^2$ and a $g^4$ term is prescription dependent: if
$\gamma_{\mathcal E}=a_1\delta g^2+hg^4$, then $c_2=a_1/b_0+h/b_0^2$ is invariant under
$\mathcal E_R\to e^{\kappa g^2}\mathcal E_R$, which shifts $a_1$ by $2\kappa$ and $h$ by $-2\kappa b_0$. Minimal
subtraction has $a_1=0$ and $h=\hMS=c_2b_0^2\approx-6.54$.

\subsection{Localized source, elliptic cover and the finite contact}\label{sec:source2}
The fourth-order logarithm is extracted with the coupling smeared by a localized source $f(w)=e^{-|w|^2}$ about the
origin and with the two energy operators at $0$ and $\infty$. Write $S_4$ for the four-spin correlator
$\langle\sigma\sigma\sigma\sigma\rangle$ and $H_4$ for the same correlator with the two energies inserted, and $S_{ij}$,
$H_{ij}$ for the corresponding two-spin functions of spins $i$ and $j$. The fourth-order term of the logarithm of the
energy two-point function has the connected kernel $\mathcal F=H_4-S_4-H_{12}H_{34}+S_{12}S_{34}$ for the GFF pairing
$(12)(34)$, and the three equal pairings give $3/4!=1/8$; its integral over the four spin positions is $A_4(\delta)$.
The local coupling counterterm $b_0g^3f^3/(2\delta)$ contributes at the same order through the second-order integral
$A_{a,b}(\delta)$, the two-spin integral with Gaussian weights $e^{-a|w_1|^2-b|w_2|^2}$: at fourth order it gives
$b_0A_{3,1}(\delta)/\delta$ (Appendix~\ref{app:cform:source}), so that
\begin{equation}
 z_{4,1}=\Res A_4+b_0B_{31},\qquad B_{31}=A_{3,1}(0)=\frac{\pi^2}{96}\bigl(12-\pi\sqrt3-9\log3\bigr)=-0.34223970692408023655\ldots.
\end{equation}
A M\"obius map returns the energies to $0$ and $\infty$ (Appendix~\ref{app:cform:ellipticmap}); the integral over the
overall scale gives $A_4(\delta)=\frac\pi8\Gamma(2\delta)J(\delta)$, where $J(\delta)$ is the remaining integral over the
four-spin cross ratio $z$ and the two energy positions, weighted by $W^{2\delta}$, with $W$ the Gaussian source factor
carried through the map. The energy positions are uniformized by the elliptic curve $E_z$, the double cover of the $t$-plane branched at
$0$, $z$, $1$ and $\infty$, on which the
Ising kernel becomes an algebraic expression in Weierstrass functions, and the torus integration defines the
elliptic finite part
\begin{equation}
 \cI=\FP\!\int_{\mathbb C}\dd^2z\,|z|^{-15/4}\frac{\mathcal A(z)^2}{4}
 \Bigl[w_+|c_+-2q|^2+w_-|c_--2q|^2-w_0|c_0-2q|^2\Bigr],
 \label{eq:I}
\end{equation}
with the torus area $\mathcal A=16\,{\rm Re}[\overline{K(z)}K(1-z)]$, the circular mean
$q=E(z)/(4K(z))-(2-z)/12-\pi\overline{K(z)}/(K(z)\mathcal A)$, the Ising block weights
$w_\pm=|1\pm\sqrt{1-z}|/(2|z|^{1/4}|1-z|^{1/4})$, $w_0=|z|^{-1/4}$, and $c_\pm=(z-2\pm3\sqrt{1-z})/6$,
$c_0=(2-z)/12$. In this notation the four-spin function at cross ratio $z$ is $S_4(z)=w_++w_-$.

The flat torus finite part is not yet the limit of the source-weighted integral $J(\delta)$. When both energies
approach the same branch point, the first derivative of the source factor $W^{2\delta}$ survives against a vanishing
zeroth-order angular residue and leaves a finite contact $A_{\rm br}$ (Appendix~\ref{app:cform:contact}). It is
related exactly to the lower Gaussian integral,
\begin{equation}
 A_{\rm br}=4B_{31}-\frac{\pi^2}{4},
\end{equation}
and at fixed modulus $z$ the regulated energy integral equals the torus integrand plus $4A_{\rm br}[S_4(z)-|z|^{-1/4}]$.
The energy-primary term of the shrinking spin pair cancels exactly in the connected logarithm, and no scalar with a
logarithmic threshold remains at the modulus boundaries (Appendix~\ref{app:cform:qualification}).

\subsection{Assembly and the matching term}\label{sec:assembly2}
Integrating the fixed-modulus comparison over $z$, with $\AC\!\int|z|^{-15/4}[S_4-|z|^{-1/4}]\dd^2z=-b_0/\pi$, gives
$J(0)=\cI-16b_0B_{31}/\pi+\pi b_0$ (Appendix~\ref{app:cform:assembly}), so $B_{31}$ cancels and
\begin{equation}
 z_{4,1}=\frac{\pi\cI+\pi^2b_0}{16},\qquad
 \boxed{c_2=-\frac{\pi\cI}{4b_0^2}-\frac{\pi^2}{4b_0}.}
 \label{eq:c2I}
\end{equation}
The finite matching term $-\pi^2/(4b_0)=-1.94527965797909034563\ldots$ is comparable to the period term
$-\pi\cI/(4b_0^2)=-2.12099049015343120711\ldots$. The study's own source-contact checks (Appendix~\ref{app:cform:numerics}) and its review showed that
it does not depend on the shape of the source (analytically, and numerically for two other shapes), and that it equals $a_1/b_0$, where $a_1=-\pi^2/4$
is the $O(\delta)$ slope of the leading-order formula of BRRZ with the $\delta$-deformed $\chi$ propagator; the finite
redefinition with $\kappa=-\pi^2/8$ maps minimal subtraction to that scheme. The cylinder calculation of
Sec.~\ref{sec:cyl} finds the same $a_1=-\pi^2/4$ independently (Appendix~\ref{app:cyl:second}). Without the matching term $c_2$ would be about $-2.12$,
which the Monte Carlo of Sec.~\ref{sec:mc} excludes by about 100 standard errors.

\subsection{Closed form of the elliptic finite part}\label{sec:closed2}
The integral $\cI$ is reduced by a Landen substitution and a single-valued period pairing to periods of $K(t)^2$,
$K(t)K(1-t)$ and $K(1-t)^2$ over $0<t<1$, then by contour reflection to symmetric elliptic moments, and finally, through summand-level rational telescoping certificates with
their nonzero endpoint terms, to three convergent hypergeometric values (Appendices~\ref{app:cform:moments} to~\ref{app:cform:telescoping}). With $R_\Gamma=\Gamma(\frac14)/\Gamma(\frac34)$,
so that $b_0=9\pi^2/(8R_\Gamma^2)$, and
\begin{equation}
 \mathsf P=\pFq{-\frac12,-\frac12,\frac32,\frac34}{1,2,\frac14}{1},\quad
 \mathsf Q=\pFq{\frac12,\frac32,\frac32,\frac94}{3,3,\frac74}{1},\quad
 \mathsf T=\pFq{\frac32,-\frac14,\frac34,\frac34}{\frac14,\frac94,\frac94}{1},
\end{equation}
the coefficient is
\begin{equation}
 c_2=\frac{16}{R_\Gamma^2}-\frac{R_\Gamma^2}{9}
 -\frac{\pi^2R_\Gamma^2}{16}\Bigl(\mathsf P-\frac{5R_\Gamma^2}{384}\mathsf Q\Bigr)
 \Bigl(\mathsf P-\frac{5R_\Gamma^2}{128}\mathsf Q-\frac{512}{25\pi R_\Gamma^2}\mathsf T\Bigr),
 \label{eq:c2closed}
\end{equation}
which evaluates to
\begin{equation}
 c_2=-4.0662701481325215527400276085728764875299011411866\ldots .
 \label{eq:c2num}
\end{equation}
Each series has $\sum b_j-\sum a_i=2$, so its terms decay as $n^{-3}$, it converges absolutely and its unaccelerated
remainder is $O(N^{-2})$. The three constants need no analytic continuation, and the expression contains no unevaluated
integral, parameter derivative or unresolved subtraction constant. No reduction to
Gamma values alone has been established (an integer-relation search found none). The constants are
\begin{center}\small
\begin{tabular}{@{}lr@{\qquad}lr@{}}
\toprule
$R_\Gamma$ & $2.95867511918863889231$ & $b_0$ & $1.26840430893914254379$\\
$\mathsf P$ & $1.62562945952539344424$ & $B_{31}$ & $-0.34223970692408023655$\\
$\mathsf Q$ & $1.28812240964200676062$ & $\cI$ & $4.34474465231888563075$\\
$\mathsf T$ & $0.76385352476101335168$ & $z_{4,1}$ & $1.63550416445718870488$\\
\bottomrule
\end{tabular}
\end{center}
\paragraph{Checks of the evaluation} (Appendix~\ref{app:cform:numerics}). The six-function and three-function expressions and a half-interval Frobenius
series of the moments agree to $\sci{2.53}{-80}$ (the 80-digit arithmetic floor). An endpoint-subtracted quadrature with
Carlson's integrals, which uses no Frobenius coefficients, agrees to $\sci{1.40}{-61}$, a difference set by the
finite upper limit of its last integral (doubling that limit lowers it to $\sci{7.4}{-77}$). Exact tail inequalities for the three
series give $-4.066270148133018<c_2<-4.066270148131205$ (floating-point propagation, not a certified rounding). The
base study of the cylinder calculation (Sec.~\ref{sec:cyl}) contains an integer-interval evaluation that certifies the
value of the closed expression to 51 decimal places (Appendix~\ref{app:cyl:closedchecks}; its code is not in the
release, Appendix~\ref{app:cannot}). In the review, $\cI$ was also evaluated directly from its defining integral, without
the reduction, to 51 digits in one implementation and to $\sci{5}{-15}$ by $z$-plane integration in another, and every
intermediate identity of the chain was checked at 30 to 80 digits.

\subsection{The step that is not proved}\label{sec:unproved}
Equation~\eqref{eq:c2I} needs one step that is not proved: exchanging the limit $\delta\to0$ with the nested
subtracted integrals of the global comparison. In each collision sector the comparison subtracts the local OPE and
Taylor terms, restores powers by their Mellin continuation and spin terms by their circular finite parts, and then
takes the regulator limit of the subtracted remainder. That exchange needs an integrable bound on the remainder,
uniform near $\delta=0$, including the degenerating shapes where sectors intersect. The study classifies the channels
and gives a general subtraction argument (Appendix~\ref{app:cform:qualification}), but this step is incomplete in two
respects: not every sectorwise
subtraction is written out, and the uniform bound is not proved. We count both as the one unproved step of the
abstract and Table~\ref{tab:main}. Its weakest point is where a same-branch contact meets the corner $z\to0$.

The review of the study (summarized in \file{docs/c2_closed_d2/REVIEWS.md}) looked for each mechanism by which an
$O(1)$ term could arise there, and found none. There is no scalar with a logarithmic threshold at dimension two (the energy
primary cancels), no dependence of the leading regulator on $\arg z$ at the corners, and no term above threshold from
subleading regulator terms combined with spin-one terms. Two numerical tests bear on the step: the Monte Carlo of
Sec.~\ref{sec:mc} integrates the full regulated integral, and the cylinder calculation of Sec.~\ref{sec:cyl} uses a
different regularization. The identification of Eq.~\eqref{eq:c2closed} with the physical coefficient is therefore a
formal perturbative result in the critically tuned analytic minimal-subtraction prescription, conditional on this
step; the evaluation of $\cI$ itself is proved.

\subsection{Independent Monte Carlo of the regulated integral}\label{sec:mc}
The review integrated the complete regulated eight-dimensional integral over the four insertion positions by Monte
Carlo, with the exact Ising correlators built by bosonization and no M\"obius map, elliptic cover, contact derivation
or sector subtraction. Because it integrates the whole regulated integral, it tests the global exchange step
numerically. Three configurations with different channel weights were run, two on the symmetry cluster and one on a
personal desktop, with $\sci{4.26}{10}$ samples in all (142 chunks, about 135 core-hours); all three enter the
result. Their inverse-variance combination gives
\begin{center}\small
\begin{tabular}{@{}lll@{}}
\toprule
quantity & Monte Carlo & closed-form route\\
\midrule
$J(0)$ & $10.559\pm0.039$ & $10.5404$\\
$z_{4,1}$ & $1.6391\pm0.0076$ & $1.635504$\\
$c_2$ & $-4.075\pm0.019$ & $-4.066270$\\
\bottomrule
\end{tabular}
\end{center}
The errors are statistical (one standard error). The recombination gives $c_2=-4.0752\pm0.0189$, quoted as
$-4.075\pm0.019$; it differs from the closed form by $-0.009$, $0.5$ standard errors. The two configurations run on
the symmetry cluster alone give $-4.0724\pm0.0197$. The three configurations agree on the integral itself ($\chi^2=1.29$
for two degrees of freedom). A source-shape control, the difference between a quartic and a Gaussian source,
reproduces its prediction when the three are combined ($16.185\pm0.053$ against $16.150$), but the three
configurations disagree among themselves on it ($\chi^2=7.6$ for two degrees of freedom). The control does not enter
$c_2$; the disagreement is unexplained and left open. The same code reproduces $b_0$, the vanishing first-order shift,
$B_{31}$, and $z_{2,1}=-b_0/2$, which gives $\Ds=\frac18+\delta$.

\subsection{Fourth-order perturbation theory on the cylinder}\label{sec:cyl}
\paragraph{Setup.} A separate calculation (its method in full in Appendix~\ref{app:cyl}) does Rayleigh--Schr\"odinger
perturbation theory for the Ising CFT times the GFF on the radial cylinder, with $H_0=D_{\rm Ising}+D_\chi$ and $V=\int_0^{2\pi}\dd\theta\,\sigma\chi$. It shares
with the closed-form route only the crossover description, the exact Ising data and $b_0$; it uses none of $\cI$,
the hypergeometric form, the branch contact or the proposed value. The observable is
$X=2E_{\mathcal E}-E_d-E_0$, with $\ket d=L_{-1}\ket{\mathcal E}$, which equals $\De-1$ at a conformal fixed point and
removes the direct fourth-order local terms (Appendix~\ref{app:cyl:counterterms}). In the scheme defined by the cylinder Hamiltonian with its locally
derived counterterms, $H=H_0+gV+\sum_{k\ge2}g^kW_k$, write $X=a_1\delta g^2+hg^4+\cdots$. With $h_0$ the value of $h$
once the regulator is removed, the invariant combination of Sec.~\ref{sec:rg2} gives
\begin{equation}
 c_2=\frac{a_1}{b_0}+\frac{h_0}{b_0^2};
\end{equation}
here $h_0\approx-3.41$, while minimal subtraction has $a_1=0$ and $\hMS\approx-6.54$. The second-order coefficient is
an exact meromorphic continuation (Appendix~\ref{app:cyl:second}),
$a(\delta)=\pi^2\Gamma(1-\frac\delta2)^2\Gamma(\delta-1)/[\Gamma(\frac\delta2)^2\Gamma(2-\delta)]=-\frac{\pi^2}4\delta+O(\delta^2)$,
so $a_1=-\pi^2/4$ and $a_1/b_0=-1.94527965797909\ldots$. The fourth-order coefficient $h$ is computed at $\delta=0$
with a family of smooth positive oscillator regulators of order $r=1,2,3$ and length $\lambda$, locally derived thermal
and spin-two counterterms, and exact chiral conformal blocks as spectral generating functions
(Appendices~\ref{app:cyl:regulator}, \ref{app:cyl:counterterms} and~\ref{app:cyl:blocks}). The spectral sums run
over intermediate states labelled by three level indices and are truncated at a cube $0\le n,m,\ell<N$, with
$N\le1024$; they are extrapolated first in $N$ at fixed $(r,\lambda)$ and then in $\lambda\to0$. No fit has the
physical $\delta$ as its variable. As a lower-order check, the renormalized magnetic gap gives
$[g^2](E_\sigma-E_0)=1.268404004672631$ at $r=3$, $\lambda=1/32$, which differs from $b_0$ by $\sci{-3.04}{-7}$; it
was not used to tune the counterterms (Appendix~\ref{app:cyl:magnetic}).

\paragraph{Continuum analysis.} The base study, with $N\le384$, gave $c_2\simeq-4.07$ (representative extrapolations
$-4.0669$, $-4.0680$ and $-4.0713$ for $r=1,2,3$). The continuum analysis uses 25 fixed-regulator points $(r,\lambda)$,
for each the record with the largest $N$. Fourteen are new production records with $N=768$ or $1024$, computed on SDSC
Expanse and the symmetry cluster ($\lambda=0.05$ to $0.175$). The other eleven are archived records of the base study,
computed by GPT-6 astra (Sec.~\ref{sec:prov}): $N=256$ at $\lambda=0.125$ to $0.25$ for $r=1$ and $0.175$ to $0.35$ for
$r=2$, and $N=384$, $256$, $256$ and $128$ at $\lambda=0.2$, $0.25$, $0.35$ and $0.5$ for $r=3$. Their spectral sums are
taken as archived and extrapolated under the same rules; their auxiliary sums were recomputed on the symmetry cluster at the common cutoffs of all points
(below). Nine of them are among the 23 points of the joint fit. The rules were all fixed in code before the final
numbers were computed and are copied into every output:
\begin{itemize}
\item \emph{Large-$N$ limit at fixed $(r,\lambda)$.} The leading large-$N$ correction, a power $N^{-7/4}$ from the faces
 of the truncation cube with an analytically derived coefficient (Appendix~\ref{app:cyl:facetail}), is subtracted; the remaining tail is fitted with
 sets of powers and exponentials fixed in advance, on windows that are fixed fractions of $N$, and the median of the
 selected fits is $h(\lambda)$.
\item \emph{Checks of that limit.} The same extrapolation is repeated on the sums truncated at $N'=N/2,\dots,N$, and
 records with different $N$ are compared. Points with $\lambda N$ below a minimum (50, 60 and 75 for $r=1,2,3$) are
 flagged, since the cutoff must resolve the regulator; the five flagged points, all archived, stay in the fits, and
 only records above the minimum enter the cross-$N$ comparison. The spectral coefficients, computed by discrete Cauchy integrals on a grid of points on a
 torus (Appendix~\ref{app:cyl:blocks}), are extracted again on a second grid with a different radius, for every record, and roundoff is estimated. Three
 auxiliary truncations (of the mixed thermal sums, of the collapsed sums and of the two-index leading-tail
 calculation; Appendices~\ref{app:cyl:collapsed} and~\ref{app:cyl:facetail}) are set to common values for all points
 (16384, 800000 and 8192), and their shifts under a change (to 400000 and 2048 for the last two) are propagated.
\item \emph{Regulator limit} (Appendix~\ref{app:cyl:tails}). Per order, $h(\lambda)=h_0+k_1\lambda^2+k_2\lambda^2\log\lambda+k_3\lambda^3$, the
 preferred model fixed in advance, on $\lambda\le0.25$ for $r=1,2$ and $\lambda\le0.35$ for $r=3$ (the smallest
 $\lambda$ is 0.05, 0.0625 and 0.1; the points $(2,0.35)$ and $(3,0.5)$ lie outside these windows). Alternative models
 and windows are admitted by fixed rules, and the regulator-model error is the largest deviation of an admissible fit.
\item \emph{Choice of estimator.} A rule fixed in advance chooses between the per-order results and a joint fit. Here
 the three orders are consistent (each pairwise difference of $h_0$, at most $\sci{1.9}{-3}$, is below the sum of
 the two errors), and the joint fit with a shared $h_0$ is accepted (23 of the 25
 fixed-regulator points enter it).
\end{itemize}
The closed form was never an input, a prior or a selection criterion. The per-order and joint results are
\begin{center}\small
\begin{tabular}{@{}lllll@{}}
\toprule
& points & $h_0$ & $c_2$ & error in $c_2$\\
\midrule
$r=1$ & 9 & $-3.414150$ & $-4.06739$ & $0.01006$\\
$r=2$ & 7 & $-3.412774$ & $-4.06653$ & $0.00417$\\
$r=3$ & 7 & $-3.412269$ & $-4.06622$ & $0.00346$\\
joint & 23 & $-3.413609\pm0.003794$ & $-4.0670498$ & $0.0023580$\\
\bottomrule
\end{tabular}
\end{center}
so that
\begin{equation}
 c_2=-4.0670498\pm0.0023580,\qquad\text{quoted }-4.0670(24).
\end{equation}
The error is a linear sum of systematic numerical diagnostics, not a standard deviation and not a rigorous bound.
Its components in $c_2$ are: the choice of regulator model $\sci{1.52}{-3}$ (64\%); the large-$N$ extrapolation of the
points, propagated through the fit, $\sci{4.04}{-4}$ (17\%); the cutoff and model spread of the mixed thermal sums
$\sci{2.14}{-4}$ (9\%); roundoff $\sci{8.7}{-5}$ (4\%); coefficient extraction, from the comparison of the two Cauchy
grids, $\sci{7.5}{-5}$ (3\%); the cutoff of the collapsed sums $\sci{6.2}{-5}$ (3\%); and the cutoff of the
leading-tail calculation $\sci{2.5}{-9}$. The difference from the closed form is $-\sci{7.8}{-4}$, or
$|{-4.0670498}+4.0662701|/4.06627=\sci{1.9}{-4}$ relative, 0.33 of the quoted error.

\paragraph{Caveats, as the paper states them.}
\begin{itemize}
\item The closed form was known before this analysis, but it enters the analysis in no way.
\item The error rule was revised after the first result had been seen. The first complete analysis gave the same
 central value with $E(c_2)=\pm2.21$ (absolute, in $c_2$). Of that, 99.9\% was the coefficient-extraction term,
 bounded by a worst case over sign patterns: the weights $\sum_n|v_n|$ of the large-$N$ fits on the raw sums, up to
 about $\sci{2.4}{5}$, times differences between two coefficient grids of $\sci{5}{-7}$ to $\sci{7}{-5}$, measured at six of the 25 fit points and transported
 to the others. The revised rule, written after that result and approved in the Codex review, measures the grid
 difference directly at all 25 points (20 new extraction-check records), carries its smooth part with its measured
 sign and keeps only the rough residual at worst case; it gives $\sci{7.5}{-5}$ in $c_2$. Every other component is
 identical in both analyses ($\sci{2.28}{-3}$ in all). The quoted error therefore fell from $\pm2.21$ to $\pm0.00236$,
 about 940-fold, entirely through the replaced extraction bound, and the central value is unchanged by construction.
 The superseded result is archived.
\item The two $N=1024$ coefficient grids were rebuilt on the symmetry cluster, and the last nine extraction-check
 computations ran there, because SDSC Expanse's file system stopped serving two small files of the grids built there
 (an author-approved change of the plan). The grids are a deterministic function of their parameters; a pilot at
 $N=512$ built on both machines agreed to $\sci{6.9}{-11}$ in the raw sums.
\item A difference between two Cauchy grids measures extraction sensitivity; it does not bound an error that both
 share. The regulator extrapolation dominates the error and is not rigorously bounded.
\end{itemize}

\subsection{Summary for \texorpdfstring{$d=2$}{d=2}}
The closed form is $c_2=-4.06627014813\ldots$, conditional on one unproved step. The Monte Carlo of the regulated
integral gives $-4.075\pm0.019$, one standard error ($0.5\%$ relative), and differs from it by $0.5$ standard errors.
The cylinder Hamiltonian gives $-4.0670(24)$, where the error is a linear sum of diagnostics with no probability
meaning, and differs from it by $\sci{1.9}{-4}$ relative, 0.33 of that error. The Monte Carlo integrates the
full regulated integral without the exchange step, and the cylinder calculation uses a different regularization; both
test the unproved step numerically. No published $O(\delta^2)$ value exists in two dimensions to
compare with.

%====================================================================
\section{Three dimensions: \texorpdfstring{$c_1$}{c1}}\label{sec:d3}
%====================================================================
\subsection{The first-order coefficient as a ratio of integrated correlators}
In $d=3$ the same crossover theory applies with $\Delta_\chi=3-\Ds^{\rm SR}-\delta$. With
$\beta(g)=-\delta g+b_0g^3$ ($b_0$ again the one-loop coefficient, now in $d=3$) and $\gamma_{\mathcal E}=a_0g^2$,
\begin{equation}
 c_1=\frac{a_0}{b_0},\qquad a_0=-2\pi I_\epsilon,\quad b_0=-2\pi I_\sigma,
\end{equation}
where, with one external state at infinity,
\begin{equation}
\begin{gathered}
 I_O=\AC\!\int_{\mathbb R^3}\dd^3y\,|e-y|^{-2(3-\Ds^{\rm SR})}\bigl[H_O(y)-|e-y|^{-2\Ds^{\rm SR}}\bigr],\\
 H_O(y)=\lim_{R\to\infty}R^{2\Delta_O}\langle O(0)\sigma(y)\sigma(e)O(Re)\rangle,\qquad O\in\{\sigma,\epsilon\},
\end{gathered}
\end{equation}
with $e$ a unit vector and $\AC$ the analytic continuation through the power-divergent short-distance terms. The
denominator comes from the protected spin dimension, $a_\sigma=b_0$ (Appendices~\ref{app:cone:cpt}
and~\ref{app:cone:functional}). Both integrals are evaluated on a folded
region with conformal-block sums of the 3d Ising CFT: 60 primaries from the spectrum of
Simmons-Duffin~\cite{Simmons-Duffin:2016wlq} (the stable entries with $\Delta<24$), with $\Ds^{\rm SR}$, $\De^{\rm SR}$
and the stress-tensor couplings from Ref.~\cite{Chang:2024whx} and the Ward identity, radial and small-$z$ block
recurrences, and analytic continuation of the power-divergent terms. Appendix~\ref{app:cone} derives the folding
(Appendix~\ref{app:cone:folding}), the blocks and the integration (Appendices~\ref{app:cone:blocks}
and~\ref{app:cone:recurrences}) and lists the retained spectrum (Table~\ref{tab:cone:spectrum}).

\subsection{Retained-spectrum value}\label{sec:retained}
With the retained 60 primaries,
\begin{equation}
 a_0=3.3474594916917857,\qquad b_0=12.254914105481044,\qquad c_1=0.27315242382601645,
\end{equation}
quoted as $0.2732$. The long strings are numerical regression targets, not physical digits: the uncertainty of $c_1$ is
set by the omitted spectrum, through the estimated range of Sec.~\ref{sec:range} (the three meanings of precision
are separated in Appendix~\ref{app:cone:precision}). The numbers agree with
BRRZ~\cite{Behan:2017dwr,Behan:2017emf} component by component: $I_\sigma=-b_0/(2\pi)=-1.950430$ against $-1.950(5)$,
$b_0=12.2549$ against $12.26(3)$, $a_0=3.3475$ against $3.3(5)$, and $c_1$ against their estimate $0.27(4)$; later
bootstrap studies found $0.269$~\cite{Behan:2018hfx} and $0.27$~\cite{Behan:2023ile}.

\emph{Implementations.} The package's full workflow, rerun from scratch, reproduces every shipped number bit for bit in
the review's run and on the Intel nodes of the symmetry cluster, and within the replay's tolerances on its AMD nodes
($\sci{5}{-10}$ for $a_0$, $b_0$ and $c_1$). The review of the package wrote two independent implementations, A and B,
each with its own conformal blocks, quadrature and analytic continuation, neither importing the package; they agree on
$c_1$ to $10^{-12}$, including all 60 per-family contributions. Both are included in \file{review_indep/} of the package
and rerun by its replay on the retained spectrum: on the symmetry cluster (Intel and AMD nodes) A differs from the
package's $c_1$ by at most $\sci{6}{-14}$, B by at most $\sci{5}{-13}$, and A from B by at most $\sci{4}{-13}$, all
absolute (the replay tolerance is $10^{-12}$). For $a_0$ and $b_0$ the agreement is only at the level of
$\sci{2}{-12}$ relative, so the claim of $10^{-12}$ is made for $c_1$ alone. The package's own two integrated forms
(\texttt{blocks} and \texttt{kernels}) agree to about $\sci{1.5}{-13}$ but share the folded domain, so they are a
consistency check, not an independent one. The blocks satisfy the Casimir equation and agree with the
Zamolodchikov-type recursion of implementation B to $\sci{6}{-13}$; the package's own block checks, against a small-$z$
series and a closed form on the diagonal, are in Appendix~\ref{app:cone:blocks}.

\subsection{Estimated range}\label{sec:range}
The omitted spectrum is controlled only under explicit hypotheses:
\begin{itemize}
\item every primary missing from the retained list and exchanged in $\sigma\times\sigma$, $\sigma\times\epsilon$ or
 $\epsilon\times\epsilon$ has dimension at least 8 if it is $\mathbb Z_2$-even and at least 7 if it is
 $\mathbb Z_2$-odd (published spectra motivate these gaps but do not prove
 them~\cite{Simmons-Duffin:2016wlq,Henriksson:2022,Zhu:2022,Fardelli:2026});
\item each retained input lies within twice its quoted error (the input box). The errors of the higher families are
 spreads over spectra at fixed derivative order, not rigorous; of the five low inputs from Ref.~\cite{Chang:2024whx},
 $\Ds^{\rm SR}$ and $\De^{\rm SR}$ have rigorous errors there, the two OPE coefficients and $c$ do not
 (Appendix~\ref{app:cone:inputs});
\item the conditional bootstrap result $\lambda_{\epsilon\epsilon\epsilon'}=1.5362(12)$ of Ref.~\cite{Reehorst:2021},
 with its sparsity assumptions, applies;
\item the margins added to the linearized crossing constraints and to the integrated ratio, which cover nonlinear terms
 and sampled remainders, and the sampled allowances for the denominator ($|\Delta b_{\rm in}|\le0.001428$), for the
 reference bounds (a factor 1.001) and for the quadrature, are large enough.
\end{itemize}
The hypotheses are discussed in Appendices~\ref{app:cone:tails} (gaps and $\lambda_{\epsilon\epsilon\epsilon'}$),
\ref{app:cone:variations} (the input box), and \ref{app:cone:optimization} and~\ref{app:cone:validation} (the margins
and their sampled tests).
Under them, an identical-channel positive remainder, a Cauchy--Schwarz bound on the even mixed remainder, a sign-aware
integration of the odd remainder and an affine linear program over crossing constraints and inputs
(Appendices~\ref{app:cone:tails} to~\ref{app:cone:budget}) give the nominal interval $[0.2697959,\,0.2752135]$ (ten
constraint updates); doubling the crossing margins gives $[0.2692738,\,0.2754104]$ (the package's recorded run, eight
updates; ten give $[0.26927,\,0.27541]$). The
outward-rounded estimated range is
\begin{equation}
 c_1\in[0.2690,\,0.2760]\qquad(\text{conditional}),
\end{equation}
to be read next to the retained value $0.2732$, about $-0.0042$ and $+0.0028$ from it. It is neither a certified
bound nor a confidence interval. A weaker-gap scenario (6 and 6) gives $[0.26400,\,0.27842]$ (recorded run, eight updates; ten give
$[0.26420,\,0.27834]$).

\subsection{How conditional the range is}\label{sec:howcond}
The review found, and the revised package now shows, that the range is set by the assumed input box and the gap
hypotheses rather than by numerical precision.
\begin{itemize}
\item \emph{Feasibility} (Appendix~\ref{app:cone:optimization}). With all inputs at their central values, the final
 linear program is infeasible: 68 of
 its 552 constraints, all in the spin four-point channel, are violated. With the nominal program held fixed, the
 smallest feasible displacement is $0.574$ of the box half-width, that is, 1.15 quoted errors for the input that moves
 most (it is neither a common displacement of all inputs nor a significance). Underestimated spectrum errors and
 a missing contribution below dimension 8 are natural explanations, and the chosen numerical relaxation may contribute
 (Appendix~\ref{app:cone:optimization}); the data cannot tell which.
\item \emph{Box size} (Appendix~\ref{app:cone:sensitivity}). Replacing each quoted error $e_k$ by $me_k$, so that the inputs range over
 $|u_k-u_k^{(0)}|\le2me_k$, and rebuilding the program (sensitivities, crossing secants and margins; the
 $\lambda_{\epsilon\epsilon\epsilon'}$ window held fixed) gives, with ten constraint updates each,
\begin{center}\small
\begin{tabular}{@{}lllll@{}}
\toprule
$m$ & $\le0.6$ & $0.8$ & $1$ (nominal) & $2$\\
interval & infeasible & $[0.27027,0.27492]$ & $[0.26980,0.27521]$ & $[0.26708,0.27684]$\\
\bottomrule
\end{tabular}
\end{center}
 The infeasibility at $m=0.6$ ($\pm1.2$ quoted errors) does not contradict the 1.15 above: that minimization kept the
 nominal program fixed, whereas this scan rebuilds it.
\item \emph{Gaps} (Appendices~\ref{app:cone:tails} and~\ref{app:cone:sensitivity}). Mild relaxations of the gap hypotheses extend the interval outside the estimated range
 $[0.2690,0.2760]$, while every interval still contains the retained value:
\begin{center}\small
\begin{tabular}{@{}lllll@{}}
\toprule
(even, odd) gaps & $(8,6.5)$ & $(7,7)$ & $(8,6)$ & $(6,5.5)$\\
interval & $[0.26890,0.27541]$ & $[0.26891,0.27629]$ & $[0.26720,0.27576]$ & $[0.25839,0.27948]$\\
\bottomrule
\end{tabular}
\end{center}
 Lowering the odd gap mainly lowers the lower end, because the odd remainder enters with a definite sign; lowering the
 even gap, as in $(7,7)$, widens both ends. Unlisted odd primaries are estimated at about $7.3$ to $7.7$, only
 0.3 to 0.7 above the assumed gap: spin six at about $7.3$~\cite{Henriksson:2022} and $7.73(5)$~\cite{Fardelli:2026},
 spin three at $7.53(6)$~\cite{Fardelli:2026}.
\end{itemize}
The package contains its review revision of 8 October 2026, which adds these feasibility and sensitivity tests, the
missing derivations (now in Appendix~\ref{app:cone}) and explicit guards, with the numerical values unchanged; its
\texttt{replay-full} reruns the feasibility diagnostics and the gap, margin and box scenarios. The three-dimensional
Monte Carlo of the accompanying paper~\cite{Paper} cannot test $c_1$: its only accepted point, at $s=1.6$
($\delta=0.18$), is too far from $s_*$ to isolate the linear term.

%====================================================================
\section{Authorship and provenance}\label{sec:prov}
%====================================================================
The author of the project chose the problems, approved plans and production runs, and decided between the options
the agents presented; the author wrote none of the computations below. They were done as follows.
\begin{itemize}
\item \emph{$\CSR_1$, $\CSR_2$}: published by Benedetti, Lauria, Maz\'a\v c and van Vliet.
\item \emph{$\CSR_3$ to $\CSR_8$}: computed by another AI system, OpenAI GPT-6 astra (via ChatGPT), and supplied by the
 author as a research compendium with its derivations, code, exact expressions, certificates and numerical records. It
 was reviewed by Claude Code (Anthropic) on 2 October 2026 (Sec.~\ref{sec:review1}). The derivations of
 Appendices~\ref{app:dualA} to~\ref{app:dualC} are GPT-6 astra's.
\item \emph{Independent re-derivation of the inputs of $\CSR_5$ to $\CSR_8$}: Claude Code sessions, 2 to 7 October
 2026, with adversarial verification by separate agents for independence, statistics and mathematics; production on
 the symmetry cluster, part of the eighth-order work on a personal desktop.
\item \emph{The two-coupling closure and prescription arguments} (Appendix~\ref{app:clos}): written by OpenAI GPT-6
 astra (via ChatGPT) and supplied by the author on 15 October 2026, the second with its checker; not refereed. Claude
 Code assessed both at the author's request, checked the finite algebra of the first independently, reran the checker
 of the second and re-derived its beta recurrence independently (Appendix~\ref{app:clos:checks};
 \file{docs/closure_d1/REVIEWS.md}).
\item \emph{$c_2$ closed form}: derived by OpenAI GPT-6 astra (via ChatGPT), supplied by the author (the final-method
 study; its derivation is Appendix~\ref{app:cform}). Reviewed by Claude Code on 8 October 2026; the review wrote and
 ran the Monte Carlo of Sec.~\ref{sec:mc}.
\item \emph{Cylinder Hamiltonian for $c_2$}: the base study was written by OpenAI GPT-6 astra (via ChatGPT) and supplied
 by the author; its method is Appendix~\ref{app:cyl}. All four archives it consolidates are GPT-6 astra's: the final-method study, the independent formula,
 elliptic and local checks (with the integer-interval certificate of the closed expression), and the two Hamiltonian
 packages (the solver). Its 17 archived raw records ($N=128$ to 384) were computed by GPT-6 astra in its own
 environment; their machine, dates and software versions are not recorded. Eleven of them are points of the
 continuum analysis (nine of the 23 in its joint fit; Sec.~\ref{sec:cyl}), with their auxiliary sums recomputed on the
 symmetry cluster. Claude Code sessions wrote the cluster
 layer for SDSC Expanse and the symmetry cluster, ran the production and wrote the arbitrary-$N$ continuum analysis.
 OpenAI Codex reviewed the analysis plan in seven rounds, approving it at revision 7, and inspected the cluster layer
 and the analysis code twice each. The fixes made after the last inspection of each were not re-inspected (for the
 analysis, four hardening findings with no numerical change, by the author's choice).
\item \emph{$c_1$ in $d=3$}: the package was written by OpenAI GPT-6 astra (via ChatGPT), supplied by the author and
 dated 7 October 2026; its derivation is Appendix~\ref{app:cone}. It was reviewed by Claude Code on 7 October 2026; the two independent re-implementations of the
 review (\file{review_indep/}) were written by Claude Code review agents. The package contains its review revision of
 8 October 2026 (UTC; 7 October local time; a full rerun with explicit guards, the missing derivations and the sensitivity tests; numbers
 unchanged), which OpenAI GPT-6 astra made after our review and the author supplied; the four earlier archives it was
 built from are also GPT-6 astra's. Claude Code checked the revision on 8 October 2026.
\item \emph{Magnetization-scheme value of $\CSR_8$} (Sec.~\ref{sec:integ}, Appendix~\ref{app:integ}): Claude Code
 sessions, 14 to 17 October 2026, on the symmetry cluster (about $\sci{2.7}{4}$ core-hours, including pilots and failed jobs); OpenAI Codex reviewed the
 plan and inspected the code (Appendix~\ref{app:integ:disclosures}). Its code and records ship as
 \file{code/integ_d1/} (commit 8bd5a22), with \file{public.mk} and \file{public/} added by the same session on 18 October
 2026.
\item \emph{This note}: written by a Claude Code session on 15 October 2026 from the sources above, and revised the
 same day after an independent readability and precision review. On 16 October 2026, at the author's request that each
 direction of the release have a single note, the derivations of the sources, which had accompanied the packages as six
 separate notes by GPT-6 astra, were brought into Appendices~\ref{app:dualA} to~\ref{app:cone} by Claude Code
 sessions, edited to the notation and status words of this note and checked against the sources (A to C from the
 compendium, D from the two notes on the closure and the prescription, E from the study, F from the base study, G from
 the $c_1$ package); short explanatory steps were added in the merge and checked, and none changes a result. The
 separate notes are not part of the release, nor are the scripts that only build them; the few TeX blocks of the
 compendium that the checks parse are kept verbatim in \file{code/dual_d1/public/compendium_extract.json}
 (Appendix~\ref{app:repro}). On 17 October 2026, at the author's request, the magnetization-scheme calculation was added
 as Sec.~\ref{sec:integ} and Appendix~\ref{app:integ}, from its records and documentation, which are not in the
 release.
\end{itemize}

%====================================================================
\section{Open problems}\label{sec:open}
%====================================================================
\begin{itemize}
\item A proof of Eq.~\eqref{eq:C8cand}: exact expressions for all inputs of $\CSR_8$ and a weight-seven certificate.
\item The cause of the low independent value of $\CSR_8$, which Sec.~\ref{sec:integ} indicates is the one in error
 and the re-derivation's own input comparison locates in $A_9[\varepsilon^0\ell^1]$: a rerun with fresh random shifts
 would show whether it is a fluctuation with an underestimated error or a systematic error.
\item Beyond the magnetization scheme of Sec.~\ref{sec:integ}, which checks the prescription argument of
 Appendix~\ref{app:clos} numerically with a second regulator and different responses but assumes the same two-coupling
 description: a check that does not assume it, for example a higher-order analytic bootstrap (a ring also needs the
 monodromy sector of the compact boson), and a source-level audit of the historical high-order integrations, whose 21
 stage programs are missing from the archive.
\item $\CSR_9$, which needs the next determinant coefficient $D_5$ (a rational-part guess for it,
 $D_5^{(0)}=\frac{141}{4}$, is untested; Appendix~\ref{app:dualC:conjectures}).
\item Interval enclosures of the compendium's high-order numerical integrals (Appendix~\ref{app:dualC:numerics}), and a
 proof of the weight bound that the closed forms through $\CSR_7$ obey (Appendix~\ref{app:dualC:conjectures}).
\item In $d=2$, the uniform remainder bound that would make Eq.~\eqref{eq:c2closed} a theorem within the prescription.
\item In $d=2$, the disagreement of the three Monte Carlo configurations on the source-shape control ($\chi^2=7.6$ for
 two degrees of freedom), which does not affect $c_2$.
\item In $d=3$, gap hypotheses with independent support, or input errors with a calibrated meaning, so that the range of
 $c_1$ stops depending on chosen margins.
\end{itemize}

\appendix
\addtocontents{toc}{\protect\setcounter{tocdepth}{1}}
%====================================================================
\section{The dual of the chain: model, regulator and finite-response renormalization}\label{app:dualA}
%====================================================================
This appendix gives the derivations of the compendium (written by OpenAI GPT-6 astra; Sec.~\ref{sec:sources})
behind Secs.~\ref{sec:dual} to~\ref{sec:inputs1}, edited to the notation of this note. It starts from the lattice
model and derives the alternating domain-wall gas and its matrix representation
(Appendices~\ref{app:dualA:model} and~\ref{app:dualA:gas}). It then defines the regulated finite-interval responses
of the matched two-coupling prescription (Appendices~\ref{app:dualA:regulator} to~\ref{app:dualA:insertions}) and
the finite-response Jacobian and subtraction recurrence that turn them into beta coefficients
(Appendix~\ref{app:dualA:recurrence}). The last part explains why the invariants $T$ and $D$ need the
higher-insertion integrals (Appendix~\ref{app:dualA:invariants}). The beta inputs of the individual orders and their
evaluation are in Appendices~\ref{app:dualB} and~\ref{app:dualC}.

Apart from the published results quoted in Appendix~\ref{app:dualA:dictionary} and the continuum derivation of the
wall gas from the lattice model (Appendices~\ref{app:dualA:model} and~\ref{app:dualA:gas}), everything here either
defines the prescription or derives exact relations within it. Two things are not established here: that the dual
describes the critical chain, which is the published duality proposal~\cite{Benedetti:2024wgx,Benedetti:2025nzp}, and
that the renormalization of the dual involves only the couplings $g$ and $b$ and the prescription extracts their local
beta functions, which is the subject of the caveat ``Two-coupling closure'' in Sec.~\ref{sec:d1caveats} and of
Appendix~\ref{app:clos}.

%--------------------------------------------------------------------
\subsection{Notation}\label{app:dualA:notation}
As in the main text, $\db=1-s$ and $x=\Ds=\db/2$; the compendium writes $\delta$ for $\db$ (record and input names
such as $N_\delta$ keep their spelling), $C_n$ for $\CSR_n$, $b_{\rm dim}$ for the resummed charge $b$ and
$D_{2m}$ for $\Xi_{2m}$. The hard core is $\varepsilon$ and
$\ell=\log\varepsilon$; the thermal operator is $\mathcal E$, with dimension $\De$. Table~\ref{tab:dualA:notation}
collects the symbols of the dual that this appendix and the two following ones use.

{\small
\renewcommand{\arraystretch}{1.15}
\begin{longtable}{@{}>{\raggedright\arraybackslash}p{.24\textwidth}>{\raggedright\arraybackslash}p{.70\textwidth}@{}}
\caption{Symbols of the dual of the chain. The calligraphic $\mathcal C_n$ and $\mathcal K_n$ are normalized
responses (archived programs call them $C_n$ and $K_n$); they are never the thermal coefficients
$\CSR_n$.}\label{tab:dualA:notation}\\
\toprule
symbol & meaning\\
\midrule\endfirsthead
\toprule
symbol & meaning\\
\midrule\endhead
\bottomrule\endfoot
$X$, $X_i$, $q_i$ & coordinate along the chain; ordered kink (wall) positions; wall charges $\pm1$, alternating,
 $q_i=(-1)^{i+1}$\\
$\varepsilon$, $\ell$, $\Omega_n(\varepsilon)$ & dimensionless hard core between consecutive kinks; $\ell=\log\varepsilon$;
 ordered domain of $n$ kinks, Eq.~\eqref{eq:dualA:domain}\\
$\kappa$, $f_{\db}$ & reference inverse length (set to 1 from Appendix~\ref{app:dualA:regulator} on); wall
 interaction function, Eq.~\eqref{eq:dualA:fdb}\\
$\phi$, $\varphi$ & Gaussian field of the dual; coarse-grained Ising field (different fields)\\
$\hat\sigma_3$, $\hat\sigma_\pm$, $\sigma$ & qubit matrices; spin primary (unit normalized)\\
$\Og$, $\Oh$, $\chi_B$, $\chi$ & kink operator and current operator of Eq.~\eqref{eq:dual}; bare current
 $\frac{\mathrm i}{\sqrt2}\partial\phi$ (written $\chi$ in Eq.~\eqref{eq:dual}); renormalized unit-normalized current
 primary\\
$g$, $h$, $b_0$, $b$ & auxiliary wall fugacity and current coupling; reference charge; resummed charge
 $b=b_0-\sqrt2h$\\
$\eta$, $\partial_\eta$ & endpoint source $\sqrt2hb$, Eq.~\eqref{eq:dualA:eta}; derivative with respect to $\eta$ at
 $\eta=0$ (the compendium's script-E operator)\\
$B$, $v$, $z$; $Y$, $u$ & essential charge, $v=B^2$, $z=v-1$; essential wall fugacity (renamed $g$ after
 matching), $u=g^2$\\
$F$, $H$ & essential beta functions divided by $2u$ and by $1+z$, Eq.~\eqref{eq:dualA:essential}\\
$T$, $D$; $T_m$, $D_m$ & trace and determinant of the stability matrix in $(u,z)$; their coefficients of $\db^m$\\
$W_n$, $\Xi_{2m}$ & wall weight of $n$ ordered kinks; total length of the reversed domains,
 Eq.~\eqref{eq:dualA:moments}\\
$A_{2m+1}$, $Z_{2m}$, $Q_{2m}$ & odd, vacuum and current moments, Eq.~\eqref{eq:dualA:moments}\\
$\mathcal C_{2m+1}$, $\mathcal K_{2m}$ & vacuum-normalized odd and current responses,
 Eq.~\eqref{eq:dualA:normalization}\\
$I_0$, $I_1$ & coefficients $[\varepsilon^0\ell^0]I$ and $[\varepsilon^0\ell^1]I$ of a response $I$; e.g.\
 $\mathcal C_{3,1}$, $\mathcal K_{4,0}$, $Q_{2,0}$\\
$S_j$ & charge-weighted sums of $\log^j$ of the kink separations, Eq.~\eqref{eq:dualA:insertions}\\
$j_n$, $\lambda_{2m}$ & entries of the finite-response Jacobian, Eq.~\eqref{eq:dualA:jacobian}\\
$b_n$ ($n$ odd), $h_n$ ($n$ even) & coefficients of the auxiliary fugacity and current beta functions,
 Eq.~\eqref{eq:dualA:auxbeta} (archived scripts write $\eta_4,\eta_6$ for $h_4,h_6$); unrelated to $b_0$, $b$, $h$ and $\eta$\\
$a,p,b,c,f,\dots$ & beta inputs: Taylor coefficients of $F$ and $H$ at $v=1$, $\db=0$ (Sec.~\ref{sec:inputs1}); the
 input $b$ is $b_5$, not the charge\\
\end{longtable}
}

%--------------------------------------------------------------------
\subsection{The classical model and the short-range end}\label{app:dualA:model}
The object of interest is the thermal dimension of the \emph{classical} one-dimensional long-range Ising critical
family, not the spectrum of the quantum transverse-field Ising chain. A microscopic representative is the chain of
spins $m_i=\pm1$ with reduced Hamiltonian (energy divided by temperature)
\begin{equation}
 \mathcal H_{\rm red}=-K_{\rm nn}\sum_i m_im_{i+1}-K_{\rm lr}\sum_{i<j}\frac{m_im_j}{|i-j|^{1+s}},
 \qquad K_{\rm lr}>0.
 \label{eq:dualA:lattice}
\end{equation}
The nearest-neighbour coupling changes the core energies of the walls and the nonuniversal critical temperature; it
need not be present separately. The long-range critical regime of interest is $\frac12<s<1$. The endpoint $s=1$ is
marginal, and for $s>1$ there is no finite-temperature transition~\cite{Benedetti:2024wgx,Benedetti:2025nzp}.

\paragraph{The two ends of the long-range window.} The familiar coarse-grained action is
\begin{equation}
 S[\varphi]=\frac{c_s}{4}\int\dd X\dd Y\,\frac{[\varphi(X)-\varphi(Y)]^2}{|X-Y|^{1+s}}
 +\int\dd X\Bigl(\frac{r_\varphi}{2}\varphi^2+\frac{\lambda_\varphi}{4}\varphi^4\Bigr),\qquad c_s>0.
\end{equation}
Its nonanalytic quadratic kernel is proportional to $|p|^s$ in momentum space and gives $\Delta_\varphi=(1-s)/2$,
which is $\Ds$. The quartic coupling has engineering dimension $1-4\Delta_\varphi=2s-1$. The expansion about the
generalized free field (GFF) is therefore weak near $s=\frac12$, whereas the dual used here is weak near $s=1$. These
are opposite ends of the same critical family, not the same perturbative
expansion~\cite{Benedetti:2025nzp,Benedetti:2020rrq}.

\paragraph{What sits at the short-range end in $d=1$.} For the nearest-neighbour chain alone, the transfer matrix is
\begin{equation}
 \mathsf T_{\rm nn}=\begin{pmatrix}e^{K_{\rm nn}}&e^{-K_{\rm nn}}\\ e^{-K_{\rm nn}}&e^{K_{\rm nn}}\end{pmatrix}
 =\mathcal N\exp(\gamma\hat\sigma_1),\qquad \tanh\gamma=e^{-2K_{\rm nn}},
\end{equation}
with a constant $\mathcal N>0$. Its transfer Hamiltonian is a single two-state system, $\hat H=-\gamma\hat\sigma_1$ up
to an additive constant, and the inverse correlation length is $\xi^{-1}=2\gamma=\log\coth K_{\rm nn}$. At zero
temperature $\gamma\to0$, and the symmetric normalized trace gives $\frac12\operatorname{tr}\hat\sigma_3^2=1$: the
spin $\hat\sigma_3$ then has dimension zero. This endpoint spin sector plays, in $d=1$, the role of the
short-range fixed point ($\Ds^{\rm SR}=0$, Table~\ref{tab:conv}); it is not an ordinary finite-temperature
one-dimensional conformal field theory. In $d=2,3$ the crossover theory couples the short-range fixed point to a
generalized free field through the spin (Sec.~\ref{sec:d2}). In $d=1$ the analogous diagonal coupling
of the dual, $h\,\Oh=h\hat\sigma_3\chi_B$, commutes with $\hat\sigma_3$ and cannot flip the spin, so operators that create kinks are
essential. The higher-dimensional crossover construction therefore has a special one-dimensional realization, the
dual of Sec.~\ref{sec:dual}~\cite{Benedetti:2025nzp,Behan:2017dwr,Behan:2017emf}.

%--------------------------------------------------------------------
\subsection{Alternating domain walls, the Gaussian field and the matrices}\label{app:dualA:gas}
\paragraph{Domain walls.} Let $m(X)=\pm1$ be piecewise constant along the chain, with walls at $X_i$ of charge
$q_i$:
\begin{equation}
 \partial_Xm(X)=2\sum_iq_i\,\delta(X-X_i),\qquad q_i=\pm1,
\end{equation}
where $\delta(X)$ is the Dirac distribution. Ising walls necessarily alternate in sign, and a configuration that
returns to its initial spin has an even number of walls and zero total charge. With a reference inverse length
$\kappa$, define
\begin{equation}
 f_{\db}(r)=\frac{r^{\db}-\kappa^{-\db}}{\db},\qquad f_0(r)=\log(\kappa r),\qquad
 f_{\db}''(r)=-(1-\db)\,r^{-2+\db}\quad(r>0).
 \label{eq:dualA:fdb}
\end{equation}
From Appendix~\ref{app:dualA:regulator} on $\kappa=1$, so that $f_{\db}(r)=(r^{\db}-1)/\db$ as in
Sec.~\ref{sec:dual}. Take the continuum tail action, relative to the uniform configuration, to be
\begin{equation}
 S_{\rm tail}[m]=\frac{\kappa_s}{4}\int\dd X\dd Y\,\frac{[m(X)-m(Y)]^2}{|X-Y|^{2-\db}},
\end{equation}
the long-range term of Eq.~\eqref{eq:dualA:lattice} with $1+s=2-\db$. Since
$\partial_X\partial_Yf_{\db}(|X-Y|)=(1-\db)|X-Y|^{-2+\db}$ away from contact points, two integrations by parts give
\begin{equation}
 S_{\rm tail}=-\frac{\kappa_s}{2(1-\db)}\int\dd X\dd Y\,m'(X)\,m'(Y)\,f_{\db}(|X-Y|)+\text{core and boundary terms}.
\end{equation}
Inserting the wall charges, removing the self-energies and absorbing the positive coefficient into
$b_{\rm w}^2=2\kappa_s/(1-\db)$ gives the wall action
\begin{equation}
 S_{\rm walls}=-2b_{\rm w}^2\sum_{i<j}q_iq_j\,f_{\db}(X_j-X_i)+N_{\rm w}E_{\rm core}+\text{boundary terms},
 \label{eq:dualA:wallaction}
\end{equation}
with $N_{\rm w}$ the number of walls and $E_{\rm core}$ their core energy. As a check of the sign, an isolated pair
of opposite charges at separation $r$ has interaction energy $2b_{\rm w}^2f_{\db}(r)$, which increases with $r$ and is
positive beyond the reference length $1/\kappa$. Short-distance details are encoded by a fugacity and by local core
counterterms.

On a finite interval the regulated neutral gas is
\begin{equation}
 Z_{\rm gas}=\sum_{m\ge0}g^{2m}\int_{\Omega_{2m}(\varepsilon)}\prod_i\dd X_i\,
 \exp\Bigl[2b_{\rm w}^2\sum_{i<j}q_iq_j\,f_{\db}(X_j-X_i)\Bigr],
\end{equation}
up to the overall two-state trace and reference-length normalizations, with the ordered hard-core domain
$\Omega_n(\varepsilon)$ of Eq.~\eqref{eq:dualA:domain}. The positions are ordered, so there is no additional divisor
$(2m)!$. The odd projected responses of Appendix~\ref{app:dualA:moments} carry a specified compensating external
charge; they are not neutral partition functions.

\paragraph{The Gaussian field and its zero mode.} The wall weight is reproduced by vertex operators of a Gaussian
field. Choose a real Gaussian field modulo constant shifts, with covariance on neutral test functions
\begin{equation}
 C_{\db}(X-Y)=-2f_{\db}(|X-Y|),
 \label{eq:dualA:covariance}
\end{equation}
the covariance of Sec.~\ref{sec:dual}. With a finite-volume momentum cutoff and the zero mode removed, it is defined
by the positive quadratic action
\begin{equation}
 S_0[\phi]=\frac{\mathcal A_{\db}}{2}\int\frac{\dd p}{2\pi}\,|p|^{1+\db}\,|\widetilde\phi(p)|^2,
 \qquad\mathcal A_{\db}>0.
\end{equation}
The constant in $C_{\db}$ is immaterial for neutral observables, and $\phi$ itself is not a shift-invariant local
primary. Its increments have nonnegative variance. The formal scaling dimension of $\phi$ is $-\db/2$, but this does
not license treating its unphysical zero mode as a negative-norm primary.

Compactifying the constant mode with period $2\pi/b_0$ implements charge neutrality for the vertices
$V_\pm=\kappa\,{:}e^{\pm\mathrm ib_0\phi}{:}$. For $\sum_iq_i=0$, Wick's rule gives
\begin{equation}
 \Bigl\langle\prod_i{:}e^{\mathrm ib_0q_i\phi(X_i)}{:}\Bigr\rangle_0
 =\exp\Bigl[-b_0^2\sum_{i<j}q_iq_j\,C_{\db}(X_i-X_j)\Bigr]
 =\exp\Bigl[2b_0^2\sum_{i<j}q_iq_j\,f_{\db}(X_j-X_i)\Bigr],
 \label{eq:dualA:vertexgas}
\end{equation}
which is the nonlocal wall weight of Eq.~\eqref{eq:dualA:wallaction} with $b_0^2=b_{\rm w}^2$. At $\db=0$ the
two-point function of opposite charges is proportional to $r^{-2b_0^2}$.

\paragraph{Matrices enforce the Ising constraint.} With these vertices the operators of Eq.~\eqref{eq:dual} are
$\Og=\hat\sigma_+V_++\hat\sigma_-V_-$ and $\Oh=\hat\sigma_3\chi_B$, with $\chi_B=\frac{\mathrm i}{\sqrt2}\partial\phi$
the bare current (the subscript distinguishes it from the renormalized primary $\chi$ of
Appendix~\ref{app:dualA:dictionary}). Their path-ordered evolution from $U$ to $V$ is
\begin{equation}
 \mathcal D(V,U)=\mathrm P\exp\int_U^V\dd X\,\bigl[g\,\Og(X)+h\,\Oh(X)\bigr].
\end{equation}
Because $\hat\sigma_+^2=\hat\sigma_-^2=0$, only alternating charges survive in its expansion, so the traced matrix
expansion matches the regulated Ising gas, whereas an unconstrained cosine interaction would also admit nonalternating
walls. Spin correlators are obtained by inserting $\hat\sigma_3$ in the ordered trace. Two insertions multiply a wall
history by $(-1)^{N_{\rm between}}$, with $N_{\rm between}$ the number of walls between them, which is exactly the
product of the two Ising spins.

The auxiliary model also contains charged objects that are not physical operators of the chain. Constant shifts of
$\phi$ compensated by rotations of $\hat\sigma_\pm$ define a $U(1)$ action, and the physical operator algebra is its
invariant sector (Sec.~\ref{sec:dual}). The equality of consistently regulated wall weights and matrix histories is a
direct algebraic statement. The identification of the critical continuum limit with the long-range Ising CFT is the
published duality proposal and its tested consequences~\cite{Benedetti:2025nzp}; this calculation supplies no
all-orders theorem on the lattice scaling limit.

%--------------------------------------------------------------------
\subsection{The regulator and exact current resummation}\label{app:dualA:regulator}
Work on an interval of length one, with $\kappa=1$. All ordered integrations use
\begin{equation}
 \Omega_n(\varepsilon)=\{0<X_1<\cdots<X_n<1:\ X_{i+1}-X_i\ge\varepsilon\},\qquad \ell=\log\varepsilon.
 \label{eq:dualA:domain}
\end{equation}
Only consecutive kink separations have a hard core; there is no separate hard core between a wall and an endpoint.
Current insertions are resummed \emph{before} this regulator is imposed. An independent hard core on unresolved
current insertions would change the finite parts of the responses, and could not be combined with the formulas below.

\paragraph{Change of frame.} The resummation is a field-dependent change of frame,
\begin{equation}
 R(X)=\exp\bigl[\mathrm ih\,\hat\sigma_3\,\phi(X)/\sqrt2\bigr],\qquad
 \mathcal D(V,U)=R(V)\,\widetilde{\mathcal D}(V,U)\,R(U)^{-1},
\end{equation}
which removes the bulk derivative interaction $h\,\Oh$. Conjugation by $R$ multiplies $\hat\sigma_\pm$ by
$e^{\mp\mathrm i\sqrt2h\phi}$ and so shifts the bulk wall charge to
\begin{equation}
 b=b_0-\sqrt2h,
\end{equation}
the effective charge of Sec.~\ref{sec:dual} (the beta input $b$ of Sec.~\ref{sec:inputs1} is a different
quantity). The endpoint matrices $R(1)$ and $R(0)^{-1}$ remain. Equivalently, one integrates by parts on the
interval $I$, $\int m\,\partial\phi=m\phi|_{\partial I}-\int m'\phi$. The endpoint charges have magnitude $h/\sqrt2$,
and their contraction with a wall produces the \emph{endpoint source}
\begin{equation}
 \eta=\sqrt2\,h\,b=\sqrt2\,b_0h-2h^2.
 \label{eq:dualA:eta}
\end{equation}
The coefficient of $h^2$ is $-2$. An intermediate higher-logarithm check in the archive used $-1$; the beta
coefficients below use only first $h$ derivatives at $h=0$ and are unaffected, but checks that involve second or
higher source derivatives are.

\paragraph{Endpoint factors.} With alternating charges $q_i=(-1)^{i+1}$, the endpoint factors of $n$ kinks are
\begin{equation}
\begin{aligned}
 U_n^{\rm odd}(\eta)&=\exp\Bigl\{\eta\sum_iq_i\bigl[f_{\db}(X_i)+f_{\db}(1-X_i)\bigr]\Bigr\}&&(n\text{ odd}),\\
 U_n^{\rm even}(\eta)&=\exp\Bigl\{\eta\sum_iq_i\bigl[f_{\db}(X_i)-f_{\db}(1-X_i)\bigr]\Bigr\}&&(n\text{ even}).
\end{aligned}
\label{eq:dualA:endpoint}
\end{equation}
We write $\partial_\eta$ for the derivative with respect to $\eta$ at $\eta=0$ (the compendium writes it as a script E).
By the reflection $X\to1-X$, a \emph{first} source derivative may be represented after integration by
$2\sum_iq_if_{\db}(X_i)$; this replacement is not valid before squaring, for a second derivative. At the reference
charge $v=b_0^2$ the coupling $h$ enters the moments through the squared charge $v=b^2=(b_0-\sqrt2h)^2$ of the wall
weight $W_n$ (Eq.~\eqref{eq:dualA:moments}; it equals the essential $B^2$ of Appendix~\ref{app:dualA:matching} at
$\db=0$) and through $\eta$, so that
\begin{equation}
 \partial_h\big|_{h=0}=\sqrt2\,b_0\,\bigl(-2\partial_v+\partial_\eta\bigr).
 \label{eq:dualA:sourcechain}
\end{equation}
For example, for a single kink at $\db=0$ and $b_0=1$ the odd endpoint factor is $X^\eta(1-X)^\eta$, so the
one-kink moment is Euler's beta function, $A_1(h)=\Gamma(1+\eta)^2/\Gamma(2+2\eta)$, and
\begin{equation}
 A_1(h)=1-2\sqrt2\,h+\bigl(12-2\zeta(2)\bigr)h^2+O(h^3).
 \label{eq:dualA:A1}
\end{equation}

%--------------------------------------------------------------------
\subsection{Essential couplings and finite-tail matching}\label{app:dualA:matching}
Normal ordering at the cutoff and the charge shift give the essential charge $B$ and the essential wall fugacity
$Y$,
\begin{equation}
 B=(b_0-\sqrt2h)\,\varepsilon^{\db/2},\qquad
 Y=g_{\rm aux}\,\kappa\varepsilon\,\exp\Bigl[\frac{b_0^2(\kappa^{-\db}-\varepsilon^{\db})}{\db}\Bigr],
 \label{eq:dualA:matching}
\end{equation}
where $g_{\rm aux}$ is the fugacity of the auxiliary (unmatched) theory. The $h$ dependence cancels from $Y$ once
both normal-ordering factors are included. At unit reference length and $b_0=1$,
\begin{equation}
 Y=g_{\rm aux}\exp\Bigl[\ell+\frac{1-e^{\db\ell}}{\db}\Bigr]
 =g_{\rm aux}\Bigl[1-\frac{\db}{2}\ell^2+\db^2\Bigl(\frac{\ell^4}{8}-\frac{\ell^3}{6}\Bigr)+\cdots\Bigr].
 \label{eq:dualA:fugmatch}
\end{equation}
Charge derivatives require keeping $b_0$ variable until they are taken; setting $b_0=1$ prematurely loses terms.
With $Y$ renamed $g$, the essential flow is written as in Sec.~\ref{sec:dual},
\begin{equation}
 u=g^2,\qquad v=B^2=1+z,\qquad \beta_u=2uF,\qquad \beta_z=(1+z)H,\qquad
 \beta=\mu\frac{\dd}{\dd\mu}=-\frac{\dd}{\dd\ell}.
 \label{eq:dualA:essential}
\end{equation}
The finite $\db$ factors of Eq.~\eqref{eq:dualA:fugmatch} must be matched together with the endpoint integrals. It
is not consistent to compute at $\db=0$ in one scheme and append the classical-scaling term in unrelated coordinates. The
finite-tail factor is needed, for example, for the input $p$ of Sec.~\ref{sec:inputs1} (Appendix~\ref{app:dualB:threeins}).

%--------------------------------------------------------------------
\subsection{Leading spectrum and operator dictionary}\label{app:dualA:dictionary}
This subsection collects the published leading-order results~\cite{Benedetti:2024wgx,Benedetti:2025nzp} in the
notation of this appendix. The published field-theory and bootstrap methods agree at their stated orders. They give
\begin{equation}
 \Delta_\pm=1\pm\sqrt{2\db}+\frac{\db}{4}+O(\db^{3/2}),\qquad \Ds=\frac{\db}{2},\qquad \Delta_\chi=1-\frac{\db}{2}.
 \label{eq:dualA:published}
\end{equation}
The coefficients $\CSR_3$ and beyond are not quoted from those papers; they belong to the computation described in
this and the following appendices and carry its conditions (Sec.~\ref{sec:dual}).

\paragraph{OPE and leading flow.} At $\db=0$ and $b_0=1$ the ordered OPE of two kink operators is
\begin{equation}
 \Og(X)\,\Og(0)=X^{-2}\,\mathbb 1+\sqrt2\,X^{-1}\,\Oh(0)+\cdots,\qquad X>0,
\end{equation}
and the two ordered mixed OPEs, of $\Og$ with $\Oh$ in either order, each contain $\sqrt2\,\Og/X$. The beta
functions of the dual through the published order are
\begin{equation}
 \beta_g=-2\sqrt2\,gh+2gh^2+g^3+\cdots,\qquad \beta_h=-\sqrt2\,g^2+2hg^2+\cdots.
 \label{eq:dualA:leadingbeta}
\end{equation}
In the essential couplings, with $z=b^2-1=-2\sqrt2h+2h^2$ at $b_0=1$ and $\db=0$, they become $\beta_g=gF$ and
$\beta_z=2b\,\partial_hb\,\beta_h=(1+z)H$ with $F=z+u+\cdots$ and $H=4u+\cdots$; at $\db\neq0$ the classical
scaling of $B=b\,\varepsilon^{\db/2}$ adds $-\db$ to $H$, so that $H=-\db+4u+\cdots$, whose fixed point $u_*=\db/4$,
$z_*=-\db/4$ is given in Sec.~\ref{sec:dual}. This explains the expansion in $\sqrt{\db}$: the fugacity itself is of
order $\sqrt{\db}$, although $u$ and $z$ are of order $\db$. In the coordinates $(g,v)$ the leading stability matrix is
\begin{equation}
 M_*=\begin{pmatrix}\db/2&g_*\\ 8(1-\db/4)\,g_*&0\end{pmatrix}+\text{higher-order terms}.
 \label{eq:dualA:Mstar}
\end{equation}
A linear beta eigenvalue $\lambda$ corresponds to the dimension $1+\lambda$; with $g_*^2=\db/4$ the eigenvalues of
$M_*$ give $\Delta_\pm$ of Eq.~\eqref{eq:dualA:published}. The relevant branch is the thermal primary
$\mathcal E=\mathcal O_-$, so $\nu^{-1}=1-\De$, and the other eigenvalue gives the leading correction exponent
$\omega=\Delta_+-1$. Only the orders shown in Eq.~\eqref{eq:dualA:published} are controlled by this truncation.
Expanding the small matrix to more orders does not compute the higher coefficients
(Appendix~\ref{app:dualA:invariants}).

\paragraph{Spin dimension.} A direct wall-gas check gives, for the spin two-point function at separation $r$,
$G_\sigma(r)-1=-4g^2\log(r/\varepsilon)+\cdots$, because a kink pair changes the spin product only when one wall lies
between the insertions. Thus $\gamma_\sigma=2g^2+\cdots$, which reproduces $\Ds=\db/2$ at the fixed point.

\paragraph{Equation of motion and the shadow pair.} The covariant derivative of an insertion is
$D_X\mathcal O=\partial_X\mathcal O-[g\,\Og+h\,\Oh,\mathcal O]$. In particular
\begin{equation}
 D_X\hat\sigma_3=2g\bigl(\hat\sigma_+V_+-\hat\sigma_-V_-\bigr),\qquad
 C_{\db}^{-1}\phi=\frac{\mathrm ib}{2}\,D_X\hat\sigma_3
 \label{eq:dualA:eom}
\end{equation}
away from contact terms, where $C_{\db}^{-1}$ is the inverse of the covariance~\eqref{eq:dualA:covariance}. The
nonanalytic Gaussian kernel and this equation of motion give the protected shadow pair of
Eq.~\eqref{eq:dualA:published}, $\Ds+\Delta_\chi=1$: the antisymmetric flip operator
$\hat\sigma_+V_+-\hat\sigma_-V_-$ becomes the descendant of $\sigma$. The bare current $\chi_B$ can mix with $\sigma$
through power-divergent counterterms and needs a finite normalization; all CFT data below refer to the renormalized,
unit-normalized primary $\chi$, not to an unsubtracted current~\cite{Benedetti:2025nzp}.

\paragraph{Published OPE coefficients.} For completeness, with $\tau=\sqrt{\db}=\sqrt{2\Ds}$, with $\pm$ labelling the
operators $\mathcal O_\pm$ of dimensions $\Delta_\pm$, and with the published operator
signs~\cite[eqs.~(5.115) and (5.120)]{Benedetti:2025nzp},
\begin{equation}
\begin{aligned}
 c_{\sigma\chi\pm}&=\frac1{\sqrt2}\mp\frac{\tau}{16}-\frac{\tau^2}{256\sqrt2}+O(\tau^3),\\
 c_{\sigma\sigma\pm}&=\pm\frac{\tau}{\sqrt2}-\frac{15}{16}\tau^2
 \pm\frac1{\sqrt2}\Bigl(\frac{543}{256}-\frac{\pi^2}{6}\Bigr)\tau^3+O(\tau^4),\\
 c_{\chi\chi\pm}&=-\frac{\pi^2}{2}\tau^2\mp\frac{11\pi^2}{16\sqrt2}\tau^3+O(\tau^4),\\
 c_{+++}&=\frac32+\frac{39}{16\sqrt2}\tau+O(\tau^2),\qquad
 c_{---}=\frac32-\frac{39}{16\sqrt2}\tau+O(\tau^2),\\
 c_{++-}&=-\frac12+\frac{\tau}{16\sqrt2}+O(\tau^2),\qquad
 c_{+--}=-\frac12-\frac{\tau}{16\sqrt2}+O(\tau^2).
\end{aligned}
\label{eq:dualA:ope}
\end{equation}
These OPE coefficients, including their $\pi^2$ terms, are published inputs, not loop results of the compendium.
The compendium takes its published inputs from Ref.~\cite{Benedetti:2025nzp}, Secs.~2--3 and App.~B, Secs.~4.1--4.4,
Eqs.~(5.48), (5.53)--(5.54) and (5.115)--(5.120), and App.~E.

%--------------------------------------------------------------------
\subsection{Ordered moments and vacuum normalization}\label{app:dualA:moments}
\paragraph{Moments.} At unit reference length and with $q_i=(-1)^{i+1}$, define the wall weight of $n$ ordered kinks
and, for an even number of kinks, the total length of the reversed domains,
\begin{equation}
 W_n(v,\db;\boldsymbol X)=\exp\Bigl[2v\sum_{i<j}q_iq_j\,\frac{(X_j-X_i)^{\db}-1}{\db}\Bigr],\qquad
 \Xi_{2m}(\boldsymbol X)=\sum_{j=1}^m(X_{2j}-X_{2j-1}),
\end{equation}
and the odd, vacuum and current moments
\begin{equation}
\begin{gathered}
 A_{2m+1}=\int_{\Omega_{2m+1}(\varepsilon)}W_{2m+1}\prod_i\dd X_i,\\
 Z_{2m}=\int_{\Omega_{2m}(\varepsilon)}W_{2m}\prod_i\dd X_i,\qquad
 Q_{2m}=\int_{\Omega_{2m}(\varepsilon)}\Xi_{2m}W_{2m}\prod_i\dd X_i.
\end{gathered}
 \label{eq:dualA:moments}
\end{equation}
When the endpoint source is switched on, the integrands also carry the endpoint factors of
Eq.~\eqref{eq:dualA:endpoint}, on which all $\partial_\eta$ derivatives act (as in the example of $A_1$,
Eq.~\eqref{eq:dualA:A1}). The $A$ moments are coefficients of charged projected responses, with their compensating
external operator understood; they do not contradict the zero-mode neutrality of $Z$. At zero endpoint source
$A_1=1$.

\paragraph{Gap variables and the first kernels.} Integrating the first position gives a translation factor $1-R$
in the gap variables $r_i=X_{i+1}-X_i\ge\varepsilon$, with $R=\sum_ir_i\le1$. At $v=1$ and $\db=0$,
$W_n=\prod_{i<j}(X_j-X_i)^{2q_iq_j}$; naming the consecutive gaps $r,t,w,y$, the first kernels are
\begin{align}
 W_3&=\frac{(r+t)^2}{r^2t^2},\label{eq:dualA:W3}\\
 W_4&=\frac{(r+t)^2(t+w)^2}{r^2t^2w^2(r+t+w)^2},\label{eq:dualA:W4}\\
 W_5&=\frac{(r+t)^2(t+w)^2(w+y)^2(r+t+w+y)^2}{r^2t^2w^2y^2(r+t+w)^2(t+w+y)^2}.\label{eq:dualA:W5}
\end{align}
(The compendium names the second gap $t'$; it is unrelated to the trace $T$.) At $v=1$, $\db=0$ the two-kink moments
are elementary, $Z_2=\int_\varepsilon^1(1-r)r^{-2}\dd r$ and $Q_2=\int_\varepsilon^1(1-r)r^{-1}\dd r$:
\begin{equation}
 Z_2=\varepsilon^{-1}+\ell-1,\qquad Q_2=-\ell-1+\varepsilon.
 \label{eq:dualA:Z2Q2}
\end{equation}
These elementary expressions fix the signs of the finite parts of the responses.

\paragraph{Vacuum normalization before coefficient extraction.} The responses are divided by the vacuum before any
coefficient is extracted. One forms generating functions in the fugacity and divides the \emph{integrated} numerator
by the integrated vacuum denominator,
\begin{equation}
 \sum_{m\ge0}\mathcal C_{2m+1}\,g^{2m}=\frac{\sum_{m\ge0}A_{2m+1}g^{2m}}{\sum_{m\ge0}Z_{2m}g^{2m}},\qquad
 \sum_{m\ge1}\mathcal K_{2m}\,g^{2m}=\frac{\sum_{m\ge1}Q_{2m}g^{2m}}{\sum_{m\ge0}Z_{2m}g^{2m}},\qquad Z_0=1,
 \label{eq:dualA:genfun}
\end{equation}
so that $\mathcal C_1=A_1$ and $\mathcal K_2=Q_2$. At zero source ($A_1=1$) this gives, through the orders used,
\begin{equation}
\begin{aligned}
 \mathcal C_3&=A_3-Z_2,\\
 \mathcal C_5&=A_5-A_3Z_2+Z_2^2-Z_4,\\
 \mathcal C_7&=A_7-A_5Z_2+A_3(Z_2^2-Z_4)-Z_2^3+2Z_2Z_4-Z_6,\\
 \mathcal K_4&=Q_4-Q_2Z_2,\\
 \mathcal K_6&=Q_6-Q_4Z_2+Q_2(Z_2^2-Z_4),\\
 \mathcal K_8&=Q_8-Q_6Z_2+Q_4(Z_2^2-Z_4)+Q_2(-Z_2^3+2Z_2Z_4-Z_6).
\end{aligned}
\label{eq:dualA:normalization}
\end{equation}
Source derivatives differentiate \emph{every} factor, including $A_1$; for example
\begin{equation}
 \partial_\eta\mathcal C_3=\partial_\eta A_3-(\partial_\eta A_1)\,Z_2-\partial_\eta Z_2.
\end{equation}
For a response $I$ define
\begin{equation}
 I_0=[\varepsilon^0\ell^0]\,I,\qquad I_1=[\varepsilon^0\ell^1]\,I,
 \label{eq:dualA:I0I1}
\end{equation}
so that, for instance, $\mathcal C_{3,1}$ is the single-logarithm coefficient of $\mathcal C_3$ and $Q_{2,0}=-1$ the
finite part of $Q_2$. The products in Eq.~\eqref{eq:dualA:normalization} must be formed, as complete series in
$\varepsilon$ and $\ell$, before these projections are applied: since $Z_2\sim\varepsilon^{-1}$, a term such as
$\varepsilon\ell$ in another factor contributes to $I_1$. Multiplying only finite parts is not a valid normalization
rule. Derivatives are likewise taken before the projection: $\partial_v\mathcal C_{n,0}$ and $\partial_\eta\mathcal
C_{n,0}$ mean the finite parts of $\partial_v\mathcal C_n$ and $\partial_\eta\mathcal C_n$.

%--------------------------------------------------------------------
\subsection{Charge and tail insertions}\label{app:dualA:insertions}
Derivatives of the wall weight in the charge $v$ and the tail parameter $\db$ insert logarithms of the kink
separations. Let
\begin{equation}
 S_j=\sum_{i<k}q_iq_k\log^j(X_k-X_i).
\end{equation}
Since $[(X_k-X_i)^{\db}-1]/\db=\log(X_k-X_i)+\frac{\db}2\log^2(X_k-X_i)+\frac{\db^2}6\log^3(X_k-X_i)+\cdots$, the
derivatives at $v=1$, $\db=0$ are
\begin{equation}
\begin{aligned}
 \partial_vW_n&=2S_1W_n,&\partial_v^2W_n&=4S_1^2W_n,\\
 \partial_{\db}W_n&=S_2W_n,&\partial_v\partial_{\db}W_n&=(S_2+2S_1S_2)W_n,\\
 [\db^2]\,W_n&=\Bigl(\frac13S_3+\frac12S_2^2\Bigr)W_n.
\end{aligned}
\label{eq:dualA:insertions}
\end{equation}
The second term of the mixed derivative and both terms of the second tail insertion are essential. Tail derivatives
also act on the endpoint factors of Eq.~\eqref{eq:dualA:endpoint}, through $f_{\db}$. A Taylor coefficient such as
$[\db^2]$ includes the factor $1/2!$; it is not the second derivative itself.

%--------------------------------------------------------------------
\subsection{The finite-response Jacobian and the subtraction recurrence}\label{app:dualA:recurrence}
The beta functions are extracted from two vacuum-normalized responses at $h=0$: the odd response, the projection of
the ordered evolution on the odd spin-flip vertex, with finite parts built from the $\mathcal C_{2m+1}$, and the
current response, its projection on the current, built from the $\mathcal K_{2m}$ (with $\mathcal K_2=Q_2$). The
regulator convention subtracts logarithmic coupling counterterms with no additional finite coupling redefinition in
these coordinates. The finite parts of the responses are nevertheless required: the single logarithm of a normalized
response is not by itself a beta coefficient, because finite parts of lower orders, multiplied by lower beta
coefficients, also contribute to it.

\paragraph{Jacobian entries.} Set
\begin{equation}
\begin{aligned}
 j_1&=-\frac{2}{1+\db},\qquad j_n=-2\partial_v\mathcal C_{n,0}+\partial_\eta\mathcal C_{n,0}\quad(n=3,5),\\
 \lambda_2&=-2Q_{2,0}-4v\,\partial_vQ_{2,0}+2v\,\partial_\eta Q_{2,0},\\
 \lambda_4&=-2\mathcal K_{4,0}-4v\,\partial_v\mathcal K_{4,0}+2v\,\partial_\eta\mathcal K_{4,0}.
\end{aligned}
\label{eq:dualA:jacobian}
\end{equation}
The integral defining $j_1$ is $\int_0^1[f_{\db}(X)+f_{\db}(1-X)]\dd X$, the source derivative $\partial_\eta A_1$ of
the one-kink moment. The same definition of $j_7$ is used at the eighth order (Appendix~\ref{app:dualC:C8setup}).

\paragraph{Auxiliary beta coefficients.} Write the beta functions of the auxiliary couplings at $h=0$ as
\begin{equation}
 \beta_g=b_3g^3+b_5g^5+b_7g^7+\cdots,\qquad
 \frac{\beta_h}{\sqrt2\,b_0}=-g^2+h_4g^4+h_6g^6+\cdots,
 \label{eq:dualA:auxbeta}
\end{equation}
so that $b_n$ ($n$ odd) are the odd fugacity coefficients and $h_n$ ($n$ even, $h_2=-1$) the normalized current
coefficients. Archived scripts call $h_4$ and $h_6$ by $\eta_4$ and $\eta_6$; these are not the endpoint source $\eta$.
The coefficients depend on the charge $v=b_0^2$ and on $\db$, and their Taylor coefficients at $v=1$, $\db=0$, after
the matching factors of Appendix~\ref{app:dualA:matching} are restored, are the beta inputs of
Sec.~\ref{sec:inputs1}. For example $b_3=1$ and $h_2=-1$ at that point give the leading terms $u$ of $F$ and $4u$ of
$H$, while $b=b_5$, $c=-4h_4$ and $k=-4h_6$ are the coefficients of $u^2$ in $F$ and of $u^2$ and $u^3$ in $H$; the complete dictionary is in Appendices~\ref{app:dualB:conv} and~\ref{app:dualC:C8setup}.

\paragraph{The finite-response Jacobian.} Through the required orders the Jacobian is
\begin{equation}
 \mathcal J=\begin{pmatrix}
 1+3\mathcal C_{3,0}g^2+5\mathcal C_{5,0}g^4 & \sqrt2\,b_0\,(j_1g+j_3g^3+j_5g^5)\\[2pt]
 \sqrt2\,b_0\,(2Q_{2,0}g+4\mathcal K_{4,0}g^3) & 1+\lambda_2g^2+\lambda_4g^4
 \end{pmatrix},
 \label{eq:dualA:jacobianmatrix}
\end{equation}
and the vector of single-logarithm coefficients is
\begin{equation}
 R_1=\begin{pmatrix}
 \mathcal C_{3,1}g^3+\mathcal C_{5,1}g^5+\mathcal C_{7,1}g^7\\[2pt]
 \sqrt2\,b_0\,(-g^2+\mathcal K_{4,1}g^4+\mathcal K_{6,1}g^6)
 \end{pmatrix}.
 \label{eq:dualA:R1}
\end{equation}
The entries have a direct reading. The first column holds the $g$ derivatives of the finite parts of the odd
response, $g+\mathcal C_{3,0}g^3+\mathcal C_{5,0}g^5$, and of the current response,
$h+\sqrt2\,b\,(Q_{2,0}g^2+\mathcal K_{4,0}g^4)$. The second column holds their $h$ derivatives at $h=0$, taken with
Eq.~\eqref{eq:dualA:sourcechain}; the terms $-2Q_{2,0}$ and $-2\mathcal K_{4,0}$ of $\lambda_2$ and $\lambda_4$ are the
derivative of the charge $b=b_0-\sqrt2h$ that multiplies the current moments. The leading $-g^2$ in the second row of
$R_1$ is $Q_{2,1}=-1$, from Eq.~\eqref{eq:dualA:Z2Q2}.

\paragraph{The subtraction recurrence.} The prescription solves
\begin{equation}
 \mathcal J\begin{pmatrix}\beta_g\\ \beta_h\end{pmatrix}=R_1
 \label{eq:dualA:system}
\end{equation}
order by order in $g$. With $v=b_0^2$ (the factors $2v$ are $(\sqrt2b_0)^2$) this gives
\begin{equation}
\begin{aligned}
 b_3&=\mathcal C_{3,1}+2vj_1,\\
 h_4&=\mathcal K_{4,1}-2Q_{2,0}b_3+\lambda_2,\\
 b_5&=\mathcal C_{5,1}-3\mathcal C_{3,0}b_3-2vj_1h_4+2vj_3,\\
 h_6&=\mathcal K_{6,1}-2Q_{2,0}b_5-4\mathcal K_{4,0}b_3-\lambda_2h_4+\lambda_4,\\
 b_7&=\mathcal C_{7,1}-3\mathcal C_{3,0}b_5-5\mathcal C_{5,0}b_3-2vj_1h_6-2vj_3h_4+2vj_5.
\end{aligned}
\label{eq:dualA:recurrence}
\end{equation}
The classical-scaling and matching contributions of Eq.~\eqref{eq:dualA:matching} are then restored to pass from the
auxiliary couplings $(g_{\rm aux},b)$ to the essential couplings $(u,z)$. This derives the recurrence; it is not an
assignment of a raw single logarithm to a physical beta function. The recurrence uses only first $h$ derivatives at
$h=0$, which is why the coefficient of $h^2$ in Eq.~\eqref{eq:dualA:eta} does not enter it. Its extension through
the ninth order ($h_8$ and $b_9$, with $\mathcal C_9$) is in Appendix~\ref{app:dualC:C8setup}, Eqs.~\eqref{eq:dualC:j7lambda6} and~\eqref{eq:dualC:b9}.

\paragraph{The lowest order as a check of signs.} The normalized three-kink response has
$\mathcal C_{3,1}=5$ (its expansion is in Appendix~\ref{app:dualB:D2T2}), and $j_1=-2$ at $\db=0$, so
$b_3=5-4=1$. With $\mathcal K_{4,1}=-8+4\zeta(2)$, $Q_{2,0}=-1$, $\partial_vQ_{2,0}=-2$ and
$\partial_\eta Q_{2,0}=-2\zeta(2)$, one has $\lambda_2=10-4\zeta(2)$ and $h_4=4$, that is $c=-4h_4=-16$
(Sec.~\ref{sec:inputs1}).

\paragraph{Scope.} Equation~\eqref{eq:dualA:system} defines the beta functions of the prescription. That its
solution is the local beta function of the dual, up to a finite change of $(g,b)$ that leaves $T$ and $D$ unchanged,
is the formal argument of Appendix~\ref{app:clos} (Appendices~\ref{app:clos:prescription} and~\ref{app:clos:uniqueness}), summarized in the caveat ``Two-coupling
closure'' of Sec.~\ref{sec:d1caveats}.

%--------------------------------------------------------------------
\subsection{Stability invariants and the limits of the leading flow}\label{app:dualA:invariants}
\paragraph{Trace and determinant.} Linearizing $(\beta_u,\beta_z)=(2uF,(1+z)H)$ about a fixed point with $F=H=0$,
the terms $2F$ and $H$ of the derivatives vanish, and the stability matrix in $(u,z)$ is
\begin{equation}
 \begin{pmatrix} 2u_*F_u & 2u_*F_z\\ (1+z_*)H_u & (1+z_*)H_z\end{pmatrix},
 \label{eq:dualA:stability}
\end{equation}
with the derivatives taken before the fixed point is substituted. Its trace and determinant are the $T$ and $D$ of
Sec.~\ref{sec:invariants}, and its eigenvalues $\lambda_\pm=T/2\pm(T^2/4-D)^{1/2}$ give the dimensions
$\Delta_\pm=1+\lambda_\pm$ of Eq.~\eqref{eq:eigen}. A regular change of couplings transforms the exact stability
matrix by similarity, so $T$ and $D$ are invariant although individual beta coefficients are not. The even
coefficients need only the trace, $\CSR_{2m}=2^{m-1}T_m$ (Eq.~\eqref{eq:conv}): a new even coefficient needs the next
trace coefficient but not the next determinant coefficient.

\paragraph{Why the leading flow cannot be extrapolated.} If one regarded the leading $F=z+u$ and $H=-\db+4u$ as
exact, the invariants would be $T=2u_*=\db/2$ and $D=-8u_*(1+z_*)=-2\db+\db^2/2$, so $D_2=\frac12$ and $T_2=0$, and
Eq.~\eqref{eq:conv} would give $\CSR_3=\frac7{16}$ and $\CSR_4=0$ instead of the values of Eq.~\eqref{eq:closed}.
Terms omitted in that truncation change both results. For instance, a term $azu$ in $F$ is of order $\db^2$ at the
fixed point, but it adds $au_*$, of order $\db$, to the derivative $F_z$ and so changes $D_2$ (indeed
$D_2=\frac12-\frac a2-\frac c8$, Sec.~\ref{sec:inputs1}). A term $zu^2$ in $H$ has a value of order $\db^3$, but its
derivative $H_z$ is of order $\db^2$. Power counting must be applied \emph{after} differentiation. This is the
distinction between a higher-insertion calculation and a formal re-expansion of an already truncated flow, and it is
why the beta inputs include derivatives in $z$ and $\db$ of integrals with more kinks than the order would suggest.
For the same reason, expanding the leading matrix $M_*$ of Eq.~\eqref{eq:dualA:Mstar} to more orders does not compute
the higher coefficients.

%====================================================================
\section{The dual of the chain: the third to sixth orders}\label{app:dualB}
%====================================================================
This appendix gives the derivations of the compendium (written by OpenAI GPT-6 astra; Sec.~\ref{sec:sources}) for
the coefficients $\CSR_3$ to $\CSR_6$, edited to the notation of this note. It builds on the model, the regulator,
the moments and the subtraction recurrence of Appendix~\ref{app:dualA}, and gives the calculations that
Secs.~\ref{sec:inputs1} and~\ref{sec:closed} summarize for these orders; the inputs that rest on archived exact
expressions are marked as such below.

Appendices~\ref{app:dualB:D2T2} to~\ref{app:dualB:C6} assemble the invariants $D_2$, $T_2$, $D_3$ and $T_3$ from the
beta inputs. The remaining subsections evaluate the inputs: the Mellin representation of the regulated moments and
the rules for finite parts, with the charge-derivative period $\mathcal R_4$ (Appendix~\ref{app:dualB:worked}); a
trilogarithm lemma and the Spence--Kummer identity on which it rests (Appendix~\ref{app:dualB:trilog}); the closed
forms of the periods $\mathcal P_4$ and $\mathcal S_5$ (Appendix~\ref{app:dualB:periods}); the counterterm
calculation that attaches the periods to the inputs (Appendix~\ref{app:dualB:counterterms}); the charge and tail
derivatives at three insertions (Appendix~\ref{app:dualB:threeins}); and the three-kink and endpoint inputs of $D_3$
(Appendix~\ref{app:dualB:endpoints}). What these derivations establish differs by order:
\begin{itemize}
\item $\CSR_3$ and $\CSR_4$. The three periods $\mathcal P_4$, $\mathcal R_4$ and $\mathcal S_5$ are evaluated in closed
 form with explicit primitives, a table of elementary finite parts and the trilogarithm lemma, without numerical
 recognition. Their attachment to the five inputs $a,p,b,c,f$ goes through response coefficients and counterterms:
 the mixed response coefficients of Eq.~\eqref{eq:dualB:mixed} are taken from the archive without derivation, the
 counterterm algebra is displayed here in full, and the compendium checks the charge derivative of the fourth-order
 counterterm and the three-insertion derivatives with its code.
\item $\CSR_5$. The inputs $A$ and $E$ and two endpoint constants are evaluated here. The inputs $G$, $q$, $k$ are
 archived exact period expressions, and the combination of them that enters $D_3$ is reduced to zeta values by the
 exact certificate of $D_3$ (Appendix~\ref{app:dualC:certificates}; Sec.~\ref{sec:closed}).
\item $\CSR_6$. The six new inputs are archived exact period expressions, and the closed form of $\CSR_6$ is
 established by the exact certificate at weights two to five (same appendix). This appendix gives the assembly, a
 convergent check of one input and checks that do not use the new single logarithms.
\end{itemize}
Everything holds within the matched two-coupling prescription of Sec.~\ref{sec:dual}, with the caveats stated there
and in Sec.~\ref{sec:d1caveats}.

%--------------------------------------------------------------------
\subsection{Conventions and the beta inputs}\label{app:dualB:conv}
The notation is that of Table~\ref{tab:dualA:notation}. In particular $\db=1-s$; $\varepsilon$ is the hard core and
$\ell=\log\varepsilon$; $\mathcal C_{2m+1}$ and $\mathcal K_{2m}$ are the vacuum-normalized responses of
Eq.~\eqref{eq:dualA:normalization}; a second subscript 0 or 1 denotes the coefficient $[\varepsilon^0\ell^0]$ or
$[\varepsilon^0\ell^1]$ (Eq.~\eqref{eq:dualA:I0I1}); $\partial_\eta$ is the derivative with respect to the endpoint
source at $\eta=0$; and $b_n$ ($n$ odd) and $h_n$ ($n$ even) are the auxiliary beta coefficients of
Eq.~\eqref{eq:dualA:auxbeta}, determined by the recurrence~\eqref{eq:dualA:recurrence}. Derivatives in $v$ and $\db$
are taken at the marginal point $v=1$, $\db=0$, before any fixed point is substituted. We write $\log^22$ for
$(\log2)^2$, and similarly for other powers.

\paragraph{The beta inputs of orders three to six.} The beta inputs are the Taylor coefficients of $F$ and $H$ at the
marginal point. In the compendium's dictionary, with the matching factors of Eq.~\eqref{eq:dualA:fugmatch} restored as
described in Appendix~\ref{app:dualA:recurrence}, and with a subscript $v$ or $\db$ denoting a derivative of the
auxiliary coefficient, the inputs used here are
\begin{equation}
\begin{gathered}
 a=b_{3,v},\qquad p=b_{3,\db},\qquad A=\tfrac12b_{3,vv},\qquad E=b_{3,v\db},\qquad V=\tfrac12b_{3,\db\db},\\
 b=b_5,\qquad G=b_{5,v},\qquad R=b_{5,\db},\qquad J=b_7,\\
 c=-4h_4,\qquad f=-4h_{4,v},\qquad q=-4h_{4,\db},\qquad L=-2h_{4,vv},\qquad M=-4h_{4,v\db},\\
 k=-4h_6,\qquad N=-4h_{6,v}.
\end{gathered}
\label{eq:dualB:dictionary}
\end{equation}
Taylor factors appear only for repeated derivatives; the mixed derivatives $E$ and $M$ carry none. The input $b=b_5$
is not the resummed charge $b_0-\sqrt2h$, and the input $q$ is not a kink charge or a Mellin variable.

\paragraph{Local symbols.} From Appendix~\ref{app:dualB:worked} on, the letters $e$, $t$, $w$, $x$, $y$ and $z$ are
integration or auxiliary variables defined where they are used; there $x$ and $z$ are never $\Ds$ or the coupling $z$.
Where the compendium uses one letter for two quantities, the local quantity is renamed; the renamed symbols are
defined at first use.

%--------------------------------------------------------------------
\subsection{Third and fourth orders: \texorpdfstring{$D_2$}{D2} and \texorpdfstring{$T_2$}{T2}}\label{app:dualB:D2T2}
\paragraph{Beta functions.} To the first new orders the matched beta functions have the form
\begin{equation}
 F=z+(1+az+p\db)u+bu^2+\cdots,\qquad H=-\db+4u+cu^2+fzu^2+\cdots.
 \label{eq:dualB:beta2}
\end{equation}
The one-pair calculation fixes the term $4u$ of $H$, with no further terms $zu$, $z^2u$, $\db u$ or $\db zu$ at these
orders in this prescription. Indeed the one-pair current moment is
\begin{equation}
 Q_2(v,\db;\varepsilon)=\int_\varepsilon^1r(1-r)\,e^{-2vf_{\db}(r)}\dd r,
\end{equation}
with $f_{\db}$ of Eq.~\eqref{eq:dualA:fdb} at $\kappa=1$. At the marginal point its single logarithm is $-\ell$
(Eq.~\eqref{eq:dualA:Z2Q2}), while derivatives in the charge or the tail parameter add powers of $\log r$ to the
integrand and so produce higher powers of $\ell$, not a new single logarithm
(Appendix~\ref{app:dualB:counterterms}). The tree coefficient of the fugacity in the essential variables is exactly
$z$.

\paragraph{Inputs and periods.} The five inputs $a,p,b,c,f$ are listed in Sec.~\ref{sec:inputs1}. They are fixed by
elementary constants and three regulator periods: $\mathcal P_4$ and $\mathcal R_4$, Laurent coefficients of the
angular Mellin function of the four-kink current moment and of its charge derivative, and $\mathcal S_5$, a Laurent
coefficient of the angular Mellin function of the five-kink odd moment (Eq.~\eqref{eq:dualB:laurent}). The
derivations of Appendices~\ref{app:dualB:counterterms} and~\ref{app:dualB:threeins}, and of
Appendix~\ref{app:dualA:recurrence} for $c$, give
\begin{equation}
\begin{gathered}
 a=-4\zeta(2),\qquad c=-16,\qquad p=8-2\zeta(2)-7\zeta(3)-8\log2,\\
 b=-21+38\zeta(2)+60\zeta(3)+32\log2-\mathcal S_5,\qquad
 f=16\mathcal P_4+4\mathcal R_4-32+16\zeta(2)+32\zeta(3),
\end{gathered}
\label{eq:dualB:inputsperiods}
\end{equation}
and Appendices~\ref{app:dualB:worked} and~\ref{app:dualB:periods} evaluate the periods,
\begin{equation}
\begin{aligned}
 \mathcal P_4&=4\zeta(2)+\tfrac83\zeta(3)-2+3\log3+\log^22+2\Li_2(-\tfrac12),\\
 \mathcal R_4&=24-20\zeta(2)-30\zeta(3)-16\log2,\\
 \mathcal S_5&=44\zeta(2)+\tfrac{152}{3}\zeta(3)-3+12\log3+4\log^22+8\Li_2(-\tfrac12).
\end{aligned}
\label{eq:dualB:periods}
\end{equation}
In particular
\begin{equation}
 \mathcal S_5-4\mathcal P_4=5+28\zeta(2)+40\zeta(3).
 \label{eq:dualB:SminusP}
\end{equation}
With Eq.~\eqref{eq:dualB:periods}, Eq.~\eqref{eq:dualB:inputsperiods} gives the values of $b$ and $f$ in
Sec.~\ref{sec:inputs1}. The regulator constants $\log2$, $\log3$, $\log^22$ and $\Li_2(-\frac12)$ are genuine parts of
the individual inputs. In $T_2$ below, $b$ and $f$ enter as
\begin{equation}
 \frac b4+\frac f{16}=-\frac{29}{4}+\frac{21}{2}\zeta(2)+17\zeta(3)+8\log2
 -\frac{\mathcal S_5-4\mathcal P_4}{4}+\frac{\mathcal R_4}{4},
\end{equation}
so the periods enter only through $\mathcal R_4$ and the combination~\eqref{eq:dualB:SminusP}, and the regulator
constants cancel. The evaluation of $\mathcal P_4$ and $\mathcal S_5$ turns on one finite part,
\begin{equation}
 J_L=\lim_{T\to\infty}\Bigl[\int_2^T\frac{\log(t-1)\log(t+1)}{t}\dd t-\frac13\log^3T\Bigr]=-\frac13\zeta(3),
 \label{eq:dualB:JLshort}
\end{equation}
which Appendix~\ref{app:dualB:trilog} proves with explicit primitives and a real-branch Spence--Kummer identity, not
by numerical recognition.

\paragraph{Normalized responses.} With $\ell=\log\varepsilon$, the normalized low-order responses at $v=1$, $\db=0$
begin
\begin{align}
 \mathcal C_3&=2\ell^2+5\ell+3-2\zeta(2)+\cdots,\label{eq:dualB:C3}\\
 \mathcal K_4&=-\frac43\ell^3-6\ell^2+\bigl[-8+4\zeta(2)\bigr]\ell+\mathcal P_4+4\zeta(2)-10+\cdots,\label{eq:dualB:K4}\\
 \mathcal C_5&=\frac{10}{3}\ell^4+\frac{62}{3}\ell^3+\Bigl(\frac{95}{2}-24\zeta(2)\Bigr)\ell^2
 +\bigl[96-40\zeta(2)-\mathcal S_5\bigr]\ell+\cdots,\label{eq:dualB:C5}
\end{align}
where the dots stand for terms that vanish as $\varepsilon\to0$ and, in $\mathcal C_5$, for its finite part, which is
not needed here. Their substitution into the subtraction recurrence~\eqref{eq:dualA:recurrence}, together with the
finite parts and the charge and source derivatives of the lower responses, yields the inputs; the raw single
logarithms alone do not. The finite-tail factor of Eq.~\eqref{eq:dualA:fugmatch} is needed for $p$
(Appendix~\ref{app:dualB:threeins}).

\paragraph{Invariant assembly.} Solving $H=0$ and $F=0$ order by order gives $u_*=\db/4-c\,\db^2/64+O(\db^3)$ and
$z_*=-\db/4+O(\db^2)$. In $D=2u(1+z)(F_uH_z-F_zH_u)$ the derivative $H_z=fu^2+\cdots$ is of order $\db^2$, so through
order $\db^2$, with $F_z=1+au+\cdots$ and $H_u=4+2cu+\cdots$,
\begin{equation}
 D=-2u(1+z)F_zH_u+O(\db^3)=-8u\Bigl(1+z+au+\frac c2u\Bigr)+O(\db^3).
\end{equation}
In $T=2uF_u+(1+z)H_z$, with $F_u=1+az+p\db+2bu+\cdots$,
\begin{equation}
 T=2u\bigl(1+az+p\db+2bu\bigr)+fu^2+O(\db^3).
\end{equation}
Inserting the fixed point gives $D=-2\db+D_2\db^2+\cdots$ and $T=\db/2+T_2\db^2+\cdots$ with
\begin{equation}
 D_2=\frac12-\frac a2-\frac c8,\qquad T_2=-\frac a8+\frac b4-\frac c{32}+\frac f{16}+\frac p2,
\end{equation}
as in Sec.~\ref{sec:inputs1}. With Eqs.~\eqref{eq:dualB:inputsperiods} to~\eqref{eq:dualB:SminusP} these are the
closed forms $D_2=\frac52+2\zeta(2)$ and $T_2=2-2\zeta(2)-4\zeta(3)$ of Sec.~\ref{sec:closed}, and Eq.~\eqref{eq:conv}
gives $\CSR_3=D_2-\frac1{16}$ and $\CSR_4=2T_2$ of Eq.~\eqref{eq:closed}. The highest insertion order needed for
$\CSR_4$ is five (the period $\mathcal S_5$); an insertion count is not a Feynman-loop count.

%--------------------------------------------------------------------
\subsection{Fifth order: \texorpdfstring{$D_3$}{D3}}\label{app:dualB:D3}
The determinant at order $\db^3$ needs the further sectors
\begin{equation}
\begin{aligned}
 F&=z+\bigl(1+az+p\db+Az^2+E\db z\bigr)u+(b+Gz)u^2+\cdots,\\
 H&=-\db+4u+(c+fz+q\db)u^2+ku^3+\cdots.
\end{aligned}
\label{eq:dualB:beta3}
\end{equation}

\paragraph{The three-kink inputs $A$ and $E$.} They have convergent representations and closed forms,
\begin{equation}
\begin{aligned}
 A&=-8-4\zeta(2)-2\zeta(3)+16\log2,\\
 E&=-16+14\zeta(2)-22\zeta(3)+12\zeta(4)+8\log2-8\log^22-16\Li_4(\tfrac12)\\
 &\quad-14\zeta(3)\log2+4\zeta(2)\log^22-\tfrac23\log^42.
\end{aligned}
\label{eq:dualB:AE}
\end{equation}
For $A$ one may evaluate
\begin{equation}
 A=-4\int_0^{1/2}\frac{\log^2(1-y)}{y^2(1-y)^2}\dd y,
\end{equation}
and for $E$ an independent convergent representation is
\begin{equation}
 E=p+4\int_0^{1/2}\frac{\log^2(1-y)\bigl[2\log y-\log(1-y)\bigr]}{y^2(1-y)^2}\dd y.
 \label{eq:dualB:Erep}
\end{equation}
Appendix~\ref{app:dualB:endpoints} reduces $E$ to the weight-four constants of Eq.~\eqref{eq:dualB:AE}; the key step is
\begin{equation}
 \int_0^1\frac{\log t\,\log^2(1+t)}{t}\dd t=\frac23\int_0^1\frac{\log^3(1+t)}{t}\dd t-\frac14\zeta(4).
\end{equation}

\paragraph{Endpoint constants.} Appendix~\ref{app:dualB:endpoints} also derives the finite parts of two source
derivatives of raw moments,
\begin{equation}
 \partial_v\partial_\eta A_{3,0}=-108+60\zeta(2)+22\zeta(3)-34\zeta(4)+16\log2,\qquad
 \partial_\eta Q_{4,0}=6-12\zeta(2)-54\zeta(3)-25\zeta(4).
 \label{eq:dualB:endpointconstants}
\end{equation}
Source derivatives of this kind enter the subtraction recurrence through the Jacobian entries $j_n$ and
$\lambda_{2m}$ of Eq.~\eqref{eq:dualA:jacobian}. The precise moment convention is given in
Appendix~\ref{app:dualB:endpoints}. In particular, the conversion from angular to radial variables must keep
$\log(\varepsilon/R)=\ell-\log R$ in full ($R$ the span of the kinks); dropping the radial logarithms changes finite parts.

\paragraph{The inputs $G$, $q$, $k$.} The remaining inputs, $G=b_{5,v}$, $q=-4h_{4,\db}$ and $k=-4h_6$, are archived
exact period expressions (\file{code/dual_d1/data/archived/C6_closed_form/input_expressions/}). Their values are
\begin{equation}
\begin{aligned}
 G&=32.7634092768376557983124741174\ldots,\\
 q&=-103.5603579270427462132101236545\ldots,\\
 k&=104.6723927432126373310860847487\ldots.
\end{aligned}
\end{equation}
The exact expressions, not these decimals, enter the assembly.

\paragraph{Fixed point and determinant.} Through the orders needed,
\begin{equation}
\begin{aligned}
 u_*&=\frac\db4-\frac c{64}\db^2+\Bigl(\frac{c^2}{512}+\frac{f-k}{256}-\frac q{64}\Bigr)\db^3+\cdots,\\
 z_*&=-\frac\db4+\Bigl(\frac{a-b}{16}+\frac c{64}-\frac p4\Bigr)\db^2+\cdots,
\end{aligned}
\label{eq:dualB:fixedpoint}
\end{equation}
and the determinant gives
\begin{equation}
 D_3=\frac{16A-32E-8G+8b+c^2+4f-4k+32p-8q}{64}.
 \label{eq:dualB:D3assembly}
\end{equation}
With $\mathcal R=2G+k+2q-2b-f$ and $c^2/64=4$, this is $D_3=\frac A4-\frac E2+4+\frac p2-\frac{\mathcal R}{16}$. The
exact certificate of $D_3$ (Appendix~\ref{app:dualC:certificates}; the table of Sec.~\ref{sec:closed}) proves
\begin{equation}
\begin{aligned}
 \mathcal R={}&140+64\zeta(2)-64\zeta(3)-272\zeta(4)-64\log2+64\log^22\\
 &+112\zeta(3)\log2-32\zeta(2)\log^22+\tfrac{16}{3}\log^42+128\Li_4(\tfrac12).
\end{aligned}
\label{eq:dualB:calR}
\end{equation}
With Eq.~\eqref{eq:dualB:AE} and the value of $p$, all regulator constants other than zeta values cancel, and
$D_3=\frac{21}{4}-13\zeta(2)+11\zeta(3)+11\zeta(4)=8.99403863655111357\ldots$, the closed form of
Sec.~\ref{sec:closed}. Equation~\eqref{eq:conv} then gives $\CSR_5$ of Eq.~\eqref{eq:closed},
$29.2378115809900396\ldots$.

%--------------------------------------------------------------------
\subsection{Sixth order: \texorpdfstring{$T_3$}{T3}}\label{app:dualB:C6}
To determine $T_3$, retain
\begin{equation}
\begin{aligned}
 F&=z+u\bigl(1+az+p\db+Az^2+E\db z+V\db^2\bigr)+u^2(b+Gz+R\db)+Ju^3+\cdots,\\
 H&=-\db+4u+u^2\bigl(c+fz+q\db+Lz^2+M\db z\bigr)+u^3(k+Nz)+\cdots.
\end{aligned}
\label{eq:dualB:beta6}
\end{equation}
The six new inputs of Eq.~\eqref{eq:dualB:dictionary},
\begin{equation}
 V=\tfrac12b_{3,\db\db},\quad R=b_{5,\db},\quad J=b_7,\quad L=-2h_{4,vv},\quad M=-4h_{4,v\db},\quad N=-4h_{6,v},
 \label{eq:dualB:C6inputs}
\end{equation}
are evaluated at the marginal point $v=1$, $\db=0$, not at the interacting fixed point. The term $Nzu^3$ is of order
$\db^4$ in $H$ at the fixed point but contributes at order $\db^3$ after a derivative in $z$, so it cannot be dropped.
\begin{center}\small
\begin{tabular}{@{}cll@{}}
\toprule
input & largest new moment sector & value\\
\midrule
$V$ & three-kink second tail & $-19.518741011574465808$\\
$R$ & five-kink tail & $-1.148356985726247710$\\
$J$ & seven-kink odd response & $32.280324973629037003$\\
$L$ & four-kink second charge & $-7.677303759056569914$\\
$M$ & four-kink charge--tail & $13.315217322864044765$\\
$N$ & six-kink charge derivative & $210.030599248618268410$\\
\bottomrule
\end{tabular}
\end{center}
These are beta inputs of the prescription, not universal exponents. Their exact expressions are archived with those
of $G$, $q$, $k$ (\file{code/dual_d1/data/archived/C6_closed_form/input_expressions/}).

\paragraph{A convergent check of $V$.} For three kinks, put $l=\log y$, $\bar l=\log(1-y)$ and
\begin{equation}
 P_V(y)=-l^2\bar l+l\bar l^2-\tfrac13\bar l^3+2l^2\bar l^2-2l\bar l^3+\tfrac12\bar l^4.
\end{equation}
The direct regulated calculation reduces to
\begin{equation}
 V=\frac23\log^32-2\int_0^{1/2}\Bigl[\frac{P_V(y)}{y^2(1-y)^2}-\frac{\log^2y}{y}\Bigr]\dd y.
 \label{eq:dualB:Vintegral}
\end{equation}
The subtraction removes the endpoint singularity. Both the cubic-logarithm term $\frac13S_3$ and the square
$\frac12S_2^2$ of the first tail insertion in Eq.~\eqref{eq:dualA:insertions} are needed. The finite hard-core
endpoint contributes $2\int_1^\infty\log^2w\,w^{-2}\dd w=4$, which cancels a $-4$ in the beta assembly (the coefficient
of $\db^2$ in $2vj_1=-4/(1+\db)$ at $v=1$, Eqs.~\eqref{eq:dualA:jacobian} and~\eqref{eq:dualA:recurrence}). This
finite piece is derived with the regulator kept; it is not an adjustable counterterm.

\paragraph{Trace assembly.} With the fixed point~\eqref{eq:dualB:fixedpoint}, the trace gives
\begin{equation}
 \CSR_6=\mathcal B^{(6)}_{\rm old}+2V+R+\frac38J-\frac18L+\frac14M+\frac1{16}N,
 \label{eq:dualB:C6assembly}
\end{equation}
where the part fixed by the inputs of the lower orders is
\begin{equation}
\begin{aligned}
 \mathcal B^{(6)}_{\rm old}={}&\frac A8-\frac E2-\frac G4+\frac{a^2-ab}{8}+\frac{ac}{16}-\frac{ap}{2}-\frac{bc}{8}
 +\frac{c^2}{64}\\
 &-\frac{cf}{32}-\frac{cp}{8}-\frac f{32}-\frac k{32}-\frac q8,
\end{aligned}
\end{equation}
with the archived value $\mathcal B^{(6)}_{\rm old}=-48.0003600840977427059\ldots$. The complete regulated expression
gives $\CSR_6=-58.66569747422530894971714612863082\ldots$. The exact certificate at weights two to five (Appendix~\ref{app:dualC:certificates}) establishes the closed form of $\CSR_6$ in Eq.~\eqref{eq:closed}, and $T_3=\CSR_6/4$ is the
expression of Sec.~\ref{sec:closed}.

\paragraph{Checks that do not use the new single logarithms.} The lower flow predicts, for the normalized seven-kink
response,
\begin{equation}
 [\ell^6]\,\mathcal C_7=\frac{244}{45},\qquad [\ell^5]\,\mathcal C_7=\frac{170}{3},\qquad
 [\ell^4]\,\mathcal C_7=247-\frac{316}{3}\zeta(2),
\end{equation}
and the archived exact integrations reproduce these coefficients without using $J$. For the finite-tail correction
to the five-kink response, the matching of Eq.~\eqref{eq:dualA:fugmatch} predicts
\begin{equation}
 [\ell^6]\,\partial_{\db}\mathcal C_5=-\frac{142}{45},\qquad [\ell^5]\,\partial_{\db}\mathcal C_5=-\frac{71}{5},\qquad
 [\ell^4]\,\partial_{\db}\mathcal C_5=\frac{14}{3}\zeta(2)-\frac{55}{3},
\end{equation}
which do not use $R$. The archive contains further higher-logarithm checks, exact derivatives of the primitives and
direct finite-cutoff numerical integrations of the original kernels. Some period residuals in the higher-logarithm
tests were checked numerically, not by exact certificates, and the verification reports keep that distinction.

Two archived full period evaluations used 80 and 100 decimal digits, matching points 100 and 120, and different
continuation steps and series orders; they agreed in all 65 output digits of the first. This is a convergence
check, not an interval bound; the closed form of $\CSR_6$ rests on the rational certificate, not on numerical
agreement.

%--------------------------------------------------------------------
\subsection{Mellin representation, finite parts and the period \texorpdfstring{$\mathcal R_4$}{R4}}\label{app:dualB:worked}
This and the following subsections give the low-insertion calculations behind Appendix~\ref{app:dualB:D2T2}, with
their subtractions and branch conventions. The normalized responses are those of
Eq.~\eqref{eq:dualA:normalization}, and the gap variables are those of Appendix~\ref{app:dualA:moments}.

\paragraph{Separating the radial integral.} Let $F$ be homogeneous of degree $d_F$ in $m$ gaps, and put $k_F=m+d_F$.
Define, with the translation factor of Appendix~\ref{app:dualA:moments},
\begin{equation}
 I_F(\varepsilon)=\int_{r_i\ge\varepsilon,\ \sum r_i\le1}\Bigl(1-\sum_ir_i\Bigr)F(\boldsymbol r)\,\dd^mr.
\end{equation}
First integrate $\varepsilon$ up to $\min_ir_i$, then write $r_i=Rx_i$ with $\sum_ix_i=1$. This gives
\begin{equation}
 \int_0^\infty\varepsilon^{q-1}I_F(\varepsilon)\dd\varepsilon=\frac{\mathcal A_F(q)}{q(q+k_F)(q+k_F+1)},\qquad
 \mathcal A_F(q)=\int_{\sum x_i=1}F(\boldsymbol x)\bigl(\min_ix_i\bigr)^q\dd^{m-1}x.
 \label{eq:dualB:mellin}
\end{equation}
Here $q$ is the Mellin variable, not the input $q$. The identity is first used in a half-plane where it converges and
is then continued after explicit subtraction; no divergent integral is evaluated naively at $q=0$. The pole of order
$n$ at $q=0$ gives the $\varepsilon^0$ terms of $I_F$ through
\begin{equation}
 q^{-n}\;\longleftrightarrow\;\frac{(-\ell)^{n-1}}{(n-1)!},
 \label{eq:dualB:poles}
\end{equation}
so the coefficient of $q^{-1}$ is the finite part; power divergences give poles away from $q=0$.

The needed Laurent coefficients of the angular functions define the regulator periods:
\begin{equation}
\begin{aligned}
 \mathcal A_{Q_4}(q)\big|_{v=1}&=\frac8{q^2}-\frac6q-6-4\zeta(2)+\mathcal P_4\,q+O(q^2),\\
 \partial_v\mathcal A_{Q_4}(q)\big|_{v=1}&=\frac{32}{q^3}-\frac4{q^2}+\frac{4+16\zeta(2)}{q}+\mathcal R_4+O(q),\\
 \mathcal A_{Z_4}(q)&=\frac{12}{q^2}-\frac3q-5-6\zeta(2)+O(q),\\
 \mathcal A_{A_5}(q)&=\frac{80}{q^3}-\frac{68}{q^2}-\frac{56+48\zeta(2)}{q}+\mathcal S_5+O(q).
\end{aligned}
\label{eq:dualB:laurent}
\end{equation}
The subscript names the moment whose kernel is meant: $\Xi_4W_4$ for $Q_4$, $W_4$ for $Z_4$ and $W_5$ for $A_5$ (the
compendium writes $G_5$ for the last). For $Q_4$, $k_F=4-4v$, so a charge derivative of Eq.~\eqref{eq:dualB:mellin}
also differentiates its radial denominator; that contribution must be included
(Appendix~\ref{app:dualB:counterterms}). For example, at $v=1$ ($k_F=0$) the first line and
Eq.~\eqref{eq:dualB:poles} give the $\varepsilon^0$ part of $Q_4$,
$-\frac43\ell^3-7\ell^2+[-8+4\zeta(2)]\ell+\mathcal P_4+4\zeta(2)-8$; subtracting the $\varepsilon^0$ part of $Q_2Z_2$,
which is $2-\ell^2$ by Eq.~\eqref{eq:dualA:Z2Q2}, gives $\mathcal K_4$ of Eq.~\eqref{eq:dualB:K4}.

\paragraph{Hard-cut angular functions.} Define
\begin{equation}
 \mathcal H_F(e)=\int_{x_i\ge e,\ \sum x_i=1}F(\boldsymbol x)\dd^{m-1}x,\qquad
 \mathcal A_F(q)=q\int_0^{1/m}e^{q-1}\mathcal H_F(e)\dd e.
 \label{eq:dualB:angularH}
\end{equation}
The constant term ($e^0\log^0e$) of $\mathcal H_{A_5}$ is $\mathcal S_5$, and that of $\partial_v\mathcal H_{Q_4}$ is
$\mathcal R_4$. Power divergences produce poles away from $q=0$ and do not change these identifications. The period
$\mathcal P_4$ is instead a finite part of $\int\mathcal H_{Q_4}(e)\dd e/e$ (Appendix~\ref{app:dualB:periods}).

\paragraph{Finite parts.} Throughout, $\FP$ denotes the constant coefficient that remains after all displayed powers
and logarithms of the endpoint cutoff are subtracted; for instance,
\begin{equation}
 \FP\!\int_2^\infty f(t)\dd t=\text{constant term of }\int_2^Tf(t)\dd t\text{ as }T\to\infty.
 \label{eq:dualB:FPdef}
\end{equation}
On the interval $(0,1)$ the cutoffs are $\vartheta$ and $1-\vartheta$. Finite parts are not invariant under a change of
cutoff variable unless the corresponding conversion terms are kept; an example occurs in the evaluation of
$\mathcal P_4$.

To extract the constant of a one-dimensional angular integral with moving endpoints, one uses three regions: the bulk,
a lower endpoint region $w=et$ and an upper endpoint region $w=1-et$. The integrand, \emph{including the Jacobian}, is
expanded in each region, and its coefficient of $e^0\log^0e$ is kept; the desired constant is the sum of the three
finite-part integrals. To justify this, split the original domain at fixed $\vartheta$ and $1-\vartheta$, take
$e\ll\vartheta$ first, and cancel the powers and logarithms of the overlaps before taking $\vartheta\to0$; the constant
that remains is the sum stated. This procedure keeps finite boundary contributions that an expansion of the bulk alone
would miss. In symbolic software, a term such as $1/[e(t+1)]$ must be recognized as $e^{-1}/(t+1)$ before a Laurent
coefficient is read off; reading a coefficient from an uncollected rational denominator can silently drop an endpoint
term.

\paragraph{Elementary finite parts.} With $l=\log w$ and $\bar l=\log(1-w)$, the nontrivial finite parts over $(0,1)$
needed below are (convergent integrals take their ordinary values):
\begin{center}\small
\renewcommand{\arraystretch}{1.25}
\begin{tabular}{@{}ll@{\qquad\qquad}ll@{}}
\toprule
integrand & $\FP\int_0^1$ & integrand & $\FP\int_0^1$\\
\midrule
$l\bar l$ & $2-\zeta(2)$ & $\bar l/w$ & $-\zeta(2)$\\
$l\bar l/w$ & $\zeta(3)$ & $\bar l/(w-1)$ & $0$\\
$l\bar l/(w-1)$ & $-\zeta(3)$ & $\bar l/(w-1)^2$ & $-1$\\
$l\bar l/w^2$, $l\bar l/(w-1)^2$ & $\zeta(2)-1$ & $\bar l/(w-1)^3$ & $\frac14$\\
$l$ & $-1$ & $l^2/w$ & $0$\\
$l/w$ & $0$ & $l^2/w^2$ & $-2$\\
$l/w^2$ & $-1$ & $l^2/(w+1)$ & $\frac32\zeta(3)$\\
$l/(w-1)$ & $\zeta(2)$ & $l^2/(w+1)^2$ & $\zeta(2)$\\
$l/(w-1)^2$ & $-1$ & $\bar l^2$, $\bar l^2/(w-1)$ & $2$, $0$\\
$l/(w+1)$ & $-\frac12\zeta(2)$ & $\bar l^2/w$, $\bar l^2/w^2$ & $2\zeta(3)$, $2\zeta(2)$\\
$l/(w+1)^2$ & $-\log2$ & $1/(w+1)$, $1/(w+1)^2$ & $\log2$, $\frac12$\\
$l/(w+1)^3$ & $-\frac12\log2-\frac14$ & $1/w$, $1/(w-1)$ & $0$, $0$\\
& & $1/w^2$, $1/(w-1)^2$ & $-1$, $-1$\\
& & $1/(w-1)^3$ & $\frac12$\\
\bottomrule
\end{tabular}
\end{center}
For example,
\begin{equation}
 \FP\!\int_0^1\frac{\log w\log(1-w)}{w^2}\dd w=\sum_{n=2}^\infty\frac1{n(n-1)^2}=\zeta(2)-1,
\end{equation}
where the term $n=1$ of the expansion of $\log(1-w)$ gives a pure logarithmic divergence with zero finite constant.
Likewise
\begin{equation}
 \int_0^1\frac{\log^2w}{1+w}\dd w=2\sum_{n=1}^\infty\frac{(-1)^{n-1}}{n^3}=\frac32\zeta(3).
\end{equation}
Repeated denominators follow by integration by parts, and pure powers and logarithms can be integrated directly; these
procedures derive the whole table.

\paragraph{The period $\mathcal R_4$.} On the angular simplex $r+t+w=1$ of three gaps, Eq.~\eqref{eq:dualA:W4}
simplifies to
\begin{equation}
 W_4=\frac1{r^2w^2}+\frac2{rtw}+\frac1{t^2},
 \label{eq:dualB:W4decomp}
\end{equation}
and the current kernel is $\Xi_4W_4=(r+w)W_4$. For a fixed outer gap $r$, let $\mathcal M_{\rm out}(r,e)$ be the
integral of $\Xi_4W_4$ over the other outer gap, and for a fixed middle gap $r$ let $\mathcal M_{\rm mid}(r,e)$ be the
integral over an outer gap. Elementary partial fractions give, with $\rho_r=e^{-1}-(1-r-e)^{-1}$ and
$L_r=\log[(1-r-e)/e]$,
\begin{equation}
\begin{aligned}
 \mathcal M_{\rm out}(r,e)&=\bigl(1+r^{-1}\bigr)\rho_r+\frac{(r+1)(r^2+1)}{r^2(1-r)}L_r,\\
 \mathcal M_{\rm mid}(r,e)&=\frac{2\rho_r}{1-r}+\frac{4(r^2-r+1)}{r(1-r)^2}L_r+\frac{(1-r)(1-r-2e)}{r^2}.
\end{aligned}
\label{eq:dualB:Mout}
\end{equation}
On the simplex $\log W_4=2\log[(1-r)/r]+2\log[(1-w)/w]-2\log t$, and $\partial_vW_4^v|_{v=1}=W_4\log W_4$. The symmetry
between the two outer gaps then gives the exact one-dimensional reduction
\begin{equation}
 \partial_v\mathcal H_{Q_4}(e,v)\big|_{v=1}=\int_e^{1-2e}\Bigl[4\mathcal M_{\rm out}(r,e)\log\frac{1-r}{r}
 -2\mathcal M_{\rm mid}(r,e)\log r\Bigr]\dd r.
 \label{eq:dualB:Hprime}
\end{equation}
Its three regions are the bulk (at fixed $r$, renamed $w$), the lower region $r=et$ with $t\ge1$, and the upper region
$r=1-et$ with $t\ge2$. With $l=\log w$ and $\bar l=\log(1-w)$, the constant integrands are
\begin{align}
 B_4(w)={}&l\bar l\Bigl(4+\frac{16}{w-1}-\frac8{(w-1)^2}-\frac{16}{w}-\frac4{w^2}\Bigr)
 +l\Bigl(-2-\frac8{w-1}+\frac4{(w-1)^2}+\frac8w-\frac2{w^2}\Bigr)\notag\\
 &+\bar l^2\Bigl(-4-\frac{16}{w-1}+\frac8w+\frac4{w^2}\Bigr)+\bar l\Bigl(\frac8{w-1}-\frac4w\Bigr),\label{eq:dualB:B4}\\
 L_4(t)={}&\Bigl(-8+\frac{12}{t}+\frac8{t^2}\Bigr)\log t-4,\qquad t\ge1,\\
 U_4(t)={}&\frac{16}{t}\log t\log(t-1)+\Bigl(8-\frac8{t-1}\Bigr)\log t+\frac8t\log(t-1)+4-\frac4{t-1},\qquad t\ge2.
\end{align}
The finite-part table gives, group by group,
\begin{center}\small
\begin{tabular}{@{}lllll@{}}
\toprule
group of $B_4$ & $l\bar l$ & $l$ & $\bar l^2$ & $\bar l$\\
\midrule
$\FP\int_0^1$ & $20-16\zeta(2)-32\zeta(3)$ & $-8\zeta(2)$ & $-8+8\zeta(2)+16\zeta(3)$ & $4\zeta(2)$\\
\bottomrule
\end{tabular}
\end{center}
so that
\begin{equation}
 \FP\!\int_0^1B_4(w)\dd w=12-12\zeta(2)-16\zeta(3),\qquad \FP\!\int_1^\infty L_4(t)\dd t=4.
\end{equation}
In the upper integral, the term with two logarithms is
\begin{equation}
 J_U=\FP\!\int_2^\infty\frac{\log t\log(t-1)}{t}\dd t=-\frac{\log^32}{3}-\log2\,\Li_2(\tfrac12)-\Li_3(\tfrac12)
 =-\frac78\zeta(3),
 \label{eq:dualB:JU}
\end{equation}
where the last step uses the values~\eqref{eq:dualB:halfvalues} at $\frac12$. The remaining terms have the primitive
\begin{equation}
 8t\log t-4t+8\Li_2(1-t)+4\log^2t+8\Li_2(1/t)-4\log(t-1),
\end{equation}
whose constant at infinity is $-8\zeta(2)$ and whose value at $t=2$ is $16\log2-8$. Hence
\begin{equation}
 \FP\!\int_2^\infty U_4(t)\dd t=8-8\zeta(2)-14\zeta(3)-16\log2.
\end{equation}
Adding the three regions gives
\begin{equation}
 \mathcal R_4=24-20\zeta(2)-30\zeta(3)-16\log2,
 \label{eq:dualB:Rclosed}
\end{equation}
with no numerical reconstruction.

%--------------------------------------------------------------------
\subsection{The trilogarithm lemma and the Spence--Kummer identity}\label{app:dualB:trilog}
\paragraph{Polylogarithms.} For $|z|<1$, $\Li_n(z)=\sum_{k\ge1}z^k/k^n$, continued by
$\dd\Li_n(z)/\dd z=\Li_{n-1}(z)/z$ with $\Li_1(z)=-\log(1-z)$. Real principal values are used for $z<1$; where an
argument exceeds one, the real part is indicated explicitly. The identities needed are
\begin{align}
 \Li_2(z)+\Li_2(1-z)&=\zeta(2)-\log z\log(1-z),\label{eq:dualB:Euler}\\
 \Li_2\Bigl(\frac{z}{z-1}\Bigr)&=-\Li_2(z)-\frac12\log^2(1-z),\label{eq:dualB:Li2Landen}\\
 \Li_n(z)+\Li_n(-z)&=2^{1-n}\Li_n(z^2),\label{eq:dualB:dup}\\
 \Li_3(z)+\Li_3(1-z)+\Li_3\Bigl(-\frac z{1-z}\Bigr)&=\zeta(3)+\zeta(2)\log(1-z)-\frac12\log z\log^2(1-z)
 +\frac16\log^3(1-z).\label{eq:dualB:Li3Landen}
\end{align}
The first two follow by differentiation and evaluation at an endpoint, the third from the power series, and the
fourth by differentiation using the first two. At $z=\frac12$ they give
\begin{equation}
 \Li_2(\tfrac12)=\frac{\zeta(2)}{2}-\frac{\log^22}{2},\qquad
 \Li_3(\tfrac12)=\frac78\zeta(3)-\frac{\zeta(2)\log2}{2}+\frac{\log^32}{6}.
 \label{eq:dualB:halfvalues}
\end{equation}

\paragraph{The real-domain Spence--Kummer identity.} For real $0<x<y<1$ define
\begin{equation}
\begin{array}{lll}
 f_1=\dfrac{x(1-y)^2}{y(1-x)^2},\quad & f_2=xy,\quad & f_3=x/y,\\[3mm]
 f_4=\dfrac{x(1-y)}{y(1-x)}, & f_5=-\dfrac{x(1-y)}{1-x}, & f_6=\dfrac{1-y}{1-x},\\[3mm]
 f_7=-\dfrac{1-y}{y(1-x)}, & f_8=x, & f_9=y.
\end{array}
\end{equation}
All positive arguments are smaller than one, and the two negative ones are on the real principal branch. The identity
is
\begin{equation}
 \Li_3(f_1)+\Li_3(f_2)+\Li_3(f_3)-2\sum_{i=4}^{9}\Li_3(f_i)+2\zeta(3)
 =\log^2y\,\log\frac{1-y}{1-x}-2\zeta(2)\log y-\frac13\log^3y.
 \label{eq:dualB:Kummerfull}
\end{equation}
The branch choices are the real-domain specialization of the path formulation of Ref.~\cite{Shiraishi:2023}.

A differential verification uses the five-term relation of the dilogarithm,
\begin{equation}
 \Li_2(a)+\Li_2(b)=\Li_2(ab)+\Li_2\Bigl(\frac{a(1-b)}{1-ab}\Bigr)+\Li_2\Bigl(\frac{b(1-a)}{1-ab}\Bigr)
 +\log\frac{1-a}{1-ab}\log\frac{1-b}{1-ab},
 \label{eq:dualB:Abel}
\end{equation}
together with Eqs.~\eqref{eq:dualB:Euler} and~\eqref{eq:dualB:Li2Landen}; Eq.~\eqref{eq:dualB:Abel} itself follows by
differentiation in $a$ and $b$ and evaluation at $a=0$. Differentiating Eq.~\eqref{eq:dualB:Kummerfull} with respect
to $x$ and multiplying by $x(1-x)$ leaves the dilogarithm equality
\begin{equation}
\begin{aligned}
 &(1+x)\Li_2(f_1)+(1-x)\bigl[\Li_2(f_2)+\Li_2(f_3)\bigr]-2\bigl[\Li_2(f_4)+\Li_2(f_5)\bigr]\\
 &\quad-2x\bigl[\Li_2(f_6)+\Li_2(f_7)\bigr]-2(1-x)\Li_2(x)=x\log^2y.
\end{aligned}
\label{eq:dualB:Kummerdiff}
\end{equation}
Substituting the rational $f_i$ into Eq.~\eqref{eq:dualB:Abel}, followed by Eqs.~\eqref{eq:dualB:Euler}
and~\eqref{eq:dualB:Li2Landen}, verifies this equality. The integration constant is fixed by $x\to0$, where
Eq.~\eqref{eq:dualB:Kummerfull} reduces exactly to the Landen relation~\eqref{eq:dualB:Li3Landen} at $1-y$. This
proves the real-domain identity without a numerically inferred special value.

At $x=\frac12$, $y=\frac23$ the nine arguments are
\begin{equation}
 (f_1,\ldots,f_9)=\Bigl(\frac13,\frac13,\frac34,\frac12,-\frac13,\frac23,-1,\frac12,\frac23\Bigr).
\end{equation}
Since $\Li_3(-1)=-\frac34\zeta(3)$, Eq.~\eqref{eq:dualB:Kummerfull} becomes
\begin{equation}
 2\Li_3(\tfrac13)+\Li_3(\tfrac34)-4\Li_3(\tfrac12)-2\Li_3(-\tfrac13)-4\Li_3(\tfrac23)+\frac72\zeta(3)
 =\frac23\log^3\tfrac23-2\zeta(2)\log\tfrac23.
 \label{eq:dualB:Kspecial}
\end{equation}
All its arguments are real and at most one, so no continuation across a branch cut of the trilogarithm is needed. The
compendium's script \file{code/dual_d1/code/archived/LRI_thermal_closed_form/code/verify_closed_forms.py} also checks
this specialization and the Landen eliminations that follow.

\paragraph{Lemma (a finite part with an exact zeta value).} With the finite-part convention
of Eq.~\eqref{eq:dualB:FPdef},
\begin{equation}
 J_L=\lim_{T\to\infty}\Bigl[\int_2^T\frac{\log(t-1)\log(t+1)}{t}\dd t-\frac13\log^3T\Bigr]=-\frac13\zeta(3).
 \label{eq:dualB:JL}
\end{equation}

\emph{Proof.} Define the convergent integral
\begin{equation}
 \Lambda(z)=\int_0^z\frac{\log(1-t)\log(1+t)}{t}\dd t.
\end{equation}
The substitution $t\mapsto1/t$ in Eq.~\eqref{eq:dualB:JL} gives the absolutely convergent representation
\begin{equation}
 J_L=\Lambda(\tfrac12)-\frac12\log2\,\Li_2(\tfrac14)-\frac14\Li_3(\tfrac14)-\frac13\log^32.
 \label{eq:dualB:JLhalf}
\end{equation}
For an explicit reduction introduce
\begin{equation}
 \Theta(z)=\int_0^z\frac{\log^2(1-t)}{t}\dd t=\log z\log^2(1-z)+2\log(1-z)\Li_2(1-z)-2\Li_3(1-z)+2\zeta(3);
 \label{eq:dualB:Theta}
\end{equation}
the equality follows by differentiation and the limit $z\to0$. Squaring $\log(1-t^2)=\log(1-t)+\log(1+t)$ gives
\begin{equation}
 \Lambda(z)=\frac14\Theta(z^2)-\frac12\Theta(z)-\frac12\Theta(-z).
\end{equation}
In $\Theta(-\frac12)$ the real parts of Eq.~\eqref{eq:dualB:Theta} are taken. Equivalently, for $a>0$ the manifestly
real formula
\begin{equation}
\begin{aligned}
 \Theta(-a)&=-\log a\log^2(1+a)-2\zeta(2)\log(1+a)-2\log(1+a)\Li_2(-a)\\
 &\quad-2\Li_3\Bigl(\frac1{1+a}\Bigr)+\frac13\log^3(1+a)+2\zeta(3)
\end{aligned}
\end{equation}
holds; it follows from Eq.~\eqref{eq:dualB:Euler} and from
$\operatorname{Re}\Li_3(\beta)=\Li_3(1/\beta)-\frac16\log^3\beta+2\zeta(2)\log\beta$ for real $\beta>1$. Let
\begin{equation}
 \Pi=\log^33-3\log^23\log2+6\log3\log^22-5\log^32-6\zeta(2)\log3+6\zeta(2)\log2.
\end{equation}
Substituting Eq.~\eqref{eq:dualB:Theta} and the formula for $\Theta(-a)$ into Eq.~\eqref{eq:dualB:JLhalf}, and using
Eqs.~\eqref{eq:dualB:Euler} to~\eqref{eq:dualB:dup}, gives
\begin{equation}
 J_L=\Li_3(\tfrac23)-\frac12\Li_3(\tfrac34)-\Li_3(-\tfrac12)-\frac32\zeta(3)-\frac\Pi6.
 \label{eq:dualB:JLreduced}
\end{equation}
The dilogarithms are eliminated in favour of $\Li_2(\frac13)$,
\begin{equation}
\begin{gathered}
 \Li_2(-\tfrac12)=-\Li_2(\tfrac13)-\tfrac12\log^2\tfrac32,\qquad
 \Li_2(\tfrac14)=\zeta(2)-\log^22-2\Li_2(\tfrac13)-\log^2\tfrac32,\\
 \Li_2(\tfrac34)=\zeta(2)-\Li_2(\tfrac14)+2\log2\,\log\tfrac34,
\end{gathered}
\end{equation}
and all coefficients of $\Li_2(\frac13)$ cancel. Next, Eq.~\eqref{eq:dualB:Li3Landen} at $z=\frac13$ and $z=\frac14$
eliminates $\Li_3(\frac13)$ and $\Li_3(-\frac13)$ from Eq.~\eqref{eq:dualB:Kspecial}. With
$\Li_3(\frac14)=4[\Li_3(\frac12)+\Li_3(-\frac12)]$ and Eq.~\eqref{eq:dualB:halfvalues} the result is
\begin{equation}
 \Li_3(\tfrac23)-\frac12\Li_3(\tfrac34)-\Li_3(-\tfrac12)=\frac76\zeta(3)+\frac\Pi6.
 \label{eq:dualB:keycombination}
\end{equation}
Substitution in Eq.~\eqref{eq:dualB:JLreduced} gives $J_L=-\zeta(3)/3$. Every functional identity is used with its
branch specified, and the final cancellation is symbolic, not an integer-relation fit.\hfill$\square$

%--------------------------------------------------------------------
\subsection{Closed forms of the four- and five-insertion periods}\label{app:dualB:periods}
\paragraph{The four-insertion period $\mathcal P_4$.} The decomposition~\eqref{eq:dualB:W4decomp} allows the angular
current kernel to be integrated directly. Put
\begin{equation}
 L_e=\log\frac{1-2e}{e},\qquad z_e=\frac e{1-e},\qquad 0<e<\frac13.
\end{equation}
The result is
\begin{equation}
 \mathcal H_{Q_4}(e)=\frac{2L_e+1}{e}-\frac{4L_e}{1-e}-2(1-e)L_e-2-3e+4\log^2z_e-4\zeta(2)+8\Li_2(z_e),
 \label{eq:dualB:HQexact}
\end{equation}
with the small-$e$ expansion
\begin{equation}
 \mathcal H_{Q_4}(e)=\frac{-2\log e+1}{e}+4\log^2e+6\log e-6-4\zeta(2)+O(e\log e).
\end{equation}
The finite part that corresponds to the first line of Eq.~\eqref{eq:dualB:laurent} is therefore
\begin{equation}
 \mathcal P_4=\lim_{\vartheta\to0}\Bigl[\int_\vartheta^{1/3}\frac{\mathcal H_{Q_4}(e)}{e}\dd e+\frac{2\log\vartheta+1}{\vartheta}
 +\frac43\log^3\vartheta+3\log^2\vartheta-\bigl(6+4\zeta(2)\bigr)\log\vartheta\Bigr].
 \label{eq:dualB:Pfinitepart}
\end{equation}

An explicit antiderivative evaluates it. Change variables to $z=e/(1-e)$, so that the upper endpoint is $z=\frac12$,
and write $\lambda(z)=\log(1-z)-\log z$. The integrand becomes
\begin{equation}
 A_P(z)\lambda(z)+B_P(z)+\frac{4\log^2z-4\zeta(2)+8\Li_2(z)}{z(1+z)},
\end{equation}
with
\begin{equation}
 A_P(z)=\frac2{1+z}+\frac2{(1+z)^2}-\frac6z+\frac2{z^2},\qquad
 B_P(z)=\frac2{1+z}-\frac3{(1+z)^2}-\frac2z+\frac1{z^2}.
\end{equation}
A primitive is $K_P=K_{P,\le2}+K_{P,3}$, with
\begin{align}
 K_{P,\le2}(z)={}&2\Bigl[(\log2-\log z)\log(1+z)-\Li_2\Bigl(\frac{1+z}{2}\Bigr)-\Li_2(-z)\Bigr]
 -2\lambda(z)\Bigl(\frac1z+\frac1{1+z}\Bigr)\notag\\
 &+\frac1z+\frac3{1+z}+6\Li_2(z)+3\log^2z-6\log z+3\log(1-z^2),\\
 K_{P,3}(z)={}&\frac43\log^3z-4\log^2z\log(1+z)-8\log z\Li_2(-z)+8\Li_3(-z)+8\Li_3(z)\notag\\
 &-8\Li_2(z)\log(1+z)-8\Lambda(z)-4\zeta(2)\log\frac z{1+z},
\end{align}
which is verified directly by differentiation, with $\Lambda'(z)=\log(1-z)\log(1+z)/z$ ($\Lambda$ as in
Appendix~\ref{app:dualB:trilog}). The lower endpoint is $z=\vartheta/(1-\vartheta)$, not $z=\vartheta$, and there
\begin{equation}
 K_P\Bigl(\frac\vartheta{1-\vartheta}\Bigr)=\frac{2\log\vartheta+1}{\vartheta}+\frac43\log^3\vartheta+3\log^2\vartheta
 -\bigl(6+4\zeta(2)\bigr)\log\vartheta
 +6-2\Li_2(\tfrac12)+o(1).
\end{equation}
Consequently
\begin{equation}
 \mathcal P_4=K_P(\tfrac12)-6+2\Li_2(\tfrac12).
\end{equation}
The representation~\eqref{eq:dualB:JLhalf}, Euler reflection and duplication reduce this to
\begin{equation}
 \mathcal P_4=4\zeta(2)-2+3\log3+\log^22+2\Li_2(-\tfrac12)-8J_L,
\end{equation}
and the lemma~\eqref{eq:dualB:JL} gives
\begin{equation}
 \mathcal P_4=4\zeta(2)+\frac83\zeta(3)-2+3\log3+\log^22+2\Li_2(-\tfrac12).
 \label{eq:dualB:Pclosed}
\end{equation}
The logarithms and the dilogarithm are allowed in this regulator-dependent period; they cancel from the stability
invariants (Appendix~\ref{app:dualB:D2T2}).

\paragraph{The five-insertion period $\mathcal S_5$: the middle vertex.} On the angular simplex of five kinks, fix the
first and last kinks at $0$ and $1$ and write the intermediate positions as $0<\alpha<\xi<\beta<1$, with every
consecutive gap at least $e$. Let $w=\beta-\alpha$ and $P_0(x)=x^2(1-x)^2$. Then
\begin{equation}
 W_5=\frac{w^2}{P_0(\alpha)P_0(\beta)}\,\frac{P_0(\xi)}{(\xi-\alpha)^2(\beta-\xi)^2}.
\end{equation}
With $\xi=\alpha+t$, the rational identity
\begin{equation}
 \frac{P_0(\alpha+t)}{t^2(w-t)^2}=1+\frac{\mu_1}{t}+\frac{\mu_2}{w-t}+\frac{\nu_1}{t^2}+\frac{\nu_2}{(w-t)^2}
\end{equation}
has the coefficients
\begin{equation}
 \nu_1=\frac{P_0(\alpha)}{w^2},\quad \nu_2=\frac{P_0(\beta)}{w^2},\quad
 \mu_1=\frac{P_0'(\alpha)}{w^2}+\frac{2P_0(\alpha)}{w^3},\quad
 \mu_2=-\frac{P_0'(\beta)}{w^2}+\frac{2P_0(\beta)}{w^3}.
\end{equation}
The $t$ integration runs from $e$ to $w-e$ and leaves
\begin{equation}
 \frac{w^2(w-2e)+\mathcal N_1\log[(w-e)/e]+\mathcal N_2\bigl[e^{-1}-(w-e)^{-1}\bigr]}{P_0(\alpha)P_0(\beta)},
 \label{eq:dualB:postcentral}
\end{equation}
where $\mathcal N_2=P_0(\alpha)+P_0(\beta)$ and $\mathcal N_1=P_0'(\alpha)-P_0'(\beta)+2\mathcal N_2/w$. The
integration over $\alpha$, from $e$ to $1-e-w$, is again elementary, and leaves
\begin{equation}
 \mathcal H_{A_5}(e)=\int_{2e}^{1-2e}I_5(w,e)\dd w.
\end{equation}
With
\begin{equation}
\begin{gathered}
 L_0=\log\frac{w-e}{e},\qquad \rho_0=\frac1e-\frac1{w-e},\qquad
 \rho_-=\frac1e-\frac1{1-e-w},\qquad \rho_+=\frac1{w+e}-\frac1{1-e},\\
 L_-=\log\frac{1-e-w}{e},\qquad L_+=\log\frac{1-e}{w+e},
\end{gathered}
\end{equation}
the exact remaining kernel is
\begin{equation}
\begin{aligned}
 I_5(w,e)={}&\rho_-\Bigl[\frac{2(w-2e)}{(w-1)^2}-\frac{4L_0}{w-1}+2\rho_0\Bigr]
 +\rho_+\Bigl[\frac{2(w-2e)}{(w+1)^2}-\frac{4L_0}{w+1}+2\rho_0\Bigr]\\
 &+L_-\Bigl[\frac{4(w^2-3w+1)(w-2e)}{w(w-1)^3}+\Bigl(-\frac8w+\frac8{w^2(w-1)^2}\Bigr)L_0+4\rho_0\Bigr]\\
 &+L_+\Bigl[\frac{4(w^2+3w+1)(w-2e)}{w(w+1)^3}+\Bigl(-\frac8w-\frac8{w^2(w+1)^2}\Bigr)L_0+4\rho_0\Bigr].
\end{aligned}
\label{eq:dualB:I5}
\end{equation}
Its equality with the $\alpha$ integral of Eq.~\eqref{eq:dualB:postcentral} can be checked by rational partial
fractions in $\alpha$; the compendium's code also compares both sides numerically at finite positive cutoffs, before
any subtraction or expansion. This makes the reduction reproducible without a sector-decomposition package.

\paragraph{The five-insertion period $\mathcal S_5$: bulk and endpoint finite parts.} The three regions are obtained by
expanding $I_5(w,e)$, $e\,I_5(et,e)$ and $e\,I_5(1-et,e)$ and keeping the coefficient of $e^0\log^0e$. The lower and
upper ranges begin at $t=2$, because the middle separation satisfies $w\ge2e$ and the outer complement $1-w\ge2e$.
With $l=\log w$ and $\bar l=\log(1-w)$, the bulk integrand is
\begin{equation}
\begin{aligned}
 B_5(w)={}&l^2\Bigl(\frac{16}{w+1}+\frac8{(w+1)^2}-\frac8w+\frac8{w^2}\Bigr)
 +l\bar l\Bigl(-\frac{16}{w-1}+\frac8{(w-1)^2}+\frac8w+\frac8{w^2}\Bigr)\\
 &+l\Bigl(\frac4{w+1}-\frac4{(w+1)^2}+\frac4{(w+1)^3}-\frac4{(w-1)^2}\Bigr)
 +\bar l\Bigl(\frac4{w-1}-\frac4{(w-1)^2}-\frac4{(w-1)^3}-\frac4w\Bigr)\\
 &-6-\frac2{w+1}+\frac4{(w+1)^2}+\frac6{w-1}-\frac4{(w-1)^2}+\frac2{(w-1)^3}-\frac4w-\frac6{w^2}.
\end{aligned}
\label{eq:dualB:B5}
\end{equation}
The lower region $w=et$ gives
\begin{equation}
 L_5(t)=-\frac8t\log(t-1)\log(t+1)+\Bigl(4-\frac4{t+1}-\frac8t\Bigr)\log(t-1)
 +\Bigl(-4+\frac4{t-1}\Bigr)\log(t+1)-4+\frac4{t-1}.
 \label{eq:dualB:L5}
\end{equation}
The term $-4\log(t-1)/(t+1)$ is essential; omitting it changes the finite part. The upper region $w=1-et$ gives
\begin{equation}
 U_5(t)=\Bigl(4-\frac{12}{t}-\frac8{t^2}\Bigr)\log(t-1)-6+\frac{10}{t-1}-\frac8t.
 \label{eq:dualB:U5}
\end{equation}
These are the coefficients of $e^0\log^0e$ after the Jacobians of the endpoint regions are included.

The finite-part table evaluates the bulk integral group by group,
\begin{center}\small
\begin{tabular}{@{}llllll@{}}
\toprule
group of $B_5$ & $l^2$ & $l\bar l$ & $l$ & $\bar l$ & rational\\
\midrule
$\FP\int_0^1$ & $24\zeta(3)+8\zeta(2)-16$ & $24\zeta(3)+16\zeta(2)-16$ & $-2\zeta(2)+2\log2+3$ & $4\zeta(2)+3$
 & $7-2\log2$\\
\bottomrule
\end{tabular}
\end{center}
so that
\begin{equation}
 \FP\!\int_0^1B_5(w)\dd w=48\zeta(3)+26\zeta(2)-19.
 \label{eq:dualB:Sbulk}
\end{equation}
For the upper region a primitive is
\begin{equation}
 (4t-2)\log(t-1)-10t+4-6\log^2t-12\Li_2(1/t)+\frac8t\log(t-1);
\end{equation}
its finite constant at infinity is zero and its value at $t=2$ is $-16-6\zeta(2)$. Thus
\begin{equation}
 \FP\!\int_2^\infty U_5(t)\dd t=16+6\zeta(2).
 \label{eq:dualB:Supper}
\end{equation}
The lower region contains the only nontrivial weight-three endpoint integral, $-8J_L$, from its first term. For the
remaining terms one primitive is
\begin{equation}
\begin{aligned}
 K_L(t)={}&4\bigl[(t-1)\log(t-1)-(t+1)\log(t+1)+2\bigr]-2\log^2(t+1)-4\Li_2\Bigl(\frac2{t+1}\Bigr)\\
 &+4\log2\log(t-1)-4\Li_2\Bigl(-\frac{t-1}{2}\Bigr)-4\log^2t-8\Li_2(1/t)-4t+4\log(t-1),
\end{aligned}
\end{equation}
whose finite part at infinity is $4\zeta(2)+2\log^22$. Its value at $t=2$, followed by Euler reflection and the
inversion formula of the dilogarithm at negative argument, gives
\begin{equation}
 \FP\!\int_2^\infty L_5(t)\dd t=-8J_L+12\zeta(2)+12\log3+4\log^22+8\Li_2(-\tfrac12).
 \label{eq:dualB:Slower}
\end{equation}
Combining Eqs.~\eqref{eq:dualB:Sbulk} to~\eqref{eq:dualB:Slower} with the lemma~\eqref{eq:dualB:JL} gives
\begin{equation}
 \mathcal S_5=44\zeta(2)+\frac{152}{3}\zeta(3)-3+12\log3+4\log^22+8\Li_2(-\tfrac12),
 \label{eq:dualB:Sclosed}
\end{equation}
and subtracting four times Eq.~\eqref{eq:dualB:Pclosed} gives Eq.~\eqref{eq:dualB:SminusP} without any numerical
inference. With Eq.~\eqref{eq:dualB:Rclosed}, this completes the closed-form evaluation of the three periods.

%--------------------------------------------------------------------
\subsection{Counterterms and the inputs \texorpdfstring{$b$}{b}, \texorpdfstring{$c$}{c} and \texorpdfstring{$f$}{f}}\label{app:dualB:counterterms}
The period evaluation is one part of the calculation of the exponents. This subsection records the response
functions, finite terms and counterterm insertions that attach the periods to the beta inputs.

\paragraph{Projected responses.} The two essential counterterms can be extracted by projecting the finite-interval
ordered exponential onto the odd spin-flip vertex and onto the even current operator. At the endpoint the odd
projection can equivalently be neutralized by a compensating vertex at infinity, with its free two-point
normalization divided out; the current projection is the corresponding normalized dipole response. At finite $\db$ the
calculation uses the regulated finite-interval responses and their kernel derivatives; it does not assume a charged
conformal limit of the auxiliary theory. Displaying the terms needed, the odd and even responses are
\begin{equation}
\begin{aligned}
 \Gamma_g&=g+\Gamma_{gh}\,gh+\Gamma_{gh^2}\,gh^2+\mathcal C_3\,g^3+\Gamma_{g^3h}\,g^3h+\mathcal C_5\,g^5+\cdots,\\
 \Gamma_h&=h+\sqrt2\,b_0\,g^2Q_2+\Gamma_{hg^2}\,hg^2+\sqrt2\,b_0\,g^4\mathcal K_4+\cdots,
\end{aligned}
\label{eq:dualB:responsegen}
\end{equation}
with the normalized responses of Eq.~\eqref{eq:dualA:normalization}, not raw moments. Contractions involving the
current coupling are generated by differentiating the exact current-resummed charges, keeping the endpoint factors of
Eq.~\eqref{eq:dualA:endpoint}. At $b_0=1$, $\db=0$ the mixed coefficients needed are
\begin{equation}
\begin{aligned}
 \Gamma_{gh}&=-2\sqrt2\,(\ell+1),\\
 \Gamma_{hg^2}&=4\ell^2+10\ell+10-4\zeta(2),\\
 \Gamma_{gh^2}&=4\ell^2+10\ell+12-2\zeta(2),\\
 \frac{\Gamma_{g^3h}}{\sqrt2}&=-\frac{20}{3}\ell^3-28\ell^2+\bigl[-48+24\zeta(2)\bigr]\ell+\Gamma_{40},\qquad
 \Gamma_{40}=-62+36\zeta(2)+30\zeta(3)+16\log2.
\end{aligned}
\label{eq:dualB:mixed}
\end{equation}
These finite constants belong to the kink-only prescription, not to one that regulates current--kink separations
independently. For example, $\Gamma_{40}$ multiplies the counterterm of $h$ at second order in the renormalized fugacity $g_R$ and contributes to the
fifth-order single logarithm. (The compendium writes $F_G$, $F_H$, $F_{gh}$, \dots\ for these responses and
coefficients.)

\paragraph{Lower counterterms.} The two-kink responses of Eq.~\eqref{eq:dualA:Z2Q2} give, at $h_R=0$,
\begin{equation}
 h_{\rm bare}=\sqrt2\,g_R^2\ell+O(g_R^4),\qquad g_{\rm bare}=g_R+g_R^3(2\ell^2-\ell)+O(g_R^5),
 \label{eq:dualB:lowerct}
\end{equation}
where $g_R$ and $h_R$ are the renormalized couplings. The odd subtraction can be checked directly:
$\mathcal C_3+\Gamma_{gh}\,(\sqrt2\ell)+(2\ell^2-\ell)$ has no logarithm. Charge differentiation of the one-pair
integral generates only higher powers of $\ell$, not a new charge-dependent single logarithm, and its finite-$\db$
derivative likewise gives the cubic logarithm of Appendix~\ref{app:dualB:threeins}. This is the one-pair origin of the
absent terms stated below Eq.~\eqref{eq:dualB:beta2}.

\paragraph{Fourth and fifth insertions.} At $h_R=0$ and $\db=0$ the counterterms are
\begin{equation}
\begin{aligned}
 h_{\rm bare}={}&\sqrt2\,g_R^2\ell+\sqrt2\,g_R^4\Bigl(\frac43\ell^3-2\ell^2-4\ell\Bigr)+\cdots,\\
 g_{\rm bare}={}&g_R+g_R^3(2\ell^2-\ell)
 +g_R^5\Bigl[\frac{10}{3}\ell^4-\frac{22}{3}\ell^3-\Bigl(\frac{13}{2}+8\zeta(2)\Bigr)\ell^2-b\,\ell\Bigr]+\cdots.
\end{aligned}
\label{eq:dualB:g5ct}
\end{equation}
To make the subtraction reproducible, set
\begin{equation}
 h^{\rm ct}_2=\sqrt2\,\ell,\qquad g^{\rm ct}_3=2\ell^2-\ell,\qquad
 h^{\rm ct}_4=\sqrt2\Bigl(\frac43\ell^3-2\ell^2-4\ell\Bigr).
\end{equation}
Then the coefficient of $g_R^5$ in the bare fugacity is
\begin{equation}
 g^{\rm ct}_5=-\Bigl[\mathcal C_5+\Gamma_{gh}\bigl(h^{\rm ct}_4+g^{\rm ct}_3h^{\rm ct}_2\bigr)+3\,\mathcal C_3\,g^{\rm ct}_3
 +\Gamma_{gh^2}\bigl(h^{\rm ct}_2\bigr)^2+\Gamma_{g^3h}\,h^{\rm ct}_2\Bigr]_{\rm logs},
 \label{eq:dualB:g5assembly}
\end{equation}
where $[\,\cdot\,]_{\rm logs}$ keeps the terms of order $\varepsilon^0$ with a positive power of $\ell$. All quantities
on the right are specified in Eqs.~\eqref{eq:dualB:C3},~\eqref{eq:dualB:C5} and~\eqref{eq:dualB:mixed}. The
coefficient of the single logarithm gives
\begin{equation}
 b=-21+38\zeta(2)+60\zeta(3)+32\log2-\mathcal S_5.
 \label{eq:dualB:bperiod}
\end{equation}
The fourth-insertion charge counterterm gives $c=-16$ (the value $h_4=4$ of Appendix~\ref{app:dualA:recurrence}).

\paragraph{The charge derivative that enters the trace.} Write the fourth-order auxiliary charge counterterm as
$h^{\rm ct}_4=\sqrt2\,b_0\,c^{\rm ct}_4(v_0)$ with $v_0=b_0^2$. Its derivative at $v_0=1$ is
\begin{equation}
 c^{\rm ct\,\prime}_4=-\frac83\ell^4+\frac{16}{3}\ell^3+\bigl(8+4\zeta(2)\bigr)\ell^2
 +\bigl[4\mathcal P_4+\mathcal R_4-8+4\zeta(2)+8\zeta(3)\bigr]\ell.
 \label{eq:dualB:C4prime}
\end{equation}
Consequently
\begin{equation}
 f=16\mathcal P_4+4\mathcal R_4-32+16\zeta(2)+32\zeta(3).
 \label{eq:dualB:fperiod}
\end{equation}
The lower endpoint integral needed in this derivative is
\begin{equation}
\begin{gathered}
 J_{\rm end}(v)=2\int_0^1r^{1-2v}\bigl[(1-r)\log(1-r)+r\log r\bigr]\dd r,\\
 J_{\rm end}(1)=-2\zeta(2),\qquad J_{\rm end}'(1)=-4\bigl(\zeta(2)+\zeta(3)\bigr),
\end{gathered}
\end{equation}
and it gives
\begin{equation}
 \partial_v\Gamma_{hg^2}\big|_{v=1}=-\frac83\ell^3+2\ell^2+24\ell+44-12\zeta(2)-8\zeta(3).
 \label{eq:dualB:Hmixedprime}
\end{equation}

Equation~\eqref{eq:dualB:C4prime} can be reconstructed algebraically. Let $Q_2\to-\ell-1$ and
$\partial_vQ_2\to\ell^2-2$ denote the $\varepsilon^0$ parts of the one-pair current moment and of its charge
derivative, and $\partial_vg^{\rm ct}_3=-\frac43\ell^3+3\ell^2+4\zeta(2)\ell$. Differentiating
Eq.~\eqref{eq:dualB:mellin} at $k_F=4-4v=0$ gives
\begin{equation}
 \partial_v\mathcal M_{Q_4}=\frac{\partial_v\mathcal A_{Q_4}}{q^2(q+1)}+\frac{4(2q+1)\,\mathcal A_{Q_4}}{q^3(q+1)^2},
 \label{eq:dualB:MQprime}
\end{equation}
where $\mathcal M_{Q_4}$ is the Mellin transform of $Q_4$. Its poles give the $\varepsilon^0$ part of $\partial_vQ_4$
through Eq.~\eqref{eq:dualB:poles}. Subtracting the derivative of $Q_2Z_2$ forms $\partial_v\mathcal K_4$; with
$Z_2=\int_\varepsilon^1(1-r)r^{-2v}\dd r$ and $Q_2=\int_\varepsilon^1(1-r)r^{1-2v}\dd r$ at $\db=0$, the
$\varepsilon^0$ part of $\partial_v(Q_2Z_2)$ at $v=1$ is $2\ell^3-8\ell$. Then
\begin{equation}
 c^{\rm ct\,\prime}_4=-\partial_v\mathcal K_4\big|_{\rm logs}-2(\partial_vQ_2)\,g^{\rm ct}_3-2Q_2\,\partial_vg^{\rm ct}_3
 -(\partial_v\Gamma_{hg^2})\,\ell+\Gamma_{hg^2}\,\ell^2.
 \label{eq:dualB:C4prassembly}
\end{equation}
Equations~\eqref{eq:dualB:laurent} and~\eqref{eq:dualB:MQprime} fix every term of this expression; the compendium
checks it symbolically in its code.

%--------------------------------------------------------------------
\subsection{Charge and finite-\texorpdfstring{$\db$}{delta} derivatives at three insertions:
\texorpdfstring{$a$}{a} and \texorpdfstring{$p$}{p}}\label{app:dualB:threeins}
\paragraph{An incomplete-beta representation.} For three kinks with gaps $r$ and $t$, take the kernel
$r^{-2p_{12}}t^{-2p_{23}}(r+t)^{2p_{13}}$ with general powers (the subscripts name the kink pairs). Its Mellin
integral is
\begin{equation}
 \mathcal M_3(q;p_{12},p_{23},p_{13})
 =\frac{\mathrm B_{1/2}(q+1-2p_{12},\,1-2p_{23})+\mathrm B_{1/2}(q+1-2p_{23},\,1-2p_{12})}{q(q+k_3)(q+k_3+1)},
 \label{eq:dualB:M3}
\end{equation}
with $k_3=2-2p_{12}-2p_{23}+2p_{13}$ and the incomplete beta function
$\mathrm B_x(a,b)=\int_0^xy^{a-1}(1-y)^{b-1}\dd y$, continued meromorphically from its domain of convergence. The
formula follows by writing $r=Ry$ and $t=R(1-y)$, splitting at $y=\frac12$ and doing the radial integral. The kernel
$W_3$ at $v=1$, $\db=0$ corresponds to $p_{12}=p_{23}=p_{13}=1$.

\paragraph{The charge derivative and $a$.} The charge derivative varies all three powers together. The normalized
three-kink response then has
\begin{equation}
 \partial_v\mathcal C_3\big|_{v=1}=-\frac83\ell^3-3\ell^2+\bigl(4-4\zeta(2)\bigr)\ell
 +18-14\zeta(2)-15\zeta(3)-8\log2+\cdots.
 \label{eq:dualB:C3v}
\end{equation}
Combining it with the lower mixed response and counterterms gives
\begin{equation}
 a=-4\zeta(2).
\end{equation}
Equation~\eqref{eq:dualB:M3} is a compact exact representation from which the logarithmic and finite terms of
Eq.~\eqref{eq:dualB:C3v} can be regenerated.

\paragraph{The finite-tail deformation of the kernel and $p$.} For separated charges the Gaussian weight expands as
\begin{equation}
 \exp\Bigl[\frac{2v}{\db}\sum_{i<j}q_iq_j\bigl(r_{ij}^{\db}-1\bigr)\Bigr]
 =W^v\Bigl[1+v\,\db\sum_{i<j}q_iq_j\log^2r_{ij}+O(\db^2)\Bigr],
\end{equation}
where $r_{ij}$ is the separation of kinks $i$ and $j$ and $W^v$ the weight at $\db=0$. Thus the three-kink Mellin
derivative is
\begin{equation}
 \partial_{\db}\mathcal M_3\big|_{\db=0}=\frac14\bigl(\partial_{p_{13}}^2-\partial_{p_{12}}^2-\partial_{p_{23}}^2\bigr)
 \mathcal M_3\big|_{p_{12}=p_{23}=p_{13}=1}.
 \label{eq:dualB:M3delta}
\end{equation}
The positive sign belongs to the like-charge outer pair, the negative signs to the two neighbouring opposite-charge
pairs. Expanding Eq.~\eqref{eq:dualB:M3delta} and normalizing by the two-kink moment gives
\begin{equation}
\begin{gathered}
 \partial_{\db}\mathcal C_3\big|_{\db=0}=-\ell^4-\ell^3-2\zeta(2)\ell^2+(p-4)\ell+O(1),\qquad
 \partial_{\db}\Gamma_{gh}\big|_{\db=0}=\sqrt2\,(2-\ell^2),\\
 p=8-2\zeta(2)-7\zeta(3)-8\log2.
\end{gathered}
\label{eq:dualB:panswer}
\end{equation}
The $\mathrm B_{1/2}$ formula, the explicit differential operator and the lower moments fix this single logarithm
independently of any fit to thermal exponents. The compendium's three-vertex residue script evaluates these
derivatives as a separate check.

The finite-$\db$ current counterterm at order $g_R^2$ is
\begin{equation}
 h_{\rm bare}\big|_{\db g_R^2}=-\frac{\sqrt2}{3}\,\db\,g_R^2\ell^3.
\end{equation}
Substitution into the odd response gives
\begin{equation}
 g_{\rm bare}\big|_{\db g_R^3}=\db\,g_R^3\Bigl[-\frac13\ell^4+\frac53\ell^3+2\zeta(2)\ell^2-p\,\ell\Bigr].
 \label{eq:dualB:auxDeltaCt}
\end{equation}
This auxiliary fugacity is not yet the essential wall fugacity. Multiplication by the factor of
Eq.~\eqref{eq:dualA:fugmatch} changes it to
\begin{equation}
 Y_{\rm bare}\big|_{\db g_R^3}=\db\,g_R^3\Bigl[-\frac43\ell^4+\frac{13}{6}\ell^3+2\zeta(2)\ell^2-p\,\ell\Bigr].
 \label{eq:dualB:physicalDeltaCt}
\end{equation}
The difference in the terms $\ell^4$ and $\ell^3$ is the finite-$\db$ matching, and it is indispensable.

%--------------------------------------------------------------------
\subsection{Analytic three-kink and endpoint inputs}\label{app:dualB:endpoints}
This subsection gives the analytic inputs of Appendix~\ref{app:dualB:D3}. These reductions of lower inputs do not
assume the higher-kink identity~\eqref{eq:dualB:calR} that the certificate proves. Here $\partial_\eta A_3$ denotes the integrated first source derivative
of the raw odd three-kink moment, not the derivative of an already normalized CFT correlator; the normalization and
derivative factors of the subtraction recurrence must still be applied. At $\db=0$ the endpoint factors of
Eq.~\eqref{eq:dualA:endpoint} are
\begin{equation}
 U_n^{\rm odd}(\eta)=\prod_i\bigl[X_i(1-X_i)\bigr]^{\eta q_i},\qquad
 U_n^{\rm even}(\eta)=\prod_i\bigl[X_i/(1-X_i)\bigr]^{\eta q_i},
\end{equation}
and
\begin{equation}
 \partial_\eta A_3=\partial_\eta\int W_3^vU_3^{\rm odd}\Big|_{\eta=0},\qquad
 \partial_\eta Q_4=\partial_\eta\int\Xi_4W_4^vU_4^{\rm even}\Big|_{\eta=0},
\end{equation}
with the integrals over $\Omega_3(\varepsilon)$ and $\Omega_4(\varepsilon)$; the endpoints of the interval are not cut.
The charge-curvature input has the convergent representation
\begin{equation}
 A=-4\int_0^{1/2}\frac{\log^2(1-t)}{t^2(1-t)^2}\dd t=-8-4\zeta(2)-2\zeta(3)+16\log2,
\end{equation}
obtained by partial fractions and integration by parts.

\paragraph{The mixed three-kink coefficient $E$.} The representation~\eqref{eq:dualB:Erep} reads
\begin{equation}
 E=p-4\,\mathsf I_3+8\,\mathsf I_{12},\qquad
 \mathsf I_3=\int_0^{1/2}\frac{\log^3(1-y)}{y^2(1-y)^2}\dd y,\qquad
 \mathsf I_{12}=\int_0^{1/2}\frac{\log y\log^2(1-y)}{y^2(1-y)^2}\dd y,
 \label{eq:dualB:Estart}
\end{equation}
and it can be integrated exactly. The result is the expression of $E$ in Eq.~\eqref{eq:dualB:AE}.

\emph{Proof.} Put $t=y/(1-y)$. Then $0<t<1$, $\log(1-y)=-\log(1+t)$, and the rational measure becomes
$(1+2/t+1/t^2)\dd t$. Define
\begin{equation}
 \mathsf B_3=\int_0^1\frac{\log^3(1+t)}{t}\dd t,\qquad \mathsf B_{12}=\int_0^1\frac{\log t\log^2(1+t)}{t}\dd t.
\end{equation}
The elementary reductions by integration by parts are
\begin{equation}
\begin{aligned}
 \int_0^1\Bigl(1+\frac2t+\frac1{t^2}\Bigr)\log^3(1+t)\dd t&=-6\log^22+12\log2-6+\frac34\zeta(3)+2\mathsf B_3,\\
 \int_0^1\Bigl(1+\frac2t+\frac1{t^2}\Bigr)\log t\log^2(1+t)\dd t&=2\mathsf B_{12}+2\zeta(2)-4\log^22+8\log2-6
 -\frac32\zeta(3).
\end{aligned}
\end{equation}
Thus
\begin{equation}
 E=-16+14\zeta(2)-22\zeta(3)-8\log^22+8\log2+16\mathsf B_{12}-8\mathsf B_3.
 \label{eq:dualB:Eintermediate}
\end{equation}
To evaluate $\mathsf B_3$, set $u=(1+t)^{-1}$ and expand $(1-u)^{-1}$. A primitive needed for the logarithmic part is
\begin{equation}
 \int\frac{\log^3u}{1-u}\dd u=-\log^3u\log(1-u)-3\log^2u\Li_2(u)+6\log u\Li_3(u)-6\Li_4(u).
\end{equation}
Evaluation between $u=\frac12$ and $1$, with Eq.~\eqref{eq:dualB:halfvalues}, gives
\begin{equation}
 \mathsf B_3=6\zeta(4)-6\Li_4(\tfrac12)-\frac{21}{4}\zeta(3)\log2+\frac32\zeta(2)\log^22-\frac14\log^42.
 \label{eq:dualB:B3}
\end{equation}
For $\mathsf B_{12}$ introduce
\begin{equation}
 \mathsf K(\alpha,\beta)=\int_0^1t^{\alpha-1}(1+t)^\beta\dd t.
\end{equation}
Splitting the complete integral over $(0,\infty)$ at one and inverting $t$ in the second half proves, first in a
common domain of convergence and then by meromorphic continuation,
\begin{equation}
 \mathsf K(\alpha,\beta)+\mathsf K(-\alpha-\beta,\beta)=\mathrm B(\alpha,-\alpha-\beta),
 \label{eq:dualB:betareflection}
\end{equation}
with $\mathrm B$ Euler's beta function. The cubic regular part of the right side, after its pole, is
\begin{equation}
 \frac74\zeta(4)\bigl(\alpha^2\beta+\alpha\beta^2\bigr)+\zeta(4)\beta^3,
\end{equation}
which follows from the expansion
\begin{equation}
 \mathrm B(\alpha,-\alpha-\beta)=\frac{\beta}{\alpha(\alpha+\beta)}\exp\sum_{n\ge2}\frac{(-1)^n\zeta(n)}{n}
 \bigl[\alpha^n+(-\alpha-\beta)^n-(-\beta)^n\bigr];
\end{equation}
the term of degree four in the exponential, including its quadratic $\zeta(2)^2$ contribution, supplies the cubic
regular part, with $\zeta(2)^2=\frac52\zeta(4)$. The coefficient of $\beta^3$ on the left of
Eq.~\eqref{eq:dualB:betareflection} is
\begin{equation}
 \frac{\mathsf B_3}{3}-\frac{\mathsf B_{12}}{2}+\frac12\int_0^1\frac{\log^2t\log(1+t)}{t}\dd t
 =\frac{\mathsf B_3}{3}-\frac{\mathsf B_{12}}{2}+\frac78\zeta(4).
\end{equation}
Therefore
\begin{equation}
 \mathsf B_{12}=\frac23\mathsf B_3-\frac14\zeta(4).
 \label{eq:dualB:B12}
\end{equation}
Substituting Eqs.~\eqref{eq:dualB:B3} and~\eqref{eq:dualB:B12} into Eq.~\eqref{eq:dualB:Eintermediate} gives the
closed form of $E$.\hfill$\square$

An independent quadrature of Eq.~\eqref{eq:dualB:Estart}, which does not use the functional identity, gives
\begin{equation}
 E=-21.663876674184610081020147908761053559049257064774\ldots,
\end{equation}
in agreement with the closed form to more than eighty decimal places. This precision refers to a convergent
one-dimensional check, not to the full quadrature of $D_3$.

\paragraph{The three-kink endpoint derivative.} For the uncut interval endpoints and the consecutive-kink hard core,
the finite part of the charge-differentiated source derivative of the odd three-kink moment is
\begin{equation}
 \partial_v\partial_\eta A_{3,0}=-108+60\zeta(2)+22\zeta(3)-34\zeta(4)+16\log2.
 \label{eq:dualB:ep3}
\end{equation}
Its reduction is explicit. Put
\begin{equation}
 j(r)=(1-r)\log(1-r)-r\log r.
\end{equation}
For the two gaps $r$ and $t$ with span $R=r+t$, the integral over the first position gives the odd endpoint factor
\begin{equation}
 \Phi_{\rm odd}(r,t)=2j(R)-j(r)-j(t)-2(1-R).
\end{equation}
Consequently
\begin{equation}
 \partial_v\partial_\eta A_3(\varepsilon)=2\int_{2\varepsilon}^1\bigl[j(R)-(1-R)\bigr]I_{R,v}(R,\varepsilon)\dd R
 -2\int_\varepsilon^{1-\varepsilon}j(r)\,I_{r,v}(r,\varepsilon)\dd r.
 \label{eq:dualB:ep3reduce}
\end{equation}
Both inner integrals, $I_{R,v}$ at fixed span and $I_{r,v}$ at fixed gap, have elementary dilogarithmic primitives.
Define
\begin{equation}
\begin{aligned}
 J_0(e)&=\frac2e-\frac2{1-e}-4\log e+4\log(1-e),\\
 \Phi_{\log}(y)&=-\frac{\log y+1}{y}+\frac{\log y}{1-y}-\log y+\log(1-y)+\log^2y-2\log y\log(1-y)-2\Li_2(y),\\
 J_1(e)&=-4\bigl[\Phi_{\log}(1-e)-\Phi_{\log}(e)\bigr].
\end{aligned}
\end{equation}
Then
\begin{equation}
 I_{R,v}=R^{-1}\bigl[-2\log R\,J_0(\varepsilon/R)+J_1(\varepsilon/R)\bigr],
\end{equation}
and indeed $\Phi_{\log}'(y)=\log y/[y^2(1-y)^2]$. For the second integral let
\begin{equation}
\begin{aligned}
 \Psi_0(t)&=t+2\log t-\frac1t,\\
 \Psi_1(t)&=t\bigl[\log(1+t)-\log t\bigr]-2\Li_2(-t)-\log^2t+\frac{\log t+1-\log(1+t)}{t}+\log t.
\end{aligned}
\end{equation}
The derivative of $\Psi_1$ is $(1+2/t+1/t^2)[\log(1+t)-\log t]$, and
\begin{equation}
 I_{r,v}=\frac2r\Bigl[\Psi_1\Bigl(\frac{1-r}{r}\Bigr)-\Psi_1\Bigl(\frac\varepsilon r\Bigr)
 -\log r\Bigl\{\Psi_0\Bigl(\frac{1-r}{r}\Bigr)-\Psi_0\Bigl(\frac\varepsilon r\Bigr)\Bigr\}\Bigr].
\end{equation}

The finite part of Eq.~\eqref{eq:dualB:ep3reduce} consists of two bulk terms and the lower-span region
$R=\varepsilon t$; the second integral has no additional finite boundary term. The pieces are
\begin{center}\small
\begin{tabular}{@{}ll@{}}
\toprule
piece & finite contribution\\
\midrule
first bulk & $-144+64\zeta(2)+32\zeta(3)-24\zeta(4)$\\
second bulk & $60-16\zeta(2)-40\zeta(3)-10\zeta(4)$\\
lower-span boundary & $-24+16\log2+12\zeta(2)+30\zeta(3)$\\
\bottomrule
\end{tabular}
\end{center}
and their sum is Eq.~\eqref{eq:dualB:ep3}. The first bulk is
\begin{equation}
 2\,\FP\!\int_0^1\frac{j(r)-(1-r)}{r}\bigl[-4\log^2r+8\log r+8+8\zeta(2)\bigr]\dd r.
\end{equation}
For the second bulk set $l=\log r$, $\bar l=\log(1-r)$, $\tau=(1-r)/r$ and
\begin{equation}
 \Upsilon=-2\tau(\bar l+l)+4\Li_2(1-r)-2\bar l^2+2l^2+\frac{2(\bar l+l+1)}{\tau}+2\bar l+2.
\end{equation}
The integrand is $-2(\tau\bar l-l)\Upsilon$; its part with logarithms only integrates to
$12+8\zeta(2)-24\zeta(3)-16\zeta(4)$, and its dilogarithm part to $48-24\zeta(2)-16\zeta(3)+6\zeta(4)$. The
lower-span contribution is
\begin{equation}
 2\,\FP\!\int_0^{1/2}\frac{-2\log e\,J_0(e)-J_1(e)}{e}\dd e.
\end{equation}
Euler's identity~\eqref{eq:dualB:Euler} reduces its integrand, with $l=\log e$ and $\bar l=\log(1-e)$, to
\begin{equation}
 \frac{8(\bar l+1)}{e^2}-\frac{8(\bar l+1)}{e(1-e)}+\frac{16(l-\bar l)}{e}+\frac{8(\bar l^2+l^2)}{e}
 -\frac{16\zeta(2)}{e}+\frac{32\Li_2(e)}{e},
\end{equation}
whose six finite integrals over $(0,\frac12)$ are, in order,
\begin{equation}
 -24+16\log2,\quad 4\zeta(2),\quad 8\zeta(2),\quad 2\zeta(3)-\tfrac{16}{3}\log^32,\quad 16\zeta(2)\log2,\quad
 28\zeta(3)-16\zeta(2)\log2+\tfrac{16}{3}\log^32.
\end{equation}
This gives the boundary entry of the table directly. The compendium checks all primitives and rational-logarithm
moment operations symbolically (\file{code/dual_d1/code/archived/LRI_1d_complete/stages/D3_exact/reproducibility/}).

\paragraph{The four-kink endpoint finite part.} The current endpoint response has the expansion
\begin{equation}
\begin{aligned}
 \partial_\eta Q_4(\varepsilon)&=-\frac{2\zeta(2)}{\varepsilon}+\ell^3-\bigl(1+8\zeta(2)\bigr)\ell^2
 -\bigl(2+32\zeta(2)+28\zeta(3)\bigr)\ell\\
 &\quad+6-12\zeta(2)-54\zeta(3)-25\zeta(4)+o(1),
\end{aligned}
\label{eq:dualB:ep4asympt}
\end{equation}
so that
\begin{equation}
 \partial_\eta Q_{4,0}=6-12\zeta(2)-54\zeta(3)-25\zeta(4).
 \label{eq:dualB:ep4}
\end{equation}
Let $\mathsf f(t)=t\log t+(1-t)\log(1-t)$, with continuous values at the endpoints. For the gaps $r,t,w$ and span
$R=r+t+w$, the integral of the endpoint logarithm over the first position is
\begin{equation}
 \Phi_4=2\mathsf f(R)+\mathsf f(r)+\mathsf f(w)-\mathsf f(R-r)-\mathsf f(R-w).
\end{equation}
The rational four-kink kernel is $W_4$ of Eq.~\eqref{eq:dualA:W4}. At fixed span put $y=r/R$ and $e=\varepsilon/R$;
the integration over the other angular gap is elementary. With
\begin{equation}
 \mathcal A(y)=\frac{(y+1)(y^2+1)}{y^2(1-y)}=-1+\frac4{1-y}+\frac2y+\frac1{y^2},\qquad
 \mathsf F_R(y)=\mathsf f(R)+\mathsf f(Ry)-\mathsf f\bigl(R(1-y)\bigr),
\end{equation}
and the function $\mathcal M_{\rm out}(y,e)=(1+1/y)[e^{-1}-(1-y-e)^{-1}]+\mathcal A(y)\log[(1-y-e)/e]$ of
Eq.~\eqref{eq:dualB:Mout}, one obtains the exact two-dimensional formula
\begin{equation}
 \partial_\eta Q_4(\varepsilon)=2\int_{3\varepsilon}^1\frac{\dd R}{R}\int_{\varepsilon/R}^{1-2\varepsilon/R}
 \mathcal M_{\rm out}(y,\varepsilon/R)\,\mathsf F_R(y)\dd y,
 \label{eq:dualB:ep4kernel}
\end{equation}
which is also integrated directly at finite cutoff as a separate check.

\emph{Angular finite part and the radial logarithms.} Let $\varkappa(R)=2\log R-\log(1-R)-1$ and
$\mathsf F_1=2\mathsf f(R)$. Near $y=0$, $\mathsf F_R(y)=Ry[\log y+\varkappa(R)]+O(y^2)$; near $y=1$,
$\mathsf F_R(y)\to\mathsf F_1$. Define
\begin{equation}
 \mathcal H(e,R)=\int_e^{1-2e}\mathcal M_{\rm out}(y,e)\,\mathsf F_R(y)\dd y,\qquad
 J_{\FP}(R)=\FP\!\int_0^1\mathcal A(y)\,\mathsf F_R(y)\dd y.
\end{equation}
Differentiating $\mathcal H$ with respect to $\log e$, including the moving boundaries, gives its part of order $e^0$,
\begin{equation}
 \partial_{\log e}\mathcal H\big|_{e^0}=\frac{3R}{2}\log^2e+\bigl(2R\varkappa+4\mathsf F_1\bigr)\log e-J_{\FP}-R
 +2\mathsf F_1.
\end{equation}
Hence
\begin{equation}
 \mathcal H\big|_{e^0}=\frac R2\log^3e+\bigl(R\varkappa+2\mathsf F_1\bigr)\log^2e+\bigl(-J_{\FP}-R+2\mathsf F_1\bigr)\log e
 +C_0(R),
\end{equation}
where the constant is
\begin{equation}
 C_0(R)=\FP\!\int_0^1\bigl[-\mathcal B(y)+\mathcal A(y)\log(1-y)\bigr]\mathsf F_R(y)\dd y+R(1-\varkappa)
 -\bigl(4+2\zeta(2)\bigr)\mathsf F_1,
\end{equation}
with $\mathcal B(y)=(1+1/y)/(1-y)$. These expressions keep both angular endpoints.

It is essential to insert $\log e=\ell-\log R$ in the \emph{whole} polynomial; integrating only $C_0(R)$ would miss
finite terms that come from powers of $\log e$. The radial moments are
\begin{equation}
\begin{aligned}
 \mathsf U(y)&=\int_0^1\mathsf F_R(y)\frac{\dd R}{R}=-2\Li_2(y)-\log y\log(1-y),\\
 \mathsf V(y)&=\int_0^1\log R\,\mathsf F_R(y)\frac{\dd R}{R}=\zeta(3)-\mathsf U(y)+\Li_3(y)-\Li_3(1-y),
\end{aligned}
\end{equation}
which follow, for example, from
\begin{equation}
 \int_0^1\mathsf f(cR)\frac{\dd R}{R}=\mathsf f(c)-\Li_2(c),\qquad
 \int_0^1\log R\,\mathsf f(cR)\frac{\dd R}{R}=-\mathsf f(c)+\Li_2(c)+\Li_3(c).
\end{equation}
The remaining bulk integral is
\begin{equation}
 2\,\FP\!\int_0^1\Bigl[-\mathcal B\,\mathsf U+\mathcal A\bigl\{\log(1-y)\,\mathsf U+\mathsf V\bigr\}\Bigr]\dd y
 =12-26\zeta(3)-29\zeta(4).
 \label{eq:dualB:ep4bulk}
\end{equation}
One integration by parts removes the trilogarithms. A primitive of $\mathcal A$ is
\begin{equation}
 \Pi_{\mathcal A}(y)=-y+2\log y-4\log(1-y)-\frac1y,
\end{equation}
and the finite endpoint term for $\mathcal D(y)=\zeta(3)+\Li_3(y)-\Li_3(1-y)$ is
\begin{equation}
 \bigl[\Pi_{\mathcal A}\mathcal D\bigr]_0^1=1+\zeta(2)-4\zeta(3).
\end{equation}
The resulting rational-logarithm and dilogarithm moments are elementary derivatives of beta integrals.

The other endpoint constants contribute $4+16\zeta(2)+20\zeta(4)$. The shift from the radial logarithms is
\begin{equation}
 \Delta_{\rm rad}=2\int_0^1\Bigl[-\frac{2\mathsf F_1}{R}\log R+\Bigl(\varkappa+\frac{2\mathsf F_1}{R}\Bigr)\log^2R
 -\frac12\log^3R\Bigr]\dd R=-10-28\zeta(2)-28\zeta(3)-16\zeta(4),
\end{equation}
which follows from
\begin{equation}
\begin{gathered}
 \int_0^1\mathsf f(R)\frac{\log R}{R}\dd R=\zeta(2)+\zeta(3),\qquad
 \int_0^1\mathsf f(R)\frac{\log^2R}{R}\dd R=-2\bigl(\zeta(2)+\zeta(3)+\zeta(4)\bigr),\\
 \int_0^1\varkappa(R)\log^2R\dd R=-8-2\zeta(2)-2\zeta(3).
\end{gathered}
\end{equation}
Adding Eq.~\eqref{eq:dualB:ep4bulk}, the other endpoint constants and $\Delta_{\rm rad}$ gives
Eq.~\eqref{eq:dualB:ep4}, and collecting the remaining logarithms gives Eq.~\eqref{eq:dualB:ep4asympt}.

%====================================================================
\section{The dual of the chain: exact certificates, the seventh and eighth orders}\label{app:dualC}
%====================================================================
This appendix gives the derivations of the compendium (written by OpenAI GPT-6 astra; Sec.~\ref{sec:sources}) behind
Secs.~\ref{sec:closed} and~\ref{sec:C8}, edited to the notation of this note. As in Appendices~\ref{app:dualA}
and~\ref{app:dualB}, everything here holds within the matched two-coupling prescription of Sec.~\ref{sec:dual}: the
certificates and the rational part of $\CSR_8$ prove reductions of the archived expressions of that prescription, and
the numerical values are evaluations within it (Sec.~\ref{sec:dual}, ``What `proved within the prescription'
covers''). Appendix~\ref{app:dualC:certificates}
sets up the exact integration and the word algebra in which the certificates of $\CSR_5$ to $\CSR_7$ are written, and
Appendix~\ref{app:dualC:C7} completes the seventh order. Appendices~\ref{app:dualC:C8setup} to~\ref{app:dualC:numerics}
treat the eighth order: its trace reduction, the exact rational part, the four computational stages of the full
value, the assembly with the zeta-value candidate, and the numerical evaluation of periods.
Appendices~\ref{app:dualC:conjectures} and~\ref{app:dualC:endpoints} record the general formulas that were tested and
the relation to the other end of the long-range window; Appendices~\ref{app:dualC:shadow} and~\ref{app:dualC:logs} give two
results that use only published data of the dual; Appendix~\ref{app:dualC:Bold} prints the polynomial $B_{\rm old}$ of
Eq.~\eqref{eq:C8assembly}.

\paragraph{Notation.} As in Sec.~\ref{sec:dual}, $\db=1-s$, $x=\Ds=\db/2$, $\varepsilon$ is the hard core and
$\ell=\log\varepsilon$; $u=g^2$ and $v=B^2=1+z$ are the essential couplings, with $\beta_u=2uF$ and $\beta_z=(1+z)H$.
The beta inputs are Taylor coefficients of $F$ and $H$ at the marginal reference point $v=1$, $\db=0$, not at the
fixed point. Their names ($a,p,b,c,f,\dots,J_9,K_v$) are the compendium's and denote coefficients only. Several of these letters
also have unrelated uses here, distinguished by subscripts or context: $\Omega_n(\varepsilon)$ is the hard-core domain,
$W_n$ a wall weight and $W_N\mathcal Z$ a span of MZVs, $Q_{2m}$ a current moment, $R$ a sum of gaps and $R_w$ a
regularized constant, $S_j$ an insertion, $P_6$, $P_8$, $P_{w,k}$ periods and polynomials, $L$ a logarithmic
variable, and $B$, $B_j$ a boundary derivative and its face integrals (Appendices~\ref{app:dualC:C8rational}
and~\ref{app:dualC:C8packages}), not the essential charge. The input $N_\delta$ keeps the
compendium's name; the $\delta$ in it is $\db$. For a regulated response $I$ we write $I_0=[\varepsilon^0\ell^0]I$
and $I_1=[\varepsilon^0\ell^1]I$, and more generally $(I)_m=[\varepsilon^0\ell^m]I$. A subscript $v$ or $\db$ after a
comma denotes a derivative at the reference point, as in $b_{3,v\db}=\partial_v\partial_{\db}b_3$, while a final
subscript $0$ or $1$ after a comma is the projection $I_0$ or $I_1$, as in $Q_{2,0}$, $\mathcal K_{8,1}$ and
$A_{9,1}=(A_9)_1$, the coefficient written $A_9[\varepsilon^0\ell^1]$ in Sec.~\ref{sec:indep}; thus $(\partial_vA_7)_1$
is the coefficient of $\varepsilon^0\ell^1$ in $\partial_vA_7$. The symbol $\partial_\eta$ is the derivative with respect to the
endpoint source $\eta$ of Sec.~\ref{sec:dual}, item~(ii), at $\eta=0$. The other symbols of the dual are those of
Appendix~\ref{app:dualA:notation}: the ordered moments $A_{2m+1}$, $Z_{2m}$ and $Q_{2m}$ with the current weight
$\Xi_{2m}$ (Eq.~\eqref{eq:dualA:moments}), the vacuum-normalized odd and current
responses $\mathcal C_{2m+1}$ and $\mathcal K_{2m}$ (Eq.~\eqref{eq:dualA:normalization}), the charge and tail
insertions $S_j=\sum_{i<k}q_iq_k\log^j(X_k-X_i)$ (Eq.~\eqref{eq:dualA:insertions}), the Jacobian entries $j_n$ and
$\lambda_{2m}$ (Eq.~\eqref{eq:dualA:jacobian}) and the subtraction recurrence (Eq.~\eqref{eq:dualA:recurrence}).
The formal projection $\pi_0$ is defined in Appendix~\ref{app:dualC:C8rational}; a superscript $(0)$ denotes it, as in
$U^{(0)}=\pi_0U$.

%--------------------------------------------------------------------
\subsection{Exact integration, regularized periods and the certificates}\label{app:dualC:certificates}
%--------------------------------------------------------------------
The numerical tables of the compendium are not the definitions of its results. The primary exact data are rational
functions of the cutoff multiplying iterated integrals. This subsection explains how they are integrated and reduced,
including the constants that a symbol calculation does not see.

\paragraph{Keeping the hard core during exact integration.} After the integral over the overall translation, use
cumulative coordinates and shift $X_j=j\varepsilon+\tilde X_j$. The remaining simplex is
$0<\tilde X_1<\cdots<\tilde X_{n-1}<1-(n-1)\varepsilon$. At $v=1$, $\db=0$ the kernel is rational, and the charge and tail
derivatives insert polynomials in logarithms. Partial fractions, integration by parts and shuffle products reduce the
successive integrations to hyperlogarithms, defined with the outermost letter first by
\begin{equation}
 G(a_1,\ldots,a_n;\Lambda)=\int_0^\Lambda\frac{\dd t'}{t'-a_1}\,G(a_2,\ldots,a_n;t'),\qquad G(;\Lambda)=1
 \label{eq:dualC:gpl}
\end{equation}
(the compendium writes $T$ for the argument $\Lambda$; here $T$ is the stability trace). For zero letters the
tangential-base convention $G(0^n;\Lambda)=\log^n\Lambda/n!$ is used. The total differential is
\begin{equation}
 \dd G(a_1,\ldots,a_n;\Lambda)=\sum_{i=1}^nG(a_1,\ldots,\widehat{a_i},\ldots,a_n;\Lambda)\,
 \dd\log\frac{a_{i-1}-a_i}{a_{i+1}-a_i},\qquad a_0=\Lambda,\quad a_{n+1}=0,
 \label{eq:dualC:totaldiff}
\end{equation}
with coincident letters taken as limits. For affine letters it gives an exact recurrence in rational kernels.
Antiderivatives are checked by differentiation, and their constants are fixed by the specified base point;
differentiation alone does not fix a constant.

For a word $w$ of nonnegative letters write $-w$ for the word with every letter negated. At infinity
\begin{equation}
 G(-w;\Lambda)=\sum_{n\ge0}\Lambda^{-n}P_{w,n}(\log\Lambda),\qquad
 P_{w,0}(L)=\sum_{j=0}^{|w|}\frac{L^j}{j!}\,R_{w_{j+1}\cdots w_{|w|}},
 \label{eq:dualC:infinity}
\end{equation}
with the constants $R$ defined below. The inverse-power coefficients follow from the differential equation. A rational
prefactor can grow as a power of $\Lambda$, so enough inverse powers must be kept before a finite coefficient is taken.
In the same way, positive powers of $\varepsilon$ are kept whenever a vacuum denominator can lower their degree.

\paragraph{The rational shuffle algebra.} For a word $w=(a_1,\ldots,a_n)$ of nonnegative letters define the
regularized constant at infinity
\begin{equation}
 R_w=\operatorname{Reg}_{\Lambda\to\infty}G(-a_1,\ldots,-a_n;\Lambda),\qquad R_\varnothing=1,
 \label{eq:dualC:Rw}
\end{equation}
where $\operatorname{Reg}$ takes the coefficient of $\Lambda^0(\log\Lambda)^0$. The hyperlogarithms themselves have no
positive powers of $\Lambda$, so the constants obey the shuffle product,
\begin{equation}
 R_wR_{w'}=\sum_{w''\in w\mathbin{\amalg}w'}R_{w''},
\end{equation}
with shuffle multiplicities. The basic values are $R_0=R_1=0$ and $R_a=-\log a$ for $a>0$. A subscript $ab$ that
denotes the single integer product satisfies $R_{ab}=R_a+R_b$; it must not be confused with the two-letter word $(a,b)$.

\emph{Lemma (repeated-letter values).} For $b>0$, $a\ge0$ and $n\ge2$,
\begin{equation}
 R_{a,b^{n-1}}=\frac{(-\log b)^n}{n!}+\Li_n(1-a/b),
 \label{eq:dualC:repeated}
\end{equation}
where $b^{n-1}$ is the letter $b$ repeated $n-1$ times.

\emph{Proof.} The inner repeated-letter integral is $\log^{n-1}(1+t'/b)/(n-1)!$. Set $\vartheta=\log(1+t'/b)$ in the outer
integration. The kernel splits into a constant $1$ and $(1-a/b)/(e^\vartheta-(1-a/b))$. The first term integrates to the
logarithmic polynomial. The remaining convergent geometric-series integral is $\Li_n(1-a/b)$, extended along the
nonsingular real domain; its integrable limit at $a=0$ is $\zeta(n)$. $\square$

Hence
\begin{equation}
 \zeta(n)=-\frac{2^{n-1}}{2^{n-1}-1}\,R_{2,1^{n-1}},\qquad
 \Li_n(r)=R_{a,b^{n-1}}-R_{b^n}\quad(1-r=a/b>0).
\end{equation}
These identities, together with the prime factorization of rational logarithms, give a deterministic lift of the exact
beta expressions to rational vectors of words. For words whose last letter is positive, uniform scaling gives
\begin{equation}
 R_{\kappa w}=\sum_{j=0}^{|w|}R_{\kappa^j}\,R_{w_{j+1}\cdots w_{|w|}},\qquad \kappa>0,
\end{equation}
where $\kappa w$ multiplies every letter by $\kappa$ and $\kappa^j$ is the word of $j$ letters $\kappa$. The restriction
on the last letter prevents silently rescaling a trailing-zero tangential logarithm.

\paragraph{Confluence fixes the full constants.} \emph{Lemma (convergent confluence).} Let $e=(e_1,\ldots,e_n)$ be a
binary word with $e_1=0$ and $e_n=1$. Choose $a_i\ge0$ for a zero bit and $a_i>0$ for a one bit, and set
$\lambda_i(\Lambda)=-a_i$ if $e_i=0$ and $\lambda_i(\Lambda)=\Lambda+a_i$ if $e_i=1$. Then
\begin{equation}
 \lim_{\Lambda\to\infty}G\bigl(\lambda_1(\Lambda),\ldots,\lambda_n(\Lambda);\Lambda\bigr)=G(e_1,\ldots,e_n;1).
 \label{eq:dualC:confluence}
\end{equation}

\emph{Proof.} Rescale every integration variable by $\Lambda$. The two kinds of kernels become
$\dd\upsilon/(\upsilon+a_i/\Lambda)$ and $-\dd\upsilon/(1-\upsilon+a_i/\Lambda)$, whose absolute values are bounded by
$\dd\upsilon/\upsilon$ and $\dd\upsilon/(1-\upsilon)$. The product of these bounds is integrable on the ordered simplex
because the binary word begins with $0$ and ends with $1$. Dominated convergence gives the limit of the actual
integral, including its constant. Fixed zero intercepts do not invalidate the bound. $\square$

Applying Eq.~\eqref{eq:dualC:totaldiff} to the affine letters, every nonconstant logarithmic derivative has a fixed
nonpositive pole, and deleting a letter produces a finite recursive expansion
\begin{equation}
 G\bigl(\lambda_1(\Lambda),\ldots,\lambda_n(\Lambda);\Lambda\bigr)=\sum_wf_w\,G(-w;\Lambda),\qquad f_w\in\mathbb Z.
\end{equation}
The original convergent word vanishes at $\Lambda\to0$, while pure zero subwords keep their tangential logarithm; these
conditions fix the constants. Taking $\operatorname{Reg}$ yields the exact relations $\sum_wf_wR_w=G(e;1)$. For example,
\begin{equation}
 -R_{a,b}-R_{b,a}+R_{a+b,b}+R_{a+b,a}=-\zeta(2)\qquad(a,b>0),
\end{equation}
and the zero-intercept case is $-R_{0,b}+R_{b,b}=-\zeta(2)$. This zero-letter extension is needed by the target of
$\CSR_6$; treating all letters as strictly positive would not suffice.

\paragraph{The weight-five constants.} Multiple zeta values (MZVs) are written with decreasing indices,
\begin{equation}
 \zeta(n_1,\ldots,n_r)=\sum_{k_1>\cdots>k_r\ge1}\frac1{k_1^{n_1}\cdots k_r^{n_r}},\qquad n_1\ge2.
\end{equation}
For convergent words the right-hand sides of Eq.~\eqref{eq:dualC:confluence} are MZVs, with a sign set by the number of
one letters. Through weight four their reduction uses shuffle, stuffle, duality and $\zeta(2)^2=\frac52\zeta(4)$. At
weight five the required height-one identity follows from
\begin{equation}
 1-\frac{\Gamma(1-u)\Gamma(1-v)}{\Gamma(1-u-v)}
 =u\int_0^1t'^{-u-1}\bigl[(1-t')^{-v}-1\bigr]\dd t'
 =\sum_{m,n\ge1}\zeta(m+1,1^{n-1})\,u^mv^n,
 \label{eq:dualC:heightone}
\end{equation}
where $u$ and $v$ are formal variables here. The Taylor expansion of $\log\Gamma$ and the coefficient of $u^3v^2$
give $\zeta(4,1)=2\zeta(5)-\zeta(2)\zeta(3)$. The convergent stuffle and shuffle products give respectively
\begin{equation}
 \zeta(2)\zeta(3)=\zeta(2,3)+\zeta(3,2)+\zeta(5),\qquad
 \zeta(2)\zeta(3)=\zeta(2,3)+3\zeta(3,2)+6\zeta(4,1),
\end{equation}
hence
\begin{equation}
 \zeta(3,2)=3\zeta(2)\zeta(3)-\frac{11}{2}\zeta(5),\qquad
 \zeta(2,3)=\frac92\zeta(5)-2\zeta(2)\zeta(3).
\end{equation}
Duality supplies the remaining binary words of weight five. No assumption of numerical independence is needed. The
iterated-integral and MZV framework is standard~\cite{Brown:2011uzk,Brown:2013depth}; the programs of the compendium
implement the explicit conventions above.

\paragraph{The certificates and their logical scope.} The proposed closed expression is subtracted from the exact
archived expression of the renormalization-group (RG) assembly, and the difference is lifted to the word algebra. The
following finite identities then hold in the free rational vector space of words ($\CSR_7$ is in
Appendix~\ref{app:dualC:C7}):
\begin{center}\small
\begin{tabular}{@{}lrrr@{}}
\toprule
target & weight & nonzero words & certificate rows\\
\midrule
$D_3-D_3^{\rm closed}$ & 2 & 14 & 10\\
 & 3 & 54 & 47\\
 & 4 & 102 & 419\\
\midrule
$\CSR_6-\CSR_6{}^{\rm closed}$ & 2 & 26 & 21\\
 & 3 & 169 & 142\\
 & 4 & 608 & 3\,879\\
 & 5 & 1\,142 & 52\,339\\
\bottomrule
\end{tabular}
\end{center}
The totals, 170 words and 476 rows for $D_3$ and 1\,945 words and 56\,381 rows for $\CSR_6$, are those of
Sec.~\ref{sec:closed}. For each weight the target $\mathsf X$ is written explicitly as
\begin{equation}
 \mathsf X=\sum_jr_j\bigl(\mathsf B_j\mathbin{\amalg}\langle v_j\rangle\bigr),\qquad r_j\in\mathbb Q,
\end{equation}
where $\langle v_j\rangle$ is a word and every $\mathsf B_j$ is one of the analytically established zero relations
above (shuffle, stuffle, duality, the confluence relations and the height-one identity). The verifier regenerates
these relations from their labels and checks the rational coefficients word by word. Since the relations evaluate to
zero, so does the target.

The electronic certificate is part of the proof of the reduction, not optional numerical support: it contains all coefficients, and
none is inferred from decimal agreement. Both certificates are plain JSON embedded in standalone Python
verifiers (for $D_3$, \file{code/dual_d1/code/archived/D3_analytic_proof/verify_D3_proof_standalone.py}). Modular elimination was
used to discover the certificates, but the final verification uses exact rational arithmetic and no numerical period.
The proof does not require the chosen relations to span all MZV identities. Both certificates replay successfully
(Sec.~\ref{sec:review1}). In addition, a standard-library checker
(\file{code/dual_d1/code/archived/C7_closed_form_proof/code/verify_C6_input_map.py}) parses the archived beta
expressions as a restricted syntax tree, lifts the individual beta coefficients directly into the shuffle algebra and
verifies the trace assembly, the archived $\CSR_6$ expression, its difference from the closed form and all 1\,945
target coefficients, so that the certificate is tied to the RG inputs. It does not reintegrate the original kink
kernels.

%--------------------------------------------------------------------
\subsection{The seventh order and the weight-six certificate}\label{app:dualC:C7}
%--------------------------------------------------------------------
By Eq.~\eqref{eq:conv}, $\CSR_7=4D_4-T_3-T_2^2+\frac12\CSR_3\CSR_5$ needs a new determinant coefficient but no new
trace coefficient. This follows from expanding the characteristic polynomial, not from extrapolating a truncated beta
function. The closed form of $D_4$ is in Sec.~\ref{sec:closed}, and
\begin{equation}
 D_4=-2.5346297995382501152606064124\ldots,\qquad \CSR_7=51.06898861292810883118\ldots .
\end{equation}

\paragraph{New derivatives.} At this order the required jets are
\begin{equation}
\begin{aligned}
 F={}&z+u\bigl(1+az+p\db+Az^2+E\db z+V\db^2+\alpha z^3+\beta z^2\db+\gamma z\db^2\bigr)\\
 &+u^2\bigl(b+Gz+R\db+\Phi z^2+\Psi z\db\bigr)+u^3(J+\Omega z)+\cdots,\\
 H={}&-\db+4u+u^2\bigl(c+fz+q\db+Lz^2+Mz\db+Q\db^2\bigr)+u^3(k+Nz+P\db)+Ku^4+\cdots,
\end{aligned}
\label{eq:dualC:C7jets}
\end{equation}
and the determinant combination is
\begin{equation}
 D_4=B^{D_4}_{\rm old}-\frac{3\alpha}{32}+\frac\beta4-\frac\gamma2+\frac\Phi{16}-\frac\Psi8-\frac\Omega{32}
 -\frac Q8-\frac P{16}-\frac{3K}{128}.
 \label{eq:dualC:D4assembly}
\end{equation}
The polynomial $B^{D_4}_{\rm old}$ involves only earlier inputs; it is regenerated by
\file{code/dual_d1/code/archived/C7_closed_form_proof/code/derive_D4.py}. Differentiation must precede the
fixed-point substitution, since a term that is small at the fixed point can contribute to a derivative at the
required order. The derivatives that define the new inputs are listed in Appendix~\ref{app:dualC:C8setup}.

The two largest new integrations were the eight-kink current response and the charge derivative of the seven-kink odd
response. The latter enters
\begin{equation}
 \Omega=\partial_vb_7\big|_{v=1,\db=0}=-234.4214552889636875265144452523\ldots .
\end{equation}
Every term of $b_7$ in the subtraction recurrence, Eq.~\eqref{eq:dualA:recurrence}, is differentiated, including the
endpoint terms and the explicit factors $v$, and so is every product in the normalized response
$\mathcal C_7=A_7-A_5Z_2+A_3(Z_2^2-Z_4)-Z_2^3+2Z_2Z_4-Z_6$; the raw integral $\partial_vA_7$ alone is not $\Omega$.
An early partial calculation of $\CSR_7$ had used an incorrect numerical value of the second-tail input; the rebuilt
value is $Q=-107.0654814955476394045726813165\ldots$, and with all other inputs fixed
$\CSR_7=21.76630670180764789\ldots-\Omega/8$. The certificate below is tied to the rebuilt input expressions.

\paragraph{Higher-logarithm checks.} The rebuilt numerical integration passed the following higher-logarithm tests,
which do not use the new single-logarithm coefficient:
\begin{equation}
\begin{aligned}
 [\ell^8]\partial_v\mathcal C_7&=0, & [\ell^7]\partial_v\mathcal C_7&=-\frac{6416}{315}, &
 [\ell^6]\partial_v\mathcal C_7&=-\frac{1622}{9}, & [\ell^5]\partial_v\mathcal C_7&=\frac{3848\zeta(2)-8824}{15},\\
 && [\ell^7]\mathcal K_8&=-\frac{1088}{315}, & [\ell^6]\mathcal K_8&=-\frac{664}{15}, &
 [\ell^5]\mathcal K_8&=96\zeta(2)-\frac{3656}{15}.
\end{aligned}
\end{equation}
A direct finite-cutoff test compares the whole of $\partial_vA_7$ at $\varepsilon=1/10$ with a separate cubature of the
original kernels, not with a fitted asymptotic polynomial.

\paragraph{What the certificate establishes.} All 25 archived beta inputs of $\CSR_3$ to $\CSR_7$ are parsed
directly from the archived expressions and assembled into $D_4$ and $\CSR_7$. Subtracting the proposed zeta-value result and lifting
logarithms and polylogarithms to the word algebra gives a target with the homogeneous components
\begin{center}\small
\begin{tabular}{@{}rrrr@{}}
\toprule
weight & target words & certificate rows & final residual words\\
\midrule
2 & 37 & 33 & 0\\
3 & 337 & 284 & 0\\
4 & 1\,810 & 2\,880 & 0\\
5 & 7\,428 & 54\,363 & 0\\
6 & 15\,600 & 531\,318 & 0\\
\midrule
total & 25\,212 & 588\,878 & 0\\
\bottomrule
\end{tabular}
\end{center}
The weight-six identity includes the products of lower-weight periods; a reduction modulo products would not suffice.
In the free rational word space the certificate has the form
\begin{equation}
 \mathsf X_6=\sum_jr_j\,\mathsf B_j\mathbin{\amalg}\mathsf M_j,\qquad r_j\in\mathbb Q,
\end{equation}
where $\mathsf X_6$ is the weight-six target, every $\mathsf B_j$ is a proved zero relation and every multiplier
$\mathsf M_j$ is expanded with its full shuffle multiplicities, so that evaluation annihilates the right-hand side. The
verifier does not use the Lyndon-polynomial representation employed during discovery; it expands ordinary words with
exact integers after clearing denominators.

All 16 convergent binary boundary values of weight six are derived from 31 rational equations, using the shuffle and
stuffle products, duality and the height-one identity~\eqref{eq:dualC:heightone}. Examples are
\begin{equation}
\begin{aligned}
 \zeta(5,1)&=\tfrac34\zeta(6)-\tfrac12\zeta(3)^2, & \zeta(4,2)&=-\tfrac43\zeta(6)+\zeta(3)^2,\\
 \zeta(3,3)&=-\tfrac12\zeta(6)+\tfrac12\zeta(3)^2, & \zeta(2,4)&=\tfrac{25}{12}\zeta(6)-\zeta(3)^2.
\end{aligned}
\end{equation}
No assumption of numerical independence of MZVs is used. The final certificates and the target suffice for
verification; the large discovery matrices and redundant intermediate certificates are not part of the proof input.

%--------------------------------------------------------------------
\subsection{The eighth order: trace reduction and recurrence}\label{app:dualC:C8setup}
%--------------------------------------------------------------------
The integer powers of $x$ in $\De$ (the even powers of $\sqrt x$) come entirely from the half-trace $T/2$, so
$\CSR_8=8T_4$ (Eq.~\eqref{eq:conv}). No
new determinant coefficient $D_5$ is needed, but $T_4$ needs new derivatives that $D_4$ does not fix. The jets
required for this trace, keeping only the sectors that can contribute after differentiation, extend
Eq.~\eqref{eq:dualC:C7jets} by ten coefficients:
\begin{equation}
\begin{aligned}
 F={}&z+u\bigl(1+az+p\db+Az^2+E\db z+V\db^2+\alpha z^3+\beta\db z^2+\gamma\db^2z+W\db^3\bigr)\\
 &+u^2\bigl(b+Gz+R\db+\Phi z^2+\Psi\db z+S\db^2\bigr)+u^3(J+\Omega z+U\db)+J_9u^4+\cdots,\\
 H={}&-\db+4u+u^2\bigl(c+fz+q\db+Lz^2+M\db z+Q\db^2+L_3z^3+M_{21}\db z^2+M_{12}\db^2z\bigr)\\
 &+u^3\bigl(k+Nz+P\db+N_2z^2+N_\delta\db z\bigr)+u^4(K+K_vz)+\cdots.
\end{aligned}
\label{eq:dualC:C8jets}
\end{equation}
These are not the exact beta functions beyond the displayed order. For example, the term $K_vzu^4$ of $H$ must be kept,
because $H_z$ contributes at order $\db^4$.

\paragraph{Derivative dictionary.} All derivatives are taken at $v=1$, $\db=0$. The normalized auxiliary current
beta function is $\beta_h/(\sqrt2b_0)=-g^2+h_4g^4+h_6g^6+h_8g^8+\cdots$ (so $h_2=-1$), and $b_3,b_5,b_7,b_9$ are the odd
fugacity coefficients of the subtraction recurrence. The lower jets are
\begin{equation}
\begin{aligned}
 a&=b_{3,v}, & p&=b_{3,\db}, & A&=\tfrac12b_{3,vv}, & E&=b_{3,v\db}, & V&=\tfrac12b_{3,\db\db},\\
 \alpha&=\tfrac16b_{3,vvv}, & \beta&=\tfrac12b_{3,vv\db}, & \gamma&=\tfrac12b_{3,v\db\db}, && &&\\
 b&=b_5, & G&=b_{5,v}, & R&=b_{5,\db}, & \Phi&=\tfrac12b_{5,vv}, & \Psi&=b_{5,v\db},\\
 J&=b_7, & \Omega&=b_{7,v}, &&&&&&\\
 c&=-4h_4, & f&=-4h_{4,v}, & q&=-4h_{4,\db}, & L&=-2h_{4,vv}, & M&=-4h_{4,v\db},\\
 Q&=-2h_{4,\db\db}, & k&=-4h_6, & N&=-4h_{6,v}, & P&=-4h_{6,\db}, & K&=-4h_8,
\end{aligned}
\end{equation}
and the ten additional inputs of $\CSR_8$ are
\begin{equation}
\begin{aligned}
 W&=\tfrac16b_{3,\db\db\db}, & S&=\tfrac12b_{5,\db\db}, & U&=b_{7,\db}, & J_9&=b_9, & L_3&=-\tfrac23h_{4,vvv},\\
 M_{21}&=-2h_{4,vv\db}, & M_{12}&=-2h_{4,v\db\db}, & N_2&=-2h_{6,vv}, & N_\delta&=-4h_{6,v\db}, & K_v&=-4h_{8,v}.
\end{aligned}
\end{equation}
A mixed first derivative carries no factor $\frac12$; Taylor factors appear only for repeated derivatives, as
displayed.

\paragraph{Finite-response recurrence through ninth order.} The eight-kink current response is
\begin{equation}
 \mathcal K_8=Q_8-Q_6Z_2+Q_4(Z_2^2-Z_4)+Q_2(-Z_2^3+2Z_2Z_4-Z_6),
 \label{eq:dualC:K8norm}
\end{equation}
as in Eq.~\eqref{eq:dualA:normalization}, and the nine-kink odd response, the coefficient of $g^8$ in the ratio of the
odd generating function to the vacuum one (Appendix~\ref{app:dualA:moments}) at zero source ($A_1=1$), is
\begin{equation}
 \mathcal C_9=A_9-A_7Z_2+A_5(Z_2^2-Z_4)+A_3(-Z_2^3+2Z_2Z_4-Z_6)+Z_2^4-3Z_2^2Z_4+Z_4^2+2Z_2Z_6-Z_8.
 \label{eq:dualC:C9norm}
\end{equation}
The Jacobian entries needed next are
\begin{equation}
 j_7=-2\partial_v\mathcal C_{7,0}+\partial_\eta\mathcal C_{7,0},\qquad
 \lambda_6=-2\mathcal K_{6,0}-4v\,\partial_v\mathcal K_{6,0}+2v\,\partial_\eta\mathcal K_{6,0},
 \label{eq:dualC:j7lambda6}
\end{equation}
with the same rule as $j_3,j_5$ and $\lambda_2,\lambda_4$ in Eq.~\eqref{eq:dualA:jacobian}. The compendium defines
$j_7$ in this form but uses $\lambda_6$ without writing it; Eq.~\eqref{eq:dualC:j7lambda6} gives $\lambda_6$ by that
rule. The new counterterm recurrences, which continue Eq.~\eqref{eq:dualA:recurrence}, are
\begin{equation}
\begin{aligned}
 h_8={}&\mathcal K_{8,1}-2Q_{2,0}b_7-4\mathcal K_{4,0}b_5-6\mathcal K_{6,0}b_3-\lambda_2h_6-\lambda_4h_4+\lambda_6,\\
 b_9={}&\mathcal C_{9,1}-3\mathcal C_{3,0}b_7-5\mathcal C_{5,0}b_5-7\mathcal C_{7,0}b_3
 -2vj_1h_8-2vj_3h_6-2vj_5h_4+2vj_7.
\end{aligned}
\label{eq:dualC:b9}
\end{equation}
Products and derivatives are formed before the finite parts are extracted. An unavailable raw scalar is an error, not a
zero coefficient, and the compendium's full scalar-aware convolution code enforces this.

Solving $F=H=0$ with Eq.~\eqref{eq:dualC:C8jets} through $\db^4$ and expanding $T=2uF_u+(1+z)H_z$ gives
Eq.~\eqref{eq:C8assembly}: $\CSR_8$ is $B_{\rm old}$ plus a fixed linear combination of the ten new inputs. The
polynomial $B_{\rm old}$ involves only the 25 earlier inputs; it is printed in Appendix~\ref{app:dualC:Bold}. Its value is
$-861.6635672354797792014971909923\ldots$, and its rational projection is $\frac{1153}{6}$. Replacing it by its
rational projection before the assembly would lose genuine terms.

%--------------------------------------------------------------------
\subsection{The rational part of the eighth order}\label{app:dualC:C8rational}
%--------------------------------------------------------------------
The formal weight-zero map $\pi_0$ acts on the graded representation of an integrated quantity: it removes the
positive-weight constant periods. It is not a numerical limit in which the real zeta values vanish, and variable
logarithms cannot be discarded before integration; for example,
\begin{equation}
 \int_0^\infty\frac{\log(1+t')}{(1+t')^2}\,\dd t'=1.
\end{equation}
Products, coefficient convolutions and power divergences must already have been treated consistently when $\pi_0$ is
applied.

\paragraph{Four raw residues.} Define
\begin{equation}
\begin{aligned}
 \mathcal R_6&=\pi_0(\partial_v\partial_{\db}Q_6)_1, &
 \mathcal R_8&=\pi_0(\partial_vQ_8)_1,\\
 \mathcal R_{\rm tail}&=\pi_0(\partial_{\db}A_7)_1-\pi_0(\partial_{\db}Z_6)_1, &
 \mathcal R_9&=\pi_0(A_9)_1-\pi_0(Z_8)_1+2\pi_0(\partial_\eta A_7)_0,
\end{aligned}
\end{equation}
the residues of the six-kink and eight-kink current moments, of the seven-kink tail and of the nine-kink fugacity.
Before any of them was inserted, the exact trace and subtraction reduction gave
\begin{equation}
 \pi_0\CSR_8=\frac{5345}{12}-\frac12\mathcal R_6-\frac18\mathcal R_8+\frac34\mathcal R_{\rm tail}+\frac14\mathcal R_9 .
\end{equation}
The completed values are
\begin{equation}
 \mathcal R_6=\frac{399343}{2},\qquad \mathcal R_8=-757993,\qquad \mathcal R_{\rm tail}=-\frac{910891}{6},\qquad
 \mathcal R_9=\frac{1423087}{3},
\end{equation}
and after the first two are inserted the known term is $-\frac{111389}{24}$, which gives the sum
$\pi_0\CSR_8=88$ of Sec.~\ref{sec:C8}. None of the raw integrals was inferred from that answer or from the earlier
conjecture $160$. For reconstruction, the last five raw scalars are
\begin{center}\small
\begin{tabular}{@{}lr@{\qquad}lr@{}}
\toprule
raw scalar & rational value & raw scalar & rational value\\
\midrule
$\pi_0(\partial_{\db}A_7)_1$ & $-\frac{468683}{3}$ & $\pi_0(Z_8)_1$ & $\frac{23294}{3}$\\[2pt]
$\pi_0(\partial_{\db}Z_6)_1$ & $-\frac{8825}{2}$ & $\pi_0(\partial_\eta A_7)_0$ & $-\frac{184413}{2}$\\[2pt]
$\pi_0(A_9)_1$ & $666540$ & & \\
\bottomrule
\end{tabular}
\end{center}

\paragraph{The rational inputs and the full-Jacobian check.} A separate finite-response reconstruction gives
\begin{equation}
 U^{(0)}=\frac{16534}{3},\qquad J_9^{(0)}=-\frac{10192}{3},\qquad N_\delta^{(0)}=-20848,\qquad K_v^{(0)}=\frac{77920}{3},
\end{equation}
and the other six new rational inputs are
\begin{equation}
 W^{(0)}=168,\quad S^{(0)}=-1751,\quad L_3^{(0)}=128,\quad M_{21}^{(0)}=-384,\quad M_{12}^{(0)}=2720,\quad
 N_2^{(0)}=3712.
\end{equation}
A separate literal full-Jacobian calculation then yields
$T^{(0)}(\db)=\frac{\db}{2}+2\db^2+\frac{13}{2}\db^3+11\db^4+O(\db^5)$, which reproduces $\pi_0\CSR_8=8\times11$ without
using the residue assembly above.

\paragraph{Boundary differentiation and its checks.} If $M(\varepsilon)=\sum_{p,m}c_{p,m}\varepsilon^p\ell^m$ and
$B=-\varepsilon\partial_\varepsilon M$, then
\begin{equation}
 [\varepsilon^p\ell^0]B=-p\,c_{p,0}-c_{p,1},\qquad [\varepsilon^0\ell^1]M=-[\varepsilon^0\ell^0]B.
 \label{eq:dualC:Euler}
\end{equation}
Differentiating a consecutive-gap lower limit removes one integration. For $n=2m$ kinks the reflection weights of the
$2m-1$ faces are $(2,\ldots,2,1)$; for odd $n$ the faces are all paired. These factors and the minus sign in
Eq.~\eqref{eq:dualC:Euler} are essential. The negative powers of $\varepsilon$ are predicted from lower moments. For
$Q_{6,v\db}$ the predictions are $(c_{-2,0},c_{-2,1},c_{-1,0},c_{-1,1})^{(0)}=(-7,-32,2193,555)$; for $Q_{8,v}$ they are
\begin{center}\small
\begin{tabular}{@{}rrr@{}}
\toprule
power $p$ & $c_{p,0}^{(0)}$ & $c_{p,1}^{(0)}$\\
\midrule
$-3$ & $\frac23$ & $2$\\[2pt]
$-2$ & $-\frac{104}{3}$ & $-\frac{52}{3}$\\[2pt]
$-1$ & $-\frac{80255}{12}$ & $-\frac{20429}{6}$\\
\bottomrule
\end{tabular}
\end{center}
All comparisons in the rational calculation are exact. They are generated from lower moments and do not use the target
finite-power residue. The tail and ninth-order groups also pass complete comparisons of their negative-power
logarithmic polynomials, and further direct finite-cutoff and alternative endpoint-representation tests are kept in
their stage records.

%--------------------------------------------------------------------
\subsection{The full eighth order: the four computational stages}\label{app:dualC:C8packages}
%--------------------------------------------------------------------
The rational result does not give the full coefficient. Restoring the constant periods requires undoing every earlier
projection, keeping the integer-offset logarithms and lower-endpoint constants, and keeping the full counterterm
contributions.

\paragraph{The six inherited inputs.} Six of the ten new inputs have recorded numerical evaluations only:
\begin{center}\small
\begin{tabular}{@{}lrr@{}}
\toprule
input & full numerical value & rational projection\\
\midrule
$W$ & $-37.708111514952555\ldots$ & $168$\\
$S$ & $878.183733747745103\ldots$ & $-1751$\\
$L_3$ & $-0.638692810583406\ldots$ & $128$\\
$M_{21}$ & $-171.235576565588699\ldots$ & $-384$\\
$M_{12}$ & $-422.715057969894520\ldots$ & $2720$\\
$N_2$ & $189.742012494084223\ldots$ & $3712$\\
\bottomrule
\end{tabular}
\end{center}
They enter the assembly as inherited numerical records: decimal values from one recorded 100-digit evaluation, with no
surviving exact expression and no second evaluation (Sec.~\ref{sec:C8}). This is one reason why a complete analytic
proof for $\CSR_8$ is a further task. The other four new inputs were computed in four stages, F1 ($N_\delta$), F2
($U$), F3 ($K_v$) and F4 ($J_9$). Their definitions and subtraction factors are fixed by Eq.~\eqref{eq:C8assembly}
and the recurrence; they are not adjustable parameters.

\paragraph{F1a: the six-kink mixed logarithm.} Define
\begin{equation}
 P_6=(Q_{6,v\db})_1\big|_{v=1,\db=0}=[\varepsilon^0\ell^1]\int_{\Omega_6(\varepsilon)}\Xi_6W_6\,(S_2+2S_1S_2)\,\dd^6X,
\end{equation}
where $\Omega_6(\varepsilon)$ is the ordered hard-core domain of six kinks, $W_6$ the six-kink weight,
$\Xi_6=\sum_{j=1}^3(X_{2j}-X_{2j-1})$ the current weight, and $(S_2+2S_1S_2)W_6=\partial_v\partial_{\db}W_6$ the
mixed insertion, with both terms and no Taylor factor. Early quadrature estimates around $1261.9448\ldots$ were
not accurate enough for an integer-relation search with many constants. The final computation instead integrated the
complete cutoff-boundary source with all constant periods kept. With the five gaps written $r_i$ and $R=\sum_ir_i$,
the face integrals are
\begin{equation}
 B_j=\varepsilon\int_{\substack{r_i\ge\varepsilon\ (i\ne j)\\ R<1}}(1-R)\,\Xi_6W_6\,(S_2+2S_1S_2)\prod_{i\ne j}\dd r_i,
 \qquad r_j=\varepsilon,
\end{equation}
with $B=2B_1+2B_2+B_3$ (the reflection weights of the five faces) and $P_6=-[\varepsilon^0\ell^0]B$ by
Eq.~\eqref{eq:dualC:Euler}. The final record includes all 21 first-stage blocks and all 64 terminal outputs; exact
collection gives 66\,404 period terms, and the rational projection is $\frac{399343}{2}=\mathcal R_6$. The
value of $P_6$ is given in Sec.~\ref{sec:C8}. The two fresh evaluations at 70 and 100 working digits differ in the
matching point, the step size, the Taylor order and the asymptotic order. A complete four-kink mixed control agrees
exactly with the archived expression; a comparison of the full stage in moving-letter and in ordinary-word form agrees
in 196\,169 rational coefficients; and all 23\,816 rational prefactors are reconstructed exactly from their partial
fractions.

On a face all separations carry a factor $\varepsilon$, and the scale is restored by
\begin{equation}
 S_1=-3\ell+s_1,\qquad S_2=-3\ell^2+2s_1\ell+s_2,
\end{equation}
where $s_1$ and $s_2$ are the insertions of the rescaled configuration (the six alternating charges give
$\sum_{i<k}q_iq_k=-3$). The logarithm-free boundary coefficient uses $s_2+2s_1s_2$, but its variable and constant
logarithms remain. The full positive-weight divergence differences at the powers $\varepsilon^{-2}$ and
$\varepsilon^{-1}$ are numerically consistent with zero at about $10^{-88}$ and $10^{-87}$; these are not exact period
certificates.

\paragraph{F1b: vacuum normalization.} The six-kink current response is
$\mathcal K_6=Q_6-Q_4Z_2+Q_2Z_2^2-Q_2Z_4$, so that
\begin{equation}
 (\partial_v\partial_{\db}\mathcal K_6)_1=P_6+D_{\rm norm},\qquad
 D_{\rm norm}=[\varepsilon^0\ell^1]\,\partial_v\partial_{\db}\bigl(-Q_4Z_2+Q_2Z_2^2-Q_2Z_4\bigr)\big|_{v=1,\db=0}.
\end{equation}
The three product contributions are
\begin{equation}
 -1300.763962200486633\ldots,\qquad -108\ \text{(exact)},\qquad -1261.404579171677298\ldots,
\end{equation}
which give
\begin{equation}
 D_{\rm norm}=-2670.168541372163931149878165554559\ldots,\qquad \pi_0D_{\rm norm}=-\frac{8719}{2}.
\end{equation}
All 13 Leibniz terms are included, with their repeated-factor multiplicities. The terms with zero cutoff powers alone
would give $-1497.3871181439\ldots$, so the products of suppressed and divergent powers are substantial. The input map
uses 32 complete scalar coefficients; the 16 two-kink ones can be rederived by elementary integration, and the
four-kink finite expressions are inherited exact inputs, not integrated afresh in this step.

\paragraph{F1c: finite-response correction.} Define
\begin{equation}
 D_{\rm ct}=\partial_v\partial_{\db}\bigl[-2Q_{2,0}b_5-4\mathcal K_{4,0}b_3-\lambda_2h_4+\lambda_4\bigr]\big|_{v=1,\db=0},
\end{equation}
the derivative of the counterterm terms of $h_6$ in the recurrence. The complete result is
\begin{equation}
 D_{\rm ct}=-134.559964098286854425679630067592\ldots,\qquad \pi_0D_{\rm ct}=-190100.
\end{equation}
With $k_4=\mathcal K_{4,0}$ and $e_4=\partial_\eta\mathcal K_{4,0}$, the explicit coupling factors of $\lambda_4$ give
\begin{equation}
 (\lambda_4)_{v\db}=-6k_{4,v\db}-4k_{4,vv\db}+2e_{4,\db}+2e_{4,v\db};
\end{equation}
the extra terms would be lost if the prefactors were treated as constants. Two new complete four-kink finite parts are
\begin{equation}
\begin{aligned}
 (Q_{4,vv\db})_0&=17290.15809730659828787662714975844\ldots,\\
 (\partial_\eta Q_{4,v\db})_0&=34463.59564045596175027158175593235\ldots .
\end{aligned}
\end{equation}
Since $N_\delta=-4h_{6,v\db}$ and $h_{6,v\db}=P_6+D_{\rm norm}+D_{\rm ct}$, the three F1 stages combine into
\begin{equation}
 N_\delta=-4P_6+11218.914021881803142302231182488606\ldots ,
 \label{eq:dualC:NdP6}
\end{equation}
where the ellipsis abbreviates an archived exact period expression, not an unknown contribution. The rational check is
$\frac{399343}{2}-\frac{8719}{2}-190100=5212$, hence $N_\delta^{(0)}=-20848$.

\paragraph{F2: the seven-kink tail coefficient.} The derivative of the $b_7$ recurrence with respect to $\db$ isolates
\begin{equation}
 U=(\partial_{\db}A_7)_1-(\partial_{\db}Z_6)_1+2(\partial_\eta\partial_{\db}A_5)_0-2(\partial_\eta\partial_{\db}Z_4)_0
 +D_{\rm lower},
\end{equation}
where $D_{\rm lower}$ collects the remaining lower-order terms of the differentiated recurrence. The lower shift
terms, the vacuum tail and the endpoint terms are kept unprojected. Collecting everything but the first term into $D_{\rm known}$,
\begin{equation}
\begin{aligned}
 (\partial_{\db}A_7)_1&=315459.451152690228448576716538300797\ldots,\\
 D_{\rm known}&=-316756.139112981754940233110311642866\ldots,
\end{aligned}
\end{equation}
and their sum is the value of $U$ in Sec.~\ref{sec:C8}, $U=-1296.687960291526491656393773342069\ldots$, with
$U^{(0)}=\frac{16534}{3}$. The seven-kink expression has 66\,476 terms and the final $U$ 81\,495; their exact scalar
relation and the rational projection are checked, and full evaluations at 80 and 110 digits cover 57\,346 periods
through weight seven. A separately derived leading logarithm,
\begin{equation}
 [\varepsilon^0\ell^8]\partial_{\db}A_7=-\frac{2362}{315},\qquad [\varepsilon^0\ell^9]\partial_{\db}A_7=0,
\end{equation}
agrees with the leading-logarithm recursion of the lower flow; this is not a replay of the whole subleading logarithmic
polynomial. The historical coverage record of F2 names 258 intermediate files that are not among the surviving files.
The complete period expression of $U$ and its numerical replay survive, but an independent replay from the full kernels
would require regenerating those source stages.

\paragraph{F3: the eight-kink logarithm and six-kink finite responses.} Define the full, unprojected scalars
\begin{equation}
 P_8=(Q_{8,v})_1,\qquad X_{vv}=(Q_{6,vv})_0,\qquad X_E=(\partial_v\partial_\eta Q_6)_0 .
\end{equation}
The finite-response calculation of $K_v=-4h_{8,v}$ gives
\begin{equation}
 K_v=-4P_8+16X_{vv}-8X_E+\mathcal S_3,\qquad \mathcal S_3=24144.925877712789692098427372845263\ldots,
\end{equation}
with
\begin{equation}
\begin{aligned}
 X_{vv}&=8059.736102186616125752727901120985\ldots,\\
 X_E&=-88077.172363552137052197675538291236\ldots,\\
 P_8&=214801.261132598641881962805253618627\ldots,\\
 K_v&=-1486.962109278823406127742917363593\ldots .
\end{aligned}
\end{equation}
The complete period expression of $P_8$ has 120\,411 terms, and all 37 original source blocks are included through the
weighted boundary calculation. Its rational projection is $-757993=\mathcal R_8$, a separate check of the rational
calculation.

For certain six-kink finite endpoint responses, scaling gives an overall factor $\varepsilon$ multiplying the last
integral; proper rational terms then grow only logarithmically after integration and vanish as
$\varepsilon\log^k\varepsilon$, so only the part with polynomial growth contributes to the finite constant. This
shortcut is not applied to $P_8$, whose scaling differs and whose proper rational terms contribute to the residue;
both endpoints and the corresponding regularized constants are kept. The complete expression of $Q_{6,vv}$ at
$\varepsilon=1/10$ equals $425.5449632473382428\ldots$ and agrees with direct quadrature of the original kernels. The
full source of $P_8$ passes the six rational divergence checks of the table in Appendix~\ref{app:dualC:C8rational}, and a
positive-weight divergence comparison is numerically below $10^{-85}$. Exact source coverage and numerical period checks
are distinct kinds of check.

\paragraph{F4: the ninth-order fugacity.} The last new combination is
\begin{equation}
 J_9=A_{9,1}-Z_{8,1}+2(\partial_\eta A_7)_0-2(\partial_\eta Z_6)_0+D_{9,\rm lower},
\end{equation}
with the full inputs
\begin{center}\small
\begin{tabular}{@{}lr@{}}
\toprule
quantity & numerical value\\
\midrule
$A_{9,1}$ & $-613817.497451118368344265\ldots$\\
$Z_{8,1}$ & $11059.024476469114411035\ldots$\\
$(\partial_\eta A_7)_0$ & $-14142.938944879185749695\ldots$\\
$(\partial_\eta Z_6)_0$ & $305.264828822013665772\ldots$\\
$D_{9,\rm lower}$ & $653499.204958990556939720\ldots$\\
\bottomrule
\end{tabular}
\end{center}
They give $J_9=-273.724515999324646514077910398647\ldots$ and $J_9^{(0)}=-\frac{10192}{3}$; the value is a cancellation
between terms of order $6\times10^5$. The complete expression has 84\,928 terms. Some base-integral stages were
restored from rational-stage files only after an exact audit of the omitted constant terms, and the endpoint-offset
logarithms were restored separately. An intermediate error in an endpoint-vacuum multiplicity was corrected, and the
affected outputs were regenerated before the accepted evaluations; only the corrected scalar enters the assembly. The
independent re-derivation of Sec.~\ref{sec:indep} differs from this record in the first entry of the table,
$A_{9,1}$, and agrees on the other four pieces of $J_9$.

%--------------------------------------------------------------------
\subsection{The full numerical assembly and the zeta-value candidate}\label{app:dualC:C8assembly}
%--------------------------------------------------------------------
The six inherited inputs combined with $B_{\rm old}$ give
\begin{equation}
 B_{\rm done}=B_{\rm old}+4W+2S+\tfrac3{32}L_3-\tfrac14M_{21}+\tfrac12M_{12}-\tfrac1{16}N_2
 =563.404066124777663118309147204616\ldots,
\end{equation}
and Eq.~\eqref{eq:C8assembly} becomes the scalar identity
\begin{equation}
 \CSR_8=B_{\rm done}+\frac{N_\delta}{8}+\frac{3U}{4}+\frac{K_v}{32}+\frac{J_9}{4}.
\end{equation}
With Eq.~\eqref{eq:dualC:NdP6} it reduces to $\CSR_8=C_{\rm fixed}-\frac12P_6$, with the $C_{\rm fixed}$ and the
record of Eq.~\eqref{eq:C8num} given in Sec.~\ref{sec:C8}. Every additive ingredient has a numerical value; there is
no missing F3 or F4 term, and the exact $P_6$ integration removes the earlier $10^{-4}$ discrepancy of the
quadrature estimates of $P_6$ (Appendix~\ref{app:dualC:numerics}). A separate literal-polynomial implementation solves
$F=H=0$ with Eq.~\eqref{eq:dualC:C8jets} through $\db^4$, differentiates the full beta functions and obtains the same
$8T_4$. The exact equality of this result, of the literal trace polynomial and of $C_{\rm fixed}-\frac12P_6$, with the
supplied decimal strings treated as fractions, checks the assembly; it does not certify the numerical integration
errors of those strings.

\paragraph{Recognition.} The first fresh evaluation, at 70 digits, was used for an integer-relation search in the
basis of Sec.~\ref{sec:C8}, $\{\zeta(2),\zeta(3),\zeta(4),\zeta(5),\zeta(2)\zeta(3),\zeta(6),\zeta(3)^2,\zeta(7),
\zeta(2)\zeta(5),\zeta(3)\zeta(4)\}$. The relation found for $(\CSR_8-88,\ \text{basis})$ is
\begin{equation}
 (3,\,534,\,-330,\,-339,\,-14676,\,7128,\,8630,\,-5616,\,36504,\,-18144,\,-4608),
\end{equation}
which is Eq.~\eqref{eq:C8cand}. The second evaluation, at 100 digits and with the candidate held fixed, differs from
it by about $\sci{3.0854}{-86}$; a separate search with the rational intercept free recovers $88$. A looser earlier
search had produced a different relation that failed at higher precision; its rejected record is kept as a warning
against overinterpreting integer relations. The highest-weight part of the candidate,
$-12168\zeta(7)+6048\zeta(2)\zeta(5)+1536\zeta(3)\zeta(4)$, is compatible with the maximal weight seven
(Appendix~\ref{app:dualC:conjectures}); compatibility proves neither this coefficient identity nor a period-class theorem
to all orders. With the even zeta values written out, the same candidate is
\begin{equation}
\begin{aligned}
 \CSR_8\overset{?}{=}{}&88-\frac{89\pi^2}{3}+\frac{113\pi^4}{90}-\frac{1726\pi^6}{567}
 +\Bigl(110-396\pi^2+\frac{256\pi^4}{15}\Bigr)\zeta(3)+(4892+1008\pi^2)\zeta(5)\\
 &+1872\zeta(3)^2-12168\zeta(7),
\end{aligned}
\end{equation}
algebraically identical to Eq.~\eqref{eq:C8cand}.

\paragraph{Status.} The record is numerical: its digits are cross-checked only through $P_6$, to about 57 decimal
places; six of its inputs are single inherited records; and it carries no error bound (Sec.~\ref{sec:C8}).
Equation~\eqref{eq:C8cand} is conjectured: it has no certificate, and the record is its only high-precision support, so its agreement
with the record is not independent evidence for the record's digits. The independent re-derivation of
Sec.~\ref{sec:indep} gives $\CSR_8=242.8\pm2.4$ (one standard error), 1.9 to 4.0 standard errors below the record
depending on the error estimator, with the difference in $A_{9,1}$. The magnetization-scheme calculation of
Appendix~\ref{app:integ}, which shares no input with the compendium, gives $247.38118\pm0.13$, $\sci{2.5}{-5}$ from the
record; it supports the leading digits of the record but not the coefficients of Eq.~\eqref{eq:C8cand}. A proof of
Eq.~\eqref{eq:C8cand} would follow the
route stated in Sec.~\ref{sec:C8}; its verifier would also have to regenerate each source relation and check the
complete rational vector, not only its primitive quotient or its numerical value. No such certificate exists.

%--------------------------------------------------------------------
\subsection{Numerical evaluation of periods and error control}\label{app:dualC:numerics}
%--------------------------------------------------------------------
The numerical value of a regulated period is not obtained by substituting a large cutoff into a divergent integral; it
is the constant term of a controlled expansion in logarithms and inverse powers. For a fixed nonnegative letter $a$
the differential system is triangular,
\begin{equation}
 \frac{\dd}{\dd\Lambda}G(-a,-w;\Lambda)=\frac{G(-w;\Lambda)}{\Lambda+a},
\end{equation}
where $(-a,-w)$ is the word $-w$ with the letter $-a$ prefixed. The empty word is one, all-zero words have the tangential value
$\log^n\Lambda/n!$, and words beginning with a zero letter need their regular and singular parts treated explicitly.

For $a>0$ and a regular Taylor expansion $G(-w;\Lambda)=\sum_nc_{w,n}\Lambda^n$, the coefficients $c_{a,w,n}$ of $G(-a,-w;\Lambda)$ near the origin obey
\begin{equation}
 c_{a,w,n}=\frac1n\sum_{j=0}^{n-1}(-1)^ja^{-j-1}c_{w,n-1-j},\qquad n\ge1.
\end{equation}
The local Taylor expansion is continued along the positive axis in steps smaller than the distance to the nearest
singularity. At a large enough matching point the expansion~\eqref{eq:dualC:infinity},
$G(-w;\Lambda)=\sum_{k\ge0}\Lambda^{-k}P_{w,k}(\log\Lambda)$, is used. For the word $(a,w)$ the leading polynomial satisfies
$P'_{a,w,0}=P_{w,0}$, with the constant fixed by matching, and the higher polynomials solve
\begin{equation}
 (\partial_L-k)P_{a,w,k}(L)=\sum_{j=0}^k(-a)^jP_{w,k-j}(L),\qquad k\ge1.
\end{equation}
The regularized period is $R_w=P_{w,0}(0)$. Ordering the suffixes depth first lowers the peak memory without changing
the recurrence. A C++ implementation with the GMP library supplies the arithmetic; a separate Python implementation with
decimal arithmetic checks sampled periods and closed controls.

Every accepted run records its working precision, matching point, Taylor order, asymptotic order, step fraction,
input hashes and backend. Repeating a run with more printed digits only is not an independent accuracy test; the
comparisons of the compendium change several settings at once. Their stability and the controls on lower integrals
support the quoted values, but no rigorous global error bound for the high-order assembly exists (Sec.~\ref{sec:C8},
``How many digits'').

\paragraph{The superseded Mellin quadrature.} For a cutoff moment $M(\varepsilon)$ the Mellin transform
$\mathcal M(\varsigma)=\int_0^{\varepsilon_{\max}}\varepsilon^{\varsigma-1}M(\varepsilon)\,\dd\varepsilon$ converges
in a right half-plane and is continued analytically; a term $\varepsilon^0\log\varepsilon$ contributes
$-\varsigma^{-2}$. Sector
decomposition with analytic power and logarithm moments evaluates the continued transform without integrating
divergent powers. The compendium's six-kink runs of this kind covered all ordered-gap sectors with compensated
summation, but its finest runs of $P_6$ still differed at about $10^{-4}$. They are kept only as historical
diagnostics and are not used in the high-precision test of the candidate; the exact boundary calculation of
Appendix~\ref{app:dualC:C8packages} replaces them.

%--------------------------------------------------------------------
\subsection{General-formula conjectures and what survived}\label{app:dualC:conjectures}
%--------------------------------------------------------------------
Matching a finite series and deriving the next coefficient from an independent integral are different tests. The
compendium used the latter to reject several attractive extrapolations; none of the general generating formulas it
tried is validated.

\paragraph{Exact recursion versus an all-orders solution.} Define
\begin{equation}
 Q_{\rm root}(\db)=\frac{T(\db)^2/4-D(\db)}{2\db}=1+\sum_{m\ge1}q_m\db^m,\qquad
 \sqrt{Q_{\rm root}}=1+\sum_{m\ge1}s_m\db^m .
\end{equation}
Then, with $T_1=\frac12$ and $s_0=1$,
\begin{equation}
\begin{gathered}
 q_m=\frac18\sum_{i=1}^mT_iT_{m+1-i}-\frac12D_{m+1},\qquad
 s_m=\frac12\Bigl(q_m-\sum_{i=1}^{m-1}s_is_{m-i}\Bigr),\\
 \CSR_{2m}=2^{m-1}T_m,\qquad \CSR_{2m+1}=-2^{m+1}s_m .
\end{gathered}
 \label{eq:dualC:allrecursion}
\end{equation}
These equations are exact algebra; they do not determine uncalculated invariants. For instance, $\CSR_9$ needs the new
determinant coefficient $D_5$, which $T_4$ does not fix.

\paragraph{Filtered MZV structure.} Give $\zeta(n_1,\ldots,n_r)$ the weight $\sum_in_i$, with weights adding under
products, and let $W_N\mathcal Z$ be the rational span of MZVs of weight at most $N$. The structural conjecture that
remains viable is
\begin{equation}
 \CSR_n\in W_{n-1}\mathcal Z .
 \label{eq:dualC:weightconj}
\end{equation}
It is not uniform transcendentality, and it does not specify the rational coefficients. The closed forms through
$\CSR_7$ obey it, and the candidate~\eqref{eq:C8cand} is consistent with it; its unproved status must not be upgraded
into a theorem to all orders. The geometry of ordered points on the line (genus zero) motivates the period
class~\cite{Brown:2009qja,Brown:2011uzk}. A conditional angular-counting argument suggests
\begin{equation}
 [u^kz^j\db^r]F\in W_{2k+j+2r-1}\mathcal Z,\qquad [u^kz^j\db^r]H\in W_{2k+j+2r-2}\mathcal Z .
\end{equation}
The individual inputs of the prescription contain the regulator constants $\log2$, $\log3$, $\Li_2(-\frac12)$ and
$\Li_4(\frac12)$, which are not written as MZVs and which cancel only from the invariants (Sec.~\ref{sec:inputs1}; $\Li_4(\frac12)$ enters through the input $E$,
Appendix~\ref{app:dualB:D3});
the compendium does not say how the bound is to be read for them. If subtraction and matching preserve these bounds and
close in the two essential couplings, the fixed-point recursion divides only by the rational determinant $-4$ of the
leading Jacobian of $(F,H)$ in $(u,z)$, and the bounds imply Eq.~\eqref{eq:dualC:weightconj}. Proving that every
required renormalization operation preserves them is the missing step. At weight eight the motivic MZV algebra admits a
new depth-two direction, representable by $\zeta(5,3)$, beyond products of single zeta values~\cite{Brown:2013depth}.
Under the conjectured bound it could first enter $\CSR_9$, and it need not have a nonzero coefficient; this is not a
claim of numerical independence of the corresponding real constants.

\paragraph{Specific guesses and direct tests.} For the inverse curve $y=1-\De$ write $x=y^2\mathcal G(y)/4$. A
specified trial function $\widehat{\mathcal G}$ predicts, by Lagrange inversion,
\begin{equation}
 \widehat C_n=-\frac{2^n}{n}[y^{n-1}]\,\widehat{\mathcal G}(y)^{-n/2}.
\end{equation}
Two rational--digamma trials, with scales $(1,2,3)$ and $(1,2,4)$, fit $\CSR_1$ to $\CSR_6$ exactly but predicted
$\CSR_7\simeq17.91784$ and $247.18113$; the completed value $51.06899\ldots$ rejects both fixed generators. A later
implicit Gamma family used
\begin{equation}
 \mathcal K_{ab}(y,t)=1-\frac{\Gamma(1+ay)^2}{\Gamma(1+ay+bt)\,\Gamma(1+ay-bt)},\qquad
 \mathcal S=y^2-4t^2+t^2\Bigl[P(y)+\sum P_{ab}(y)\,\mathcal K_{ab}\Bigr],
\end{equation}
with $t=\sqrt x$. It fits all seven supplied coefficients but predicts the wrong rational part of $\CSR_8$. Quadratic
log-Gamma and Borel--Laplace variants also fit their inputs but fail the same direct test. The rational-invariant
Pad\'e prescription $\widehat C_8=\CSR_6{}^2/\CSR_4$ fails after clearing denominators and taking weight zero:
\begin{center}\small
\begin{tabular}{@{}>{\raggedright\arraybackslash}p{.38\textwidth}>{\raggedright\arraybackslash}p{.29\textwidth}>{\raggedright\arraybackslash}p{.22\textwidth}@{}}
\toprule
fixed hypothesis & prediction tested & outcome\\
\midrule
geometric rational trace & $\pi_0\CSR_8=160$ & false: $88$\\
implicit Gamma fit & $\pi_0\CSR_8=-\frac{55777}{64}$ & false\\
quadratic log-Gamma and Borel--Gamma fits & $\pi_0\CSR_8=-\frac{78669187}{917504}$ & false\\
rational-invariant Pad\'e & $\pi_0\CSR_8=26^2/4=169$ & false\\
rational--digamma fits & $\CSR_7\simeq17.92$ or $247.18$ & false\\
filtered MZV class and weight bound & constants allowed at each order & unproved; not contradicted\\
determinant-only rational guess & $D_m^{(0)}=(3^m+2^m+7)/8$ ($m\ge2$) & next value $D_5^{(0)}=\frac{141}{4}$ untested\\
\bottomrule
\end{tabular}
\end{center}
The rational trace guess $T_m^{(0)}=(3^m-1)/4$ is excluded by $T_4^{(0)}=11\ne20$; this alone does not exclude its
determinant-only companion. Adjusting a new parameter after learning $\CSR_8$ is a new fit, not a successful prior
prediction.

\paragraph{Gamma-product and Euler-kernel obstructions.} The logarithm of a finite product of Gamma functions with
rational linear arguments consists of single-zeta terms, whereas the inverse series contains
\begin{equation}
 [y^5]\log\mathcal G(y)\big|_{\text{weight }5}=-\frac{17}{2}\zeta(5)+2\zeta(2)\zeta(3).
\end{equation}
This excludes the simple Gamma-product ansatz as a formal (motivic) identity, but not every nonlinear implicit Gamma
equation. A broader single-product trial,
\begin{equation}
 K_{\rho,e}(y)=\sum_{n\ge1}\frac1{n(n+ey)}\prod_{k<n}\rho(y/k),\qquad
 \mathcal G_{\max}(y)=1+\frac{y^2}{2}K_{\rho,e}(y),
\end{equation}
was tested with an arbitrary analytic $\rho$, $\rho(0)=1$. Formal matching of the product terms requires both
\begin{equation}
 8e^2-16e+23=0,\qquad 864e^4-4176e^3+16504e^2-21884e+14095=0,
\end{equation}
whose resultant is $401375222784\ne0$. This exact exclusion applies to the stated family, not to all possible spectral
kernels; the code derives the elimination equations rather than loading them as empirical numerical constraints.

%--------------------------------------------------------------------
\subsection{The two ends are not a coefficient-by-coefficient duality}\label{app:dualC:endpoints}
%--------------------------------------------------------------------
Both perturbative descriptions concern the same long-range family. The expansion of Sec.~\ref{sec:d1} is about $s=1$, with
$x\to0$; the expansion about the generalized free field (GFF) with a quartic coupling is about $s=\frac12$, with
\begin{equation}
 \varepsilon_{\rm G}=2s-1=1-4x\to0 .
\end{equation}
The strong--weak duality compares two descriptions at the same $s$; it is not a symmetry that exchanges the two
ends~\cite{Benedetti:2025nzp}. A finite GFF expansion is not controlled at $\varepsilon_{\rm G}=1$ and cannot be
substituted there to compute a coefficient $\CSR_n$.

In one dimension $\De=1-\nu^{-1}$, and the GFF expansion gives
\begin{equation}
 \De=\frac12-\frac{1-4x}{6}+\frac{\pi+4\log2}{9}(1-4x)^2+O\bigl((1-4x)^3\bigr),
\end{equation}
where the three-loop diagram must be the corrected tetrahedron value, not the superseded
one~\cite{Benedetti:2020rrq,Benedetti:2024mqx}. With $t=\sqrt x$ and $f(t)=\De(t^2)$, this gives the local derivatives
at the GFF end
\begin{equation}
 f(\tfrac12)=\frac12,\qquad f'(\tfrac12)=\frac23,\qquad f''(\tfrac12)=\frac43+\frac{32}{9}(\pi+4\log2).
\end{equation}
They are not relations derived from $\CSR_1$ to $\CSR_7$. Only under extra convergence and termwise-continuation
assumptions could one write, for the GFF coefficients $a_m$ of $\varepsilon_{\rm G}^m$,
\begin{equation}
 a_m=(-1)^m\sum_{n\ge1}\frac{\CSR_n}{2^n}\binom{n/2}{m},\qquad m\ge1,
\end{equation}
and even then each $a_m$ mixes infinitely many orders of the expansion about $s=1$. For an asymptotic series alone the
identity is not justified: a flat term $e^{-1/t^2}h(t)$ changes the data at the GFF end without changing any formal
coefficient $\CSR_n$. This is a logical example, not a proposed physical correction.

\paragraph{An independent global exclusion.} A specified global Gamma trial fails the value at the GFF end outright.
Its one-parameter repair enforces the value and the slope, and matching the curvature fixes its last parameter; the
resulting third derivative is less than $927$. With $A_{\rm G}=\pi+4\log2$, the corrected cubic coefficient
$b_{\rm G}$ of the GFF expansion of $\nu^{-1}=1-\De$ (the coefficient of $\varepsilon_{\rm G}^3$; the compendium calls it
the cubic mass coefficient) and the third derivative are
\begin{equation}
 b_{\rm G}=\frac{\mathfrak a}{81}-\frac{A_{\rm G}^2+\pi^2/3}{108},\qquad f'''(\tfrac12)=\frac{64A_{\rm G}}{3}+384b_{\rm G},
\end{equation}
where, in the corrected integral normalization~\cite{Benedetti:2024mqx},
\begin{equation}
 \mathfrak a=2\int_{\mathbb R^2}\frac{\dd u\,\dd v}{\sqrt{|u(1-u)v(1-v)(u-v)|}}
\end{equation}
($u$ and $v$ are integration variables here). Reflection and inversion divide the integral into three equal diagonal
and six equal off-diagonal sectors; with their integrals $\mathcal D$ and $\mathcal O$, $\mathfrak a=6\mathcal
D+12\mathcal O$. The off-diagonal substitution $v\to1/v$ gives $\mathcal O\ge\pi^2$. For independent arcsine variables
$U_1,U_2$ on $[0,1]$, $\mathcal D=\pi^2\,\mathbb E|U_1-U_2|^{-1/2}$ and $\mathbb E|U_1-U_2|=4/\pi^2$, so Jensen's
inequality gives
\begin{equation}
 \mathcal D\ge\frac{\pi^3}{2},\qquad \mathfrak a\ge3\pi^3+12\pi^2>211 .
\end{equation}
This implies $f'''(\frac12)>992$, which contradicts the repaired global trial; the archived proof gives exact rational
interval enclosures for the strict inequalities. The argument excludes that global function family; without further
analyticity or uniqueness assumptions it does not by itself exclude a formal rule for the coefficients $\CSR_n$. The
rational tests of $\CSR_8$ in Appendix~\ref{app:dualC:conjectures} do not need that global assumption.

%--------------------------------------------------------------------
\subsection{Further consequences of the shadow relation}\label{app:dualC:shadow}
%--------------------------------------------------------------------
The results of this subsection are separate from the thermal series. They use the published shadow identity and
leading data of operator product expansion (OPE) coefficients of the dual~\cite{Benedetti:2025nzp}, not the integrals
of many kinks. Unit-normalized primaries obey $\langle\mathcal O_i(X)\mathcal O_j(0)\rangle=\delta_{ij}|X|^{-2\Delta_i}$,
and $\chi$ denotes the renormalized, unit-normalized primary of dimension $\Delta_\chi=1-x$ obtained from the current
of Eq.~\eqref{eq:dual}, not the unsubtracted current. Here $x=\Ds=\db/2$ as throughout.

\paragraph{The gamma-function factor.} A scalar shadow transform is proportional to
\begin{equation}
 \mathsf S_xf(X)=\mathcal N_x\int\frac{f(Y)\,\dd Y}{|X-Y|^{2(1-x)}},
\end{equation}
defined by continuation of homogeneous distributions with contact subtractions. For three-point functions at separated
points the factor that depends on the pair follows from
\begin{equation}
 \int\dd X\,e^{-\mathrm ipX}|X|^{-\alpha}=2^{1-\alpha}\sqrt\pi\,\frac{\Gamma((1-\alpha)/2)}{\Gamma(\alpha/2)}\,|p|^{\alpha-1}.
\end{equation}
Putting one spectator at zero and one at infinity and convolving the remaining power law gives
\begin{equation}
 q(x;\Delta_i,\Delta_j)=\frac{\Gamma\bigl((1-x+\Delta_i-\Delta_j)/2\bigr)\Gamma\bigl((1-x-\Delta_i+\Delta_j)/2\bigr)}
 {\Gamma\bigl((x+\Delta_i-\Delta_j)/2\bigr)\Gamma\bigl((x-\Delta_i+\Delta_j)/2\bigr)}.
\end{equation}
The published relation has the form $c_{\chi ij}/c_{\sigma ij}=\mathcal N(\db)\,q(x;\Delta_i,\Delta_j)$, with a
normalization independent of the spectator pair (Ref.~\cite{Benedetti:2025nzp}, Eq.~(5.116)). For the leading
operator $\mathcal O_+$, of dimension $\Delta_+$ (Eq.~\eqref{eq:eigen}), the published data give
$c_{\sigma\chi+}/c_{\sigma\sigma+}\sim1/\sqrt{\db}$, while $q(x;x,\Delta_+)\sim1/(\sqrt2\pi\sqrt{\db})$; hence
$\mathcal N(0)=\sqrt2\pi$. Applying the relation twice to an even primary $\mathcal O$ of dimension $\Delta$ gives
\begin{equation}
 \frac{c_{\chi\chi\mathcal O}}{c_{\sigma\sigma\mathcal O}}=\mathcal N(\db)^2P(x,\Delta),\qquad
 P(x,\Delta)=\frac{\Gamma\bigl((1-2x+\Delta)/2\bigr)\Gamma\bigl((2-2x-\Delta)/2\bigr)}
 {\Gamma\bigl((2x-\Delta)/2\bigr)\Gamma\bigl((2x-1+\Delta)/2\bigr)}.
 \label{eq:dualC:shadowP}
\end{equation}

\paragraph{The isolated dimension-two primary.} At the endpoint the unique primary at dimension two that is even under
parity and under the Ising symmetry is $\rho\to{:}\chi^2{:}/\sqrt2$. Its published leading data are
$c_{\chi\chi\rho}=\sqrt2+\cdots$ and $c_{\sigma\sigma\rho}=\db/\sqrt2+\cdots$, so the double-shadow ratio diverges as
$2/\db$. A generic displacement $\Delta=2+a'\db+\cdots$ leaves Eq.~\eqref{eq:dualC:shadowP} finite; the required pole
lies at $\Delta=2-\db+e$ with $e=o(\db)$, where
\begin{equation}
 P\bigl(\tfrac{\db}{2},2-\db+e\bigr)=\frac{\Gamma(\frac32-\db+\frac e2)\,\Gamma(-\frac e2)}
 {\Gamma(-1+\db-\frac e2)\,\Gamma(\frac12+\frac e2)}\sim\frac{\db}{e} .
\end{equation}
Matching $2/\db\sim2\pi^2\db/e$ gives
\begin{equation}
 \Delta_\rho=2-\db+\pi^2\db^2+O(\db^3)=2-2x+4\pi^2x^2+O(x^3).
\end{equation}
For the same solution $q(x;x,\Delta_\rho)=\pi\db+O(\db^2)$, hence
\begin{equation}
 c_{\sigma\chi\rho}=\pi^2\db^2+O(\db^3).
\end{equation}
The sign assumes $c_{\chi\chi\rho}>0$ and the stated sign of $\chi$. The integer-power remainders assume analytic
continuation of the isolated even branch; without that regularity the last justified term is followed by $o(\db^2)$.
Conjugation by $\hat\sigma_3$ implements $g\mapsto-g$, which explains why an isolated even analytic branch has no odd
powers of $g$. No new fully renormalized four-point integral was used.

\paragraph{A weighted statement, not a resolved higher spectrum.} Let $N=2,4,6,\ldots$, and let $\alpha$ label the
branches that approach dimension $N$, with $v_\alpha$ the value of $c_{\chi\chi\alpha}$ at $\db=0$ and
$u_\alpha=\lim_{\db\to0}c_{\sigma\sigma\alpha}/\db$. For regular branches visible to the leading $\chi\chi$ OPE write
$\Delta_\alpha=N-\db+\gamma^{[2]}_\alpha\db^2+o(\db^2)$. The shadow residue implies
\begin{equation}
 v_\alpha\gamma^{[2]}_\alpha=\pi^2N(N-1)\,u_\alpha,
\end{equation}
with a zero $u_\alpha$ interpreted by continuity. The published mixed-correlator weight and the free $\chi$ block
coefficient are
\begin{equation}
 A_N=\sum_\alpha u_\alpha v_\alpha=\frac{4(N-1)[(N-2)!]^2}{N(2N-2)!},\qquad
 a_N=\sum_\alpha v_\alpha^2=\frac{2N(N-1)[(N-1)!]^2}{(2N-2)!};
\end{equation}
the second follows by matching $1+z_c^2+[z_c/(1-z_c)]^2$ to $1+\sum_Na_Nz_c^N{}_2F_1(N,N;2N;z_c)$, with $z_c$ the cross
ratio. Therefore
\begin{equation}
 \frac{\sum_\alpha v_\alpha^2\gamma^{[2]}_\alpha}{\sum_\alpha v_\alpha^2}=\frac{2\pi^2}{N(N-1)} .
\end{equation}
For $N=2$ the primary is unique. For higher $N$ this is only an average weighted by the OPE: it does not assign the
same correction to every degenerate branch, and it does not determine branches invisible to the leading OPE.

%--------------------------------------------------------------------
\subsection{The logarithmic endpoint at \texorpdfstring{$s=1$}{s=1}}\label{app:dualC:logs}
%--------------------------------------------------------------------
At $\db=0$, with $b_0=1$, the published beta functions of the dual are those of
Eq.~\eqref{eq:dualA:leadingbeta}~\cite{Benedetti:2025nzp}.
Choose the marginally irrelevant branch $g>0$, $h<0$, and write its critical separatrix as
$h=-g/\sqrt2+a_2g^2+O(g^3)$. The cubic beta functions give $\beta_h-h'(g)\beta_g=-6a_2g^3+O(g^4)$, hence $a_2=0$.
With $L=\log(r/r_0)$ (not the input $L$),
\begin{equation}
 \frac{\dd g}{\dd L}=-2g^2-2g^3+O(g^4),\qquad
 g(L)=\frac1{2L}-\frac{\log L}{4L^2}+\frac{C}{L^2}+O\Bigl(\frac{\log^2L}{L^3}\Bigr).
\end{equation}
Let $M_\sigma^2=\lim_{r\to\infty}\langle\sigma(r)\sigma(0)\rangle$ on this trajectory. The wall gas gives
$G_\sigma(r)-1=-4g^2\log(r/\varepsilon)+\cdots$, because a kink pair changes the spin product only when one wall lies
between the insertions; so $\gamma_\sigma=2g^2+\cdots$, and
\begin{equation}
 \log\frac{G_\sigma(r)}{M_\sigma^2}=2\int_L^\infty\gamma_\sigma\bigl(g(u'),h(u')\bigr)\dd u'+\log F_\sigma\bigl(g(L),h(L)\bigr).
\end{equation}
The finite matching function satisfies $F_\sigma=1+O(g^2)$: at $g=0$ the spin is conserved for every $h$, and a nonzero
spin-flip correction needs two flips, so it cannot change the term in $\log L$. Inserting the running coupling and
exponentiating gives
\begin{equation}
 \frac{G_\sigma(r)}{M_\sigma^2}=1+\frac1L-\frac{\log L}{2L^2}+\frac{K_{\rm match}}{L^2}+O\Bigl(\frac{\log^2L}{L^3}\Bigr).
\end{equation}
The coefficient $-\frac12$ is fixed by the cubic flow; $K_{\rm match}$ depends on the matching convention and on the
definition of $r_0$. This extends the leading endpoint logarithm of Ref.~\cite{Benedetti:2025nzp} (its Appendix~E).
The difference $G_\sigma(r)-M_\sigma^2$ is best called the plateau-subtracted correlator: in the symmetric mixture
$\langle\sigma\rangle=0$ while the plateau is nonzero, so it need not equal the usual connected subtraction. Neither
the results of Appendices~\ref{app:dualC:shadow} and~\ref{app:dualC:logs} nor the thermal coefficients rely on a
stronger four-point statement that the compendium's earliest record had argued from a formal
principal-value cancellation in a current integral; such an identity alone does not determine composite-operator
counterterms or the fully renormalized four-point function.

%--------------------------------------------------------------------
\subsection{The polynomial \texorpdfstring{$B_{\rm old}$}{B-old} of the eighth order}\label{app:dualC:Bold}
%--------------------------------------------------------------------
The coefficient names are those of the derivative dictionary of Appendix~\ref{app:dualC:C8setup}. All terms are fixed by
the 25 inputs of $\CSR_3$ to $\CSR_7$; the ten new inputs enter only through the linear terms of
Eq.~\eqref{eq:C8assembly}. This is the unprojected expression: its rational projection is $\frac{1153}{6}$, and its
value is $-861.6635672354797792014971909923\ldots$.
\begin{align}
 B_{\rm old}={}&- \gamma - \frac{5 c^{3}}{1024} - \frac{3 \Omega}{16} - \frac{\Psi}{2} - \frac{Q}{4} - \frac{M}{16} - \frac{P}{16}\notag\\
 &- \frac{\alpha}{16} - \frac{a^{3}}{16} - \frac{K}{64} - \frac{N}{64} + \frac{\beta}{4} + \frac{\Phi}{8} + \frac{f^{2}}{64}\notag\\
 &+ \frac{3 L}{64} + \frac{A p}{2} + \frac{E a}{2} - E p - V a - \frac{9 J c}{64} - \frac{5 a c^{2}}{256}\notag\\
 &- \frac{3 A a}{16} - \frac{3 A c}{64} - \frac{3 a^{2} c}{64} - \frac{3 N c}{128} - \frac{G p}{2} - \frac{E b}{4} - \frac{L p}{4}\notag\\
 &- \frac{R a}{4} - \frac{R c}{4} - \frac{V c}{4} - \frac{b q}{4} - \frac{p q}{4} - \frac{G b}{8} - \frac{J a}{16}\notag\\
 &- \frac{L b}{16} - \frac{M c}{16} - \frac{b k}{16} - \frac{f q}{16} - \frac{k p}{16} - \frac{a f}{64} - \frac{f k}{64}\notag\\
 &+ \frac{a^{2} p}{4} + \frac{A b}{8} + \frac{E c}{8} + \frac{a q}{8} + \frac{L a}{16} + \frac{a^{2} b}{16} + \frac{c q}{16}\notag\\
 &+ \frac{a k}{32} + \frac{c^{2} p}{32} + \frac{c f}{256} + \frac{3 G a}{16} + \frac{3 G c}{32} + \frac{3 L c}{64} + \frac{3 b f}{64}\notag\\
 &+ \frac{5 b c^{2}}{128} + \frac{5 c k}{256} + \frac{5 c^{2} f}{512} + \frac{a c p}{8} + \frac{3 a b c}{64}.\label{eq:dualC:Bold}
\end{align}
The machine-readable version is the field \texttt{old\_part} of
\file{code/dual_d1/data/archived/C8_full_work/full_trace.json}. The separate full-Jacobian program recomputes this
expression from the beta functions; it does not read an expected value of $\CSR_8$ to obtain it.

\section{The two-coupling closure and the prescription}\label{app:clos}
%====================================================================
This appendix gives two formal arguments about the dual of Sec.~\ref{sec:dual}, both written by OpenAI GPT-6 astra
(via ChatGPT) and supplied by the author on 15 October 2026 (Secs.~\ref{sec:sources} and~\ref{sec:prov}), merged and edited to the
notation of this note: the \emph{closure argument}, on the bulk counterterms of the dual, and the \emph{prescription
argument}, on what the matched two-coupling prescription computes. Neither has been refereed. This note counts both as
formal, unrefereed arguments (Secs.~\ref{sec:dual} and~\ref{sec:d1caveats}). Where an argument states a proposition
or a theorem, it is given below as the statement that the argument establishes formally, coefficientwise in $g$, $h$
and $\db$; it is not a proved result of this note. Appendix~\ref{app:clos:checks} lists what the review checked.

Throughout, $\db=1-s$ and $x=\Ds=\db/2$. The dual, its operators $\Og$ and $\Oh$ and the charge $b=b_0-\sqrt2h$
are those of Eq.~\eqref{eq:dual}. The bare current $\chi_B=\frac{\mathrm i}{\sqrt2}\partial\phi$ is the $\chi$ of
Eq.~\eqref{eq:dual}; as in Appendix~\ref{app:dualA:gas}, the subscript distinguishes it from the renormalized current
primary, and neither is the generalized free field $\chi$ of Secs.~\ref{sec:d2} and~\ref{sec:d3}. The regulator, the
current resummation, the finite-tail matching, the vacuum normalization, and the finite-response Jacobian and
recurrence are given as the compendium states them in Appendices~\ref{app:dualA:regulator}, \ref{app:dualA:matching},
\ref{app:dualA:moments} and~\ref{app:dualA:recurrence}, whose notation is used here (in particular $\partial_\eta$ is
the derivative with respect to the endpoint source $\eta$ at $\eta=0$). Where the prescription argument re-derives
one of these objects, its derivation is given here and the formula is referred to there.

\subsection{The question, three claims and the common framework}\label{app:clos:question}
The thermal series of Sec.~\ref{sec:d1} rests on two assumptions (Sec.~\ref{sec:dual}): the renormalization of the dual
involves only the two couplings $(g,b)$, and the matched prescription extracts their $\beta$ functions. Three claims
must be kept apart.
\begin{enumerate}[label=(\arabic*),leftmargin=2em]
\item \emph{The full Wilsonian action at every finite cutoff contains only two interactions.} This is false in an
 ordinary fixed-field cutoff scheme (Appendix~\ref{app:clos:counterexample}).
\item \emph{Only two nontrivial independent bulk counterterm structures are needed in the formal perturbation theory
 about the endpoint.} This is the restricted statement of the closure argument (Appendix~\ref{app:clos:closure}).
\item \emph{The archived subtraction algorithm implements that renormalization, so that its computed coefficients are
 those of the dual.} This follows neither from operator counting nor from the exact period certificates of
 Sec.~\ref{sec:closed}. The prescription argument addresses the analytic recipe behind it
 (Appendix~\ref{app:clos:prescription}); whether the historical integrations follow that recipe is a separate question
 (Sec.~\ref{sec:indep} and Appendix~\ref{app:cannot}).
\end{enumerate}
The distinction between (1) and (2) is not semantic: finite irrelevant interactions can contribute to higher-order
marginal $\beta$ functions even when no independent irrelevant counterterm is required
(Appendix~\ref{app:clos:notmean}).

Claim (2) does not imply claim (3). The argument ``there are two counterterms, and two response equations determine
two numbers'' is insufficient, because for arbitrary responses the two numbers can absorb an endpoint or insertion
anomalous dimension. Schematically, the single-logarithm equation of two responses could be
\begin{equation}
 R_1=\mathcal J\beta+\Gamma_\partial R_0+\mathcal A_{\rm probe},
 \label{eq:clos:anomaly}
\end{equation}
where $R_0$ and $R_1$ are the finite and single-logarithm parts of the responses, $\mathcal J=D_cR_0$ is the
Jacobian of their finite parts (Appendix~\ref{app:clos:jacobian}; its tree-level value is the identity), $\Gamma_\partial$ an endpoint anomalous dimension and $\mathcal A_{\rm probe}$ a renormalization of the
probe. Inverting $\mathcal J$ does not show that the last two terms vanish. The prescription argument argues that both
are absent for the specific normalized, separated-probe prescription of Sec.~\ref{sec:dual}, and only then inverts
$\mathcal J$.

\paragraph{Common framework.} Both arguments work coefficientwise in a joint formal series in $g$, $h$ and $\db$ about
$g=h=\db=0$ (and, where needed, in the reference charge about $b_0=1$), at zero magnetic field. They use a fixed compact
target coordinate and a reference normalization of the Wick products. Distances in the Gaussian formulas are
dimensionless; a reference inverse length $\kappa$ is restored by replacing $r$ by $\kappa r$. An infrared regulator,
or a smooth compactly supported switching function of the interaction, is kept throughout. The subject is ultraviolet
collisions, not an infrared limit and not the renormalization of arbitrary composite insertions.

The published analysis identifies two symmetry-invariant marginal perturbations at $\db=0$ but distinguishes this
explicitly from an all-orders renormalizability result for its non-scaling perturbation at $\db>0$
\cite[Secs.~3.2.3, 4 and~6]{Benedetti:2025nzp}; the compendium likewise qualifies its higher coefficients by the
two-coupling and matching assumptions.

\subsection{The closure statement}\label{app:clos:closure}
The closure argument, a formal and unrefereed argument written by OpenAI GPT-6 astra (see the head of this
appendix; its scope is in Appendix~\ref{app:clos:scope} and its checks in Appendix~\ref{app:clos:checks}), establishes
formally the following statement, which it states as a proposition.
\begin{quote}
\emph{Formal bulk counterterm class.} For the matrix-valued Gaussian model of Eq.~\eqref{eq:dual}, impose its diagonal
shift symmetry, internal reflection and parity (Appendix~\ref{app:clos:reference}). Work coefficientwise in $g$, $h$
and $\db$ about the marginal endpoint, using the Wick products of the deformed covariance
(Appendix~\ref{app:clos:power}), with a fixed infrared cutoff. Then local symmetry-preserving subtractions can be
constructed recursively such that every required bulk action counterterm for constant couplings lies in
\begin{equation}
 \mathrm{span}\{\mathbf 1,\Og,\Oh\},
 \label{eq:clos:span}
\end{equation}
modulo total derivatives. Only two nontrivial independent bulk coupling counterterms are therefore required.
Field-insertion and endpoint counterterms are not included in this count.
\end{quote}
The argument has three parts: the exact selection rules and the classification of local operators
(Appendix~\ref{app:clos:classification}), power counting (Appendix~\ref{app:clos:power}), and the recursive subtraction
of all partial coincidences (Appendix~\ref{app:clos:forest}). Since $h$ and $b=b_0-\sqrt2h$ are two coordinates on one
coupling direction (Appendix~\ref{app:clos:current}), the statement supplies a counterterm-class argument for the two
essential coordinates $(g,b)$. The closure argument sums it up as follows: in the zero-field, symmetry-preserving formal
expansion about the marginal endpoint, two nontrivial independent bulk ultraviolet counterterm directions suffice, in
addition to vacuum and external-insertion or endpoint renormalization; a finite-cutoff Wilsonian action nevertheless
contains irrelevant terms, and their finite contributions must be incorporated in any correct high-order calculation.

The statement uses the standard local extension construction explicitly. It is not a claim that any procedure that
cancels two chosen scalar divergences implements that construction; that is the subject of the prescription argument.

\subsection{The prescription statement}\label{app:clos:prescription}
The prescription argument, a formal and unrefereed argument written by OpenAI GPT-6 astra (see the head of this
appendix; its scope is in Appendix~\ref{app:clos:scope} and its checks in Appendix~\ref{app:clos:checks}), establishes
formally the following statement, which it states as a theorem. Its hypotheses are:
\begin{enumerate}[label=(\alph*),leftmargin=2em]
\item the reference Gaussian and matrix operator algebra of Appendix~\ref{app:clos:reference}, at zero magnetic field;
\item a coefficientwise expansion in $g$, $h$ and $\db$ and, when needed, in the reference charge about $b_0=1$;
\item the diagonal $O(2)$ symmetry (the $U(1)$ of Sec.~\ref{sec:dual} and the internal reflection) is preserved;
\item the interaction is restricted to a finite interval, with a hard core on successive kink positions only;
\item current products are defined by exact Gaussian current resummation, retaining the specified Wick-ordering and
 endpoint factors;
\item every response is normalized by the same vacuum expectation of that interval operator;
\item the reference-theory probe of each response lies outside the interaction support, and any projection to an
 infinitely distant probe is taken only afterwards.
\end{enumerate}
\begin{quote}
\emph{Formal normalized-interval prescription.} Within local, scaling-degree-preserving position-space
renormalization, the bulk counterterms are $\mathbf 1$, $\Og$ and $\Oh$, and all ultraviolet counterterms localized at
either endpoint of the interval are scalar multiples of $\mathbf 1$. Consequently the normalized responses require no
additional endpoint or probe counterterm. If their tree-level two-by-two response matrix is invertible, their
single-logarithm equation determines the local logarithmic $\beta$ function. Relative to any other local normalization
in the same class, the result is related by a regular finite change of the two nontrivial bulk couplings, so the
fixed-point stability eigenvalues are identical.
\end{quote}
As defined in Sec.~\ref{sec:dual}, the matched two-coupling prescription has the form of hypotheses (d) to (f): (d)
is its item (i), (e) its items (ii) and (iii), and (f) its item (iv). The prescription argument identifies its two
responses with the far-field projections of separated reference probes, as hypothesis (g) requires
(Appendix~\ref{app:clos:endpoint}); that projection is stated rather than shown in detail
(Appendix~\ref{app:clos:checks}). Identical
eigenvalues mean identical invariants $T$ and $D$ (Sec.~\ref{sec:invariants}), and so an identical series
$\De=1+\sum_n\CSR_nx^{n/2}$.

The statement concerns the prescription as an analytic recipe, not whether a particular computer output implemented
every part of it without error. The complete normalized moments, the first endpoint-source derivatives and the exact
normal-ordering map are essential hypotheses, not optional approximations; the argument compares them with the
compendium's formulas (Appendices~\ref{app:clos:current} and~\ref{app:clos:jacobian}). It uses standard local
distribution-extension and finite-renormalization results~\cite{Brunetti:2000,Brunetti:2009}; the model-specific
steps are the collision bounds, the endpoint classification, the spectator factorization and the audit of the
current and normal-ordering factors. These steps do not use Lorentzian locality of the nonlocal Gaussian action: they
concern the explicit Euclidean Wick kernels, which are smooth away from coincidence diagonals.

The order of the steps is what makes the argument non-circular: the local renormalization is built for arbitrary
spectators first (Appendix~\ref{app:clos:forest}), the endpoint cancellation and the absence of probe renormalization
follow (Appendix~\ref{app:clos:endpoint}), and the invertibility of the response matrix is used only at the end, for
uniqueness (Appendix~\ref{app:clos:uniqueness}).

\subsection{The reference theory, its symmetries and the regulated responses}\label{app:clos:reference}
\paragraph{Reference theory.} At $\db=0$ the field normalization and compactification are chosen so that
\begin{equation}
 \langle\phi(X)\phi(0)\rangle_0=-2\log(\kappa|X|),\qquad
 V_\pm=\kappa\,{:}e^{\pm\mathrm i b_0\phi}{:},\qquad
 \chi_B=\frac{\mathrm i}{\sqrt2}\partial\phi,\qquad b_0=1.
 \label{eq:clos:reference}
\end{equation}
Thus $V_\pm$ and $\chi_B$ have reference dimension one. The deformed covariance, before any expansion in $\db$, is
\begin{equation}
 C_{\db}(r)=\frac{2\bigl(\kappa^{-\db}-|r|^{\db}\bigr)}{\db}=-2f_{\db}(|r|),\qquad
 f_{\db}(r)=\frac{r^{\db}-\kappa^{-\db}}{\db},
 \label{eq:clos:cov}
\end{equation}
which at $\kappa=1$ is the covariance of Sec.~\ref{sec:dual}. The zero mode is removed, and complete charged
observables carry the appropriate neutralizing probe. With $\hat\sigma_\pm=(\hat\sigma_1\pm\mathrm i\hat\sigma_2)/2$,
the interactions are $\Og=\hat\sigma_+V_++\hat\sigma_-V_-$ and $\Oh=\hat\sigma_3\chi_B$, and the evolution over an
interval $I$ is
\begin{equation}
 \mathcal D_I(c)=\mathrm P\exp\int_I\dd X\,\bigl[g\Og+h\Oh\bigr],\qquad c=(g,h).
 \label{eq:clos:interaction}
\end{equation}

\paragraph{Exact symmetries.} With $A$ a matrix of the qubit, the global symmetry includes
\begin{equation}
 \phi\mapsto\phi+\alpha/b_0,\quad A\mapsto e^{-\mathrm i\alpha\hat\sigma_3/2}Ae^{\mathrm i\alpha\hat\sigma_3/2};
 \qquad\qquad
 \phi\mapsto-\phi,\quad A\mapsto\hat\sigma_1A\hat\sigma_1.
 \label{eq:clos:symmetry}
\end{equation}
The first is the $U(1)$ of Sec.~\ref{sec:dual}, the second the internal reflection; together they form the diagonal
$O(2)$. Parity combines $X\mapsto-X$, $\phi(X)\mapsto-\phi(-X)$, the reversal of the path ordering and matrix
transposition. These transformations preserve the regulated matrix expansion
\cite[Eqs.~(3.31)--(3.35)]{Benedetti:2025nzp}. At zero magnetic field the reflection is not broken. The internal
reflection, not spatial parity alone, is what excludes a local $\hat\sigma_3$ counterterm at a single endpoint
(Appendix~\ref{app:clos:power}).

\paragraph{The regulated responses.} The interaction is restricted to $I=[0,L]$, and the distinct-kink integrations
run over
\begin{equation}
 0<X_1<\cdots<X_n<L,\qquad X_{i+1}-X_i\ge\varepsilon,
 \label{eq:clos:cutoff}
\end{equation}
with no hard core between a kink and an endpoint; $\ell=\log\varepsilon$ (or $\log(\varepsilon\mu)$ at a reference
scale $\mu$). The compendium works at $L=\kappa^{-1}=1$ (Eq.~\eqref{eq:dualA:domain}). At a finite regulator it
is useful first to keep a smoothed Gaussian covariance, for which the matrix change-of-frame identity of
Appendix~\ref{app:clos:current} is exact, and then to pass to normal-ordered products. Changes of the smoothing at
separations of order $\varepsilon$ are local regulator changes, handled by the forest construction of
Appendix~\ref{app:clos:forest}. The coincident Wick convention used in the archive is $C(0)\to C_{\db}(\varepsilon)$.

For a reference probe $\mathcal X$ separated from $I$, the normalized response is
\begin{equation}
 \mathcal R_{\mathcal X}(\varepsilon;c)
 =\frac{\frac12\operatorname{tr}\langle\mathcal X\,\mathcal D_{I,\varepsilon}(c)\rangle_0}
 {\frac12\operatorname{tr}\langle\mathcal D_{I,\varepsilon}(c)\rangle_0},
 \label{eq:clos:ratio}
\end{equation}
with the same interval and regulator in numerator and denominator. The probes used for the two coordinates are the
reference vertex-and-matrix and current-and-matrix probes, normalized by their free two-point functions. They are
\emph{not} the fully interacting, renormalized primaries inserted inside the interacting interval.

\subsection{Matrix algebra and the local operators below threshold}\label{app:clos:classification}
\paragraph{Matrix algebra excludes higher-charge vertices.} The exact relations
\begin{equation}
 \hat\sigma_+^2=\hat\sigma_-^2=0,\qquad
 \hat\sigma_+\hat\sigma_-=\frac{\mathbf 1+\hat\sigma_3}{2},\qquad
 \hat\sigma_-\hat\sigma_+=\frac{\mathbf 1-\hat\sigma_3}{2}
\end{equation}
force the sequence of nonzero flips to alternate. Inserting diagonal $\hat\sigma_3$ factors, as the current
interaction does, does not change this, since $\hat\sigma_\pm\hat\sigma_3^k\hat\sigma_\pm=0$. Consequently every
nonzero contiguous matrix word $M$ has vertex charge $Q\in\{0,+1,-1\}$ and satisfies $[\hat\sigma_3,M]=2QM$: an even
number of flips gives $Q=0$ and a diagonal matrix, an odd number gives $Q=\pm1$ and a matrix proportional to
$\hat\sigma_\pm$. This holds for words of every length, not only for a two-vertex product. In particular, no ordering
of the interactions generates a bulk counterterm of charge two.

Equivalently, the complete shift-invariant local Wick algebra generated by these fields consists of
\begin{equation}
 P_{\mathbf 1}(\partial\phi,\partial^2\phi,\ldots)\,\mathbf 1,\qquad
 P_3(\partial\phi,\partial^2\phi,\ldots)\,\hat\sigma_3,\qquad
 P_+(\partial\phi,\ldots)\,\hat\sigma_+V_+,\qquad
 P_-(\partial\phi,\ldots)\,\hat\sigma_-V_-,
 \label{eq:clos:algebra}
\end{equation}
with $P_{\mathbf 1},P_3,P_\pm$ polynomials in the derivatives of $\phi$. This is the perturbative local algebra:
disorder or boundary-condition-changing fields outside it are not adjoined to the theory. An undifferentiated
polynomial in $\phi$ is not an extra shift-invariant local operator ($\phi$ itself is the unobservable zero-mode
field); expanding a vertex in powers of $\phi$ and giving those powers independent couplings would break the exact
diagonal shift symmetry.

\paragraph{Classification.} The reference dimension of a derivative monomial is
\begin{equation}
 d\Bigl(\prod_{j\ge1}(\partial^j\phi)^{n_j}\Bigr)=\sum_{j\ge1}jn_j,
\end{equation}
and a charged vertex sector adds one. These are dimensions at $\db=0$; they do not claim that the deformed vertex is a
scaling operator at $\db\ne0$. Through dimension one the complete list of candidates is that of
Table~\ref{tab:clos:operators}.
\begin{table}[htbp]\centering\small
\renewcommand{\arraystretch}{1.15}
\begin{tabular}{@{}lll@{}}
\toprule
candidate & reference dimension & internal reflection\\
\midrule
$\mathbf 1$ & $0$ & even\\
$\hat\sigma_3$ & $0$ & odd\\
$\partial\phi$ & $1$ & odd\\
$\hat\sigma_3\partial\phi$ & $1$ & even\\
$\hat\sigma_+V_++\hat\sigma_-V_-$ & $1$ & even\\
$\hat\sigma_+V_+-\hat\sigma_-V_-$ & $1$ & odd\\
\bottomrule
\end{tabular}
\caption{The $U(1)$-invariant local operators of the algebra~\eqref{eq:clos:algebra} through reference dimension one,
with their behaviour under the internal reflection of Eq.~\eqref{eq:clos:symmetry}.}
\label{tab:clos:operators}
\end{table}
The two even entries of dimension one also have the required parity. The only nontrivial symmetry-invariant local
action densities at dimension one are therefore $\Og$ and $\Oh=\frac{\mathrm i}{\sqrt2}\hat\sigma_3\partial\phi$. The
identity changes the vacuum normalization. A source for $\hat\sigma_3$ (a magnetic field) would be an additional
relevant coupling; it is excluded by the zero-field symmetry of the problem considered.

At dimension zero the invariant sector is $\mathbb R\,\mathbf 1$: a neutral constant matrix is diagonal, and the
internal reflection equates its two entries ($\hat\sigma_3$ is odd, and the bare $\hat\sigma_\pm$ are charged). This
is the input of the endpoint classification of Appendix~\ref{app:clos:power}.

\subsection{The product \texorpdfstring{$\Og\Og$}{OgOg}: no literal two-operator closure}\label{app:clos:counterexample}
At a symmetric pair of points with $r>0$, Wick's theorem with Eq.~\eqref{eq:clos:reference} gives
\begin{align}
 \Og(X+\tfrac r2)\,\Og(X-\tfrac r2)
 &=r^{-2}\Bigl[\mathbf 1+\mathrm i r\hat\sigma_3\partial\phi-\frac{r^2}{2}{:}(\partial\phi)^2{:}+O(r^3)\Bigr]
 \notag\\
 &=\frac{\mathbf 1}{r^2}+\frac{\sqrt2}{r}\Oh(X)+{:}\chi_B^2{:}(X)+O(r).
 \label{eq:clos:ope}
\end{align}
Only $\hat\sigma_+\hat\sigma_-$ and $\hat\sigma_-\hat\sigma_+$ contribute; they multiply
$r^{-2}{:}e^{\pm\mathrm i(\phi(X+r/2)-\phi(X-r/2))}{:}$, and $\phi(X+\frac r2)-\phi(X-\frac r2)=r\partial\phi+O(r^3)$.
The midpoint convention removes a spurious first derivative of the expansion point. The coefficient of ${:}\chi_B^2{:}$
is one, not zero. Integrating a short-distance shell $\varepsilon<r<\varepsilon e^{\dd\ell}$ in a standard real-space
shell prescription therefore produces, besides the two lower-dimensional terms, an irrelevant term proportional to
\begin{equation}
 g^2\varepsilon\,\dd\ell\int\dd X\,{:}\chi_B^2{:}(X).
\end{equation}
Its detailed Wilsonian coefficient depends on the scheme; its appearance as an allowed term of the operator product
expansion does not.

The current-squared channel is finite as $\varepsilon\to0$ in this two-insertion integral (its radial integrand is
$r^0$). This is why its generation does not refute the closure of the ultraviolet counterterm class
(Appendix~\ref{app:clos:closure}). It does refute the unqualified statement that no other operator coefficient is
generated at finite cutoff, that is, claim (1) of Appendix~\ref{app:clos:question}. Some such terms can be removed by
redundant transformations in particular schemes, but that is not the same as their absence from a fixed-field
Wilsonian action. How these finite terms enter the marginal $\beta$ functions is the subject of
Appendix~\ref{app:clos:notmean}.

\subsection{Power counting in the bulk and at the endpoints}\label{app:clos:power}
\paragraph{Bulk clusters.} Consider $N\ge2$ insertions of reference dimension one collapsing toward one interior
point. Removing the centre coordinate leaves $N-1$ relative coordinates, written as a scale $\rho$ and shape
variables, with radial measure $\rho^{N-2}\dd\rho$. If the product produces a local field $\mathcal A$ of reference
dimension $d_{\mathcal A}$, its coefficient scales as $\rho^{d_{\mathcal A}-N}$, up to logarithms introduced by
subtractions, and the radial integral has the form
\begin{equation}
 \int_0^{\rho_0}\dd\rho\,\rho^{N-2}\rho^{d_{\mathcal A}-N}(\log\rho)^M
 =\int_0^{\rho_0}\dd\rho\,\rho^{d_{\mathcal A}-2}(\log\rho)^M .
 \label{eq:clos:bulk}
\end{equation}
The threshold does not depend on $N$:
\begin{equation}
 d_{\mathcal A}=0:\ \text{power divergence};\qquad
 d_{\mathcal A}=1:\ \text{logarithmic divergence};\qquad
 d_{\mathcal A}\ge2:\ \text{integrable},
 \label{eq:clos:threshold}
\end{equation}
and for any finite $M$ the logarithms do not change it. By Table~\ref{tab:clos:operators}, symmetry leaves
$\mathbf 1$, $\Og$ and $\Oh$ below the threshold.

\paragraph{Alternating vertices.} The same follows directly from Wick's theorem. For $n$ vertex insertions with charges
$q_i=\pm1$ and total charge $Q$, at $b_0=\kappa=1$,
\begin{equation}
 \prod_iV_{q_i}(X_i)=\prod_{i<j}|X_i-X_j|^{2q_iq_j}\,{:}\exp\Bigl(\mathrm i\sum_iq_i\phi(X_i)\Bigr){:},
 \qquad 2\sum_{i<j}q_iq_j=Q^2-n .
\end{equation}
For a cluster $X_i=X+\rho r_i$ the normal product is organized exactly as
\begin{equation}
 \prod_iV_{q_i}(X_i)=\mathcal W(\rho,\boldsymbol r)\,
 {:}\exp\Bigl[\mathrm i b_0Q\phi(X)+\mathrm i b_0\sum_{k\ge1}\frac{\rho^km_k}{k!}\partial^k\phi(X)\Bigr]{:},\qquad
 m_k=\sum_iq_ir_i^k,
 \label{eq:clos:wick}
\end{equation}
where the prefactor $\mathcal W$ has radial degree $Q^2-n$ at the reference point. A descendant of derivative level $k$
adds a factor $\rho^k$, so the radial power is
\begin{equation}
 \rho^{n-2}\rho^{Q^2-n+k}=\rho^{Q^2+k-2}.
\end{equation}
For even (neutral) clusters, $Q=0$, only $k\le1$ diverges; for odd clusters, $Q=\pm1$, only $k=0$; all higher
derivatives are integrable. The operator content follows without reference to any probe: summing the two orientations
of an even cluster gives the identity at level zero and a multiple of $\hat\sigma_3\partial\phi$ at level one, while an
odd cluster gives $\Og$ at level zero and its first-derivative descendants already have dimension two. Current
insertions obey the same bound, because $\chi_B$ has reference dimension one and its contractions use one or two
derivatives of the logarithmic covariance. Simultaneous collapse is not the whole argument, since the shape variables
still have partial-collision singularities; these are treated in Appendix~\ref{app:clos:forest}.

\paragraph{Endpoint clusters: the degree differs by one.} At a fixed endpoint of $I$ the centre of the cluster also
approaches the endpoint. There are then $N$, not $N-1$, scaled coordinates, and
\begin{equation}
 \int_0\dd\rho\,\rho^{N-1}\rho^{d_{\mathcal A}-N}(\log\rho)^M=\int_0\dd\rho\,\rho^{d_{\mathcal A}-1}(\log\rho)^M .
 \label{eq:clos:edge}
\end{equation}
An endpoint divergence therefore requires $d_{\mathcal A}\le0$. The invariant dimension-zero sector is
$\mathbb R\,\mathbf 1$ (Appendix~\ref{app:clos:classification}), hence
\begin{equation}
 \text{every primitive ultraviolet counterterm at an endpoint is a scalar multiple of }\mathbf 1 .
 \label{eq:clos:scalaredge}
\end{equation}
This holds separately at the left and at the right endpoint; their scalar coefficients need not be assumed equal. The
classification is made in the original invariant $(g,h)$ frame. In the frame in which the current is resummed, the
endpoint vertices are dressed and can carry non-integer auxiliary charges; they must be transformed together with the
bulk charge and the Wick factors, and treating them as unrelated new endpoint interactions would not describe the same
theory. The estimates~\eqref{eq:clos:bulk} and~\eqref{eq:clos:edge} apply after proper collision subgraphs have been
subtracted; they must not be applied separately to unprepared products of a power-divergent counterterm and an omitted
remainder.

\paragraph{The long-range deformation does not enlarge the threshold.} At fixed reference scale ($\kappa=1$) the
deformation of the covariance~\eqref{eq:clos:cov} is
\begin{equation}
 f_{\db}(r)=\frac{r^{\db}-1}{\db}=\sum_{j\ge0}\frac{\db^{\,j}\log^{j+1}r}{(j+1)!}
 =\log r+\frac{\db}{2}\log^2r+\frac{\db^2}{6}\log^3r+\cdots.
 \label{eq:clos:deform}
\end{equation}
At each fixed order in $\db$, a Gaussian vertex coefficient is the power-law coefficient at $\db=0$ times a finite
polynomial in logarithms. Spatial derivatives of $\log^jr$ have the same inverse powers of the distance as the same
derivatives of $\log r$, multiplied by logarithms and constants, so they do not raise the superficial degree of
divergence at that order. This is visible before the outgoing operator is expanded: the Wick normal product has the
exact form
\begin{equation}
 {:}\exp\Bigl(\mathrm i\sum_iq_i\phi(X_i)\Bigr){:}_{\db}
 ={:}\exp\Bigl(\mathrm iQ\phi(X)+\mathrm i\sum_{k\ge1}\frac{\rho^km_k}{k!}\partial^k\phi(X)\Bigr){:}_{\db},
\end{equation}
with the moments $m_k$ of Eq.~\eqref{eq:clos:wick}. The deformation changes the contraction coefficients and the
normal-ordering convention, but not the list of possible outgoing charges or derivative levels. The normal ordering is
varied consistently with the contractions, rather than differentiating the contraction and forgetting the definition of
the local field. Consequently, for every fixed coefficient of $g^{n_1}h^{n_2}\db^{\,n_3}$, the coefficient
$t_{\mathcal A}$ of an outgoing field obeys the reference scaling-degree bound
\begin{equation}
 \operatorname{sd}t_{\mathcal A}\le N-d_{\mathcal A},
 \label{eq:clos:sd}
\end{equation}
with possibly higher powers of logarithms, so that the degree of divergence is
$\omega_{\mathcal A}=\operatorname{sd}t_{\mathcal A}-(N-1)\le1-d_{\mathcal A}$ in the bulk and
$\omega_{\mathcal A}\le-d_{\mathcal A}$ at an endpoint. Derivatives in the charge, that is an expansion in the resummed
charge $b^2-1$, likewise insert only logarithms, provided the expansion is taken order by order about the reference
value and is not a statement uniform in an arbitrary fixed $b$.

\paragraph{Order of limits.} The bound~\eqref{eq:clos:sd} concerns a formal coefficient before all powers of $\db$ are
summed. It does not justify interchanging that infinite sum with $\varepsilon\to0$. At fixed $\db>0$, $r^{\db}$ tends to
zero as $r\to0$, while each fixed Taylor coefficient in Eq.~\eqref{eq:clos:deform} grows logarithmically; these are
different ultraviolet organizations. The existence of coefficientwise counterterms establishes neither convergence,
nor Borel summability, nor a unique identification of the formal series with a theory at fixed $\db$; neither argument
asserts such an interchange.

\subsection{The local forest construction}\label{app:clos:forest}
This construction supplies the renormalized family \emph{before} any response condition is used, which is what keeps
the prescription argument from being circular. It is the recursive position-space (Epstein--Glaser) construction,
with the matrix multiplication kept in its original order.

\paragraph{The extension principle.} A distribution in $m$ relative variables, defined away from a coincidence point and
of finite scaling degree $\operatorname{sd}$, has an extension of the same scaling degree. If $\operatorname{sd}<m$ the
extension is unique; otherwise its local ambiguity is a finite sum of derivatives, of order at most
$\lfloor\operatorname{sd}-m\rfloor$, of the Dirac distribution at the coincidence point. A finite logarithmic scaling
anomaly does not change that bound~\cite{Brunetti:2000,Dutsch:2014}; Euclidean recursive versions are also
available~\cite{Keller:2009}. The arguments apply these results to the explicit finite-order Gaussian and
operator-product coefficient distributions of the dual; the cited works do not establish the nonperturbative impurity
theory.

\paragraph{Structure at separated points.} At separated points every fixed perturbative coefficient is a finite sum of
products of rational powers of differences and finite powers of logarithms, multiplying multilocal normal-ordered
exponentials and derivative polynomials with Pauli matrices. Only coincidence loci can be singular, and all
contractions with spectators a fixed distance away are smooth functions of the cluster coordinates.

\paragraph{The expansion is operator valued and made before projection.} The collision of a cluster is expanded with
Eq.~\eqref{eq:clos:wick} into the local operators of Eq.~\eqref{eq:clos:algebra}, as an operator statement, with no
reference to the two probes used in the computation. To see that its coefficients do not depend on distant spectators,
take external charges $Q_\alpha$ at points $y_\alpha$ a fixed positive distance from the cluster. In a chamber where the
signed distances are fixed, their Wick factors expand as
\begin{equation}
 \prod_{i,\alpha}|y_\alpha-X-\rho r_i|^{2q_iQ_\alpha}
 =\prod_\alpha|y_\alpha-X|^{2QQ_\alpha}
 \exp\Bigl[-2\sum_{k\ge1}\frac{\rho^km_k}{k}\sum_\alpha\frac{Q_\alpha}{(y_\alpha-X)^k}\Bigr].
 \label{eq:clos:spectator}
\end{equation}
The level-zero and level-one terms are exactly the Wick correlators of the local operators of
Eq.~\eqref{eq:clos:wick} with those spectators; derivative spectators follow by differentiation. Spectator dependence
therefore multiplies universal local collision coefficients and does not enter their definition.

\paragraph{A finite Taylor jet suffices.} The expansion into local output fields is in general infinite, but only a
finite Taylor jet is needed: neutral clusters need the constant and first-derivative jets, and charged clusters the
leading vertex jet. The next Taylor remainder gains at least two powers of $\rho$ in the neutral channel and one in the
charged channel. Tested against separated smooth spectators, it obeys the improved bound in the remaining shape
sectors, so the argument does not interchange an unproved infinite operator-product sum with a divergent integral.

\paragraph{Induction over clusters.} Suppose that the products of fewer than $N$ insertions have been extended at all
their coincidence loci. The $N$-point configuration space is subdivided by relative cluster scales. In one dimension a
collapsing subset of an ordered chamber is consecutive, and compatible scale hierarchies consist of nested or disjoint
clusters, that is, forests. On sectors where some clusters have already collapsed, the lower-order local subtractions
are inserted. This defines the renormalized product away from the total diagonal. Overlapping candidate subclusters
occur in different hierarchy sectors, so they are never subtracted twice in one forest, nor treated as independent
simultaneous subtractions.

For the remaining total collision of a cluster $\gamma$, let $t_{\gamma,\mathcal A}$ be the prepared coefficient
distribution of its output field $\mathcal A$, in $m$ relative coordinates. By Eq.~\eqref{eq:clos:sd} its excess of
scaling degree obeys
\begin{equation}
 \omega_{\gamma,\mathcal A}\le
 \begin{cases}1-d_{\mathcal A}&\text{(bulk)},\\ -d_{\mathcal A}&\text{(fixed endpoint)}.\end{cases}
\end{equation}
Choose a smooth compactly supported function $w$ equal to one near zero, at a reference scale $\mu$. For
$\omega\ge0$ a test function $f$ is replaced by
\begin{equation}
 (\mathcal T_\omega f)(\xi)=f(\xi)-w(\xi)\sum_{|\alpha|\le\omega}\frac{\xi^\alpha}{\alpha!}\,\partial^\alpha f(0).
 \label{eq:clos:taylor}
\end{equation}
The extra vanishing order makes $t_{\gamma,\mathcal A}(\mathcal T_\omega f)$ integrable; for $\omega<0$ nothing is
subtracted. Finite local normalization constants are then added to specify the chosen extension. The subtractions are
local by Eq.~\eqref{eq:clos:taylor}, and their coefficients are independent of spectators by
Eq.~\eqref{eq:clos:spectator}. In channels with $d_{\mathcal A}\ge2$ the scaling-degree-preserving extension is unique
once the subdivergences have been removed, so no new irrelevant bulk counterterm is needed. Only fields of dimension at
most one can carry a counterterm ambiguity, and Table~\ref{tab:clos:operators} leaves $\mathbf 1$, $\Og$ and $\Oh$.
Derivatives of the Dirac distribution in a dimension-zero channel act on switching functions or produce total
derivatives; they are not extra bulk interactions with constant couplings. Counterterms of the identity give the
vacuum normalization.

\paragraph{Stability under counterterm insertions.} An inserted counterterm is an effective vertex of dimension at most
one. With $N'$ effective vertices of dimensions $d_i\le1$, the new degree of divergence is bounded by
\begin{equation}
 \omega_{\mathcal A}\le\sum_{i=1}^{N'}d_i-d_{\mathcal A}-(N'-1)\le1-d_{\mathcal A}
\end{equation}
in the bulk, with one further unit of improvement at an endpoint. All proper collisions therefore remain subtractable
within the same class. Finite powers of logarithms from $\db$ derivatives and from earlier subtractions do not change
the power bound. This completes the induction at every finite order, rather than assuming only simultaneous collapse.

\paragraph{Symmetric extensions.} Symmetric extensions can be chosen: the original kernels are invariant, and averaging
an extension over the compact internal symmetry and the finite group of reflection and parity preserves its scaling
bound and its factorization, once the lower extensions have been chosen invariantly. No anomalous symmetry breaking is
forced by this finite-cutoff matrix and Gaussian construction.

\paragraph{The hard-core regulator.} For a coefficient regulated by the hard core, the difference between its integral
and the Taylor-subtracted one is a finite sum of local Taylor jets,
\begin{equation}
 \langle t_{\gamma,\mathcal A;\varepsilon},f\rangle-\langle t_{\gamma,\mathcal A;\varepsilon},\mathcal T_\omega f\rangle
 =\sum_{|\alpha|\le\omega}c_{\gamma,\mathcal A,\alpha}(\varepsilon,\mu)\,\partial^\alpha f(0),
 \label{eq:clos:cutoffdiff}
\end{equation}
where the coefficients are regulated moments of $t_{\gamma,\mathcal A;\varepsilon}\,w(\xi)\,\xi^\alpha/\alpha!$. Their
divergent and finite parts define the cutoff counterterms recursively. Replacing the cutoff for an already collapsed
effective vertex changes only a shrinking collision region, and the next step of the same Taylor construction absorbs
its local difference. Thus no exact finite-cutoff factorization of the exclusion region of a collapsed kink is assumed.

\paragraph{Identity terms and sharp endpoints.} The identity and its first test-function derivatives are vacuum and
switching-function counterterms. For a constant coupling on the interval they give a bulk vacuum energy and scalar
endpoint factors; the endpoint Taylor ambiguities have order zero and multiply only $\mathbf 1$. The construction can
first be formulated with a smooth switching function and then taken to the sharp-edge limit. The renormalization is
local in the coupling profile, and at finite smoothing width a profile is not assumed equal to its square. In the sharp
limit all nontrivial bulk terms have the common plateau coupling, while the additional ultraviolet terms at either
edge obey Eq.~\eqref{eq:clos:edge} and are scalar. Finite higher-dimensional regulator remainders are kept in the
prepared integrals; dropping them by hand would not be this construction.

\paragraph{No new nonlocal Gaussian counterterm.} A primitive ultraviolet counterterm is supported on a collision
diagonal. In momentum space it is therefore polynomial in the external momenta, not a new nonanalytic kernel such as
$|p|^{1+\db}$ or one of its logarithmic Taylor coefficients. The only local quadratic derivative term that could occur
has reference dimension at least two and is excluded from the required counterterms by the same power counting. This
concerns divergent counterterms; finite self-energy effects and anomalous dimensions are not excluded.

\subsection{What the closure statement does not mean}\label{app:clos:notmean}
\paragraph{Irrelevant operators feed the marginal \texorpdfstring{$\beta$}{beta} functions.} A simple formal example has
a marginal coupling $g$ and an irrelevant dimensionless coupling $\vartheta$, with (in the $\mu\partial_\mu$
convention)
\begin{equation}
 \beta_g=e_1g^3+e_2g\vartheta+\cdots,\qquad \beta_\vartheta=\vartheta-e_3g^2+\cdots .
\end{equation}
An invariant perturbative trajectory has $\vartheta=e_3g^2+O(g^4)$, so that
\begin{equation}
 \beta_g^{\rm reduced}=(e_1+e_2e_3)g^3+O(g^5).
\end{equation}
Setting $\vartheta=0$ blindly loses a cubic contribution, although $\vartheta$ need not be an independent continuum
parameter. In the dual such finite effects, among them the ${:}\chi_B^2{:}$ term of Eq.~\eqref{eq:clos:ope}, occur inside
the fully integrated multi-kink kernels and their finite subtractions. The absence of an independent counterterm does
not license truncating every intermediate operator product to its singular terms when finite higher-order coefficients
are computed. The coefficient of ${:}\chi_B^2{:}$ at level two is locally integrable in the expansion about the
endpoint; all such finite terms are retained in the full cutoff integrals and finite responses of the prescription.
The prescription argument would fail for an algorithm that replaced every cluster by its two marginal channels and
discarded the integrable remainder before the higher integrations.

\paragraph{The stability matrix is tangent to the renormalized family.} Suppose a full Wilsonian description has
coordinates $(q^a,\vartheta^i)$, with $q$ the two retained couplings, and a two-parameter invariant surface
$\vartheta^i=\Theta^i(q)$ with slope $\Theta'=\partial\Theta/\partial q$. At its fixed point write the full
linearization in blocks as
\begin{equation}
 M_{\rm full}=\begin{pmatrix}M_{qq}&M_{q\vartheta}\\ M_{\vartheta q}&M_{\vartheta\vartheta}\end{pmatrix}.
\end{equation}
The reduced matrix and the invariance condition are
\begin{equation}
 M_{\rm red}=M_{qq}+M_{q\vartheta}\Theta',\qquad
 M_{\rm full}\begin{pmatrix}\mathbf 1\\ \Theta'\end{pmatrix}=\begin{pmatrix}\mathbf 1\\ \Theta'\end{pmatrix}M_{\rm red}.
\end{equation}
The two eigenvalues of $M_{\rm red}$ are genuine tangent eigenvalues of the full flow; the block $M_{qq}$ alone is in
general not the correct matrix. This identity explains what a correctly renormalized two-coupling description
computes; it does not construct an infinite-dimensional Wilsonian invariant manifold.

\paragraph{The statement is formal, not uniform in a finite charge displacement.} Even at $\db=0$, resumming a charge
displacement $b^2=1+z$ before expanding can change the ultraviolet power counting. For an even cluster of $n$ flips the
Wick prefactor scales as $\rho^{(1+z)(Q^2-n)}=\rho^{-n(1+z)}$, so the level-two channel has radial behaviour
$\rho^{n-2}\rho^{-n(1+z)}\rho^2=\rho^{-nz}$. It is integrable for $nz<1$, while each of its fixed Taylor coefficients in
$z$ is an integrable polynomial in $\log\rho$. A formal all-order statement about coefficients is therefore not a
uniform statement for an arbitrary fixed $b$ and an arbitrarily large number of insertions. The analogous
non-uniformity in $\db$ is the order of limits of Appendix~\ref{app:clos:power}.

\paragraph{Identity, endpoints and composite insertions.} The identity counterterm cancels vacuum divergences and may
depend on the couplings. It drops out of normalized bulk correlators, but it cannot be ignored before the normalization
is formed correctly. Sharp interval endpoints and external operator insertions can require their own local
renormalizations. Such constants are not new translation-invariant couplings of the action, but they can affect an
extraction prescription based on integrated responses on a finite interval; the endpoint sources created by the
current resummation belong to this separate bookkeeping. This is the gap that the prescription argument addresses
(Appendices~\ref{app:clos:endpoint} to~\ref{app:clos:uniqueness}).

\subsection{Endpoint cancellation and the separated probes}\label{app:clos:endpoint}
\paragraph{Factorization and cancellation.} In correlators with separated probes, the operator-valued forest
construction gives the asymptotic factorization
\begin{equation}
 \mathcal D_{I,\varepsilon}(c_B)=\mathcal V(\varepsilon,c_B,L)\,\mathcal Z_0(\varepsilon,c_B)\,
 \mathcal Z_L(\varepsilon,c_B)\,\mathcal D_I^{\rm loc}(c_R)+o(1),
 \label{eq:clos:factor}
\end{equation}
after the same nontrivial local bulk redefinition $c_B=\mathcal B(c_R,\varepsilon\mu,\db)$ at each order. Here $c_B$
and $c_R$ are the bare and renormalized couplings, $\mathcal D_I^{\rm loc}$ is the locally renormalized evolution,
$\mathcal V$ is the bulk identity normalization, and $\mathcal Z_0$, $\mathcal Z_L$ are the scalar factors of the two
endpoints (Eq.~\eqref{eq:clos:scalaredge}). A finite scalar factor can be moved to either side without changing the
conclusion. Equation~\eqref{eq:clos:factor} holds coefficientwise, after the proper subdivergences are prepared; it is
not a uniform nonperturbative limit in operator norm.

For any separated reference probe in the numerator and the vacuum in the denominator, the three scalar factors cancel:
\begin{equation}
 \mathcal R_{\mathcal X,\varepsilon}\bigl(\mathcal B(c_R)\bigr)\longrightarrow
 \frac{\operatorname{tr}\langle\mathcal X\,\mathcal D_I^{\rm loc}(c_R)\rangle_0}
 {\operatorname{tr}\langle\mathcal D_I^{\rm loc}(c_R)\rangle_0}.
 \label{eq:clos:cancel}
\end{equation}
The cancellation happens before any differentiation with respect to parameters, so derivatives of coupling-dependent
scalar endpoint normalizations cancel as well.

Only the ultraviolet normalization factors cancel. The finite endpoint-source interactions created by the current
resummation do not; they remain in the finite responses and in the Jacobian of Appendix~\ref{app:clos:jacobian}.
Ignoring them would be a different, unjustified prescription.

\paragraph{The first endpoint logarithm appears at two kinks.} On an interval of length $L$ with hard core
$\varepsilon$, the two-kink vacuum and current integrals at $\db=0$, $b_0=1$ are
\begin{align}
 Z_2(L,\varepsilon)&=\int_\varepsilon^L\frac{L-r}{r^2}\dd r=\frac L\varepsilon+\log\frac\varepsilon L-1,
 \label{eq:clos:Z2}\\
 Q_2(L,\varepsilon)&=\int_\varepsilon^L\frac{L-r}{r}\dd r=L\log\frac L\varepsilon-L+\varepsilon,
 \label{eq:clos:Q2}
\end{align}
which at $L=1$ are the moments of Appendix~\ref{app:dualA:moments}. In $Z_2$ the term $L/\varepsilon$ is the bulk
identity normalization and the logarithm, which does not scale with $L$, is an endpoint identity contribution; it is
not zero. Both multiply any distant numerator by the same identity coefficient and cancel when it is divided by the
vacuum moment. In $Q_2$ the logarithm is proportional to $L$: it multiplies the bulk operator $\Oh$ and survives the
normalization. With the operator product $\Og(X+r)\Og(X)=r^{-2}\mathbf 1+\sqrt2r^{-1}\Oh(X)+O(1)$ it gives the
standard leading $\beta_h=-\sqrt2g^2$. The two integrals exhibit, rather than assume, the difference between an
endpoint logarithm and a bulk one.

\paragraph{Separated reference probes need no ultraviolet counterterm.} Fix a probe at $y$, a positive distance from
$I$. There is no collision diagonal involving $y$ and an integrated point, so its fixed free normal ordering suffices.
Counterterms inside $I$ act on any spectator through the universal local expansion~\eqref{eq:clos:spectator}. A qubit
matrix attached to the probe may be transported freely through the free exterior; the colliding interval product still
multiplies it by the same scalar endpoint coefficient, and transporting the matrix creates no new Gaussian
contraction.

At $\db>0$ the projection to a distant probe is normalized by the free two-point function, as in
\cite[Eq.~(4.18)]{Benedetti:2025nzp}. For a vertex probe the remaining position-dependent external Wick factors contain
$C(y-X)-C(y)$. On a fixed compact interval these tend to zero uniformly, like $O(y^{\db-1})$ for $\db<1$, or
coefficientwise like $y^{-1}$ times finite polynomials in logarithms. For a neutral current probe the monopole term
cancels and the first moment survives after the free normalization: it is the total length $\Xi_{2m}$ of the reversed domains, which
weights the current moments $Q_{2m}$ (Appendix~\ref{app:dualA:moments}), plus the exact endpoint current term. This yields the
archived odd responses $A_{2m+1}$ and current moments $Q_{2m}$. The limit is taken after the local ultraviolet
subtractions and introduces no further cutoff renormalization of the probe.

An interacting primary \emph{inside} the interval is different: it needs its own insertion renormalization, and it
cannot replace the exterior reference probe in the same two-response formula.

\subsection{The current resummation and the normal-ordering factors}\label{app:clos:current}
\paragraph{From the longitudinal coupling to the charge.} For smooth regulated fields, the change of frame with
$R_h(X)=\exp[\mathrm ih\hat\sigma_3\phi(X)/\sqrt2]$, $\mathcal D_h(L,0)=R_h(L)\widetilde{\mathcal D}(L,0)R_h(0)^{-1}$
(Appendix~\ref{app:dualA:regulator}), removes the derivative interaction from the interior and shifts the kink charge
to $b=b_0-\sqrt2h$. It follows either by differentiating the identity or by integrating the piecewise-constant spin
history $m(X)=\pm1$ by parts, $\int m\,\partial\phi=m\phi|_{\partial I}-\int m'\phi$. It is not a change of the
Gaussian integration variable and needs no path-integral Jacobian. The endpoint matrices remain and must not be
omitted \cite[Sec.~3.2.1]{Benedetti:2025nzp}. With the reference normalization fixed, $h$ and $b$ are therefore two
coordinates on the same coupling direction, not independent third and fourth parameters; at $\db\ne0$ the classical
scaling must also be included in passing to a dimensionless charge. This is why the closure statement in $(g,h)$ is a
statement about the two essential coordinates $(g,b)$.

\paragraph{Wick-ordering factors.} Wick ordering supplies a local factor at each kink
\cite[Eqs.~(3.20)--(3.23)]{Benedetti:2025nzp},
\begin{equation}
 g_h=g\exp\bigl[(\sqrt2b_0h-h^2)\,C(\varepsilon)\bigr],
 \label{eq:clos:gh}
\end{equation}
with $C=C_{\db}$ of Eq.~\eqref{eq:clos:cov}. The endpoint charges have magnitude $h/\sqrt2$. Their two self-contractions
give the common scalar $e^{-h^2C(\varepsilon)/2}$, identical in numerator and denominator. The finite endpoint-to-endpoint
and endpoint-to-kink contractions must be kept. The endpoint-to-kink source parameter is the $\eta=\sqrt2hb
=\sqrt2b_0h-2h^2$ of Eq.~\eqref{eq:dualA:eta}, and for alternating histories with $q_i=(-1)^{i+1}$ the finite endpoint
factors are the $U_n^{\rm odd}$ and $U_n^{\rm even}$ of Eq.~\eqref{eq:dualA:endpoint}, with $f_{\db}(L-X_i)$ in place of
$f_{\db}(1-X_i)$ on an interval of length $L$. At
$L=\kappa^{-1}=1$ the finite endpoint-pair covariance vanishes ($C_{\db}(\kappa^{-1})=0$); a general $L$ retains it as
part of the finite response, not as a new bulk coupling. The first source derivative may be combined by reflection, as
in the archive; that replacement cannot be made before squaring for a second source derivative.

\paragraph{The essential charge and fugacity.} The dimensionless essential charge and fugacity are
\begin{equation}
 B=(b_0-\sqrt2h)\,\varepsilon^{\db/2},\qquad
 Y=g_h\,\kappa\varepsilon\,e^{b^2C(\varepsilon)/2}=g\,\kappa\varepsilon\,e^{b_0^2C(\varepsilon)/2},
 \label{eq:clos:essential}
\end{equation}
where the second equality is the exact cancellation
\begin{equation}
 (\sqrt2b_0h-h^2)+\tfrac12(b_0-\sqrt2h)^2=\tfrac12b_0^2 .
 \label{eq:clos:wickcancel}
\end{equation}
With Eq.~\eqref{eq:clos:cov}, $Y=g\kappa\varepsilon\exp[b_0^2(\kappa^{-\db}-\varepsilon^{\db})/\db]$; this is precisely
the finite-tail matching~\eqref{eq:dualA:matching} of the compendium (expanded in Eq.~\eqref{eq:dualA:fugmatch}), where $g$ is called $g_{\rm aux}$ and $Y$ is
renamed $g$ after matching (Sec.~\ref{sec:dual}). The known explicit scale factors must be differentiated when the essential $\beta$ functions are
formed. On the reference section $h=0$, with the couplings identified at the reference matching scale and $v=B^2$,
the chain rule applied to Eq.~\eqref{eq:clos:essential} gives
\begin{equation}
 \frac{\beta_Y}{Y}=v-1+\frac{\beta_g^{\log}}{g},\qquad
 \frac{\beta_v}{v}=-\db-\frac{4\beta_h^{\log}}{\sqrt2b_0},
 \label{eq:clos:classical}
\end{equation}
where $\beta^{\log}$ denotes the logarithmic part extracted from the responses. This explains the classical-scaling terms
$v-1$ and $-\db$, rather than appending them in an unrelated scheme. With $u=Y^2$ and $z=v-1$ these are $F$ and $H$ of
Sec.~\ref{sec:dual} (Eq.~\eqref{eq:dualA:essential}); for instance $\beta_h^{\log}/(\sqrt2b_0)=-g^2+\cdots$ gives
$H=-\db+4u+\cdots$.

\paragraph{An exact check at \texorpdfstring{$g=0$}{g=0}.} With $J_I=\int_I\chi_B$,
\begin{equation}
 \frac{\operatorname{tr}\langle\hat\sigma_3\chi_B(y)\,e^{h\hat\sigma_3J_I}\rangle_0}
 {\operatorname{tr}\langle e^{h\hat\sigma_3J_I}\rangle_0}=h\,\langle\chi_B(y)J_I\rangle_0 ,
\end{equation}
since the qubit traces give $e^{hJ_I}\mp e^{-hJ_I}$, and for a Gaussian field
\begin{equation}
 \langle\chi_B(y)\,e^{\pm hJ_I}\rangle_0=\pm h\,\langle\chi_B(y)J_I\rangle_0\,\langle e^{hJ_I}\rangle_0 .
\end{equation}
Every Gaussian endpoint normalization factor cancels. The normalized far current response is exactly $hL$, so an
endpoint anomalous factor cannot generate a spurious pure-$h$ quantum $\beta$ function.

\subsection{The response Jacobian}\label{app:clos:jacobian}
Dividing the integrated odd and current numerators by the same integrated vacuum series gives, at $h=0$ and on the unit
interval, the two normalized responses
\begin{equation}
 R_g=g+\mathcal C_3g^3+\mathcal C_5g^5+\cdots,\qquad
 R_h=\sqrt2b_0\bigl(Q_2g^2+\mathcal K_4g^4+\mathcal K_6g^6+\cdots\bigr),
\end{equation}
and away from that section $R_h$ also contains the tree term $h$. The coefficients $\mathcal C_{2m+1}$ and
$\mathcal K_{2m}$ are the vacuum-normalized moments of Eq.~\eqref{eq:dualA:normalization}, for example
$\mathcal C_3=A_3-Z_2$ and $\mathcal K_4=Q_4-Q_2Z_2$. They are the coefficients of the normalized
correlators~\eqref{eq:clos:ratio}, not two functions selected after integration. The products are formed before
$[\varepsilon^0\ell^0]$ or $[\varepsilon^0\ell^1]$ is extracted, since a term of order $\varepsilon$ can multiply one of
order $1/\varepsilon$; the common-scalar cancellation~\eqref{eq:clos:cancel} requires this, and it is not a dispensable
numerical detail.

The projections $\mathcal C_{n,0}$, $\mathcal C_{n,1}$, $\mathcal K_{2m,0}$, $\mathcal K_{2m,1}$ and $Q_{2,0}$ (finite
part and single-logarithm coefficient) are those of Appendix~\ref{app:dualA:moments}, and $v=b_0^2$ on the reference
section. Exact resummation gives, on the normalized moment part at $h=0$, the source chain rule
$\partial_h=\sqrt2b_0(-2\partial_v+\partial_\eta)$ of Eq.~\eqref{eq:dualA:sourcechain}, because $h$ enters only through
the charge $b^2=v$ and the endpoint source $\eta$. The local factor~\eqref{eq:clos:gh} has no
finite constant at $\kappa\varepsilon=1$ (there $C(\varepsilon)=0$); its logarithms are accounted for by the matching.
Differentiating also the explicit current prefactor $\sqrt2b$, with $\partial_h(\sqrt2b)=-2$, gives the finite
source and current entries
\begin{equation}
\begin{aligned}
 j_1&=-\frac2{1+\db},\qquad
 j_{2m+1}=-2\partial_v\mathcal C_{2m+1,0}+\partial_\eta\mathcal C_{2m+1,0},\\
 \lambda_{2m}&=-2\mathcal K_{2m,0}-4v\partial_v\mathcal K_{2m,0}+2v\partial_\eta\mathcal K_{2m,0},
\end{aligned}
 \label{eq:clos:jlambda}
\end{equation}
with $\mathcal K_2=Q_2$; for $m=1,2$ these are the entries of Eq.~\eqref{eq:dualA:jacobian}. Here $j_1$ is the case
$m=0$, with $\mathcal C_1=A_1$: the one-kink response has no charge
dependence, and $j_1$ is its endpoint integral $\int_0^1[f_{\db}(X)+f_{\db}(1-X)]\dd X=2(\frac1{1+\db}-1)/\db$. The
coupling and source derivatives thus produce exactly the compendium's finite-response matrix $\mathcal J=D_cR_0$ of
Eq.~\eqref{eq:dualA:jacobianmatrix}, continued to the orders needed: $\partial_g$ of the finite parts of $R_g$ and $R_h$
gives its first column, and $\partial_h$, through Eq.~\eqref{eq:clos:jlambda}, its second. Its free value is the
identity, so it is formally invertible; the invertibility is used only \emph{after} locality and the endpoint cancellation have
been established.

\subsection{Uniqueness of the two-response solution, the recurrence and scheme independence}\label{app:clos:uniqueness}
\paragraph{Uniqueness.} Keep the known engineering and normal-ordering factors separate for the moment, and write the
normalized bare responses, in the common reference coordinates $c$, as
\begin{equation}
 R(c;\varepsilon)=R_0(c)+\ell R_1(c)+\ell^2R_2(c)+\cdots+o(1),\qquad \mathcal J=D_cR_0 .
 \label{eq:clos:Rexp}
\end{equation}
The absence of residual power divergences follows from the operator-valued local construction and the scalar
normalization, not from dropping such terms in the raw moments. The independently constructed local scheme has a bare
map
\begin{equation}
 c_B=\mathcal B(\tilde c,\ell),\qquad \mathcal B(\tilde c,0)=\mathcal G(\tilde c),\qquad
 \mathcal G(\tilde c)=\tilde c+O(\tilde c^2),
\end{equation}
whose coefficients are finite formal coefficients: they are the finite local Taylor moments of
Eq.~\eqref{eq:clos:cutoffdiff}, recursively including forests. Local finite normalizations can be chosen so that
$\mathcal G$ is the identity, and no response condition is needed for that choice; for an arbitrary local normalization
$\mathcal G$ is kept explicit. At fixed bare couplings, with $\ell=\log(\varepsilon\mu)$, the local logarithmic
renormalization-group identity is
\begin{equation}
 \partial_\ell\mathcal B+D_{\tilde c}\mathcal B\,\tilde\beta^{\log}=0 .
\end{equation}
Every normalized spectator response is finite in the renormalized couplings, so its $\ell$ derivative at fixed
$\tilde c$ vanishes. At first order in $\ell$ this gives
\begin{equation}
 R_1\bigl(\mathcal G(\tilde c)\bigr)=DR_0\bigl(\mathcal G(\tilde c)\bigr)\,D\mathcal G(\tilde c)\,
 \tilde\beta^{\log}(\tilde c),
 \label{eq:clos:localtoR}
\end{equation}
and therefore
\begin{equation}
 \beta^{\log}_{\rm archive}\bigl(\mathcal G(\tilde c)\bigr)=D\mathcal G(\tilde c)\,\tilde\beta^{\log}(\tilde c),
 \qquad \beta^{\log}_{\rm archive}=\mathcal J^{-1}R_1 .
 \label{eq:clos:equivalence}
\end{equation}
The known explicit classical-scaling and normal-ordering maps are then applied on both sides, including their scale
derivatives as in Eq.~\eqref{eq:clos:classical}; this gives the corresponding relation for the full essential vector
fields. The charge is not held fixed when a change of scale physically changes $B$.

The argument is not circular. It constructs the local renormalization and the finiteness for all spectators in
Appendices~\ref{app:clos:power} to~\ref{app:clos:endpoint}, independently of the response equations;
Equation~\eqref{eq:clos:localtoR} then makes the local $\beta$ function a solution, and the invertibility of
$\mathcal J$ makes it the \emph{unique} solution in the specified coordinates. No
independent anomalous term of Eq.~\eqref{eq:clos:anomaly} remains: the endpoint factors cancel before
differentiation, and the exterior reference probe does not collide with an integrated operator. If either fact were
false, the uniqueness argument would not apply.

\paragraph{The recurrence at every required power.} With the auxiliary coefficients $b_{2m+1}$ and $h_{2m}$ of
Appendix~\ref{app:dualA:recurrence} ($\beta_g=b_3g^3+b_5g^5+\cdots$, $\beta_h/(\sqrt2b_0)=-g^2+h_4g^4+\cdots$),
solving Eq.~\eqref{eq:clos:equivalence} on the section $h=0$ means solving the system $\mathcal J\beta=R_1$ of
Eq.~\eqref{eq:dualA:system}, with $R_1$ of Eq.~\eqref{eq:dualA:R1} continued to the orders needed. Order by order in
$g$ it is triangular: the coefficient of $g^{2m+1}$ in the first
row contains $b_{2m+1}$ with coefficient one and otherwise only lower coefficients, and the coefficient of $g^{2m}$ in
the second row does the same for $h_{2m}$ (the factor $\sqrt2b_0$ of $\mathcal J_{12}$ times the factor $\sqrt2b_0$ of $\beta_h$ gives the $2v$ of the
recurrence; in the second row the factor $\sqrt2b_0$ of $\mathcal J_{21}$ cancels against that of $R_1$). The solution is the compendium's recurrence: Eq.~\eqref{eq:dualA:recurrence} through $b_7$, and its continuation to
$h_8$ and $b_9$ in Appendix~\ref{app:dualC:C8setup} (Eq.~\eqref{eq:dualC:b9}). All finite and source terms are
retained; it is not a new fit to the final $\CSR_n$. It was reproduced by the prescription
argument's symbolic checker and re-derived independently by the review (Appendix~\ref{app:clos:checks}).

The lowest order, also the check of signs of Appendix~\ref{app:dualA:recurrence}, shows what the finite endpoint
term does. At $\db=0$ and $v=1$ the archived exact normalization gives
\begin{equation}
 \mathcal C_3=2\ell^2+5\ell+3-2\zeta(2)+o(1),\qquad j_1=-2,
\end{equation}
so the first line of the recurrence, $b_3=\mathcal C_{3,1}+2vj_1$, gives $b_3=5-4=1$, while $h_2=-1$ follows from
Eq.~\eqref{eq:clos:Q2}. This is the leading $F=z+u+\cdots$ of Sec.~\ref{sec:dual}. Discarding the finite endpoint
derivative $j_1$ would give $b_3=5$, a wrong cubic coefficient, even though the ultraviolet endpoint factor itself has
cancelled.

\paragraph{Scheme independence.} For the full essential vector fields, Eq.~\eqref{eq:clos:equivalence} is the usual
transformation under a finite change of coordinates. At a fixed point $c_*$,
\begin{equation}
 M_{\rm archive}\bigl(\mathcal G(c_*)\bigr)=D\mathcal G(c_*)\,M_{\rm local}(c_*)\,D\mathcal G(c_*)^{-1};
\end{equation}
the term with second derivatives of $\mathcal G$ vanishes because the $\beta$ function itself vanishes there. The
trace, the determinant and the two scaling eigenvalues therefore coincide coefficientwise. The finite map is regular in
the formal $\db$ expansion, so substituting the fixed-point series in $\sqrt{\db}$ does not introduce a hidden change of
the physical parameter $s$. This gives the formal scheme statement at any finite computed order. It needs no separately
evaluated integral on a ring, and none is claimed: a ring in the compact boson theory also has a monodromy sector that
must be treated correctly, and deleting the endpoints of an existing interval integral is not automatically a second
prescription~\cite{Benedetti:2025nzp}.

\paragraph{Other observables.} For any other separated observable the same local forest construction and the same bare
map apply: Eq.~\eqref{eq:clos:spectator} identifies its local collision coefficients before projection, and
Eq.~\eqref{eq:clos:cancel} removes its scalar endpoint factors. There is therefore no spectator-dependent bulk
counterterm. Physical operator insertions inside the support are renormalized separately as composite operators, and
their mixing is not a third bulk coupling.

\subsection{What the arguments establish and what they do not}\label{app:clos:scope}
\paragraph{For the operator content.} The closure argument improves the status of the operator-content assumption: at
the finite orders used for the expansion about the endpoint, no third symmetry-preserving bulk operator that needs an
ultraviolet counterterm is expected or required in the formal local subtraction construction, and the same argument
applies at each finite order rather than stopping at a two-insertion test. It does not certify the archived numerical
$\beta$ coefficients. The exact certificates of Sec.~\ref{sec:closed} prove identities between supplied integral
expressions; they do not show that the finite-response prescription assembled the correct local counterterms for all
collision configurations.

\paragraph{For the prescription.} To connect the archive with the closure statement, the closure argument named a
concrete audit in three steps: (i) construct the collision subtractions of the regulated, path-ordered kernels as
operator-valued products, including every proper cluster and not only the projections that define the two response
coordinates; (ii) show that the archive's vacuum normalization, current resummation, endpoint prescription and
finite-$\db$ matching implement those subtractions, or differ from them by an explicit finite redefinition of the two
essential couplings and separately specified operator and endpoint normalizations; (iii) check an independent family of
non-projected correlators after inserting the same counterterms, whose coefficients in the $\mathbf 1$, $\Og$ and $\Oh$
channels must not depend on the distant spectator configuration, with locally integrable remainders of derivative order
two and higher. In this note's assessment (Appendix~\ref{app:clos:checks}), the prescription argument supplies steps
(i) and (ii) at the formal level
(Appendices~\ref{app:clos:forest} to~\ref{app:clos:uniqueness}). The prescription argument addresses the analytic content of step (iii)
through the spectator expansion~\eqref{eq:clos:spectator} and the paragraph ``Other observables'' of
Appendix~\ref{app:clos:uniqueness}; the computation itself has not been done. A numerical check of the prescription
with a second regulator and different responses, within the same two-coupling description, is the magnetization-scheme
calculation of Sec.~\ref{sec:integ}, which reproduces $\CSR_3$ to $\CSR_7$ and the record of $\CSR_8$ (Sec.~\ref{sec:open}). Given (i) and (ii), regular
finite coupling redefinitions transform the stability matrix by similarity, so its eigenvalues do not depend on the
chosen two-coupling coordinates.

\paragraph{What is not established.} Neither argument proves the lattice duality, the identification of the dual and
of its thermal branch with the critical chain, or the infinite-volume limit; these keep the scope of the published
proposal~\cite{Benedetti:2024wgx,Benedetti:2025nzp}. Neither proves convergence or resummability at fixed $\db$, nor
a non-perturbative renormalizability statement. Neither recomputes $\CSR_3$ to $\CSR_7$ in a second regulator or
audits the arithmetic of the historical high-order integrations, whose 21 stage programs are missing
(Appendix~\ref{app:cannot}). Their fidelity to the recipe is tested only numerically: by the independent
re-derivation (Sec.~\ref{sec:indep}) and, through the invariants they produce, by the magnetization-scheme calculation
in a second regulator (Sec.~\ref{sec:integ}); the exact certificates (Sec.~\ref{sec:closed}) concern the later reduction of the archived
expressions to zeta values. Neither bears on the conjectured zeta-value form of
$\CSR_8$ (Eq.~\eqref{eq:C8cand}) or on the statistical claims of the NRG or Monte Carlo comparisons, for which no raw
data were supplied to them. The statement of the prescription argument
concerns the mathematically specified recipe: implementations must still be checked to follow it, and its formula audit
covers the normalization, the current and endpoint identity, the Wick factors, the finite-response Jacobian and its
recurrence. The arguments are also narrower than a proof of the fixed-$s$ duality, and they avoid the opposite mistake
of treating every generated irrelevant Wilsonian coefficient as a new independent tuning parameter.

\paragraph{Status in this note.} Both arguments are formal and unrefereed; this note uses them as such. The status
word ``proved closed form within the prescription'' (Sec.~\ref{sec:words}) is unchanged by them: they bear on whether
the prescription computes the coefficients of the dual, not on the reduction of the archived expressions to zeta
values. The evidence that does not use the dual at all is the NRG test of $\CSR_3$ to $\CSR_5$
(Sec.~\ref{sec:d1caveats}).

\subsection{Assessment and checks by the review}\label{app:clos:checks}
Claude Code assessed both arguments at the author's request on 15 October 2026 (Sec.~\ref{sec:prov}); the caveat
``Two-coupling closure'' of Sec.~\ref{sec:d1caveats} is the summary that the author approved afterwards.

\paragraph{The closure argument.} Verdict: narrowed; the operator content is closed formally, and the matching of the
archived prescription was left open. The argument is sound as a standard
perturbative power-counting (Epstein--Glaser) argument, and no gap was found that changes its conclusion. It is a
careful sketch: the uniform scaling bounds in the shape sectors and the sufficiency of a finite Taylor jet are
asserted, not estimated in detail, and the extension step is invoked from the cited theorems. Checked independently:
the classification through dimension one under the shift symmetry, the internal reflection and parity
(Table~\ref{tab:clos:operators}); the charge selection rule; the radial-power identity; the product
$\Og\Og$~\eqref{eq:clos:ope}; the coefficientwise logarithmic deformation~\eqref{eq:clos:deform}; symmetric extensions by
group averaging; and the frame change between $h$ and $b$, identical to the compendium's. The assessment found that
the archive's prescription did not obviously implement the local subtraction, and named the risks: (i) the $\beta$
functions are solved from exactly two projected responses, a square system that would absorb a mis-assigned endpoint
or insertion logarithm silently; (ii) all responses live on one unit interval, with the current resummed into endpoint
sources and no endpoint hard core, so endpoint and bulk logarithms are not separated by construction; (iii) the
finite-$\db$ matching is the exact normal-ordering factor, coefficientwise a regular redefinition, but its equivalence
to the local scheme was not shown; (iv) the current resummation and (v) the vacuum normalization were judged low risk.
The prescription argument answers this open link: (i) by the scalar endpoint factors, the separated probes and the
uniqueness argument (Appendices~\ref{app:clos:endpoint} and~\ref{app:clos:uniqueness}); (ii) by the endpoint power
counting and the two-kink integrals~\eqref{eq:clos:Z2} and~\eqref{eq:clos:Q2}; (iii) by the normal-ordering audit
(Appendix~\ref{app:clos:current}).

\paragraph{The prescription argument.} Verdict: the link to the prescription is closed at the formal level, at the same
level of rigour as the closure argument. Checked: the endpoint power counting (one fewer free coordinate, so endpoint
divergences only at reference dimension zero or below); the endpoint classification (the invariant dimension-zero
sector is the identity, so endpoint ultraviolet factors are scalars that cancel in the normalized responses); that the
finite endpoint-source terms survive and are the compendium's $j$ and $\lambda$ terms ($b_3=1$, against $5$ without
$j_1$); the cancellation~\eqref{eq:clos:wickcancel} and $\eta=\sqrt2b_0h-2h^2$; the two-kink integrals at general $L$
and $\varepsilon$; and the recurrence for $b_3$ to $b_9$ and $h_4$ to $h_8$, re-derived from $\mathcal J\beta=R_1$ and
equal to the compendium's formulas. The logic is not circular: local renormalization for arbitrary spectators is
constructed first, and the invertibility of the response system is used only for uniqueness. Stated rather than shown
in detail: the sharp-edge limit from smooth switching, the recursion relating the hard-core regulator to the local
subtraction when exclusion regions do not factorize (Eq.~\eqref{eq:clos:cutoffdiff}), and the projection to a distant
probe.

\paragraph{Code.} \file{code/closure_d1/public/check_closure.py}, written by the review and run by the package's
\texttt{test} and \texttt{replay} targets (standard library and sympy, a few seconds), checks the finite algebra:
\begin{itemize}
\item all $3+3^2+\cdots+3^9=29\,523$ words of length one to nine in $\hat\sigma_+$, $\hat\sigma_-$, $\hat\sigma_3$, in
 exact integer arithmetic: the 2\,035 nonzero words all have $Q\in\{0,\pm1\}$ and $[\hat\sigma_3,M]=2QM$, with no
 violation;
\item the radial-power identity $2\sum_{i<j}q_iq_j=Q^2-n$ for alternating charges, $n=1$ to $11$;
\item the coefficients $1$, $\sqrt2$, $1$ of the midpoint product~\eqref{eq:clos:ope} (to $10^{-15}$);
\item a rerun of the prescription argument's own checker,
 \file{code/closure_d1/prescription/code/check_prescription.py} (sympy): all 22 of its checks are true and equal to its
 shipped record. They cover the matrix words, the scalar invariant matrix at dimension zero, the bulk and endpoint
 radial powers, the cancellation~\eqref{eq:clos:wickcancel} and $\eta$, the source chain rule of $\lambda_{2m}$, the
 integral $j_1$, every coefficient equation of the response system through $g^9$, both two-kink integrals at general
 $L$ and $\varepsilon$, the spectator expansion at zeroth and first order, a toy response-generator identity, and the
 cancellation of a common scalar factor;
\item independently of that checker, the recurrence $b_3,\dots,b_9$, $h_4,\dots,h_8$ solved order by order from
 $\mathcal J\beta=R_1$ (sympy), which equals the compendium's formulas exactly and gives $b_3=1$ at $v=1$,
 $\mathcal C_{3,1}=5$, $j_1=-2$.
\end{itemize}
The export tests on the Intel and AMD nodes of the symmetry cluster give the same table.

\paragraph{What the checks establish.} They establish the finite algebra on which the arguments rest: the selection
rule for words up to length nine (the general rule is argued for every length independently of the enumeration), the
pair product, the radial powers, the normal-ordering identities, the two-kink integrals, and that the compendium's
recurrence is exactly what the response system $\mathcal J\beta=R_1$ implies. They do not check the analytic
induction, the uniform bounds in the shape sectors, the sharp-edge limit, the hard-core recursion or the far-field
projection; they recompute no archived $\beta$ coefficient; and they are not a numerical test of the prescription. The
closure argument describes its own script of finite checks (the same algebra, a bound on the logarithmic degree
through order nine and an exact example of the tangent-stability identity); the script was not supplied and is not part
of the release, and the review's checker above replaces its first three items.

%====================================================================
\section{The closed form of \texorpdfstring{$c_2$}{c2}: derivation}\label{app:cform}
%====================================================================
This appendix gives the derivation of the closed form of $c_2$ in the closed-form study (``the study'', written by
OpenAI GPT-6 astra; Sec.~\ref{sec:sources}), edited to the notation of this note.

Section~\ref{sec:d2} states the steps and their results; here each step is derived, in the order in which the
results depend on one another. Appendices~\ref{app:cform:conv} to~\ref{app:cform:assembly} reduce the coefficient to
the elliptic finite part $\cI$ of Eq.~\eqref{eq:I} and give Eq.~\eqref{eq:c2I};
Appendices~\ref{app:cform:moments} to~\ref{app:cform:telescoping} evaluate $\cI$ in closed form; and
Appendix~\ref{app:cform:numerics} describes the numerical checks of that evaluation. The status is that of
Table~\ref{tab:main}: the evaluation of $\cI$ is proved, and $c_2$ is a closed form conditional on one unproved step,
the uniform-remainder exchange (Appendix~\ref{app:cform:qualification} and Sec.~\ref{sec:unproved}) on which the
identification of Eq.~\eqref{eq:c2closed} with the physical coefficient rests.

\paragraph{Notation.} We write $c_2$ for the study's $C_2$, $\hMS$ for the fourth-order coefficient of the anomalous
dimension in minimal subtraction (the study's $h_0$), $\kappa$ for the parameter of a finite operator redefinition (the
study's $q$), and $\mathsf P,\mathsf Q,\mathsf T$ for its three ${}_4F_3(1)$ constants. Auxiliary symbols are
renamed where the study's symbol would clash with another one in this note: the M\"obius scale $\Lambda$ (study: $a$), the factor $\Upsilon$
($Y$), the kernel combination $\Sigma(x,y;z)$ ($S$), the weights $\Omega_1,\Omega_0$ ($U,W$), the constant matrices
$\mathsf J,\mathsf L,\mathsf C_\pm,\mathsf B_-,\mathsf H,\mathsf D$ ($J,L,C_\pm,B_-,H,D$), the charge factors
$\mathcal W_{\boldsymbol e},\mathcal P_{\boldsymbol e}$ ($W_{\boldsymbol e},P_{\boldsymbol e}$), the moment combinations
$m_U,m_V$ ($f,g$) and $d_O$ ($d$), the polynomial $\varpi$ ($p$), the normalized integral $\hat K=2K/\pi$ ($F$), the
hypergeometric functions $\Phi_A,\Phi_D$ ($A,D_h$), the period vectors $\mathbf X,\mathbf Y,\mathbf Z,\mathbf W$
($X_v,Y_v,Z_v,W_v$), the derivative combinations $\mathcal D_1,\mathcal D_0$ ($D_1,D_0$), and the Mellin integrals
$j_0,j_1,j_2$ ($J_0,J_1,J_2$). The complex beta integral is written throughout as
\begin{equation}
 \mathcal B(a,b)=\AC\!\int_{\mathbb C}|z|^{2a-2}|1-z|^{2b-2}\dd^2z
 =\pi\frac{\Gamma(a)\Gamma(b)\Gamma(1-a-b)}{\Gamma(1-a)\Gamma(1-b)\Gamma(a+b)},
 \label{eq:cform:beta}
\end{equation}
which the study writes as $J(a-1,b-1)$ in its lower-order steps.

%--------------------------------------------------------------------
\subsection{Observable, regulator and continuation}\label{app:cform:conv}
The crossover theory is that of Sec.~\ref{sec:d2}: the Ising CFT times a generalized free scalar $\chi$, with
unit-normalized fields, $\langle\sigma(x)\sigma(0)\rangle=|x|^{-1/4}$ and $\langle\epsilon(x)\epsilon(0)\rangle=|x|^{-2}$,
and the Ising data listed there. An insertion at infinity means
$O(\infty)=\lim_{|x|\to\infty}|x|^{2\Delta_O}O(x)$, so that $\langle O(0)O(\infty)\rangle=1$. The fusion rules are
$\sigma\times\sigma=\mathbb 1+\epsilon$ and $\epsilon\times\epsilon=\mathbb 1$, each family including all its Virasoro
descendants. The generalized free field is fixed by Wick's theorem and
\begin{equation}
 \langle\chi(x)\chi(y)\rangle=|x-y|^{-15/4+2\delta},
 \label{eq:cform:gff}
\end{equation}
so that $\Delta_\chi=15/8-\delta$ and the perturbing operator $\sigma\chi$ has dimension $2-\delta$. The bare action is
\begin{equation}
 S=S_{\rm Ising}+S_\chi+\int g_0(x)\,\sigma\chi(x)\dd^2x+\int t_0(x)\,\epsilon(x)\dd^2x,
\end{equation}
where $g_0$ has mass dimension $\delta$ and the thermal source $t_0$ is tuned to the critical surface. The thermal
source cannot be left out of the definition of the prescription, because the short-distance product contains
\begin{equation}
 \sigma\chi(x)\,\sigma\chi(0)\supset|x|^{-4+2\delta}\,\mathbb 1+\tfrac12|x|^{-3+2\delta}\epsilon(0).
\end{equation}
The second term produces a thermal power divergence. It is retained by analytic continuation, as appropriate to the
critically tuned theory.

\paragraph{Observable.} The observable is the dimension $\De$ of the critical thermal eigenoperator continued from
$\epsilon$,
\begin{equation}
 \De=1+c_2\delta^2+O(\delta^3);
\end{equation}
the absence of a linear term is derived in Appendix~\ref{app:cform:lower}. In terms of $s$,
$\De=1-1.01656753703313038819\ldots\,(7/4-s)^2+O\bigl((7/4-s)^3\bigr)$, the coefficient being $c_2/4$
(Table~\ref{tab:conv}). Since $1/\nu=2-\De$, also $\nu=1+c_2\delta^2+O(\delta^3)$.

\paragraph{Regulator and continuation.} All complex integrals use $\dd^2z=\dd(\mathrm{Re}\,z)\,\dd(\mathrm{Im}\,z)$.
The renormalized coupling is localized by the source profile $f(w)=e^{-|w|^2}$; its width and the renormalization
scale $\mu$ are set to one after the renormalization identities have been derived. The Gaussian controls the
infrared. Ultraviolet integrals are first defined in a domain of convergence and then continued meromorphically in
$\delta$ and, where needed, in auxiliary exponents; $\AC$ denotes this continuation. In particular,
\begin{equation}
 \AC\!\int_0^R r^{p-1+b\delta}\dd r=\frac{R^{p+b\delta}}{p+b\delta}.
 \label{eq:cform:mellin}
\end{equation}
For $p\ne0$ the value at $\delta=0$ need not vanish; for $p=0$ the integral has a logarithmic pole. A circular finite
part, $\FP$, is defined with rotationally symmetric exclusions about each specified singularity, and the subtracted
local terms are restored in the same continuation. This convention fixes the angular contacts used below.

%--------------------------------------------------------------------
\subsection{Renormalization-group reduction}\label{app:cform:rg}
This subsection explains why a fourth-order residue determines the quadratic coefficient. In pole-only analytic
minimal subtraction let
\begin{equation}
 \beta(g,\delta)=\mu\frac{\dd g}{\dd\mu}=-\delta g+b_0g^3+O(g^5),\qquad
 \gamma_{\mathcal E}(g)=\hMS g^4+O(g^6).
\end{equation}
The interacting fixed point has $g_*^2=\delta/b_0+O(\delta^2)$, and the dimension there is
$\De=1+\gamma_{\mathcal E}(g_*)$, so $c_2=\hMS/b_0^2$. The thermal operator renormalizes multiplicatively,
$\mathcal E_0=Z_{\mathcal E}\mathcal E_R$, with $\gamma_{\mathcal E}=\beta\,\partial_g\log Z_{\mathcal E}$. There is no
energy pole at order $g^2$: the second-order integral is finite at $\delta=0$ because the first-order shift vanishes
(Appendices~\ref{app:cform:lower} and~\ref{app:cform:source}). Therefore
\begin{equation}
 \log Z_{\mathcal E}=\frac{z_{4,1}}{\delta}g^4+O(g^6),\qquad
 \gamma_{\mathcal E}=\Bigl(-\delta g+\cdots\Bigr)\frac{4z_{4,1}g^3}{\delta}=-4z_{4,1}g^4+O(g^6),
\end{equation}
which gives $\hMS=-4z_{4,1}$ and the relation $c_2=-4z_{4,1}/b_0^2$ of Sec.~\ref{sec:rg2}.

The factor four relies on there being a single renormalization factor. The source is localized at the origin, so in
the correlator with the energies at $0$ and $\infty$ the only local ultraviolet factor is the $Z_{\mathcal E}$ at $0$;
since $f(\infty)=0$, the energy at infinity supplies no second source-dependent ultraviolet factor.

The restriction to pole-only subtraction is substantive. The split of $c_2$ between a $\delta g^2$ term and a $g^4$
term of $\gamma_{\mathcal E}$ depends on the scheme (Sec.~\ref{sec:rg2}): the finite redefinition
$\mathcal E_R\to e^{\kappa g^2}\mathcal E_R$ moves weight from one to the other. The finite lower-order source contact
and the fourth-order pole must therefore be matched in one and the same convention, which is what
Appendices~\ref{app:cform:source} to~\ref{app:cform:assembly} do.

%--------------------------------------------------------------------
\subsection{Exact Ising correlators}\label{app:cform:correlators}
For $2n=2,4$ distinct spin positions $z_i$, sum over the neutral charges $e_i=\pm1$ with $\sum_ie_i=0$, and define
\begin{equation}
 \mathcal W_{\boldsymbol e}=\prod_{i<j}|z_i-z_j|^{e_ie_j/2},\qquad
 \mathcal P_{\boldsymbol e}=1-\frac14\Bigl(\sum_ie_iz_i\Bigr)\Bigl(\sum_i\frac{e_i}{z_i}\Bigr).
\end{equation}
The Euclidean correlators needed are~\cite{Ardonne:2010zz}
\begin{equation}
 S_{2n}=\langle\sigma_1\cdots\sigma_{2n}\rangle=\Bigl[2^{-n}\sum_{\boldsymbol e}\mathcal W_{\boldsymbol e}\Bigr]^{1/2},\qquad
 H_{2n}=\langle\epsilon(0)\epsilon(\infty)\sigma_1\cdots\sigma_{2n}\rangle
 =\Bigl[2^{-n}\sum_{\boldsymbol e}\mathcal W_{\boldsymbol e}|\mathcal P_{\boldsymbol e}|^2\Bigr]^{1/2}.
 \label{eq:cform:neutral}
\end{equation}
The positive Euclidean square root is fixed by the identity OPE. These formulas follow by combining two independent
Ising copies, $\sigma^{(1)}\sigma^{(2)}=\sqrt2\cos(\varphi/2)$ with $\langle\varphi(z)\varphi(w)\rangle=-2\log|z-w|$,
and representing the product of energies by currents. Charge neutrality gives the displayed sum; taking the square
root removes the second copy. For energies at finite positions $x,y$, $\mathcal P_{\boldsymbol e}$ is replaced by
\begin{equation}
 \frac1{(x-y)^2}+\frac14\Bigl(\sum_i\frac{e_i}{x-z_i}\Bigr)\Bigl(\sum_i\frac{e_i}{y-z_i}\Bigr).
 \label{eq:cform:finiteenergy}
\end{equation}
For two spins (energies at $0$ and $\infty$) this gives
\begin{equation}
 S_2(a,b)=|a-b|^{-1/4},\qquad H_2(a,b)=\frac{|a+b|^2}{4|ab|\,|a-b|^{1/4}}.
 \label{eq:cform:twospin}
\end{equation}

%--------------------------------------------------------------------
\subsection{The vanishing first-order shift and the exact \texorpdfstring{$b_0$}{b0}}\label{app:cform:lower}
\paragraph{First order.} The first-order energy shift is proportional to the integral $I_2(\delta)$ of
Sec.~\ref{sec:d2}, whose integrand is $|1-z|^{-4+2\delta}$ times $H_2(z,1)/S_2(z,1)-1$ from
Eq.~\eqref{eq:cform:twospin}. The combined integrand is absolutely integrable for $0<\mathrm{Re}\,\delta<1/2$. The term
$-1$ (the subtracted spin two-point function) gives a scaleless integral, which vanishes in the continuation. With
$|1+z|^2=2(1+|z|^2)-|1-z|^2$, the rest is three complex beta integrals~\eqref{eq:cform:beta}:
\begin{equation}
 I_2(\delta)=\tfrac12\mathcal B(\tfrac12,\delta-1)+\tfrac12\mathcal B(\tfrac32,\delta-1)-\tfrac14\mathcal B(\tfrac12,\delta)
 =\frac{\pi\delta^2}{4(1-\delta)^2}\frac{\Gamma(\delta)\Gamma(\frac12-\delta)}{\Gamma(1-\delta)\Gamma(\frac12+\delta)}.
\end{equation}
Expanding the Gamma ratios gives $I_2(\delta)=\frac\pi4[\delta+(2+4\log2)\delta^2+O(\delta^3)]$. Thus $I_2(0)=0$, which
reproduces the leading cancellation ($c_1=0$), while $I_2'(0)=\pi/4$. The finite derivative is retained in the source
matching of Appendix~\ref{app:cform:source}; it is not a separate contribution to be added to $c_2=\hMS/b_0^2$.

\paragraph{The exact normalization of $b_0$.} In the same field normalization the leading coefficient of the beta
function is
\begin{equation}
 b_0=-\pi I_\beta,\qquad
 I_\beta=\FP\!\int_{\mathbb C}|z|^{-15/4}\bigl[S_4(0,z,1,\infty)-|z|^{-1/4}\bigr]\dd^2z,
 \label{eq:cform:Ibeta}
\end{equation}
with $S_4$ from Eq.~\eqref{eq:cform:neutral}. The quadratic map $z=4w/(1+w)^2$, followed by $u=w^2$, with its
multiplicity and Jacobian, reduces the integral to
\begin{equation}
 I_\beta=\frac1{64}\AC\!\int_{\mathbb C}\bigl(|u|^{-3}+|u|^{-5/2}\bigr)|1-u|^{3/2}\dd^2u.
\end{equation}
Inversion, $u\mapsto1/u$, interchanges the two terms. Equation~\eqref{eq:cform:beta} therefore gives
\begin{equation}
 I_\beta=\frac1{32}\mathcal B(-\tfrac12,\tfrac74)=-\frac{9\pi}{8}\Bigl[\frac{\Gamma(3/4)}{\Gamma(1/4)}\Bigr]^2,
\end{equation}
which is the exact $b_0=\frac{9\pi^2}{8}[\Gamma(\frac34)/\Gamma(\frac14)]^2$ of Sec.~\ref{sec:rg2}. All power terms in
this reduction use Eq.~\eqref{eq:cform:mellin}.

%--------------------------------------------------------------------
\subsection{The localized source and the fourth-order logarithm}\label{app:cform:source}
\paragraph{The connected kernel.} With the energies at $0$ and $\infty$, write $S_{ij}=S_2(w_i,w_j)$ and
$H_{ij}=H_2(w_i,w_j)$ from Eq.~\eqref{eq:cform:twospin}, and $S_4,H_4$ for the four-spin functions of
Eq.~\eqref{eq:cform:neutral}. The fourth-order coefficient of a logarithm is its fourth-order moment minus one half
the square of its second-order moment. Applied to the logarithm of the vacuum-normalized energy two-point function,
for the fixed generalized-free pairing $(12)(34)$, this gives the connected kernel
$\mathcal F=H_4-S_4-H_{12}H_{34}+S_{12}S_{34}$ of Sec.~\ref{sec:d2}; its last two terms are fixed by the logarithm.
The three Wick pairings have equal integrals by permutation symmetry, so $3/4!=1/8$ and
\begin{equation}
 A_4(\delta)=\frac18\AC\!\int_{\mathbb C^4}\prod_{i=1}^4\dd^2w_i\,e^{-\sum_i|w_i|^2}\,
 |w_{12}|^{-15/4+2\delta}|w_{34}|^{-15/4+2\delta}\,\mathcal F(w_1,w_2,w_3,w_4),
 \label{eq:cform:A4}
\end{equation}
with $w_{ij}=w_i-w_j$. Connectedness alone does not make Eq.~\eqref{eq:cform:A4} an ordinary ultraviolet-convergent
integral.

\paragraph{Why the finite Gaussian two-point term is needed.} Locality of the coupling counterterm gives
\begin{equation}
 g_0(w)=\mu^\delta\Bigl[g\,f(w)+\frac{b_0}{2\delta}g^3f(w)^3+O(g^5)\Bigr].
 \label{eq:cform:ct}
\end{equation}
The power $f^3$ follows from the three perturbing insertions in the local coupling logarithm. It cannot be replaced by
$f$ inside a finite lower-order integral. Define the second-order integral with two Gaussian weights,
\begin{equation}
 A_{a,b}(\delta)=\frac12\AC\!\int_{\mathbb C^2}\dd^2w_1\dd^2w_2\,e^{-a|w_1|^2-b|w_2|^2}\,
 |w_{12}|^{-15/4+2\delta}\bigl(H_{12}-S_{12}\bigr).
 \label{eq:cform:Aab}
\end{equation}
Inserting Eq.~\eqref{eq:cform:ct} into the second-order term gives two equivalent placements of the cubic
counterterm, whose total fourth-order contribution is $b_0A_{3,1}(\delta)/\delta$. Hence $z_{4,1}=\Res A_4+b_0B_{31}$
with $B_{31}=A_{3,1}(0)$, as in Sec.~\ref{sec:d2}. There is one ultraviolet external-operator factor: the localized
source produces collisions at the finite energy insertion and is suppressed near infinity, so there is no extra
factor of two from the second energy.

\paragraph{Evaluation of $B_{31}$.} Set $w_1=\Lambda z$ and $w_2=\Lambda$ in Eq.~\eqref{eq:cform:Aab}, with $\Lambda$
the overall complex scale. The integral over $\Lambda$ gives
\begin{equation}
 A_{a,b}(\delta)=\frac\pi2\Gamma(\delta)\AC\!\int_{\mathbb C}\dd^2z\,(a|z|^2+b)^{-\delta}F_\delta(z),\qquad
 F_\delta(z)=|1-z|^{-4+2\delta}\Bigl(\frac{|1+z|^2}{4|z|}-1\Bigr).
 \label{eq:cform:Fdelta}
\end{equation}
By Appendix~\ref{app:cform:lower}, $\AC\!\int F_\delta(z)\dd^2z=I_2(\delta)=\frac\pi4\delta+O(\delta^2)$, so the
vanishing two-point residue leaves a finite contribution when it is multiplied by the scale factor $\Gamma(\delta)$.
For smooth radial test functions $h$ with the required decay, and by continuation for the logarithmic weights below,
the circular finite part about $z=1$ obeys the distributional identity
\begin{equation}
 \FP\!\int_{\mathbb C}\dd^2z\,F_0(z)\,h(|z|)=\frac\pi2\int_0^\infty\frac{|r-1|}{(1+r)^3}h(r)\dd r-\frac\pi4h(1).
 \label{eq:cform:radialfp}
\end{equation}
For $r\ne1$ the first term follows from the angular integral. The second is the local term from the circular exclusion
about $z=1$. Its coefficient is also fixed by $I_2(0)=0$, since the bulk radial integral for $h=1$ is $\pi/4$. It must
be retained when the source weight is differentiated. Expanding Eq.~\eqref{eq:cform:Fdelta} at $\delta=0$ gives
\begin{equation}
 A_{a,b}(0)=\frac{\pi^2}{8}-\frac\pi2\FP\!\int_{\mathbb C}\dd^2z\,F_0(z)\log(a|z|^2+b).
 \label{eq:cform:Aab0}
\end{equation}
At $(a,b)=(3,1)$, inversion of the part $r>1$ and an integration by parts reduce the remaining finite part to an
ordinary integral:
\begin{gather}
 B_{31}=\frac{\pi^2}{8}-\frac{\pi^2}{2}\int_0^1\frac{\dd r}{(1+r)^2}\Bigl(\frac1{1+3r^2}-\frac{r^2}{r^2+3}\Bigr),\notag\\
 \int_0^1\frac{\dd r}{(1+r)^2}\Bigl(\frac1{1+3r^2}-\frac{r^2}{r^2+3}\Bigr)=\frac{\pi\sqrt3}{48}+\frac{3\log3}{16},
\end{gather}
which gives $B_{31}=\frac{\pi^2}{96}(12-\pi\sqrt3-9\log3)$, the value quoted in Sec.~\ref{sec:d2}. It is the
combination $\Res A_4+b_0B_{31}$, not this source-dependent finite number alone, that determines the exponent.

%--------------------------------------------------------------------
\subsection{The regulated conformal map and the elliptic energy integral}\label{app:cform:ellipticmap}
\paragraph{Keeping the source factor through the coordinate change.} Put the four spins at the reference positions
$(0,z,1,\infty)$ and the energies at $(x,y)$. The M\"obius map
\begin{equation}
 M(t)=\Lambda\frac{t-x}{t-y},\qquad (w_1,w_2,w_3,w_4)=\bigl(M(0),M(z),M(1),M(\infty)\bigr),
 \label{eq:cform:mobius}
\end{equation}
returns the energies to $(0,\infty)$; $\Lambda$ is again an overall complex scale. Its complex Jacobian determinant is
\begin{equation}
 \det\frac{\partial(w_1,w_2,w_3,w_4)}{\partial(\Lambda,x,y,z)}=-\frac{\Lambda^3(x-y)^2}{y^2(y-z)^2(y-1)^2}.
\end{equation}
Hats denote correlators with the spins at these reference positions. In particular
$\widehat H_4=\langle\epsilon(x)\epsilon(y)\sigma(0)\sigma(z)\sigma(1)\sigma(\infty)\rangle$, and
$\widehat H_{12},\widehat H_{34}$ each include both energy insertions; their energy two-point function is not divided
out. The transformed connected kernel is
\begin{equation}
 \mathcal G(x,y;z)=\frac{\widehat H_4(x,y;z)}{|x-y|^2}-\frac{\widehat S_4(z)}{|x-y|^4}
 -\widehat H_{12}(x,y)\widehat H_{34}(x,y)+\frac{|z|^{-1/4}}{|x-y|^4},
 \label{eq:cform:G}
\end{equation}
where, from Eqs.~\eqref{eq:cform:neutral} and~\eqref{eq:cform:finiteenergy},
\begin{equation}
 \widehat H_{12}(x,y)=|z|^{-1/4}\Bigl|\frac1{(x-y)^2}+\frac{z^2}{4xy(x-z)(y-z)}\Bigr|,\qquad
 \widehat H_{34}(x,y)=\Bigl|\frac1{(x-y)^2}+\frac1{4(x-1)(y-1)}\Bigr|.
\end{equation}
Define the exact Gaussian factors
\begin{equation}
\begin{gathered}
 R(x,y,z)=\Bigl|\frac xy\Bigr|^2+\Bigl|\frac{z-x}{z-y}\Bigr|^2+\Bigl|\frac{1-x}{1-y}\Bigr|^2+1,\\
 \Upsilon(y,z)=|y(y-z)(y-1)|,\qquad
 W(x,y,z)=\frac{|x-y|^2}{\Upsilon(y,z)R(x,y,z)}.
\end{gathered}
 \label{eq:cform:W}
\end{equation}
Combining the conformal weights, the Wick kernels and the Jacobian gives the full density, apart from the factor $1/8$,
\begin{equation}
 |\Lambda|^{-2+4\delta}e^{-|\Lambda|^2R}\,|z|^{-15/4+2\delta}\Bigl(\frac{|x-y|^2}{\Upsilon}\Bigr)^{2\delta}\mathcal G.
\end{equation}
The scale integral $\int_{\mathbb C}|\Lambda|^{-2+4\delta}e^{-R|\Lambda|^2}\dd^2\Lambda=\pi\Gamma(2\delta)R^{-2\delta}$
therefore yields
\begin{equation}
 A_4(\delta)=\frac\pi8\Gamma(2\delta)J(\delta),\qquad
 J(\delta)=\AC\!\int_{\mathbb C}\dd^2z\,|z|^{-15/4+2\delta}\int_{\mathbb C^2}\dd^2x\dd^2y\,
 W(x,y,z)^{2\delta}\,\mathcal G(x,y;z).
 \label{eq:cform:J}
\end{equation}
This is an equality before the endpoint $\delta\to0$ is taken, and hence an equality of the prescribed meromorphic
continuations. The source factor $W^{2\delta}$ (the $W$ of Sec.~\ref{sec:d2}) is responsible for the matching contact
of Appendix~\ref{app:cform:contact}.

\paragraph{An explicit algebraic form of the lifted Ising kernel.} The energy positions are uniformized by the
elliptic curve $E_z$, the double cover of the $t$-plane branched at the zeros of $f_z(t)=t(t-z)(t-1)$ and at infinity,
with Abel coordinate $\dd u=\dd t/\sqrt{f_z(t)}$. The complete elliptic integrals use the parameter convention
\begin{equation}
 K(z)=\int_0^{\pi/2}\frac{\dd\theta}{\sqrt{1-z\sin^2\theta}},\qquad
 E(z)=\int_0^{\pi/2}\sqrt{1-z\sin^2\theta}\,\dd\theta,
\end{equation}
continued consistently away from the real unit interval. The periods of $E_z$ are $4K(z)$ and $4\mathrm iK(1-z)$; the
torus area $\mathcal A(z)$, the circular mean parameter $q(z)$, the Ising block weights $w_\pm,w_0$ and the constants
$c_\pm,c_0$ are those defined after Eq.~\eqref{eq:I}. In particular $\widehat S_4=w_++w_-$; changing the sign of
$\sqrt{1-z}$ interchanges the two block labels and leaves the sum invariant.

The passage from the exact Ising correlators to the elliptic kernel is an algebraic identity. Set $s_z=(1+z)/3$ and
\begin{equation}
 \Sigma(x,y;z)=-\frac{x+y-2s_z}{2}+\frac{f_z(x)+f_z(y)}{2(x-y)^2},\qquad
 D^2(x,y;z)=\frac{|f_z(x)f_z(y)|}{|x-y|^4}.
\end{equation}
Substituting Eq.~\eqref{eq:cform:neutral} and its finite-energy variant~\eqref{eq:cform:finiteenergy} into
Eq.~\eqref{eq:cform:G} gives
\begin{equation}
 \mathcal G=\frac{w_+|\Sigma+c_+|^2+w_-|\Sigma+c_-|^2-w_0|\Sigma+c_0|^2-(w_++w_--w_0)D^2}{|f_z(x)f_z(y)|}.
 \label{eq:cform:cover}
\end{equation}
Every quantity in this identity is explicit; no unknown conformal-block coefficient remains.

The uniformization is $x=4\wp(u)+s_z$, $y=4\wp(v)+s_z$, with $\wp$ the Weierstrass function of $E_z$. The addition
formulas give
\begin{equation}
 \Sigma(x,y;z)=\wp(u-v)+\wp(u+v),\qquad D^2(x,y;z)=|\wp(u-v)-\wp(u+v)|^2;
 \label{eq:cform:addition}
\end{equation}
for example $f_z(x)=16\wp'(u)^2$ and $\wp(u-v)-\wp(u+v)=\wp'(u)\wp'(v)/(\wp(u)-\wp(v))^2$. The denominator of
Eq.~\eqref{eq:cform:cover} is cancelled by $\dd^2x\dd^2y=|f_z(x)f_z(y)|\dd^2u\dd^2v$. Each energy plane has two
preimages on the curve, so its integral is one half of the integral over the cover, and the two energies give an
overall factor $1/4$.

\paragraph{Circular torus integration.} In the circular torus prescription
\begin{equation}
 \FP\!\int_{E_z}\wp(u)\dd^2u=-\mathcal A\,q.
 \label{eq:cform:wpmean}
\end{equation}
One way to obtain Eq.~\eqref{eq:cform:wpmean} is from the periodic completion
$\zeta_{\rm W}(u)-qu-(\pi/\mathcal A)\bar u$ of the Weierstrass zeta function $\zeta_{\rm W}$. For the half-period
$\omega_1=2K(z)$, with quasi-period $\eta_1=\zeta_{\rm W}(\omega_1)$, its periodicity gives
$q=\eta_1/\omega_1-\pi\bar\omega_1/(\omega_1\mathcal A)$, and the elliptic-integral expression for $\eta_1/\omega_1$
gives the formula for $q$ after Eq.~\eqref{eq:I}. Integrating the $u$ derivative of the completion over a torus with a
circular hole gives Eq.~\eqref{eq:cform:wpmean}, because the spin-two principal part has zero circular boundary
average.

Write $\wp_\mp=\wp(u\mp v)$. Subtracting the squared kernels before integrating gives
\begin{equation}
 |\wp_-+\wp_++c|^2-|\wp_--\wp_+|^2=4\,\mathrm{Re}(\wp_-\bar\wp_+)+2\,\mathrm{Re}\bigl[c(\bar\wp_-+\bar\wp_+)\bigr]+|c|^2.
\end{equation}
The map $(u,v)\mapsto(u-v,u+v)$ is a finite covering of $E_z^2$ and preserves the normalized Haar measure. With the
circular finite parts understood,
\begin{equation}
 \FP\!\int_{E_z^2}\wp_-\bar\wp_+\,\dd^2u\dd^2v=\mathcal A^2|q|^2,\qquad
 \FP\!\int_{E_z^2}\wp_-\,\dd^2u\dd^2v=\FP\!\int_{E_z^2}\wp_+\,\dd^2u\dd^2v=-\mathcal A^2q,
\end{equation}
and therefore
\begin{equation}
 \frac14\FP\!\int_{E_z^2}\dd^2u\dd^2v\,\bigl(|\wp_-+\wp_++c|^2-|\wp_--\wp_+|^2\bigr)=\frac{\mathcal A^2}{4}|c-2q|^2.
 \label{eq:cform:torus}
\end{equation}
Applying Eq.~\eqref{eq:cform:torus} to Eq.~\eqref{eq:cform:cover}, with $c=c_\pm,c_0$, gives
\begin{equation}
 \mathcal T(z)=\frac{\mathcal A(z)^2}{4}\Bigl[w_+|c_+-2q|^2+w_-|c_--2q|^2-w_0|c_0-2q|^2\Bigr],
\end{equation}
and $\cI=\FP\!\int\dd^2z\,|z|^{-15/4}\mathcal T(z)$ is the reduced period integral of Eq.~\eqref{eq:I}. This operation
defines the flat torus finite part. It is not yet the limit of the source-weighted integral $J(\delta)$ of
Eq.~\eqref{eq:cform:J}; the comparison follows.

%--------------------------------------------------------------------
\subsection{The finite contact}\label{app:cform:contact}
\paragraph{The four same-branch contacts.} Fix a nondegenerate modulus $z$. A single energy approaching a branch point
is locally integrable, and so are energies approaching distinct branch points. Coincident energies away from a branch
point leave a spin-two term with a radial leading regulator and zero scalar contact. The remaining case has both
energies approaching the same branch point. In Abel coordinates $u,v$ that vanish at that point, including the chart
at infinity, the exact source factor~\eqref{eq:cform:W} has the leading form
\begin{equation}
 W=\frac{|u-v|^2|u+v|^2|v|^2}{4(|u|^4+3|v|^4)}\bigl[1+O(|u|+|v|)\bigr].
 \label{eq:cform:branchW}
\end{equation}
With $\alpha_1=u-v$ and $\alpha_2=u+v$, the leading density, including the cover and linear-coordinate factors, is
\begin{equation}
 \frac{\widehat S_4-w_0}{4}\,\mathrm{Re}\frac1{\alpha_1^2\bar\alpha_2^2}\,\dd^2\alpha_1\dd^2\alpha_2.
\end{equation}
Now put $\alpha_1=\alpha_2\zeta$ and $\alpha_2=re^{\mathrm i\theta}$. Then
\begin{equation}
 W=r^2|\zeta|^2\frac{|1-\zeta|^2}{|1+\zeta|^4+3|1-\zeta|^4}+o(r^2).
\end{equation}
The radial integral up to an auxiliary radius $\rho$ is $\rho^{4\delta}/(4\delta)$. Its zeroth-order angular
coefficient is zero, but the first derivative of $W^{2\delta}$ survives and produces, per branch point, the contact
$A_{\rm br}[\widehat S_4-w_0]$ with
\begin{equation}
 A_{\rm br}=\frac\pi4\FP\!\int_{\mathbb C}\dd^2\zeta\,\mathrm{Re}(\zeta^{-2})
 \log\frac{|1-\zeta|^2}{|1+\zeta|^4+3|1-\zeta|^4}.
 \label{eq:cform:Abr}
\end{equation}
The term $2\log|\zeta|$ of $\log W$ has zero angular contribution. The auxiliary radius $\rho$ also drops out,
because it multiplies a vanishing zeroth-order residue. There are four branch points ($0$, $z$, $1$ and $\infty$); the
sheet factors are already included, so there is no further multiplicity.

\paragraph{The contact and the lower Gaussian integral.} The contact is related exactly to the integral
$A_{a,b}$ of Eq.~\eqref{eq:cform:Aab}. Let $B_{M1}=A_{M,1}(0)$ (so that $B_{31}$ is the case $M=3$), and let
$A_{\rm br}(M)$ be Eq.~\eqref{eq:cform:Abr} with
3 replaced by $M\ge0$. The distributional change of variables $\tilde w=(1+\zeta)/(1-\zeta)$ gives
\begin{equation}
 \FP\!\int_{\mathbb C}\dd^2\zeta\,\mathrm{Re}(\zeta^{-2})\,h\Bigl(\Bigl|\frac{1+\zeta}{1-\zeta}\Bigr|^2\Bigr)
 =8\,\FP\!\int_{\mathbb C}\dd^2z\,F_0(z)\,h(|z|).
 \label{eq:cform:contactdist}
\end{equation}
Indeed, the transformed bulk kernel is $4\,\mathrm{Re}[(\tilde w-1)^{-2}(\bar{\tilde w}+1)^{-2}]$. Its angular
integral, in the variable $\varrho=|\tilde w|^2$, is $4\pi|\varrho-1|/(1+\varrho)^3$, with local term $-2\pi h(1)$; this
is eight times Eq.~\eqref{eq:cform:radialfp}. Differentiating $A_{\rm br}(M)$ and $4B_{M1}$ in $M$, and using
Eq.~\eqref{eq:cform:contactdist} and inversion, gives equal derivatives. The endpoint values $A_{\rm br}(0)=\pi^2/4$
and $B_{01}=\pi^2/8$ (from Eq.~\eqref{eq:cform:Aab0}) fix the constant, $A_{\rm br}(M)=4B_{M1}-\pi^2/4$, which at $M=3$ is
the relation $A_{\rm br}=4B_{31}-\pi^2/4$ of Sec.~\ref{sec:d2}. These $M=0$ values are the prescribed continued limits;
no ordinary infrared convergence is asserted for the unweighted limit.

\paragraph{The fixed-modulus comparison.} For $z\notin\{0,1,\infty\}$ the comparison is therefore
\begin{equation}
 \lim_{\delta\to0}\AC\!\int\dd^2x\dd^2y\,W^{2\delta}\mathcal G=\mathcal T(z)+4A_{\rm br}\bigl[\widehat S_4(z)-w_0(z)\bigr],
 \label{eq:cform:fixedmod}
\end{equation}
with $\mathcal T(z)$ the torus integrand of Appendix~\ref{app:cform:ellipticmap}. The source-dependent first derivative
of the regulator is necessary even though the corresponding pointwise zeroth-order angular residue vanishes.

\paragraph{The external-primary cancellation.} Extending Eq.~\eqref{eq:cform:fixedmod} to the modulus boundaries uses a
stronger fact than the integrated equality $I_2(0)=0$. For fixed noncoincident $a,b,c,d$ (energies at $0$ and
$\infty$), the exact simultaneous OPE gives
\begin{equation}
\begin{aligned}
 H_4(\lambda a,\lambda b,c,d)&=\lambda^{-5/4}\frac{|a-b|^{3/4}|c-d|^{3/4}}{4|ab|}
 +\lambda^{-1/4}H_2(a,b)H_2(c,d)+O(\lambda^{3/4}),\\
 S_4(\lambda a,\lambda b,c,d)&=\lambda^{-1/4}S_2(a,b)S_2(c,d)+O(\lambda^{3/4}).
\end{aligned}
\end{equation}
Consequently
\begin{equation}
 \mathcal F(\lambda a,\lambda b,c,d)=\lambda^{-5/4}\frac{|a-b|^{3/4}|c-d|^{3/4}}{4|ab|}+O(\lambda^{3/4}):
 \label{eq:cform:external}
\end{equation}
the entire $\lambda^{-1/4}$ energy-primary term cancels in the connected logarithm. Multiplying by the internally
contracted generalized-free pair changes the leading degree to $-5+2\delta$ in four real coordinates, which produces
$\dd\lambda\,\lambda^{-2+2\delta}$. Its Mellin denominator is $-1+2\delta$, which does not vanish at the endpoint. The
absent term would have given $\dd\lambda\,\lambda^{-1+2\delta}$, a logarithm. The cancellation persists for nonzero
$\delta$, since the same generalized-free kernel multiplies all terms of $\mathcal F$. If the shrinking pair is not
internally contracted in the fixed pairing, the remaining generalized-free fields raise the output dimension beyond
this threshold.

%--------------------------------------------------------------------
\subsection{Modulus boundaries, subtractions and the qualification}\label{app:cform:qualification}
\paragraph{Thresholds at the modulus boundaries.} In the fixed pairing the contracted degeneration is $z\to0$; the
pair degenerations at $z\to1$ and $z\to\infty$ lack this internal generalized-free contraction. At $r=|z|\to0$, let
$\Delta$ be the exchanged Ising dimension. Including the measures of all energy coordinates that shrink with the pair
gives the scalar thresholds
\begin{center}\small
\begin{tabular}{@{}lll@{}}
\toprule
shrinking cluster & radial measure & logarithmic threshold\\
\midrule
two spins, neither energy & $\dd r\,r^{\Delta-3+2\delta}$ & $\Delta=2$\\
two spins and one energy & $\dd r\,r^{\Delta-2+2\delta}$ & $\Delta=1$\\
two spins and both energies & $\dd r\,r^{\Delta-3+2\delta}$ & $\Delta=2$\\
\bottomrule
\end{tabular}
\end{center}
The disconnected identity terms cancel in the connected expression; an identity-channel power term with denominator
$-1+2\delta$ is retained and analytically continued. The energy primary of the middle row cancels by
Eq.~\eqref{eq:cform:external}. At dimension two, the even Ising identity and energy families contain descendants of
nonzero spin but no scalar. The remaining thermal powers are restored by Eq.~\eqref{eq:cform:mellin}; their regular
denominators must not be dropped.

\paragraph{Spin residues and coupling logarithms.} The regulator also has to preserve the vanishing spin residues. In
the joint corner $x=z\xi$, $y=z\eta$ its leading term is
\begin{equation}
 W\longrightarrow\frac{|\xi-\eta|^2}{|\eta(\eta-1)|\bigl(|\xi/\eta|^2+|(1-\xi)/(1-\eta)|^2+2\bigr)}.
\end{equation}
This is independent of $\arg z$, so a leading term of nonzero spin keeps zero angular residue for nonzero $\delta$. The
corners with zero and with one energy have the same property. The only proper scalar coupling logarithm left is a
cluster of three insertions of $\sigma\chi$, which is already represented by the counterterm~\eqref{eq:cform:ct}. Two
distinct triples among four perturbing fields overlap, so they cannot both be separate proper nested or disjoint
members of one forest. Thermal pairs inside a triple are included in the continued lower-order coupling projection.

\paragraph{The global comparison and its unproved step.} The global comparison uses local OPE and Taylor subtractions
in each nested sector, in the distribution-extension framework of Ref.~\cite{Brunetti:2000}: subtract the
nonintegrable powers and spin terms, restore the powers by their Mellin continuation and the spin terms by their
circular finite parts, and then interchange the regulator limit with the subtracted remainder integral. This last
step requires an integrable dominating bound, uniform near $\delta=0$, including the degenerating tangent shapes at
intersections of sectors. The study gives the channel classification above and a general subtraction argument for
this step. An audit stage that is part of the study itself (same authorship; not the review of
Sec.~\ref{sec:sources}) checked the exact algebra, the local contacts and representative
simultaneous limits, but did not reconstruct every sectorwise subtraction or prove the uniform remainder bound. This
is the one unproved step of Sec.~\ref{sec:unproved}, and the identification of Eq.~\eqref{eq:c2closed} with the
physical coefficient is a formal perturbative result in the critically tuned analytic minimal-subtraction prescription,
conditional on it. The study adds that the qualification is not evidence for a missing contact, and that neither a
fixed-modulus check nor a high-precision evaluation of $\cI$ can replace the bound: they establish different parts of
the method.

%--------------------------------------------------------------------
\subsection{Assembly}\label{app:cform:assembly}
Under the global comparison just described, Eq.~\eqref{eq:cform:fixedmod} is integrated over $z$ with the weight
$|z|^{-15/4}$. The term $4A_{\rm br}[\widehat S_4-w_0]$ then gives $4A_{\rm br}I_\beta$, with the coupling integral of
Eq.~\eqref{eq:cform:Ibeta},
$I_\beta=\AC\!\int\dd^2z\,|z|^{-15/4}[\widehat S_4(z)-|z|^{-1/4}]=-b_0/\pi$. With Eq.~\eqref{eq:cform:J}, in which
$\Gamma(2\delta)=1/(2\delta)+O(1)$, and the contact relation of Appendix~\ref{app:cform:contact},
\begin{equation}
\begin{aligned}
 J(0)&=\cI+4A_{\rm br}I_\beta=\cI-\frac{16b_0B_{31}}{\pi}+\pi b_0,\\
 \Res A_4&=\frac\pi{16}J(0)=\frac\pi{16}\cI-b_0B_{31}+\frac{\pi^2b_0}{16},
\end{aligned}
\end{equation}
and $z_{4,1}=\Res A_4+b_0B_{31}$ is Eq.~\eqref{eq:c2I}. The finite source constant $B_{31}$ cancels between the bare
residue and its source counterterm. The remaining $-\pi^2/(4b_0)$ in $c_2$ is the finite matching contribution; its
size is comparable to that of the period term (Sec.~\ref{sec:d2}). The evaluation that follows concerns $\cI$ alone;
it supplies an explicit value once this renormalization comparison has been specified.

%--------------------------------------------------------------------
\subsection{Complex moments and their reduction to interval periods}\label{app:cform:moments}
From here to Appendix~\ref{app:cform:telescoping} the target is $\cI$ in
$c_2=-\pi\cI/(4b_0^2)-\pi^2/(4b_0)$, with $b_0=9\pi^2/(8R_\Gamma^2)$ and $R_\Gamma=\Gamma(\frac14)/\Gamma(\frac34)$. The
continuation and finite-part prescriptions remain those of Eq.~\eqref{eq:I}. In particular, an integral with a negative
endpoint exponent denotes its meromorphic continuation, not an ordinary improper integral.

\paragraph{Interval periods and complex moments.} With $K(t)$ of parameter $t$, put
\begin{equation}
\begin{gathered}
 \mathbf k(t)=\begin{pmatrix}K(t)^2\\K(t)K(1-t)\\K(1-t)^2\end{pmatrix},\qquad
 \mathsf J=\begin{pmatrix}0&0&1\\0&2&0\\1&0&0\end{pmatrix},\\
 \mathbf P(a,b)=\AC\!\int_0^1t^{a-1}(1-t)^{b-1}\mathbf k(t)\dd t,\qquad
 \mathbf p(a,b)=\frac4{\pi^2}\mathbf P(a,b).
\end{gathered}
 \label{eq:cform:periods}
\end{equation}
The vector $\mathbf k$ reverses under $t\mapsto1-t$; in particular the third component of $\mathbf p(a,b)$ is the first
component of $\mathbf p(b,a)$. Write
\begin{equation}
 \mathbf p(a,b)=\bigl(\mathcal U(a,b),\,\mathcal V(a,b),\,\mathcal U(b,a)\bigr)^{\mathsf T}.
\end{equation}
The complex area moment is
\begin{equation}
 \mathcal M(a,b)=\AC\!\int_{\mathbb C}\dd^2t\,|t|^{2a-2}|1-t|^{2b-2}\mathcal A(t)^2,
\end{equation}
and $\mathcal B(a,b)$ is Eq.~\eqref{eq:cform:beta}. The single-valued period pairing uses the elliptic connection
formulas~\cite{NIST:DLMF} and the bilinear contour construction of Ref.~\cite{Duhr:2023}. In the present
normalization it reads
\begin{equation}
 \mathcal M(a,b)=64\,\mathbf P(a,b)^{\mathsf T}\mathsf L(b)\,\mathbf P(2-a-b,b),\qquad
 \mathsf L(b)=\begin{pmatrix}0&0&\sin\pi b\\0&2\sin\pi b&-2\cos\pi b\\\sin\pi b&-2\cos\pi b&-\sin\pi b\end{pmatrix}.
 \label{eq:cform:pairing}
\end{equation}

\paragraph{Contour conventions.} On $0<t<1$ both $K(t)$ and $K(1-t)$ are real and positive. They are continued from this
interval to the slit plane, with the upper and lower boundary values used separately. For real $t>1$ they obey
\begin{equation}
 K(t\pm\mathrm i0)=t^{-1/2}\bigl[K(1/t)\pm\mathrm iK(1-1/t)\bigr],\qquad K(1-t)=t^{-1/2}K(1-1/t),
\end{equation}
so that
\begin{equation}
 \mathbf k(t\pm\mathrm i0)=t^{-1}\mathsf C_\pm\mathbf k(1/t),\qquad
 \mathsf C_\pm=\begin{pmatrix}1&\pm2\mathrm i&-1\\0&1&\pm\mathrm i\\0&0&1\end{pmatrix}.
 \label{eq:cform:boundary}
\end{equation}
These are the symmetric-square connection matrices of the elliptic equation. Moreover
$\mathcal A(t)^2=64\,\mathbf k(t)^{\mathsf T}\mathsf J\,\overline{\mathbf k(t)}$. To obtain
Eq.~\eqref{eq:cform:pairing}, take a primitive of the holomorphic one-form $t^{a-1}(1-t)^{b-1}\mathbf k(t)\dd t$ in
this single-valued pairing and apply Stokes' theorem on the cut plane. The two banks of the cut $[1,\infty)$ give
\begin{equation}
 \frac{e^{\mathrm i\pi b}\mathsf J\mathsf C_--e^{-\mathrm i\pi b}\mathsf J\mathsf C_+}{2\mathrm i}=\mathsf L(b).
\end{equation}
The substitution $t=1/\tau$ on that cut changes the power to $\tau^{1-a-b}(1-\tau)^{b-1}$ and so gives the second
interval period $\mathbf P(2-a-b,b)$. This accounts for the parameter $2-a-b$ and for every entry of $\mathsf L(b)$.
The contour identity is first applied with convergent endpoint regulators and then continued meromorphically;
small-circle terms at nonintegrable powers are continued with the same prescription, not discarded. This is the
standard single-valued period-pairing step of the calculation, not a separate numerical integration of the four-point
correlator.

%--------------------------------------------------------------------
\subsection{Landen substitution and the six complex moments}\label{app:cform:landen}
The quadratic substitution and the area Wronskian are
\begin{equation}
 z=\frac{4w}{(1+w)^2},\quad t=w^2,\qquad
 \mathcal A\Bigl(\frac{4w}{(1+w)^2}\Bigr)=\frac{|1+w|^2}{2}\mathcal A(w^2),\qquad
 |\partial_t\mathcal A|^2=\frac14\partial_t\partial_{\bar t}(\mathcal A^2)+\frac{2\pi^2}{|t(1-t)|^2}.
 \label{eq:cform:landen}
\end{equation}
They are inserted in the three-term integrand of Eq.~\eqref{eq:I} as follows. The derivative of the area and
Legendre's relation give
\begin{equation}
 q(z)=\frac{z(1-z)}{2}\frac{\partial_z\mathcal A(z)}{\mathcal A(z)}+\frac{1-2z}{12}.
\end{equation}
For $t=w^2$, abbreviate (in this subsection only) $A=\mathcal A(t)$ and $A_t=\partial_tA$. Using
$\sqrt{1-z}=(1-w)/(1+w)$ on a consistent sheet, the three factors of the integrand become
\begin{equation}
\begin{aligned}
 \mathcal A(z)(c_+-2q)&=-t(1-w)(1+\bar w)A_t,\\
 \mathcal A(z)(c_--2q)&=-(1-w)(1+\bar w)(A/2+tA_t),\\
 \mathcal A(z)(c_0-2q)&=\frac{t(1+\bar w)}{1+w}\bigl[A/2-(1-t)A_t\bigr].
\end{aligned}
\end{equation}
The exchange of the two square-root sheets interchanges the first two and leaves their sum unchanged. The Jacobian
$|\dd z/\dd w|^2=16|1-w|^2/|1+w|^6$ and the multiplicities of the degree-two maps $w\mapsto z$ ($w$ and $1/w$ have the
same image) and $w\mapsto t$ give
\begin{equation}
 \cI=\frac1{256}\AC\!\int_{\mathbb C}\dd^2t\,\biggl\{\frac{|1-t|^{7/2}}{|t|}|A_t|^2
 +\frac{|1-t|^{7/2}}{|t|^{5/2}}|A/2+tA_t|^2-\frac{|1-t|^2}{|t|}|A/2-(1-t)A_t|^2\biggr\}.
 \label{eq:cform:landenI}
\end{equation}
Under inversion $\mathcal A(1/t)=|t|\mathcal A(t)$, so the first two integrals are equal. Put
\begin{equation}
 \Omega_1=|t|^{-1}|1-t|^{7/2},\qquad\Omega_0=|t|^{-1}|1-t|^2.
\end{equation}
The two remaining integrals are evaluated by integration by parts. For the first, the Wronskian identity gives
\begin{equation}
 \AC\!\int\Omega_1|A_t|^2\dd^2t=\AC\!\int\Bigl[\frac14(\partial_t\partial_{\bar t}\Omega_1)A^2
 +\frac{2\pi^2\Omega_1}{|t(1-t)|^2}\Bigr]\dd^2t,\qquad
 \frac14\partial_t\partial_{\bar t}\Omega_1=\frac{|t|^{-3}|1-t|^{3/2}}{64}|2+5t|^2,
\end{equation}
with $|2+5t|^2=14+35|t|^2-10|1-t|^2$. In the third integral the cross term is also integrated by parts. The
resulting coefficient of $A^2$ is
\begin{equation}
 \frac{\Omega_0}4+\frac14\Bigl\{\partial_t\bigl[\Omega_0(1-t)\bigr]+\partial_{\bar t}\bigl[\Omega_0(1-\bar t)\bigr]\Bigr\}
 +\frac14\partial_t\partial_{\bar t}\bigl[\Omega_0|1-t|^2\bigr]
 =\frac{|t|^{-3}|1-t|^2}{16}|1+t|^2,
\end{equation}
with $|1+t|^2=2+2|t|^2-|1-t|^2$, and its Wronskian term is $2\pi^2\mathcal B(-\frac12,2)=0$, a zero of the reciprocal
Gamma function $1/\Gamma(1-b)$ in Eq.~\eqref{eq:cform:beta}. These integrations by parts are justified by first
replacing the coefficient of $|A_t|^2$ by an auxiliary weight $|t|^{2a}|1-t|^{2b}$ with $\mathrm{Re}\,a>0$,
$\mathrm{Re}\,b>0$ and $\mathrm{Re}(a+b)<1$, where the boundary terms vanish, and then continuing to the stated
exponents. They therefore preserve the finite-part boundary prescription. The result, written so as to keep every
endpoint contribution, is
\begin{equation}
\begin{aligned}
 \cI&=\frac{\mathcal D_1}{128}-\frac{\mathcal D_0}{1024},\\
 \mathcal D_1&=-\frac5{32}\mathcal M(-\tfrac12,\tfrac{11}4)+\frac{35}{64}\mathcal M(\tfrac12,\tfrac74)
 +\frac7{32}\mathcal M(-\tfrac12,\tfrac74)+2\pi^2\mathcal B(-\tfrac12,\tfrac74),\\
 \mathcal D_0&=\frac{\mathcal M(\frac12,2)+\mathcal M(-\frac12,2)}2-\frac14\mathcal M(-\tfrac12,3).
\end{aligned}
 \label{eq:cform:six}
\end{equation}
The derivative reduction thus gives six complex area moments and one elementary complex beta function, which
evaluates exactly to $\mathcal B(-\frac12,\frac74)=-32b_0/\pi$.

\paragraph{The vanishing of $\mathcal D_0$.} Behind $\mathcal D_0=0$ is an endpoint identity. The symmetric-square
elliptic equation gives
\begin{equation}
 \mathbf P(-\tfrac12,2)+\mathbf P(\tfrac12,2)=(-4,0,0)^{\mathsf T},
 \label{eq:cform:boundaryperiod}
\end{equation}
an identity of continued periods that includes the Mellin boundary term at $t=0$. It has an explicit differential
certificate. For any component $y$ of $\mathbf k(t)$ define
\begin{equation}
 \mathcal C[y]=-8t^{3/2}(1-t)^2y''+4\sqrt t(1-t)(5t-3)y'+\frac{2(1-t)(3t-1)}{\sqrt t}y.
\end{equation}
The symmetric-square differential equation,
$y'''+3\rho_1y''+(\rho_1'+2\rho_1^2+4\rho_0)y'+(2\rho_0'+4\rho_1\rho_0)y=0$ with $\rho_1=(1-2t)/[t(1-t)]$ and
$\rho_0=-1/[4t(1-t)]$, gives, by direct substitution,
\begin{equation}
 \frac{\dd}{\dd t}\mathcal C[y]=t^{-3/2}(1-t^2)\,y.
\end{equation}
At $t=0$ the continued constant boundary term is zero. At $t=1$ the three boundary values of $\mathcal C[y]$ are
$-4,0,0$, from $K(t)\sim\frac12\log[16/(1-t)]$ and $K(1-t)\sim\pi/2$. This proves Eq.~\eqref{eq:cform:boundaryperiod}
with its finite endpoint term. To see how it enters, let $\mathbf u=\mathbf P(-\frac12,2)$ and
$\mathbf v=\mathbf P(\frac12,2)$. Since $\mathsf L(2)(\mathbf u+\mathbf v)=0$, $\mathsf L(3)=-\mathsf L(2)$ and
$\mathbf P(-\frac12,3)=\mathbf u-\mathbf v$ (by Eq.~\eqref{eq:cform:contiguous} below), the pairing gives
\begin{equation}
 \mathcal M(-\tfrac12,3)=4\mathcal M(\tfrac12,2)=4\mathcal M(-\tfrac12,2),\qquad\mathcal D_0=0.
\end{equation}

\paragraph{Reduction to four interval periods.} Abbreviate
\begin{equation}
 \mathbf X=\mathbf P(-\tfrac12,\tfrac74),\quad\mathbf Y=\mathbf P(\tfrac12,\tfrac74),\quad
 \mathbf Z=\mathbf P(-\tfrac14,\tfrac74),\quad\mathbf W=\mathbf P(\tfrac34,\tfrac74).
\end{equation}
The contiguous identity
\begin{equation}
 \mathbf P(a,b+1)=\mathbf P(a,b)-\mathbf P(a+1,b)
 \label{eq:cform:contiguous}
\end{equation}
and $\mathsf L(\frac{11}4)=-\mathsf L(\frac74)$ turn the part of $\mathcal D_1$ without the beta function into
\begin{equation}
 10(\mathbf X-\mathbf Y)^{\mathsf T}\mathsf L(\tfrac74)(\mathbf Z-\mathbf W)+35\,\mathbf Y^{\mathsf T}\mathsf L(\tfrac74)\mathbf Z
 +14\,\mathbf X^{\mathsf T}\mathsf L(\tfrac74)\mathbf W
 =(2\mathbf X+5\mathbf Y)^{\mathsf T}\mathsf L(\tfrac74)(5\mathbf Z+2\mathbf W).
\end{equation}
Consequently
\begin{equation}
 \cI=\frac1{128}(2\mathbf X+5\mathbf Y)^{\mathsf T}\mathsf L(\tfrac74)(5\mathbf Z+2\mathbf W)-\frac{\pi b_0}{2}.
 \label{eq:cform:Iperiod}
\end{equation}
The constant here is the complex-beta term of $\mathcal D_1$. It is separate from the finite regulator-matching term
$-\pi^2/(4b_0)$ of Eq.~\eqref{eq:c2I}.

%--------------------------------------------------------------------
\subsection{Reduction to symmetric elliptic moments}\label{app:cform:symmetric}
Use normalized period vectors and set
\begin{equation}
 \mathbf a=2\mathbf p(-\tfrac12,\tfrac74)+5\mathbf p(\tfrac12,\tfrac74),\quad
 \mathbf b=5\mathbf p(-\tfrac14,\tfrac74)+2\mathbf p(\tfrac34,\tfrac74),\quad
 \mathsf H=\begin{pmatrix}0&0&1\\0&2&2\\1&2&-1\end{pmatrix},\quad
 \mathsf L(\tfrac74)=-\frac{\mathsf H}{\sqrt2}.
\end{equation}
Equations~\eqref{eq:c2I} and~\eqref{eq:cform:Iperiod} give
\begin{equation}
 c_2=-\frac{R_\Gamma^2}{9}+\frac{\pi R_\Gamma^4}{10368\sqrt2}\,\mathbf a^{\mathsf T}\mathsf H\,\mathbf b.
 \label{eq:cform:c2matrix}
\end{equation}

\paragraph{Contour reflection.} In the same boundary convention, contour reflection of these particular interval
periods gives
\begin{equation}
 \mathbf p(-\tfrac12,\tfrac74)=\mathsf D\,\mathbf p(\tfrac34,\tfrac74),\qquad
 \mathbf p(\tfrac12,\tfrac74)=\mathsf D\,\mathbf p(-\tfrac14,\tfrac74),\qquad
 \mathsf D=\frac1{\sqrt2}\begin{pmatrix}1&-2&-1\\-1&-1&0\\0&-2&-1\end{pmatrix}.
 \label{eq:cform:reflection}
\end{equation}
These relations use the same connection matrices and continued endpoint contours as Eq.~\eqref{eq:cform:pairing};
they are not identities obtained by naive substitution in divergent real integrals. They can be derived without
assuming $\mathsf D$. For real $a,b>0$ with $a'=2-a-b>0$, integrate $t^{a-1}(1-t)^{b-1}\mathbf k(t)\dd t$ along the
real axis, closed in the upper half-plane; the small and large arcs vanish in this domain of convergence. On the
negative interval set $t=-\tau/(1-\tau)$; its connection matrix is
\begin{equation}
 \mathsf B_-=\begin{pmatrix}1&0&0\\-\mathrm i&1&0\\-1&-2\mathrm i&1\end{pmatrix}.
\end{equation}
The three real intervals give the vector identity
\begin{equation}
 \mathbf P(a,b)=e^{\mathrm i\pi a}\mathsf B_-\mathbf P(a,a')+e^{-\mathrm i\pi b}\mathsf C_+\mathbf P(a',b),
\end{equation}
which continues to the parameters of Eq.~\eqref{eq:cform:reflection}. Let $\mathsf A_c=e^{\mathrm i\pi a}\mathsf B_-$ and
$\mathsf E_c=e^{-\mathrm i\pi b}\mathsf C_+$. The left side and the two interval vectors are real, so eliminating
$\mathbf P(a,a')$ between the imaginary and real parts gives
\begin{equation}
 \mathbf P(a,b)=\bigl[\mathrm{Re}\,\mathsf E_c-\mathrm{Re}\,\mathsf A_c(\mathrm{Im}\,\mathsf A_c)^{-1}\mathrm{Im}\,\mathsf E_c\bigr]
 \mathbf P(a',b).
\end{equation}
For $a=\pm\frac12$ and $b=\frac74$, $e^{\mathrm i\pi a}=\pm\mathrm i$ and $e^{-\mathrm i\pi b}=(1+\mathrm i)/\sqrt2$, and the
bracket evaluates exactly to $\mathsf D$ for either sign. This proves Eq.~\eqref{eq:cform:reflection}, and hence
$\mathbf a=\mathsf D\,\mathbf b$.

\paragraph{Symmetric moments.} Define
\begin{equation}
 U(a)=\mathcal U(a,a),\qquad V(a)=\mathcal V(a,a),\qquad
 m_U=\tfrac52U(-\tfrac14)-4U(\tfrac34),\qquad m_V=\tfrac52V(-\tfrac14)-4V(\tfrac34).
 \label{eq:cform:mUV}
\end{equation}
For $a=-\frac14$ the weight in $\mathbf b$ is $[t(1-t)]^{a-1}\varpi(t)$, where
\begin{equation}
 \varpi(t)=(5+2t)(1-t)^2,\quad\frac{\varpi(t)+\varpi(1-t)}2=\frac52-4t(1-t),\quad
 \varpi(t)-\varpi(1-t)=(1-2t)\{5+2t(1-t)\}.
\end{equation}
For the components $[\mathbf b]_i$ of $\mathbf b$ it follows directly that $([\mathbf b]_1+[\mathbf b]_3)/2=m_U$ and
$[\mathbf b]_2=m_V$.

\paragraph{The antisymmetric moment.} The remaining antisymmetric moment has an exact Gamma evaluation:
\begin{equation}
 O(a):=\mathcal U(a,a+1)-\mathcal U(a+1,a)=-\frac{2^{1-2a}}{\sqrt\pi}\Bigl[\frac{\Gamma(a)}{\Gamma(a+\frac12)}\Bigr]^3.
 \label{eq:cform:odd}
\end{equation}
The following derivation also fixes its sign. Set $\hat K(t)=2K(t)/\pi$ and $x=4t(1-t)$ on $0<t<\frac12$, and
introduce the Meijer function
\begin{equation}
 G_0(x)=G^{3,0}_{3,3}\!\left(x\,\middle|\,\begin{matrix}\frac12,\frac12,\frac12\\0,0,0\end{matrix}\right),\qquad
 \int_0^1x^{s-1}G_0(x)\dd x=\frac{\Gamma(s)^3}{\Gamma(s+\frac12)^3}.
 \label{eq:cform:G0}
\end{equation}
The elliptic symmetric-square equation $[\vartheta^3-x(\vartheta+\frac12)^3]Y=0$, with $\vartheta=x\partial_x$, together
with its endpoint normalization, gives
\begin{equation}
 \hat K(1-t)^2-\hat K(t)^2=\frac2{\sqrt\pi}G_0\bigl(4t(1-t)\bigr).
 \label{eq:cform:compdiff}
\end{equation}
The normalization can be checked directly, without an unspecified constant. The triple pole at $s=0$ of the Mellin
transform~\eqref{eq:cform:G0} gives
\begin{equation}
 \frac2{\sqrt\pi}G_0(x)=\frac1{\pi^2}\log^2\frac{64}{x}-1+O(x\log^2x),
\end{equation}
which agrees with $\hat K(1-t)^2-\hat K(t)^2$ because $t=x/4+O(x^2)$, $\hat K(t)=1+O(t)$ and
$\hat K(1-t)=\pi^{-1}\log(16/t)+O(t\log t)$. The coefficients of the squared logarithm, the logarithm and the constant
fix the three local solutions of the third-order equation. Folding the odd moment about $\frac12$ gives
\begin{equation}
 O(a)=\int_0^{1/2}[t(1-t)]^{a-1}(1-2t)\bigl\{\hat K(t)^2-\hat K(1-t)^2\bigr\}\dd t
 =-\frac{2\cdot4^{-a}}{\sqrt\pi}\int_0^1x^{a-1}G_0(x)\dd x,
\end{equation}
which proves Eq.~\eqref{eq:cform:odd} where the integrals converge, and then by meromorphic continuation. With the
decomposition of $\varpi$,
\begin{equation}
 d_O:=[\mathbf b]_1-[\mathbf b]_3=5O(-\tfrac14)+2O(\tfrac34)=\frac{576\sqrt2}{\sqrt\pi R_\Gamma^3},\qquad
 \mathbf b=\begin{pmatrix}m_U+d_O/2\\m_V\\m_U-d_O/2\end{pmatrix}.
\end{equation}
Direct matrix multiplication, with no further integral identity, gives
\begin{equation}
 \frac{(\mathsf D\mathbf b)^{\mathsf T}\mathsf H\,\mathbf b}{\sqrt2}=\frac{d_O^2}{4}-(m_U+m_V)(m_U+3m_V),
\end{equation}
and therefore
\begin{equation}
 c_2=\frac{16}{R_\Gamma^2}-\frac{R_\Gamma^2}{9}-\frac{\pi R_\Gamma^4}{10368}(m_U+m_V)(m_U+3m_V).
 \label{eq:cform:c2fg}
\end{equation}
The term $16/R_\Gamma^2$ comes from the evaluated antisymmetric period. Its origin differs from that of both constants
of the residue and complex-beta reductions.

%--------------------------------------------------------------------
\subsection{Ordinary hypergeometric evaluation of the moments}\label{app:cform:hypergeometric}
The generalized hypergeometric series is~\cite{NIST:DLMF}
\begin{equation}
 {}_pF_q\!\left(\begin{matrix}a_1,\ldots,a_p\\b_1,\ldots,b_q\end{matrix};x\right)
 =\sum_{n=0}^\infty\frac{(a_1)_n\cdots(a_p)_n}{(b_1)_n\cdots(b_q)_n}\frac{x^n}{n!},\qquad(a)_n=\frac{\Gamma(a+n)}{\Gamma(a)},
\end{equation}
the convention of $\mathsf P,\mathsf Q,\mathsf T$ in Sec.~\ref{sec:closed2}. Define
\begin{equation}
 \Phi_A(a)=\pFq{\frac12,\frac12,\frac12,a}{1,1,a+\frac12}{1},\qquad
 \Phi_D(a)=\pFq{\frac12,a,a,a}{a+\frac12,a+\frac12,a+\frac12}{1}.
\end{equation}
Both series have parametric excess one at every nonresonant parameter used below. The symmetric moments are
\begin{align}
 U(a)&=2^{1-2a}\biggl[\sqrt\pi\frac{\Gamma(a)}{\Gamma(a+\frac12)}\Phi_A(a)
 +\frac1{\sqrt\pi}\Bigl(\frac{\Gamma(a)}{\Gamma(a+\frac12)}\Bigr)^3\Phi_D(a)\biggr],
 \label{eq:cform:U}\\
 V(a)&=2^{1-2a}\biggl[\sqrt\pi\frac{\Gamma(a-\frac12)}{\Gamma(a)}\Phi_A(1-a)
 +\frac{\tan\pi a}{\sqrt\pi}\Bigl(\frac{\Gamma(a)}{\Gamma(a+\frac12)}\Bigr)^3\Phi_D(a)\biggr].
 \label{eq:cform:V}
\end{align}
These are meromorphic identities in $a$. At $a=-\frac14$ and $a=\frac34$ no limiting cancellation of Gamma poles is
needed in the displayed summands.

\paragraph{The square moment.} The quadratic transformation and Clausen's identity give
\begin{equation}
 \hat K(t)^2={}_3F_2\!\left(\begin{matrix}\frac12,\frac12,\frac12\\1,1\end{matrix};x\right),\qquad x=4t(1-t),\quad0<t<\tfrac12.
\end{equation}
Combine this with Eq.~\eqref{eq:cform:compdiff} and fold the interval. Since
$[t(1-t)]^{a-1}\dd t=4^{-a}x^{a-1}(1-x)^{-1/2}\dd x$,
\begin{equation}
 U(a)=2\cdot4^{-a}\biggl[\int_0^1x^{a-1}(1-x)^{-1/2}\,{}_3F_2\!\left(\begin{matrix}\frac12,\frac12,\frac12\\1,1\end{matrix};x\right)\dd x
 +\frac1{\sqrt\pi}\int_0^1x^{a-1}(1-x)^{-1/2}G_0(x)\dd x\biggr].
\end{equation}
The first integral is the Euler beta integral for $\Phi_A(a)$. Expanding $(1-x)^{-1/2}$ and applying the Mellin
transform~\eqref{eq:cform:G0} to each term gives $\Phi_D(a)$ for the second. This proves Eq.~\eqref{eq:cform:U}.

\paragraph{The mixed moment.} Its Mellin--Barnes representation follows from the inverse Mellin formula, on $0<x<1$,
\begin{equation}
 \hat K(t_x)\hat K(1-t_x)=\frac1{2\pi\mathrm i\sqrt\pi}\int_{c-\mathrm i\infty}^{c+\mathrm i\infty}
 \frac{\Gamma(s)^2\Gamma(\frac12-s)}{\Gamma(s+\frac12)^2\Gamma(1-s)}x^{-s}\dd s,\qquad
 t_x=\frac{1-\sqrt{1-x}}2,\quad0<c<\tfrac12.
 \label{eq:cform:mixedmellin}
\end{equation}
One way to verify it: its Mellin kernel $m(s)$ satisfies $m(s+1)/m(s)=s^3/(s+\frac12)^3$, so it solves the same
symmetric-square equation. The double pole at zero has principal part
\begin{equation}
 \frac{\Gamma(s)^2\Gamma(\frac12-s)}{\sqrt\pi\,\Gamma(s+\frac12)^2\Gamma(1-s)}
 =\frac1\pi\Bigl[\frac1{s^2}+\frac{6\log2}{s}+O(1)\Bigr].
\end{equation}
The resulting local behaviour is $\pi^{-1}\log(64/x)$ with no squared logarithm, exactly that of the left side, and
these Frobenius data determine the identity. Folding the interval in $V(a)$ and inserting
Eq.~\eqref{eq:cform:mixedmellin} adds the Euler beta factor $B(a-s,\frac12)=\sqrt\pi\,\Gamma(a-s)/\Gamma(a+\frac12-s)$.
Consequently
\begin{equation}
 V(a)=2\cdot4^{-a}\frac1{2\pi\mathrm i}\int_{c-\mathrm i\infty}^{c+\mathrm i\infty}
 \frac{\Gamma(s)^2\Gamma(\frac12-s)\Gamma(a-s)}{\Gamma(s+\frac12)^2\Gamma(1-s)\Gamma(a+\frac12-s)}\dd s,\qquad
 0<c<\mathrm{Re}\,a<\tfrac12.
\end{equation}
Close the contour to the right, with its clockwise orientation. The two pole families, $s=\frac12+n$ and $s=a+n$, give
respectively
\begin{equation}
 \sqrt\pi\frac{\Gamma(a-\frac12)}{\Gamma(a)}\Phi_A(1-a),\qquad
 \frac{\Gamma(a)^2\Gamma(\frac12-a)}{\sqrt\pi\,\Gamma(a+\frac12)^2\Gamma(1-a)}\Phi_D(a).
\end{equation}
The residues of $\Gamma(\frac12-s)$ and $\Gamma(a-s)$ have the opposite sign to ordinary left-pole residues; the
clockwise closure cancels that sign. Gamma reflection reduces the second prefactor to
$(\tan\pi a/\sqrt\pi)[\Gamma(a)/\Gamma(a+\frac12)]^3$. Restoring $2\cdot4^{-a}$ proves Eq.~\eqref{eq:cform:V} and
allows its continuation to the required values of $a$. Both pole families are part of the derivation.

\paragraph{The six-function form.} At the two values of $a$ in Eq.~\eqref{eq:cform:mUV} exactly six ordinary
${}_4F_3(1)$ values occur. They enter only through
\begin{equation}
 X=5\Phi_A(-\tfrac14)+2\Phi_A(\tfrac34),\qquad Y=5\Phi_A(\tfrac54)-6\Phi_A(\tfrac14),\qquad
 Z=5\Phi_D(-\tfrac14)+2\Phi_D(\tfrac34).
 \label{eq:cform:XYZ}
\end{equation}
Gamma recursion in Eqs.~\eqref{eq:cform:U} and~\eqref{eq:cform:V} gives
\begin{equation}
 m_U=-\frac{4\sqrt{2\pi}}{R_\Gamma}X-\frac{64\sqrt2}{\sqrt\pi R_\Gamma^3}Z,\qquad
 m_V=\frac{\sqrt{2\pi}R_\Gamma}{3}Y+\frac{64\sqrt2}{\sqrt\pi R_\Gamma^3}Z.
\end{equation}
In particular $Z$ cancels in $m_U+m_V$, but it contributes to $m_U+3m_V$. Equivalently,
\begin{equation}
 c_2=\frac{16}{R_\Gamma^2}-\frac{R_\Gamma^2}{9}-\frac{\pi^2}{15552}\bigl(12X-R_\Gamma^2Y\bigr)
 \Bigl[R_\Gamma^2\bigl(4X-R_\Gamma^2Y\bigr)-\frac{128Z}{\pi}\Bigr].
 \label{eq:cform:sixfunction}
\end{equation}
This six-function expression is an intermediate identity. The next step eliminates three independent combinations by
exact telescoping, including their boundary values.

%--------------------------------------------------------------------
\subsection{Telescoping proof of the three-function closed form}\label{app:cform:telescoping}
Let $\mathsf P,\mathsf Q,\mathsf T$ be the convergent series of Sec.~\ref{sec:closed2}. The closed
form~\eqref{eq:c2closed} follows from Eq.~\eqref{eq:cform:sixfunction} by the three exact reductions
\begin{equation}
 X=\frac92\mathsf P-\frac{2R_\Gamma}{\pi^{3/2}},\qquad
 Y=\frac{45}{64}\mathsf Q-\frac{24}{\pi^{3/2}R_\Gamma},\qquad
 Z=\frac{72}{25}\mathsf T+\frac{R_\Gamma^3}{8\sqrt\pi},
 \label{eq:cform:XYZfinal}
\end{equation}
which are proved here with every telescoping endpoint retained. Take the base summands
\begin{equation}
 a_n=\frac{(\frac12)_n^3(\frac34)_n}{(n!)^3(\frac54)_n},\qquad
 b_n=\frac{(\frac12)_n^3(\frac14)_n}{(n!)^3(\frac34)_n},\qquad
 d_n=\frac{(\frac12)_n(\frac34)_n^3}{n!\,(\frac54)_n^3}.
\end{equation}
The summands of $X,Y,Z$ in Eq.~\eqref{eq:cform:XYZ} are respectively
\begin{equation}
 -\frac{12n+7}{4n-1}a_n,\qquad\frac{3(12n-1)}{4n+3}b_n,\qquad
 \Bigl[2-5\Bigl(\frac{4n+1}{4n-1}\Bigr)^3\Bigr]d_n.
\end{equation}
Let $p_n,q_n,t_n$ denote the summands of $\mathsf P,\mathsf Q,\mathsf T$, and define the certificates
\begin{equation}
\begin{aligned}
 H_X(n)&=-\frac{4n^2(2n-3)(4n+1)^2}{(2n-1)^2(4n-1)},\qquad
 H_Y(n)=-\frac{3n(16n^3+44n^2+34n+15)}{2(n+1)^2},\\
 H_Z(n)&=\frac{2n(4n+1)(64n^3+16n^2-20n+43)}{(4n-1)^3}.
\end{aligned}
\end{equation}
Direct rational simplification proves
\begin{equation}
\begin{aligned}
 -\frac{12n+7}{4n-1}a_n-\frac92p_n&=H_X(n+1)a_{n+1}-H_X(n)a_n,\\
 \frac{3(12n-1)}{4n+3}b_n-\frac{45}{64}q_n&=H_Y(n+1)b_{n+1}-H_Y(n)b_n,\\
 \Bigl[2-5\Bigl(\frac{4n+1}{4n-1}\Bigr)^3\Bigr]d_n-\frac{72}{25}t_n&=H_Z(n+1)d_{n+1}-H_Z(n)d_n.
\end{aligned}
 \label{eq:cform:telescoping}
\end{equation}
For example, divide the first identity by $a_n$ and use
\begin{equation}
 \frac{a_{n+1}}{a_n}=\frac{(n+\frac12)^3(n+\frac34)}{(n+1)^3(n+\frac54)},\qquad
 \frac{p_n}{a_n}=\frac{(2n+1)(4n+1)}{(2n-1)^2(n+1)}.
\end{equation}
It becomes an identity of rational functions of $n$, so checking it depends neither on numerical precision nor on
summing the series. The other two identities are checked in the same way.

Every $H(0)$ vanishes, but the upper boundary terms do not:
\begin{equation}
\begin{gathered}
 \lim_{n\to\infty}\frac{(H_X,H_Y,H_Z)}{n^2}=(-8,-24,8),\\
 \lim_{n\to\infty}n^2a_n=\frac{R_\Gamma}{4\pi^{3/2}},\qquad
 \lim_{n\to\infty}n^2b_n=\frac1{\pi^{3/2}R_\Gamma},\qquad
 \lim_{n\to\infty}n^2d_n=\frac{R_\Gamma^3}{64\sqrt\pi}.
\end{gathered}
\end{equation}
Summing Eq.~\eqref{eq:cform:telescoping} from $n=0$ to $N$ and letting $N\to\infty$ gives exactly the three Gamma
terms of Eq.~\eqref{eq:cform:XYZfinal}. Substituting Eq.~\eqref{eq:cform:XYZfinal} into
Eq.~\eqref{eq:cform:sixfunction}, the boundary constants cancel from both factors, and the result is exactly the closed
form~\eqref{eq:c2closed}.

%--------------------------------------------------------------------
\subsection{Numerical evaluation and the scope of the checks}\label{app:cform:numerics}
\paragraph{Default evaluation.} The study evaluates Eq.~\eqref{eq:c2closed} at 80 working digits, and checks it against
the six-function form~\eqref{eq:cform:sixfunction} and against the endpoint series below, with 400 terms in each
half-interval series; Sec.~\ref{sec:closed2} quotes their agreement. With $N_{\rm h}$ terms in each half-interval
series, the absolute difference $|\Delta c_2|$ between the endpoint-series value and the three-function value is
\begin{center}\small
\begin{tabular}{@{}rl@{\qquad}rl@{\qquad}rl@{\qquad}rl@{}}
\toprule
$N_{\rm h}$ & $|\Delta c_2|$ & $N_{\rm h}$ & $|\Delta c_2|$ & $N_{\rm h}$ & $|\Delta c_2|$ & $N_{\rm h}$ & $|\Delta c_2|$\\
\midrule
24 & $\sci{5.2}{-7}$ & 64 & $\sci{2.8}{-19}$ & 144 & $\sci{1.5}{-43}$ & 320 & $\sci{2.5}{-80}$\\
40 & $\sci{6.0}{-12}$ & 96 & $\sci{5.1}{-29}$ & 224 & $\sci{9.4}{-68}$ & 400 & $\sci{2.5}{-80}$\\
\bottomrule
\end{tabular}
\end{center}
and it reaches the 80-digit arithmetic floor at 320 terms. This is numerical consistency, not an error bound. The values
are in \file{code/c2_closed_d2/outputs/raw/three_4f3_80dps.json} and \file{code/c2_closed_d2/outputs/raw/convergence.json}.

\paragraph{Endpoint-series evaluation.}\label{app:cform:endpoint} An evaluation independent of the
${}_4F_3$ identities uses the Frobenius expansions
\begin{equation}
\begin{gathered}
 \hat K(t)=\sum_{n\ge0}\phi_nt^n,\qquad \phi_n=\Bigl[\frac{(\frac12)_n}{n!}\Bigr]^2,\\
 G(t)=4\sum_{n\ge0}\phi_n(\mathsf h_{2n}-\mathsf h_n)t^n,\qquad
 K(1-t)=\tfrac12\bigl\{\log(16/t)\hat K(t)-G(t)\bigr\},
\end{gathered}
 \label{eq:cform:frobenius}
\end{equation}
with the harmonic numbers $\mathsf h_0=0$, $\mathsf h_n=\sum_{k=1}^n1/k$. Each period is split at $t=\frac12$, its second
half is reflected, and $(1-t)^{b-1}\hat K^2$, $(1-t)^{b-1}\hat KG$ and $(1-t)^{b-1}G^2$ are expanded in power series of
$t$, with coefficient arrays $\alpha_n,\beta_n,\gamma_n$. For each coefficient the only Mellin integrals required are,
with $p=a+n$ and $\ell_{32}=\log32$,
\begin{equation}
\begin{aligned}
 j_0(p)&=\AC\!\int_0^{1/2}t^{p-1}\dd t=\frac{2^{-p}}p,\qquad
 j_1(p)=\AC\!\int_0^{1/2}t^{p-1}\log(16/t)\dd t=2^{-p}\Bigl(\frac{\ell_{32}}p+\frac1{p^2}\Bigr),\\
 j_2(p)&=\AC\!\int_0^{1/2}t^{p-1}\log^2(16/t)\dd t=2^{-p}\Bigl(\frac{\ell_{32}^2}p+\frac{2\ell_{32}}{p^2}+\frac2{p^3}\Bigr).
\end{aligned}
\end{equation}
These hold first at $\mathrm{Re}\,p>0$ and continue analytically to the nonzero negative $p$ needed here. The three
unnormalized half-period sums (for $K(t)^2$, $K(t)K(1-t)$ and $K(1-t)^2$) are
\begin{equation}
 \sum_n\frac{\pi^2}{4}\alpha_nj_0,\qquad\sum_n\frac\pi4\bigl(\alpha_nj_1-\beta_nj_0\bigr),\qquad
 \sum_n\frac14\bigl(\alpha_nj_2-2\beta_nj_1+\gamma_nj_0\bigr),
\end{equation}
and adding the reflected half gives the full vector. Because the expansion is restricted to $0\le t\le\frac12$, this is
an exponentially convergent numerical check of the continued periods that uses neither the six- nor the
three-function formula. Agreement is a reproducible numerical consistency check, not an interval enclosure of the
final digits.

\paragraph{A separate endpoint-subtracted quadrature.} A second implementation checks the same moments by quadrature,
with the elliptic integrals evaluated through Carlson's symmetric integral,
\begin{equation}
 K(t)=R_F(0,1-t,1),\qquad K(1-t)=R_F(0,t,1),\qquad
 R_F(x,y,z)=\frac12\int_0^\infty\frac{\dd v}{\sqrt{(v+x)(v+y)(v+z)}},
\end{equation}
for positive arguments, continued to the endpoint limits shown. This representation does not use the Frobenius
coefficients. To make the negative moment at $a=-\frac14$ integrable, each symmetric moment is folded onto
$0<t<\frac12$ and its leading endpoint terms are subtracted. With $u=-\log t$ and $\ell=\log16+u$, define
\begin{equation}
\begin{aligned}
 R_U(u;a)&=(1-e^{-u})^{a-1}\bigl\{K(e^{-u})^2+K(1-e^{-u})^2\bigr\}-\frac{\pi^2+\ell^2}{4},\\
 R_V(u;a)&=(1-e^{-u})^{a-1}K(e^{-u})K(1-e^{-u})-\frac{\pi\ell}{4}.
\end{aligned}
\end{equation}
The continued Mellin integrals $j_k(a)$ then give
\begin{equation}
\begin{aligned}
 U(a)&=\frac4{\pi^2}\Bigl[\int_{\log2}^\infty e^{-au}R_U(u;a)\dd u+\frac{\pi^2j_0(a)+j_2(a)}{4}\Bigr],\\
 V(a)&=\frac8{\pi^2}\Bigl[\int_{\log2}^\infty e^{-au}R_V(u;a)\dd u+\frac{\pi j_1(a)}{4}\Bigr].
\end{aligned}
 \label{eq:cform:carlson}
\end{equation}
The remainders are $O(e^{-u}u^2)$, so the integrals converge for $\mathrm{Re}\,a>-1$, including $a=-\frac14$ and
$a=\frac34$. The implementation uses Gauss--Legendre quadrature on the intervals with endpoints
$\log2,2,5,10,20,40,80,120,200$, evaluates the integrand at 100 decimal digits, and restores the analytic leading terms
exactly as in Eq.~\eqref{eq:cform:carlson}. The finite upper limit $u=200$ leaves a rapidly decaying tail. Inserted in
Eq.~\eqref{eq:cform:c2fg}, these moments reproduce the closed form to $\sci{1.40}{-61}$ in $c_2$ (record value
$\sci{1.4045}{-61}$; Sec.~\ref{sec:closed2}). This is an observed difference; the script supplies no rigorous bound on
the quadrature error or on the cutoff (record: \file{code/c2_closed_d2/outputs/raw/independent_c2_80dps.json}).

\paragraph{Exact bounds on the series tails.} The three final series also admit direct remainder inequalities. For any
of them let $v_n$ be the absolute value of its $n$th summand, $n\ge1$. Its positive ratio is the rational function
\begin{equation}
 r_n=\frac{v_{n+1}}{v_n}=\frac{\prod_{j=1}^4(n+u_j)}{(n+1)\prod_{j=1}^3(n+l_j)},
\end{equation}
with the upper parameters $u_j$ and lower parameters $l_j$ of that series. Set $\phi_c(n)=n/2+c$ and
$E_c(n)=\phi_c(n)-1-\phi_c(n+1)r_n$. For $\mathsf P$ and $\mathsf T$ take $(c_-,c_+)=(0,1)$, and for $\mathsf Q$ take
$(1,2)$. Then $E_{c_-}(n)\le0$ and $E_{c_+}(n)\ge0$ for all integers $n\ge1$; these signs are verified exactly by
replacing $n$ with $m+1$ and checking that the numerator and denominator of the suitably signed rational expression
have nonnegative coefficients. Multiplying by $v_n$ and summing telescopically gives
\begin{equation}
 \Bigl(\frac N2+c_-\Bigr)v_N\le\sum_{n=N}^\infty v_n\le\Bigl(\frac N2+c_+\Bigr)v_N,
\end{equation}
where the upper boundary term vanishes because $v_n=O(n^{-3})$. All terms of $\mathsf P$ and $\mathsf Q$ are positive;
for $\mathsf T$ the terms with $n\ge1$ are negative. These signs decide which tail is added to or subtracted from a
partial sum. With $N=20000$, propagation of these bounds gives the envelope of $c_2$ quoted in Sec.~\ref{sec:closed2}.
The inequalities for the mathematical series are exact, but the arithmetic that evaluates their endpoints is ordinary
arbitrary-precision floating point, so the envelope is an independent check, not a certified rounding interval.

\paragraph{Scope of the checks.} The study's exact checks cover the three rational telescoping identities with their
nonzero endpoint terms, the complex coordinate Jacobians, the lower-order Gamma reductions and the scheme covariance
of $c_2$ under $\mathcal E_R\to e^{\kappa g^2}\mathcal E_R$. A separate check of the analytic bridges verifies the
Landen coefficient reduction, the cut-jump and reflection matrices, and the differential certificate used for
$\mathcal D_0=0$. The study's numerical checks test the elliptic-cover identity~\eqref{eq:cform:cover} at noncoincident
points,
the Gaussian finite term $B_{31}$, the contact relation, and representative simultaneous OPE limits; the source-contact checks
test the dependence of the finite contact on the smooth infrared cutoff and its cancellation against the local
counterterm. None of these checks evaluates every global collision sector. The discrepancies quoted above are
neither statistical error bars nor statements about the physical regulator comparison of
Appendix~\ref{app:cform:qualification}. The checks made by the review of the study are those of
Secs.~\ref{sec:closed2} and~\ref{sec:mc}.

%--------------------------------------------------------------------
\subsection{Further conditions stated by the study}\label{app:cform:status}
Beyond Secs.~\ref{sec:unproved} and~\ref{sec:open}, the study states three conditions on what it establishes.
\begin{itemize}
\item \emph{Endpoint terms are part of the value.} The evaluation requires the finite source contact, the
 interval-period pairing with its continued endpoint terms, and the nonzero endpoint terms of the telescoping proof;
 every integral with a negative endpoint exponent is a meromorphic continuation, not an ordinary improper integral.
 The endpoint values $A_{\rm br}(0)=\pi^2/4$ and $B_{01}=\pi^2/8$ of the contact relation are prescribed continued
 limits; no ordinary infrared convergence is asserted for them.
\item \emph{The digits are checked, not certified, within the study.} Its long decimal value rests on
 high-precision evaluations by different routes (the six- and three-function forms, the endpoint series and the
 Carlson quadrature); it is not presented as a certified rounding theorem. A certified enclosure would
 combine interval arithmetic with the analytic tail bounds above, or bound both the Carlson quadrature error and its
 tail beyond $u=200$. (The integer-interval evaluation to 51 decimal places mentioned in Sec.~\ref{sec:closed2} is part
 of the cylinder base study (Appendix~\ref{app:cyl:closedchecks}), not of this study.)
\item \emph{The local checks do not bear on the global step.} The fixed-modulus contact calculation and the
 high-precision agreement of the special-function evaluations establish different parts of the method; neither
 supplies the uniform remainder bound of Appendix~\ref{app:cform:qualification}.
\end{itemize}

%====================================================================
\section{\texorpdfstring{$c_2$}{c2} from fourth-order perturbation theory on the cylinder: method}\label{app:cyl}
%====================================================================
This appendix gives the method of the base study of $c_2$ with the truncated cylinder Hamiltonian (written by OpenAI
GPT-6 astra; Sec.~\ref{sec:sources}), edited to the notation of this note.

Section~\ref{sec:cyl} summarises the setup and gives the continuum analysis that produced the quoted value
$c_2=-4.0670(24)$, whose error is a linear sum of systematic diagnostics, not a standard deviation. That analysis was
done later, with $N\le1024$ and with fit and selection rules fixed in code before the final numbers were computed
(its error rule was revised once, after the first result had been seen; Sec.~\ref{sec:cyl}); it evaluates the
quantities defined here. The appendix gives
the normalization and the exact second-order coefficient (Appendices~\ref{app:cyl:setup} to~\ref{app:cyl:anomaly}), the
fourth-order Rayleigh--Schr\"odinger formula, its Hilbert space, regulator and counterterms
(Appendices~\ref{app:cyl:rs} to~\ref{app:cyl:counterterms}), the spectral generating functions that evaluate it
(Appendices~\ref{app:cyl:blocks} and~\ref{app:cyl:collapsed}), the large-$N$ tails and the regulator limit
(Appendix~\ref{app:cyl:tails}), the checks (Appendices~\ref{app:cyl:checks} to~\ref{app:cyl:closedchecks}), and the explicit
formulas behind them (Appendices~\ref{app:cyl:formfactors} to~\ref{app:cyl:algorithm}). The calculation is numerical: no
step below bounds the remaining spectral or regulator extrapolation. The base study's checks of the closed form
(Appendix~\ref{app:cyl:closedchecks}) certify the value of the closed expression, Eq.~\eqref{eq:c2closed}, not its
identification with the physical coefficient, which remains conditional on the step of Sec.~\ref{sec:unproved}.

Notation: $\delta$ is that of Table~\ref{tab:conv} ($d=2$), and $c_2$ is the coefficient the base study writes $C_2$.
As in Sec.~\ref{sec:cyl}, $\ket{\mathcal E}$ is the state of the Ising energy field $\epsilon$ at the ultraviolet
point, $E_{\mathcal E}$ its energy, and a subscript $\mathcal E$ on a state-dependent quantity refers to it.

%--------------------------------------------------------------------
\subsection{Continuum setup and normalization}\label{app:cyl:setup}
\paragraph{The short-range description.} The theory is the crossover description of Sec.~\ref{sec:d2}: the diagonal
Ising CFT, with central charge $c=\frac12$, tensored with a generalized free scalar $\chi$, with
\begin{equation}
 \Delta_\sigma=\frac18,\qquad \Delta_\epsilon=1,\qquad \Delta_\chi=\frac{15}{8}-\delta,\qquad
 C_{\sigma\sigma\epsilon}=\frac12,\qquad C_{\epsilon\epsilon\epsilon}=0,
\end{equation}
and unit-normalized plane two-point functions,
$\langle\sigma(x)\sigma(0)\rangle=|x|^{-1/4}$, $\langle\epsilon(x)\epsilon(0)\rangle=|x|^{-2}$ and
$\langle\chi(x)\chi(0)\rangle=|x|^{-2\Delta_\chi}$. The higher correlators of $\chi$ are Wick-factorized. The
perturbation is $\sigma\chi$, of dimension $p=2-\delta$. Critical thermal tuning is understood: the relevant energy
coupling generated by the perturbation is removed along the critical trajectory. The limit is $\delta\to0^+$, from the
long-range side.

\paragraph{The cylinder Hamiltonian.} On the radial cylinder of unit radius, with the unperturbed vacuum energy
subtracted,
\begin{equation}
 H_0=D_{\rm Ising}+D_\chi,\qquad V=\int_0^{2\pi}\dd\theta\,\sigma(\theta)\chi(\theta),\qquad
 V_\epsilon=\int_0^{2\pi}\dd\theta\,\epsilon(\theta),
 \label{eq:cyl:H0}
\end{equation}
where $D_{\rm Ising}$ and $D_\chi$ are the dilatation operators and $V_\epsilon$ is the integrated thermal operator
(the base study writes $\mathcal E$ for it). The angular integral projects onto equal total spin and contributes a
factor $2\pi$ at every vertex, so a chain of four vertices carries $(2\pi)^4$. These factors are part of the coupling
convention and must not be changed independently of $b_0$.

\paragraph{Fixed point and coefficient extraction.} The beta function is $\beta(g)=-\delta g+b_0g^3+O(g^5)$, with the
exact $b_0$ of Sec.~\ref{sec:rg2}. This normalization agrees with the lower-order crossover
calculation~\cite{Behan:2017emf}, and it is checked independently by the magnetic gap (Appendix~\ref{app:cyl:magnetic}).
The infrared fixed point has $g_*^2=\delta/b_0+O(\delta^2)$. Let $\ket d=L_{-1}\ket{\mathcal E}$ be the left-derivative
state at the ultraviolet point. Since the energy field has holomorphic weight $\frac12$, its norm is $2\cdot\frac12=1$;
its unperturbed dimension is two and its spin one. Following the energy eigenstates continuously connected to these
states, define
\begin{equation}
 X=2(E_{\mathcal E}-E_0)-(E_d-E_0)=2E_{\mathcal E}-E_d-E_0.
 \label{eq:cyl:X}
\end{equation}
At a conformal fixed point $E_{\mathcal E}-E_0=\De$ and $E_d-E_0=\De+1$, hence $X=\De-1$; at the ultraviolet point
$X=0$. In a consistently matched Hamiltonian scheme,
\begin{equation}
 X=a_1\delta g^2+hg^4+O(g^6,\delta^2g^2,\delta g^4),
\end{equation}
and substituting $g_*$ gives the relation $c_2=a_1/b_0+h_0/b_0^2$ of Sec.~\ref{sec:cyl}, with $h_0$ the value of
$h$ once the regulator is removed (Eq.~\eqref{eq:cyl:limits} below). Only the leading fixed-point coupling is needed at
this order, because there is no term $\delta^0g^2$. The split between the two terms depends on the prescription
(Sec.~\ref{sec:rg2}): the Hamiltonian $h$ is not the pole-only coefficient $\hMS$ of minimal subtraction, and dropping
$a_1/b_0$ would be wrong in this convention. A shift $\Delta h_0$ moves $c_2$ by $\Delta h_0/b_0^2\simeq0.62156\,\Delta
h_0$ ($b_0^2\simeq1.60885$); this is the conversion by which errors in $h_0$ are quoted in $c_2$.

\paragraph{Why the coefficients are computed at $\delta=0$.} Inferring a curvature from dimensions computed at finite
$\delta$ would magnify an error in a dimension by $\delta^{-2}$. The method instead computes the $g^4$ coefficient at
$\delta=0$ and the analytic $\delta$ derivative of the $g^2$ coefficient. The remaining fits are numerical fits in
auxiliary cutoffs and in the regulator length, never in the physical $\delta$. What the calculation shares with the
closed-form route, and what it does not use, is listed in Sec.~\ref{sec:cyl}.

\paragraph{Notation and the three separate limits.}
\begin{center}\small
\begin{tabular}{@{}l>{\raggedright\arraybackslash}p{.74\textwidth}@{}}
\toprule
symbol & meaning\\
\midrule
$\delta$ & physical distance from the short-range end; the fourth-order sums are evaluated at $\delta=0$\\
$\Delta_{\chi,0}=\frac{15}{8}$ & endpoint dimension of $\chi$ (not an oscillator index)\\
$r=1,2,3$ & order of the positive, polynomial-improved oscillator regulator (Appendix~\ref{app:cyl:regulator})\\
$\lambda>0$ & length of that smooth ultraviolet regulator on the unit cylinder\\
$N$ & cubical coefficient cutoff, $0\le n,m,\ell<N$; it is not the number of Hilbert-space states\\
$N_{\rm mix}$ & two-index cutoff of the mixed thermal sums (4096 in the base study)\\
$N_{\rm col}$ & cutoff of the collapsed sums ($200\,000$ and $800\,000$ in the base study)\\
\bottomrule
\end{tabular}
\end{center}
A third auxiliary cutoff, besides $N_{\rm mix}$ and $N_{\rm col}$, is the two-index cutoff of the leading-tail calculation (Appendix~\ref{app:cyl:facetail}). The
auxiliary cutoffs are not interchangeable with $N$. The numerical limits are taken in the order
\begin{equation}
 h_0=\lim_{\lambda\to0^+}\Bigl[\lim_{N\to\infty}h_N(\lambda)\Bigr],
 \label{eq:cyl:limits}
\end{equation}
with $N_{\rm mix}$, $N_{\rm col}$ and the tail cutoff removed together with $N$ at fixed $\lambda$ (in the analysis of
Sec.~\ref{sec:cyl} they are held at common values for all points and their shifts under a change are propagated); $h_N(\lambda)$ is
the assembled fourth-order quantity of Eq.~\eqref{eq:cyl:assembly} with the chain sums truncated at $N$. None of these
limits samples the physical $\delta$.

%--------------------------------------------------------------------
\subsection{Exact second-order spectral continuation}\label{app:cyl:second}
\paragraph{Vacuum and thermal sums.} Set $p=2-\delta$ and
\begin{equation}
 A_n=\frac{(p)_n}{n!},\qquad A_{-1}=A_{-2}=0.
\end{equation}
Expanding the exact cylinder correlators and projecting on angles gives
\begin{align}
 E_0^{(2)}&=-4\pi^2\sum_{n=0}^\infty\frac{A_n^2}{p+2n},\label{eq:cyl:E02}\\
 E_{\mathcal E}^{(2)}&=-\pi^2\sum_{n=0}^\infty\frac{(A_n+A_{n-1})^2}{p-1+2n}.\label{eq:cyl:EE2}
\end{align}
The denominators are measured relative to the external state. The sums converge in an ordinary domain, for example
$\frac12<p<1$ away from level crossings, and are then continued meromorphically; they are not convergent positive sums
at the physical endpoint. With $\tilde A_n=(p-1)_n/n!$ (the base study's $c_n$),
\begin{equation}
 A_n+A_{n-1}=\frac{2n+p-1}{p-1}\tilde A_n,
\end{equation}
so Eq.~\eqref{eq:cyl:EE2} reduces to a linear combination of two Gauss sums,
\begin{equation}
 \sum_{n\ge0}\tilde A_n^2=\frac{\Gamma(3-2p)}{\Gamma(2-p)^2},\qquad
 \sum_{n\ge0}n\tilde A_n^2=(p-1)^2\frac{\Gamma(2-2p)}{\Gamma(2-p)^2}.
\end{equation}
With $\Gamma(3-2p)=(2-2p)\Gamma(2-2p)$ the required combination vanishes exactly:
\begin{equation}
 E_{\mathcal E}^{(2)}=0\qquad\text{as a meromorphically continued function.}
\end{equation}
The vacuum sum is a well-poised hypergeometric sum. Written as $p^{-1}\,{}_3F_2(p,p,\frac p2;1,1+\frac p2;1)$, it
reduces to the standard well-poised summation, followed by Gamma reflection and recursion:
\begin{equation}
 -E_0^{(2)}=\pi^2\frac{\Gamma(p/2)^2\,\Gamma(1-p)}{\Gamma(1-p/2)^2\,\Gamma(p)}.
 \label{eq:cyl:vac}
\end{equation}
This fixes the analytic continuation used here; it does not assign a finite part to an endpoint cutoff sum.

\paragraph{The derivative state and the slope of the observable.} Define
\begin{equation}
 D_n=A_{n+1}+3A_n+3A_{n-1}+A_{n-2}.
\end{equation}
The absolute second-order correction of the derivative state is
\begin{equation}
 E_d^{(2)}=-\frac{\pi^2}{4(p-1)}\sum_{n\ge0}\tilde A_nD_n.
\end{equation}
For $\Sigma_j=\sum_{n\ge0}\tilde A_nA_{n+j}$ and $\mathcal G=\Gamma(2-2p)/\Gamma(2-p)^2$ the continued shifted
sums are
\begin{equation}
 \frac{\Sigma_1}{\mathcal G}=\frac{p(1-p)}{2-p},\qquad
 \frac{\Sigma_0}{\mathcal G}=1-p,\qquad
 \frac{\Sigma_{-1}}{\mathcal G}=p-1,\qquad
 \frac{\Sigma_{-2}}{\mathcal G}=\frac{p(p-1)}{2-p}.
\end{equation}
Consequently $\Sigma_1+3\Sigma_0+3\Sigma_{-1}+\Sigma_{-2}=0$, and $E_d^{(2)}=0$ as well. The
coefficient of $g^2$ in $X$ is therefore $-E_0^{(2)}$ at $p=2-\delta$:
\begin{equation}
 a(\delta)=\pi^2\frac{\Gamma(1-\delta/2)^2\,\Gamma(\delta-1)}{\Gamma(\delta/2)^2\,\Gamma(2-\delta)}
 =-\frac{\pi^2}{4}\delta-\frac{\pi^2}{2}\delta^2+O(\delta^3),
 \label{eq:cyl:adelta}
\end{equation}
so that
\begin{equation}
 a_1=-\frac{\pi^2}{4},\qquad \frac{a_1}{b_0}=-1.94527965797909034563\ldots,
\end{equation}
which is the finite matching term of Eq.~\eqref{eq:c2I}. A complementary check
(\file{code/c2_cylinder_d2/code/hamiltonian/slope_check.py}) differentiates the regulated Pochhammer coefficients,
the regulator factors and the resolvent denominators analytically in $\delta$; it does not finite-difference dimensions
at several values of $\delta$.

%--------------------------------------------------------------------
\subsection{A finite endpoint-regulator anomaly}\label{app:cyl:anomaly}
Setting $p=2$ first and multiplying each second-order resolvent term by $e^{-\tau(E_n-E_a)}$, where $E_n$ is the energy
of the intermediate state, $E_a$ that of the external state and $\tau>0$ a Euclidean time cutoff, direct geometric
summation gives
\begin{equation}
 E_0^{(2)}(\tau)=-\frac{\pi^2}{2\sinh^2\tau},\qquad
 E_{\mathcal E}^{(2)}(\tau)=-\frac{\pi^2\cosh\tau}{2\sinh^2\tau},
\end{equation}
hence
\begin{equation}
 E_{\mathcal E}^{(2)}(\tau)-E_0^{(2)}(\tau)=-\frac{\pi^2}{4\cosh^2(\tau/2)}\ \xrightarrow{\ \tau\to0\ }\ -\frac{\pi^2}{4}.
\end{equation}
The analytically continued endpoint gap, however, is $a(0)=0$. Even a finite difference of cutoff sums can therefore
have the wrong finite part. This is the failure of an attempted endpoint prescription, not an error in the closed-form
study; it motivates the explicitly matched Hermitian regulator of Appendix~\ref{app:cyl:regulator}. General fourth-order
difficulties of Hamiltonian cutoffs, including nonlocal cutoff effects and disconnected cancellations, are discussed in
Refs.~\cite{EliasMiro:2021,EliasMiro:2022}. An earlier explicit-basis calculation that fixed the counterterms by trial
matching to finite spectral gaps gave strongly cutoff-dependent fourth-order values; it is not used.

%--------------------------------------------------------------------
\subsection{Fourth-order Rayleigh--Schr\"odinger perturbation theory}\label{app:cyl:rs}
\paragraph{Reduced resolvents and the normalization subtraction.} For an external state $\ket a$, work in its
conserved-spin sector and use
\begin{equation}
 Q_a=1-\ket a\langle a\rvert,\qquad R_a=\frac{Q_a}{E_a^{(0)}-H_0},\qquad \ket{w_a}=R_aV\ket a,\qquad
 \ket{z_a}=Q_aV\ket{w_a}.
\end{equation}
The relevant external levels are nondegenerate within their symmetry sectors at the endpoint. Since $V_{aa}=0$, the
second- and fourth-order energy coefficients are
\begin{align}
 E_a^{(2)}&=\langle a\rvert VR_aV\ket a,\\
 E_a^{(4)}&=\langle z_a\rvert R_a\ket{z_a}-E_a^{(2)}\langle w_a\vert w_a\rangle.
 \label{eq:cyl:RS4}
\end{align}
The first term, $\langle a\rvert VR_aVR_aVR_aV\ket a$, is the chain. The second comes from the perturbative
normalization of the state and is essential. It must use the same regulator as the chain; assigning unrelated finite
parts to the two separately can destroy their cancellation.

Because $V$ is linear in $\chi$, four insertions that begin and end in the $\chi$ vacuum pass through the $\chi$
occupations $0\to1\to0\to1\to0$ or $0\to1\to2\to1\to0$. The explicit-basis implementation includes the factors
$\sqrt2$ for double occupation of a mode; the generating-function implementation (Appendix~\ref{app:cyl:blocks}) includes
the same multiplicities through the three Wick pairings. Finite-box tests compare the two
(Appendix~\ref{app:cyl:checks}).

\paragraph{Counterterms enter at the same perturbative order.} Write the matched Hamiltonian as
\begin{equation}
 H=H_0+gV+g^2W_2+g^3W_3+g^4W_4+\cdots.
\end{equation}
For any insertion $W$ of order $g^2$ define its mixed and quadratic fourth-order contributions
\begin{align}
 L_a(W)&=\langle w_a\rvert(W-W_{aa})\ket{w_a}+2\,{\rm Re}\,\langle Q_aWa\rvert R_a\ket{z_a},\notag\\
 \mathcal Q_a(W)&=\langle Q_aWa\rvert R_a\ket{Q_aWa},
 \label{eq:cyl:mixed}
\end{align}
where $\ket{Q_aWa}$ means $Q_aW\ket a$ (not an independent state). For every state-dependent quantity $Y$ write
$Y_X=2Y_{\mathcal E}-Y_d-Y_0$. For $W_2=tV_\epsilon+kD_{\rm Ising}$ the quadratic contribution of the thermal term to $X$
is exactly
\begin{equation}
 \mathcal Q_X(tV_\epsilon)=\frac{32\pi^2}{3}\,t^2:
 \label{eq:cyl:thermalsq}
\end{equation}
a free Majorana thermal deformation gives the second-order gap coefficients $8\pi^2$ and $16\pi^2/3$ for $\ket{\mathcal
E}$ and $\ket d$, and $2\cdot8\pi^2-16\pi^2/3=32\pi^2/3$. Quadratic terms involving $D_{\rm Ising}$ vanish because it is
diagonal; its \emph{mixed} fourth-order contribution $L_X(D_{\rm Ising})$ does not. The observable $X$ was chosen so
that the direct local fourth-order terms cancel (Appendix~\ref{app:cyl:counterterms}); the mixed terms of
Eq.~\eqref{eq:cyl:mixed} cannot be omitted on that basis.

%--------------------------------------------------------------------
\subsection{Hilbert space, matrix elements and oscillators}\label{app:cyl:hilbert}
\paragraph{Ising states.} The validation implementation uses an orthonormal Majorana basis of Neveu--Schwarz (NS) and
Ramond (R) states. In the NS sector the occupied chiral modes are positive half-integers, in the R sector positive
integers. Chiral R energies carry the vacuum shift $\frac1{16}$, so the full R shift is $\frac18$, and the diagonal
Ising theory retains NS left and right occupations of equal parity. The vacuum, energy, derivative and spin states are
\begin{equation}
 \ket0=(\varnothing,\varnothing)_{\rm NS},\qquad \ket{\mathcal E}=(\tfrac12,\tfrac12)_{\rm NS},\qquad
 \ket d=(\tfrac32,\tfrac12)_{\rm NS},\qquad \ket\sigma=(\varnothing,\varnothing)_{\rm R}.
\end{equation}
Products of chiral spin-field form factors give the full Ising $\sigma$ matrix elements, and the angular integration
then enforces conservation of total spin. Appendix~\ref{app:cyl:formfactors} gives the product formula and the fermion
signs; the representation is also available in Ref.~\cite{Liu:2026}.

\paragraph{Generalized-free oscillators.} For $\tau>0$ the unit-cylinder correlator is
\begin{equation}
 \langle\chi(\tau,\theta)\chi(0,0)\rangle=\bigl[2\cosh\tau-2\cos\theta\bigr]^{-\Delta_\chi}.
\end{equation}
Expanding its two chiral factors gives modes $m,n\ge0$ with
\begin{equation}
 \omega_{mn}=\Delta_\chi+m+n,\qquad J_{mn}=m-n,\qquad
 \gamma_{mn}=\sqrt{\frac{(\Delta_\chi)_m(\Delta_\chi)_n}{m!\,n!}},
\end{equation}
energy $\omega_{mn}$, spin $J_{mn}$ and oscillator coefficient $\gamma_{mn}$ (the base study's $c_{mn}$), so that
$D_\chi=\sum_{m,n}\omega_{mn}a^\dagger_{mn}a_{mn}$. These modes implement the radial Hamiltonian of the isotropic theory;
no identification with a spatially long-range quantum spin chain is assumed.

%--------------------------------------------------------------------
\subsection{The regulator family}\label{app:cyl:regulator}
At the endpoint set $\Delta_{\chi,0}=\frac{15}{8}$ and define
\begin{equation}
 P_r(x)=\sum_{j=0}^r\frac{x^j}{j!},\qquad q=e^{-\lambda}
\end{equation}
($q$ is unrelated to the circular mean of Eq.~\eqref{eq:I}). The oscillator coefficient is replaced by
\begin{equation}
 \tilde\gamma_{mn}=q^{(\Delta_{\chi,0}+m+n)/2}\sqrt{\frac{(\Delta_{\chi,0})_m(\Delta_{\chi,0})_n}{m!\,n!}\,
 P_r\bigl(\lambda(m+\tfrac{\Delta_{\chi,0}}{2})\bigr)\,P_r\bigl(\lambda(n+\tfrac{\Delta_{\chi,0}}{2})\bigr)}.
 \label{eq:cyl:regc}
\end{equation}
All factors under the square root are positive. Applying the same real coefficient to the creation and annihilation
parts gives a Hermitian interaction at fixed $\lambda>0$. The un-improved regulator is $r=0$; the analysis uses
$r=1,2,3$. Excluding a common factor $q^{\Delta_{\chi,0}/2}$, the chiral contraction series is
\begin{equation}
 H(z)=\sum_{n\ge0}H_nz^n,\qquad
 H_n=\frac{(\Delta_{\chi,0})_n}{n!}\,q^n\,P_r\bigl(\lambda(n+\tfrac{\Delta_{\chi,0}}{2})\bigr).
 \label{eq:cyl:Hseries}
\end{equation}
Since $e^{-x}P_r(x)=1+O(x^{r+1})$, low-frequency regulator artifacts are suppressed. This improvement does not remove
the need for local counterterms. At fixed $\lambda$ the $\chi$ modes are exponentially damped and the remaining Ising
collision singularities of the sums are integrable. The cubical spectral cutoff $N$ is therefore a numerical truncation
of a regulated theory, not the physical ultraviolet prescription.

%--------------------------------------------------------------------
\subsection{Local counterterms and the improved observable}\label{app:cyl:counterterms}
\paragraph{Thermal and spin-two terms from the short-distance kernel.} The Ising operator product expansion contains
\begin{equation}
 \sigma(z)\sigma(0)=|z|^{-1/4}\Bigl[1+\frac12|z|\,\epsilon(0)+\frac14z^2T(0)+\frac14\bar z^2\bar T(0)+\cdots\Bigr],
\end{equation}
with $T$, $\bar T$ the stress tensor. The $\epsilon$ term gives a thermal power subtraction. The regulator distinguishes
cylinder time from angle at short distances, so the spin-two terms generate a finite anisotropic subtraction. With the
normalization of Eq.~\eqref{eq:cyl:H0} the counterterm is
\begin{equation}
 W_2=t_r(\lambda)\,V_\epsilon-\kappa_r\,D_{\rm Ising}.
 \label{eq:cyl:W2}
\end{equation}
It is derived from the regulated kernel, not by setting the fourth-order answer or by fitting an energy gap to a
desired value. For a local coordinate $z=\tau+\mathrm ix$ and unit regulator length define
\begin{equation}
 \mathfrak h_{\alpha,r}(z)=\sum_{j=0}^r\frac{(\alpha)_j}{j!}(z+1)^{-\alpha-j}.
\end{equation}
The planar regulated $\chi$ kernel is $|\mathfrak h_{\Delta_{\chi,0},r}(z)|^2$. Its thermal moment is
\begin{equation}
 c_{t,r}=\frac12\int_{\tau>0}\dd\tau\dd x\,(\tau^2+x^2)^{3/8}\,|\mathfrak h_{\Delta_{\chi,0},r}(\tau+\mathrm ix)|^2,
 \qquad t_r(\lambda)=\frac{c_{t,r}}{\lambda}+\cdots.
\end{equation}
For the un-improved regulator
\begin{equation}
\begin{aligned}
 c_{t,0}&={}_3F_2\!\left(\begin{matrix}-\frac38,1,1\\[1pt]\frac32,\frac32\end{matrix};1\right)=0.76212879108984270790\ldots,\\
 \kappa_0&=\frac{\pi^2}{4}\,\frac{\Delta_{\chi,0}-\rho_\chi}{\Delta_{\chi,0}-1},\qquad
 \rho_\chi=\frac{\Gamma(\Delta_{\chi,0}+\frac12)}{\Gamma(\frac32)\Gamma(\Delta_{\chi,0})},\qquad
 \kappa_0=1.20829489911072968942\ldots.
\end{aligned}
\end{equation}
For $r>0$ the difference of the spin-two moments of the improved and un-improved kernels converges and determines
$\kappa_r-\kappa_0$. Appendix~\ref{app:cyl:quadratures} gives the compact finite-domain integrals used for these
constants, so that they can be recomputed without any saved number.

\paragraph{Removing a leading cylinder artifact.} The holomorphic cylinder kernel obeys
\begin{equation}
 \bigl[2\sinh(z/2)\bigr]^{-\Delta_{\chi,0}}=z^{-\Delta_{\chi,0}}-\frac{\Delta_{\chi,0}}{24}z^{2-\Delta_{\chi,0}}+\cdots.
\end{equation}
Applying the same regulator to this expansion produces a thermal term of order $\lambda$. For $r\ge1$ the subtracted
moment converges:
\begin{equation}
 c_{{\rm cyl},r}=-\frac{\Delta_{\chi,0}}{24}\int_{\tau>0}\dd^2z\,|z|^{3/4}\Bigl[{\rm Re}\bigl(\mathfrak
 h_{\Delta_{\chi,0}-2,r}(z)\,\overline{\mathfrak h_{\Delta_{\chi,0},r}(z)}\bigr)-{\rm Re}(z^2)\,|z|^{-2\Delta_{\chi,0}}\Bigr].
\end{equation}
The calculation therefore uses
\begin{equation}
 \boxed{t_r(\lambda)=\frac{c_{t,r}}{\lambda}+c_{{\rm cyl},r}\,\lambda.}
 \label{eq:cyl:tfinal}
\end{equation}
The second term vanishes in the continuum. It is an explicitly derived regulator improvement, not an adjustable finite
fourth-order constant, and it is checked separately by the protected magnetic coefficient
(Appendix~\ref{app:cyl:magnetic}). The constants are
\begin{center}\small
\begin{tabular}{@{}rrrr@{}}
\toprule
$r$ & $c_{t,r}$ & $\kappa_r$ & $c_{{\rm cyl},r}$\\
\midrule
1 & $1.96478977995214$ & $1.78455951314437$ & $-0.11068091195498$\\
2 & $3.31217153117521$ & $2.20281564903759$ & $-0.05396644709415$\\
3 & $4.74456159798846$ & $2.54264875945082$ & $-0.03541592348370$\\
\bottomrule
\end{tabular}
\end{center}
They are ordinary quadrature values, refined as described in Appendix~\ref{app:cyl:quadratures}, not interval enclosures.

\paragraph{Why the direct fourth-order local terms cancel in $X$.} After the auxiliary spectral cutoff has been removed
at fixed $\lambda$, short-distance power counting and the separate Ising and $\chi$ parity selection rules restrict the
nonvanishing relevant or marginal local counterterms of even order to the identity, the thermal operator and the
parity-even Ising spin-two term. Odd-spin total derivatives do not contribute. In particular $\chi^2$ has dimension
$\frac{15}{4}>2$ at the endpoint and is not an additional relevant or marginal subtraction. The direct fourth-order
expectations of these operators in $X$ vanish:
\begin{equation}
\begin{array}{c|c}
 \text{operator in }W_4 & 2(W_4)_{\mathcal E\mathcal E}-(W_4)_{dd}-(W_4)_{00}\\\hline
 \mathbb 1 & 2-1-1=0\\
 V_\epsilon & 0\quad\text{(each diagonal element vanishes)}\\
 D_{\rm Ising} & 2\cdot1-2-0=0
\end{array}
\label{eq:cyl:W4}
\end{equation}
The diagonal thermal elements vanish by Ising fusion ($C_{\epsilon\epsilon\epsilon}=0$) and remain zero for the
descendant. This cancellation does not remove the mixed effects of $W_2$, which are included explicitly.

At third order the coupling counterterm $W_3=\mu_3(\lambda)V$ contributes
$2\,{\rm Re}\,\langle a\rvert W_3R_aV\ket a=2\mu_3(\lambda)E_a^{(2)}$, and thus $2\mu_3(\lambda)X_2(\lambda)$ to $X$,
where $X_2(\lambda)$ is the second-order observable at finite $\lambda$ and $\delta=0$
(Appendix~\ref{app:cyl:collapsed}). The coefficient $\mu_3$ grows logarithmically, whereas $X_2(\lambda)$ tends to zero as a
positive power of $\lambda$, so this product supplies no unknown finite constant to the limiting $h$. The analysis uses
the convenient choice
\begin{equation}
 \mu_3(\lambda)=-b_0\log\lambda,\qquad \Delta X_4=-2b_0\log\lambda\,X_2(\lambda).
 \label{eq:cyl:couplinglog}
\end{equation}
Omitting this term, or shifting the reference scale of the logarithm, is a sensitivity test at finite regulator.

This is an argument about the local short-distance expansion of a specified, smoothly regulated theory, made
\emph{after} the numerical spectral limit. It does not license arbitrary finite counterterms at a fixed hard
Hilbert-space cutoff, and the numerical calculation does not supply a uniform treatment of every subleading
remainder.

\paragraph{The assembled fourth-order quantity.} Combining all terms,
\begin{equation}
\begin{aligned}
 h(\lambda)={}&X_{4,\rm chain}+X_{4,\rm norm}+t_r(\lambda)\,L_X(V_\epsilon)-\kappa_r\,L_X(D_{\rm Ising})\\
 &+\frac{32\pi^2}{3}t_r(\lambda)^2-2b_0\log\lambda\,X_2(\lambda),
\end{aligned}
\label{eq:cyl:assembly}
\end{equation}
where every sum on the right must first be converged at fixed $\lambda$, and $h_0=\lim_{\lambda\to0}h(\lambda)$. Here
$X_{4,\rm chain}$ is the chain term of Eq.~\eqref{eq:cyl:RS4} combined into $X$ (Appendix~\ref{app:cyl:blocks}),
$X_{4,\rm norm}$ its normalization term and $L_X(D_{\rm Ising})$ the mixed spin-two term (both
Appendix~\ref{app:cyl:collapsed}), and $L_X(V_\epsilon)$ the mixed thermal term (Appendix~\ref{app:cyl:mixed}). Large terms
cancel in Eq.~\eqref{eq:cyl:assembly}, and its individual pieces are not physical coefficients; this is one reason the
collapsed and mixed sums need their own convergence checks. The continuum analysis of Sec.~\ref{sec:cyl} evaluates
this assembly at each $(r,\lambda)$.

%--------------------------------------------------------------------
\subsection{Exact chiral blocks as spectral generating functions}\label{app:cyl:blocks}
\paragraph{Replacing state enumeration.} The explicit Majorana basis serves for validation but grows rapidly. The
production implementation sums the level degeneracies analytically with chiral blocks. A Taylor coefficient of a
radially ordered four-point block is the sum, over intermediate states at specified levels, of products of spin-field
matrix elements. Dividing these coefficients by the corresponding three energy denominators therefore evaluates the
same Rayleigh--Schr\"odinger sum as Eq.~\eqref{eq:cyl:RS4}. This is not a return to the elliptic-integral calculation
of the closed-form study: the blocks supply discrete spectral weights, and the operations that follow are coefficient
extraction, angular projection and Hamiltonian denominators. The underlying Ising correlators are known
independently~\cite{Ardonne:2010zz}.

\paragraph{Vacuum, thermal and derivative blocks.} Place four chiral spin insertions at
\begin{equation}
 z_1=1,\qquad z_2=u,\qquad z_3=uv,\qquad z_4=uvw,
\end{equation}
and define, for $|u|,|v|,|w|<1$,
\begin{equation}
\begin{gathered}
 S=\sqrt{1-uv}\sqrt{1-vw},\qquad T=\sqrt{1-v}\sqrt{1-uvw},\qquad
 \mathcal D=\prod_{x\in\{u,v,w,uv,vw,uvw\}}(1-x)^{-1/8},\\
 \mathcal R=\sqrt{(S+T)/2},\qquad \mathcal U=\sqrt{(1-u)(1-w)}
\end{gathered}
\end{equation}
(the base study writes $\mathcal Q$ for $\mathcal U$). All roots are continued analytically from the origin, with
positive values for small positive real arguments. After removal of simple radial monomials the two middle-fusion
channels ($\eta=0,1$) of the vacuum and thermal blocks are
\begin{equation}
\begin{aligned}
 g_0&=\mathcal D\mathcal R, & g_1&=\frac{\mathcal D\mathcal U}{2\mathcal R},\\
 e_0&=\frac{\mathcal D}{4\mathcal R}\bigl[S(1+uw)+T(u+w)\bigr], &
 e_1&=\frac{\mathcal D\mathcal R}{2\mathcal U}\bigl[S(1+uw)-T(u+w)\bigr].
\end{aligned}
\label{eq:cyl:blocksge}
\end{equation}
For a partition $(A,B)$ of the four position labels set
\begin{equation}
 K_{AB}=uvw\Bigl[1-\frac14\Bigl(\sum_{i\in B}z_i-\sum_{i\in A}z_i\Bigr)\Bigl(\sum_{i\in B}z_i^{-1}-\sum_{i\in A}z_i^{-1}\Bigr)\Bigr],
\end{equation}
and let $K_{13}=K_{(1,3),(2,4)}$ and $K_{14}=K_{(1,4),(2,3)}$. The derivative blocks are
\begin{equation}
 d_0=\frac{\mathcal D}{4\mathcal R}\bigl[S(1+uw)K_{13}+T(u+w)K_{14}\bigr],\qquad
 d_1=\frac{\mathcal D\mathcal R}{2\mathcal U}\bigl[S(1+uw)K_{13}-T(u+w)K_{14}\bigr].
 \label{eq:cyl:blocksd}
\end{equation}
The apparent inverse powers in $K_{13}$ and $K_{14}$ cancel against the displayed $uvw$ and the block factors, so the
series are analytic at the origin. Their analytic branches are tested at complex arguments against explicit NS and R
form factors (Appendix~\ref{app:cyl:checks}).

\paragraph{Wick pairings, angular projection and denominators.} For each fusion label $\eta$, the block is multiplied
by the chiral $\chi$ factor of one of the three Wick pairings, with the series $H$ of Eq.~\eqref{eq:cyl:Hseries}:
\begin{center}\small
\begin{tabular}{@{}lll@{}}
\toprule
pairing & chiral factor & middle occupation $o$\\
\midrule
adjacent & $H(u)H(w)$ & 0\\
crossed & $H(uv)H(vw)$ & 1\\
nested & $H(uvw)H(v)$ & 1\\
\bottomrule
\end{tabular}
\end{center}
Here $o=1$ means a middle state with two $\chi$ quanta and $o=0$ one with none. The real Taylor coefficients of the
products are written $G_{nm\ell}$, $E_{nm\ell}$ and $D_{nm\ell}$ for the vacuum, thermal and derivative blocks. The
angular projection gives squares for the vacuum and thermal states; for the spin-one derivative state it gives the
shifted product $D_{n+1,m+1,\ell+1}E_{nm\ell}$. For a given pairing and fusion label the chain energies are
\begin{equation}
\begin{aligned}
 E_{0,\eta,o}^{(4),\rm chain}&=-(2\pi)^4q^{2\Delta_{\chi,0}}\sideset{}{'}\sum_{n,m,\ell\ge0}
 \frac{G_{nm\ell}^2}{(2n+2)(2m+\eta+2\Delta_{\chi,0}o)(2\ell+2)},\\
 E_{\mathcal E,\eta,o}^{(4),\rm chain}&=-(2\pi)^4q^{2\Delta_{\chi,0}}\sideset{}{'}\sum_{n,m,\ell\ge0}
 \frac{E_{nm\ell}^2}{(2n+1)(2m-\eta+2\Delta_{\chi,0}o)(2\ell+1)},\\
 E_{d,\eta,o}^{(4),\rm chain}&=-(2\pi)^4q^{2\Delta_{\chi,0}}\sideset{}{'}\sum_{n,m,\ell\ge0}
 \frac{D_{n+1,m+1,\ell+1}E_{nm\ell}}{(2n+1)(2m-\eta+2\Delta_{\chi,0}o)(2\ell+1)}.
\end{aligned}
\label{eq:cyl:chain}
\end{equation}
The prime excludes only a zero denominator, which corresponds to the external-state projector. A negative denominator,
such as that of the lower vacuum intermediate state in a thermal channel, must be retained. Summing over both fusion
labels and all three pairings and forming $2E_{\mathcal E}-E_d-E_0$ gives $X_{4,\rm chain}$. The common negative sign
reflects three resolvents written with positive excitation-energy denominators. At cutoff $N$ the indices run from 0 to
$N-1$; because the derivative block is shifted by one, the coefficient array has shape $(6,N+1,N+1,N+1)$, with the six
channels ordered $(g_0,g_1,e_0,e_1,d_0,d_1)$.

\paragraph{Coefficient extraction and convolution.} The coefficients of the unregulated Ising blocks are extracted by
a discrete Cauchy (Fourier) transform on a three-torus. With $M$ angular points per direction and radius
$\varrho=e^{-\varsigma/(N+1)}$,
\begin{equation}
 [f]_{nm\ell}\simeq\frac{\varrho^{-n-m-\ell}}{M^3}\sum_{j_1,j_2,j_3=0}^{M-1}
 f\bigl(\varrho e^{2\pi\mathrm ij_1/M},\varrho e^{2\pi\mathrm ij_2/M},\varrho e^{2\pi\mathrm ij_3/M}\bigr)\,
 e^{-2\pi\mathrm i(j_1n+j_2m+j_3\ell)/M}.
\end{equation}
$M$ is an even FFT length of about the oversampling factor times $N+1$, and $\varsigma$ is the radius cost.
Oversampling controls aliasing; decreasing the radius reduces aliasing but amplifies roundoff through
$\varrho^{-n-m-\ell}$. The main grids of the base study use oversampling five and radius cost four at $N=256$, and
oversampling four and radius cost $4.5$ at $N=384$. These are numerical choices, not physical regulators; the
extraction check of Sec.~\ref{sec:cyl} compares two grids that differ in the radius. The $\chi$ factors are then
applied by one-dimensional convolutions along lines parallel to the monomials $u$, $w$, $uv$, $vw$, $uvw$ or $v$, and
prefix sums of the weighted three-index arrays return every cubical partial sum at once. The prefix test of
Sec.~\ref{sec:cyl}, which repeats the large-$N$ extrapolation on the sums truncated at $N'<N$, uses these partial sums.
Independent versions of the convolution and of the coefficient grid are kept for validation.

%--------------------------------------------------------------------
\subsection{Collapsed and mixed contributions}\label{app:cyl:collapsed}
\paragraph{One-dimensional two-point contractions.} The second-order energies, the norms of the normalization
subtraction and the mixed spin-two term reduce to one-index sums. Let
\begin{equation}
 p_0(z)=1,\qquad p_{\mathcal E}(z)=\frac{1+z}{2},\qquad p_d(z)=\frac{(1+z)^3}{8},
\end{equation}
and $a_f(j)=[z^j](1-z)^{-1/8}p_f(z)$ for $f\in\{0,\mathcal E,d\}$, with coefficients of negative index zero. With the
regulated sequence $H_j$ of Eq.~\eqref{eq:cyl:Hseries}, set
\begin{equation}
 B_f(n)=\sum_{j=0}^na_f(j)H_{n-j},\qquad \hat B_f(n)=\sum_{j=0}^nj\,a_f(j)H_{n-j}
\end{equation}
(the base study writes $W_f$ for $\hat B_f$). The external states are labelled by
\begin{center}\small
\begin{tabular}{@{}ccccc@{}}
\toprule
state $a$ & left block $l$ & right block $l'$ & shift $s_a$ & $\alpha_a$\\
\midrule
$0$ & $0$ & $0$ & 0 & 2\\
$\mathcal E$ & $\mathcal E$ & $\mathcal E$ & 0 & 1\\
$d$ & $d$ & $\mathcal E$ & 1 & 1\\
\bottomrule
\end{tabular}
\end{center}
With $F_2=(2\pi)^2q^{\Delta_{\chi,0}}$ and $\nu_n=2n+\alpha_a$,
\begin{equation}
\begin{aligned}
 E_a^{(2)}&=-F_2\sum_{n\ge0}\frac{B_l(n+s_a)B_{l'}(n)}{\nu_n},\qquad
 \mathcal N_a:=\langle w_a\vert w_a\rangle=F_2\sum_{n\ge0}\frac{B_l(n+s_a)B_{l'}(n)}{\nu_n^2},\\
 L_a(D_{\rm Ising})&=F_2\sum_{n\ge0}\frac1{\nu_n^2}\Bigl[\bigl(\tfrac18-E_a^{(0)}\bigr)B_l(n+s_a)B_{l'}(n)
 +\hat B_l(n+s_a)B_{l'}(n)+B_l(n+s_a)\hat B_{l'}(n)\Bigr].
\end{aligned}
\label{eq:cyl:collapsed}
\end{equation}
The finite-regulator second-order observable is $X_2(\lambda)=2E_{\mathcal E}^{(2)}-E_d^{(2)}-E_0^{(2)}$ evaluated with
these sums, and the normalization contribution is
\begin{equation}
 X_{4,\rm norm}=-2E_{\mathcal E}^{(2)}\mathcal N_{\mathcal E}+E_d^{(2)}\mathcal N_d+E_0^{(2)}\mathcal N_0.
\end{equation}
These sums are evaluated at much larger cutoffs than the three-index chain: $N_{\rm col}=200\,000$ in the base study's
$N=256$ records and $800\,000$ in its larger-cutoff records, with the small change between them recorded separately.

\paragraph{Two-dimensional thermal insertions.} The mixed thermal term $L_X(V_\epsilon)$ has the insertion $V_\epsilon$
either between the two $V$ vertices (central placement) or outside them (side placement). The basic chiral generating
functions are
\begin{equation}
\begin{aligned}
 G_c(u,v)&=(1-u)^{-1/2}(1-v)^{-1/2}(1-uv)^{3/8}H(uv),\\
 G_s(u,v)&=(1-u)^{-1/2}(1-v)^{3/8}H(v)(1-uv)^{-1/2}.
\end{aligned}
\end{equation}
For either placement the vacuum, energy and derivative factors are $G/\sqrt2$, $GP_{\mathcal E}/(2\sqrt2)$ and
$GP_d/(2\sqrt2)$; the polynomials $P_{\mathcal E,b}$ and $P_{d,b}$, the denominators and the shifted derivative projection are in
Appendix~\ref{app:cyl:mixed}. The common factor is $(2\pi)^3q^{\Delta_{\chi,0}}$, and
\begin{equation}
 L_X(V_\epsilon)=X_c+2X_s,
\end{equation}
where the factor two counts the two side placements. All square partial sums are obtained through $N_{\rm mix}$ (4096
in the base study), and the large-level limit is extrapolated separately (Appendix~\ref{app:cyl:tails}).

%--------------------------------------------------------------------
\subsection{High-state tails and the regulator limit}\label{app:cyl:tails}
\paragraph{The independently determined leading tail.} At fixed $\lambda$ the leading remaining Ising singularity on a
face of the coefficient box is $(1-u)^{-1/8}$. Its coefficients $a_0(n)=(\frac18)_n/n!$ (Appendix~\ref{app:cyl:collapsed})
behave as $a_0(n)\sim n^{-7/8}/\Gamma(\frac18)$, so
\begin{equation}
 \sum_{n\ge N}\frac{a_0(n)^2}{2n}\sim\frac{2}{7\Gamma(\frac18)^2}N^{-7/4};
 \label{eq:cyl:tailpower}
\end{equation}
the factor $\frac27$ follows by integrating $n^{-11/4}/2$. Exact limits of the generating functions on the faces of the
coefficient box determine the coefficient of this tail, with the sign convention
\begin{equation}
 X_{4,\rm chain}(\infty)-X_{4,\rm chain}(N)=A_{\rm face}(\lambda,r)\,N^{-7/4}+\text{subleading terms}.
\end{equation}
Appendix~\ref{app:cyl:facetail} gives the two-index construction of $A_{\rm face}$. It uses the block singularities and the
collapsed norms, not the observed $N$ sequence or any value of $c_2$. This is the analytically derived coefficient of
the face tail that the continuum analysis of Sec.~\ref{sec:cyl} subtracts.

\paragraph{What is still fitted at fixed regulator.} The subleading powers include the one-face corrections
$N^{-11/4},N^{-15/4},\ldots$ and overlapping-face contributions such as $N^{-7/2}$. The regulated $\chi$ kernel has its
next singularities at $e^{\lambda}$, which generate terms exponentially small in $N\lambda$; on finite windows they can
compete with the power terms. A reasonably large $\lambda N$ is therefore desirable, but this is a planning criterion,
not a remainder bound (the analysis of Sec.~\ref{sec:cyl} flags points with $\lambda N$ below a minimum). Let $h_n(\lambda)$
be the assembled quantity, Eq.~\eqref{eq:cyl:assembly}, with the chain truncated to the cube $[0,n)^3$. With the leading
tail held fixed, the corrected sequence is fitted on $n_{\min}\le n\le N$ as
\begin{equation}
 h_n(\lambda)+A_{\rm face}\,n^{-7/4}=h(\lambda)+\sum_j\xi_j\Bigl(\frac nN\Bigr)^{-p_j}+\text{exponential terms},
 \label{eq:cyl:spectralfit}
\end{equation}
with the power sets
\begin{equation}
 \{2.75,3.75\},\quad\{2.75,3.75,4.75\},\quad\{2.75,3.5,3.75,4.75\},\quad\{2.75,3.5,3.75,4.75,5.5\}
\end{equation}
and, for exponential order $n_e=0,1,2,3$, the optional columns
\begin{equation}
 e^{-\lambda(n-n_{\min})}\Bigl(\frac nN\Bigr)^{r-1-j},\qquad e^{-2\lambda(n-n_{\min})}\Bigl(\frac nN\Bigr)^{2r+3/4-j},
 \qquad j=0,\ldots,n_e-1.
\end{equation}
The base study's selected ensemble uses at least three power terms (the analysis of Sec.~\ref{sec:cyl}: three or
four), $n_e=1,2$, and the two lower endpoints $n_{\min}=240,288$ at $N=384$
and $160,192$ at $N=256$, that is $n_{\min}=5N/8$ and $3N/4$. The fixed-regulator estimate $h(\lambda)$ is the median
of the selected limits, and half the difference between the largest and the smallest is the \emph{spectral-model
half-spread}, a model diagnostic, neither a standard deviation nor a truncation bound. Small residuals alone do not
prove the tail model correct. The continuum analysis of Sec.~\ref{sec:cyl} uses these power sets, exponential columns
and selection, with the windows $n_{\min}=5N/8$ and $3N/4$ at every $N$. That these windows reproduce the base
study's relative windows does not by itself show the same tail model valid at smaller $\lambda$; the prefix and
cross-$N$ checks of Sec.~\ref{sec:cyl} test it.

The mixed thermal sums have their own tail fits in $N_{\rm mix}$, with powers drawn from $\{1,2,\frac{11}4,3,4\}$; the
chosen mixed limit is the median of the models with at least three powers and $n_{\min}\ge N_{\rm mix}/2$ (2048 in the
base study). Its spread is recorded separately and not identified with the error of the chain.

\paragraph{Removing the smooth regulator.} After the spectral limit, the regulator model is
$h(\lambda)=h_0+k_1\lambda^2+k_2\lambda^2\log\lambda+k_3\lambda^3$, the preferred model of Sec.~\ref{sec:cyl}. The
improved low-frequency filter and the removal of the leading cylinder thermal artifact motivate this form; higher terms
are not rigorously bounded. The base study varied the regulator order, the $\lambda$ window, the spectral-tail model,
the weighting and the last power ($\lambda^4$ instead of $\lambda^3$), and it tested omission and scale shifts of the
term~\eqref{eq:cyl:couplinglog}. No fit has the physical $\delta$ as its variable, but the uncertainty remains an
extrapolation uncertainty. A four-parameter model fitted to four regulator points is saturated, and a near-zero residual
is then no evidence of accuracy; two of the base study's three regulator families had only four points, which is why more
regulator points were needed. The base study's own extrapolations ($N\le384$), quoted in Sec.~\ref{sec:cyl}, were
superseded by the analysis given there; its archived records were not: eleven of them are points of that analysis.

For a fixed linear extrapolation $\hat h_0=\sum_iw_i\hat h(\lambda_i)$ with input diagnostics $\varepsilon_i$ and a
regulator-model term $\varepsilon_{\rm model}$, the error in $c_2$ is propagated as
\begin{equation}
 \frac1{b_0^2}\Bigl(\sum_i|w_i|\varepsilon_i+\varepsilon_{\rm model}\Bigr).
 \label{eq:cyl:budget}
\end{equation}
This is the form of the linear sum of Sec.~\ref{sec:cyl}. It would be a bound only if the $\varepsilon_i$ and
$\varepsilon_{\rm model}$ were valid bounds, which the half-spreads and model spreads are not. It also shows why the
fixed-$\lambda$ inputs need better precision than the desired final error. As $\lambda$ decreases, cancellations between
large raw terms and the counterterms make roundoff more important.

%--------------------------------------------------------------------
\subsection{Algebraic and finite-sum checks}\label{app:cyl:checks}
The validation suite tests different failure modes:
\begin{center}\small
\begin{tabular}{@{}>{\raggedright\arraybackslash}p{.62\textwidth}r@{}}
\toprule
check & largest reported discrepancy\\
\midrule
chiral Virasoro Ward identities (252 identities) & $\sci{4.22}{-16}$\\
exact second-order shell weights (vacuum, energy and derivative correlators) & below $\sci{5.4}{-16}$ relative\\
explicit basis against blocks, fourth-order finite sums & below $\sci{2}{-13}$\\
blocks at complex arguments against finite-level form factors & about $\sci{1.16}{-10}$ (scaled)\\
alternative finite convolution and sequence implementations & below $\sci{2}{-11}$\\
\bottomrule
\end{tabular}
\end{center}
The suite also checks the four-spin block sums and the fourth-order formula~\eqref{eq:cyl:RS4} against an independent
finite-matrix recursion for the perturbed eigenvector. The discrepancy of the complex-argument blocks includes the
truncation of the explicit form factors. The entries are precisions of cross-checks on finite problems. None of these
identities bounds the omitted continuum tail; they establish that signs, mode occupations, derivative shifts and
normalization subtractions agree between independent implementations on the tested finite problems.

%--------------------------------------------------------------------
\subsection{The protected magnetic dimension as a matching check}\label{app:cyl:magnetic}
At the fixed point the magnetic dimension is protected, $\Ds=\frac18+\delta$ in this
convention~\cite{Behan:2017emf}. With $g_*^2=\delta/b_0$, the leading renormalized Hamiltonian gap coefficient must
therefore satisfy
\begin{equation}
 [g^2](E_\sigma-E_0)=b_0.
\end{equation}
The independently derived local counterterms, including the cylinder improvement of Eq.~\eqref{eq:cyl:tfinal}, give at
$r=3$ and $\lambda=1/32$ the value $1.268404004672631$ quoted in Sec.~\ref{sec:cyl}, which differs from $b_0$ by
$\sci{-3.04}{-7}$. This value was not used to tune the counterterms. It is a lower-order check of the normalization,
not a fourth-order error bound.

%--------------------------------------------------------------------
\subsection{What is exact and what is numerical}\label{app:cyl:status}
The second-order meromorphic continuation (Appendix~\ref{app:cyl:second}) and the finite-basis identities are exact
analytic or algebraic results. The counterterm classification and the cancellation~\eqref{eq:cyl:W4} remove the
unspecified direct local fourth-order constants within the stated smoothly regulated prescription; this is an
argument, not a proof. The fourth-order continuum number is a numerical inference, supported by independent regulator
families, the analytically derived leading tail and the lower-order normalization checks; its remaining uncertainty is
systematic, not statistical. What is not supplied is a rigorous
uniform bound on every nested spectral and regulator remainder, or an independent many-digit enclosure of $h_0$.

%--------------------------------------------------------------------
\subsection{Checks of the closed form made in the base study}\label{app:cyl:closedchecks}
The base study also consolidates a set of formula, elliptic and local checks of the closed-form study
(Secs.~\ref{sec:closed2} and~\ref{sec:unproved}), likewise written by GPT-6 astra. Their code is not in the release
(Appendix~\ref{app:cannot}); the value of the closed expression is reproduced, not as an interval enclosure, by
\file{c2_closed_d2}. None of these quantities is an input
to the Hamiltonian sums. They certify the value of the closed expression~\eqref{eq:c2closed} and test local steps of
the route to Eq.~\eqref{eq:c2I}; they do not prove the identification of the expression with the physical coefficient,
which rests on the step of Sec.~\ref{sec:unproved}.

\paragraph{The certified expression.} The expression is Eq.~\eqref{eq:c2closed}, with
${}_pF_q(a_1,\ldots,a_p;b_1,\ldots,b_q;z)=\sum_n\prod_i(a_i)_n/\prod_j(b_j)_n\,z^n/n!$. All three series have excess
two and absolutely convergent summands of order $n^{-3}$. An integer-interval implementation that uses no Gamma or
generalized-hypergeometric evaluator and imports no code from the closed-form study encloses the value, with the tail
construction below at $n_0=4000$, in an interval of width about $\sci{2.69823801745}{-53}$, and so proves the rounding
\begin{equation}
 c_2^{\rm closed}=-4.066270148132521552740027608572876487529901141186618
 \label{eq:cyl:certified}
\end{equation}
to 51 decimal places. This is a computer-assisted proof about the specified expression. It is not a physical error
interval for the Hamiltonian extrapolation, nor a proof of the regulator-identification step of the closed-form study.

\paragraph{Exact tail inequalities and elementary constant bounds.} For the absolute value $v_n$ of a term of one of
the three series, the ratio is rational,
\begin{equation}
 \rho_n=\frac{v_{n+1}}{v_n}=\frac{\prod_{j=1}^4(n+a_j)}{(n+1)\prod_{j=1}^3(n+b_j)}.
\end{equation}
Rational functions $\phi(n)=\frac n2+\phi_0+\phi_1/n+\cdots+\phi_{12}/n^{12}$ are constructed by exact cancellation of
the leading powers in $\Xi(n)=\phi(n)-1-\rho_n\phi(n+1)$; changing the last coefficient by $-1$ and $+1$ gives $\phi_-$
and $\phi_+$. After a positive denominator is cleared and $n=4000+m$ substituted, nonnegative polynomial coefficients
prove
\begin{equation}
 \phi_-(n)-1-\rho_n\phi_-(n+1)\le0,\qquad \phi_+(n)-1-\rho_n\phi_+(n+1)\ge0\qquad\text{for all }m\ge0.
\end{equation}
A second exact polynomial inequality bounds $v_n\le v_{n_0}(n_0/n)^2$ and controls the telescoping boundary at
infinity. Hence
\begin{equation}
 \phi_-(n_0)\,v_{n_0}\le\sum_{n=n_0}^\infty v_n\le\phi_+(n_0)\,v_{n_0}.
\end{equation}
The terms of $\mathsf T$ after index zero are negative, so its enclosed tail is subtracted with the endpoints reversed.
For $\pi$ the implementation uses the alternating rational sums of Machin's identity
$\pi=16\arctan\frac15-4\arctan\frac1{239}$. For $R_\Gamma^2$ it uses $B(\frac14,\frac12)^2/\pi$ and the convergent split
beta integral
\begin{equation}
 B\bigl(\tfrac14,\tfrac12\bigr)=2^{-1/4}\sum_{n\ge0}\frac{(\frac12)_n}{n!}\frac{2^{-n}}{n+\frac14}
 +2^{-1/2}\sum_{n\ge0}\frac{(\frac34)_n}{n!}\frac{2^{-n}}{n+\frac12}.
\end{equation}
Both positive term ratios are below $\frac12$, so after 420 terms each remainder lies between its first omitted term and
twice that term. Integer root inequalities enclose the powers of two. All arithmetic uses outward-rounded integer
endpoints at scale $10^{110}$, and assertions check the exact sign certificates. The closed-form study's own tail check behind the envelope quoted in Sec.~\ref{sec:closed2} (Appendix~\ref{app:cform:numerics}) is of the same
type, with $\phi=\frac n2+{\rm const}$ and floating-point endpoints; the higher-order $\phi$ and the integer arithmetic
are what make this enclosure a certified rounding.

\paragraph{Elliptic integration from the period integrands.} This check starts from the period integrands, not from the final series.
With $K(t)=\int_0^{\pi/2}\dd\theta\,(1-t\sin^2\theta)^{-1/2}$ (parameter convention) and
$\mathbf k(t)=\bigl(K(t)^2,\,K(t)K(1-t),\,K(1-t)^2\bigr)^{\mathsf T}$, define
\begin{equation}
 \mathbf P(a,b)=\AC\!\int_0^1t^{a-1}(1-t)^{b-1}\mathbf k(t)\dd t,\qquad \mathbf p(a,b)=\frac4{\pi^2}\mathbf P(a,b),
\end{equation}
where negative exponents mean meromorphic continuation, not a divergent improper integral. Near zero
\begin{equation}
 \mathbf k(t)\sim\Bigl(\frac{\pi^2}4,\ \frac\pi4\log\frac{16}t,\ \frac14\log^2\frac{16}t\Bigr)^{\mathsf T}.
\end{equation}
The quadrature splits at $\frac12$, subtracts these endpoint terms, restores their elementary continued Mellin
integrals and integrates the convergent remainder after the substitution $t=u^{12}/2$; Carlson's $R_F$ supplies $K$
independently of the Frobenius arrays of the closed-form study. One nonsymmetric assembly is
\begin{equation}
\begin{gathered}
 \mathbf a=2\mathbf p(-\tfrac12,\tfrac74)+5\mathbf p(\tfrac12,\tfrac74),\qquad
 \mathbf b=5\mathbf p(-\tfrac14,\tfrac74)+2\mathbf p(\tfrac34,\tfrac74),\qquad
 \mathsf H=\begin{pmatrix}0&0&1\\0&2&2\\1&2&-1\end{pmatrix},\\
 c_2=-\frac{R_\Gamma^2}9+\frac{\pi R_\Gamma^4}{10368\sqrt2}\,\mathbf a^{\mathsf T}\mathsf H\,\mathbf b,
\end{gathered}
\end{equation}
and a separate assembly uses symmetric moments. The observed differences from the enclosed value are
\begin{center}\small
\begin{tabular}{@{}rll@{}}
\toprule
quadrature order & nonsymmetric periods & symmetric moments\\
\midrule
96 & $\sci{1.74}{-19}$ & $\sci{4.78}{-29}$\\
192 & $\sci{5.24}{-23}$ & $\sci{2.28}{-34}$\\
384 & $\sci{1.51}{-26}$ & $\sci{1.04}{-39}$\\
\bottomrule
\end{tabular}
\end{center}
These are observed agreements under refinement, not rigorous quadrature-error bounds. The closed-form study's own
endpoint-subtracted quadrature with Carlson's integrals (Sec.~\ref{sec:closed2}, Appendix~\ref{app:cform:numerics}) is a check of the same kind,
implemented separately.

\paragraph{Complex-plane normalization.} With the single-valued area pairing
$\mathcal A(t)=8\bigl[K(t)\overline{K(1-t)}+\overline{K(t)}K(1-t)\bigr]$ (the torus area of Eq.~\eqref{eq:I}) and the
continued complex moment
\begin{equation}
 \mathcal M(a,b)=\AC\!\int_{\mathbb C}\dd^2t\,|t|^{2a-2}|1-t|^{2b-2}\mathcal A(t)^2,
\end{equation}
the tested period relation is
\begin{equation}
 \mathcal M(a,b)=64\,\mathbf P(a,b)^{\mathsf T}\mathsf L(b)\,\mathbf P(2-a-b,b),\qquad
 \mathsf L(b)=\begin{pmatrix}0&0&\sin\pi b\\0&2\sin\pi b&-2\cos\pi b\\ \sin\pi b&-2\cos\pi b&-\sin\pi b\end{pmatrix}.
\end{equation}
The ordinary convergent moment $\mathcal M(\frac23,\frac23)$ was evaluated directly on the plane and through the
real-period pairing, with the values $24181.76549557689$ and $24181.76549557647$, a relative difference of
$\sci{1.72}{-14}$. Generic convergent moments also agreed. Symbolic checks gave zero residual for the endpoint
differential identity, the cut-jump matrix, the reflection reduction and the final symmetric matrix reduction. These
tests support the period identity and the normalization in the regimes checked; they do not supply a uniform proof for
the physical collisions.

\paragraph{Local contacts.} The finite localized-source term $B_{31}$ of Sec.~\ref{sec:d2} was checked with an
elementary-integral residual of about $\sci{3.17}{-71}$ at 70 working digits. Direct polar quadrature of the local
branch contact, at the source parameter $M=3$, agreed with
\begin{equation}
 A_{\rm br}=4B_{31}-\frac{\pi^2}4=-3.8363599279686606009\ldots
\end{equation}
to about $\sci{1.05}{-14}$; other source parameters gave comparable local agreement. Fixed-shape neutral-charge checks
supported the cancellation of the energy-primary term of the shrinking spin pair. In the minimal-subtraction assembly,
Eq.~\eqref{eq:c2I}, the local calculation supports the factor $3/4!=\frac18$ and the two placements of the coupling
counterterm. The two contributions to $c_2$ there are $-2.12099049015343\ldots$ and $-1.94527965797909\ldots$; the
agreement of the second with $a_1/b_0$ of Appendix~\ref{app:cyl:second} is a comparison made after both calculations, not
an input to Eq.~\eqref{eq:cyl:assembly}.

\paragraph{The global qualification remains separate.} The closed-form study states that its audit does not
reconstruct every sectorwise subtraction or prove a regulator-uniform remainder bound across nested collisions
(Sec.~\ref{sec:unproved}). Agreement of local contact terms at fixed nondegenerate modulus does not by itself control
the degeneration regions at $0$, $1$ and $\infty$, or a simultaneous shrinking of relative scales. The base study's own review of
these checks (GPT-6 astra; not the review of Sec.~\ref{sec:sources}) reached the same boundary: no error that changes $c_2$ was found in the tested expressions, but the global regulator comparison
was not proved independently. It also noted a point of exposition: the statement that identity projections cancel
should distinguish disconnected identity terms from an identity-channel power term with denominator $-1+2\delta$, which
is retained and analytically continued; this is not a numerical error. The Hamiltonian result is a different physical
consistency check on Eq.~\eqref{eq:c2I}; it does not turn these local checks into a proof of the global comparison.

%--------------------------------------------------------------------
\subsection{Explicit Majorana form factors}\label{app:cyl:formfactors}
Let $\mathbf m=(m_1,\ldots,m_{n_{\rm R}})$ be increasing positive integer R occupations and
$\mathbf s=(s_1,\ldots,s_{n_{\rm NS}})$ increasing positive half-integer NS occupations, and
\begin{equation}
 f_{\rm R}(n)=\binom{2n}{n}4^{-n},\qquad f_{\rm NS}(s)=\binom{2s-1}{s-\frac12}4^{-(s-1/2)}.
\end{equation}
In the real chiral gauge of the implementation (\file{code/c2_cylinder_d2/code/hamiltonian/ising_basis.py}) the chiral
spin-field element between these states is
\begin{equation}
\begin{aligned}
 \mathcal S(\mathbf m,\mathbf s)={}&(-1)^{(n_{\rm NS}\bmod2)\,n_{\rm R}}\,2^{-(n_{\rm R}+n_{\rm NS})/2}
 \prod_{i=1}^{n_{\rm R}}f_{\rm R}(m_i)\prod_{j=1}^{n_{\rm NS}}f_{\rm NS}(s_j)\\
 &\times\prod_{i<j}\frac{m_i-m_j}{m_i+m_j}\prod_{i<j}\frac{s_i-s_j}{s_i+s_j}\prod_{i,j}\frac{m_i+s_j}{m_i-s_j},
\end{aligned}
\label{eq:cyl:formfactors}
\end{equation}
with empty products equal to one. The full spin-field element is the product of the left- and right-chiral factors.
This representation, including its phase convention, is checked by chiral Virasoro Ward identities; it gives
$\langle\sigma\rvert\sigma\ket0=1$ and $C_{\sigma\sigma\epsilon}=\frac12$ in the chosen normalization. Related
scaling-limit product formulas are in Ref.~\cite{Liu:2026}.

For a fermion mode that creates or removes a positive occupied mode $m$, the sign is $(-1)^{\#\{k<m\}}$, where
$\#\{k<m\}$ counts the occupied modes below $m$. The diagonal element of the Ramond zero mode is $(-1)^{\#{\rm occupied}}/\sqrt2$. Products of left
and right fermion elements give the thermal field. Angular integration forces equal signed energy changes on the two
sides, so the integrated thermal matrix connects pair creation and annihilation and the diagonal of the Ramond zero
mode; it is implemented as a sparse matrix, whose angular factor $2\pi$ enters only when the Hamiltonian insertion is
assembled. A basis state with left and right occupations $\boldsymbol\ell$, $\bar{\boldsymbol\ell}$ has energy and spin
\begin{equation}
 E=\sum_i\ell_i+\sum_j\bar\ell_j+\begin{cases}0&\text{NS},\\ \frac18&\text{R},\end{cases}\qquad
 J=\sum_i\ell_i-\sum_j\bar\ell_j,
\end{equation}
and in NS only equal left and right parity is retained. The finite-basis checks built on these elements are those of
Appendix~\ref{app:cyl:checks}; they detect basis-sign and projection errors without assuming the continuum answer.

%--------------------------------------------------------------------
\subsection{Compact quadratures for the local counterterms}\label{app:cyl:quadratures}
The constants of Appendix~\ref{app:cyl:counterterms} are computed from the following finite-domain integrals
(\file{code/c2_cylinder_d2/code/hamiltonian/improved_regulator.py}). Let
\begin{equation}
 \mathcal P_{\alpha,r}(z)=\sum_{j=0}^r\frac{(\alpha)_j}{j!}z^j.
\end{equation}
For $a,b\in(0,1)$ put $Z=(1-a)\bigl[(1-b^2)+\mathrm ib\sqrt{1-b^2}\bigr]$. The thermal coefficient is
\begin{equation}
 c_{t,r}=\int_0^1\dd a\int_0^1\dd b\,(a^2+b^2-a^2b^2)^{3/8}\,|\mathcal P_{\Delta_{\chi,0},r}(Z)|^2.
 \label{eq:cyl:ctcompact}
\end{equation}
For $t\in(0,1)$ and $\varphi\in(0,\pi/2)$ put
\begin{equation}
 D_t=t^2+2t(1-t)\cos\varphi+(1-t)^2,\qquad Z_t=\frac{1-t}{1-t+te^{\mathrm i\varphi}}.
\end{equation}
The finite spin-two change is
\begin{equation}
 \kappa_r-\kappa_0=-\pi\int_0^1\dd t\int_0^{\pi/2}\dd\varphi\,\frac{\cos(2\varphi)\,t^{11/4}}{(1-t)\,D_t^{\Delta_{\chi,0}}}
 \Bigl(|\mathcal P_{\Delta_{\chi,0},r}(Z_t)|^2-1\Bigr).
 \label{eq:cyl:kcompact}
\end{equation}
For the cylinder term put $y=(1-t)/t$ and
\begin{equation}
 F(t,\varphi)=\frac{{\rm Re}\bigl[(e^{\mathrm i\varphi}+y)^2\,\mathcal P_{\Delta_{\chi,0}-2,r}(Z_t)\,
 \overline{\mathcal P_{\Delta_{\chi,0},r}(Z_t)}\bigr]}{(1+2y\cos\varphi+y^2)^{\Delta_{\chi,0}}};
\end{equation}
then
\begin{equation}
 c_{{\rm cyl},r}=-\frac{\Delta_{\chi,0}}{12}\int_0^1\dd t\int_0^{\pi/2}\dd\varphi\,\frac{F(t,\varphi)-\cos(2\varphi)}{(1-t)^2},
 \qquad r\ge1.
 \label{eq:cyl:cylcompact}
\end{equation}
The factor two relative to the half-plane expression comes from reflection in $x$. Tensor Gauss--Legendre rules
evaluate these integrals after their domains are mapped to the unit interval; orders 128, 256 and 512 are used for
refinement, and the constants are evaluated with 512 points in each direction. The subtraction stays inside the
integrand: numerically subtracting separately divergent moments would not define the same calculation.

%--------------------------------------------------------------------
\subsection{Mixed thermal polynomials and denominators}\label{app:cyl:mixed}
For the central and side kernels of Appendix~\ref{app:cyl:collapsed}, the energy polynomials are
\begin{equation}
 P_{\mathcal E,c}=1+uv-2u-2v,\qquad P_{\mathcal E,s}=-2+u+uv-2u^2v,
\end{equation}
and the derivative polynomials are
\begin{equation}
\begin{aligned}
 P_{d,c}&=\frac{uv+1}{4}\bigl(u^2v^2+2u^2v-4u^2+2uv^2-6uv+2u-4v^2+2v+1\bigr),\\
 P_{d,s}&=-\frac{v+1}{4}\bigl(4u^4v^2-2u^3v^2-2u^3v-u^2v^2+6u^2v-u^2-2uv-2u+4\bigr).
\end{aligned}
\label{eq:cyl:mixedpolys}
\end{equation}
For placement $b=c,s$ the coefficient arrays are
\begin{equation}
 V^b_{nm}=[u^nv^m]\frac{G_b}{\sqrt2},\qquad E^b_{nm}=[u^nv^m]\frac{G_bP_{\mathcal E,b}}{2\sqrt2},\qquad
 D^b_{nm}=[u^nv^m]\frac{G_bP_{d,b}}{2\sqrt2},
\end{equation}
and the derivative angular projection uses $D^b_{n+1,m+1}E^b_{nm}$. The denominators are
\begin{center}\small
\begin{tabular}{@{}lll@{}}
\toprule
placement & vacuum, $d^b_0(n,m)$ & energy and derivative, $d^b_{\mathcal E}(n,m)$\\
\midrule
central & $(2n+2)(2m+2)$ & $(2n+1)(2m+1)$\\
side & $(2n+1)(2m+2)$ & $(2n-1)(2m+1)$\\
\bottomrule
\end{tabular}
\end{center}
so that
\begin{equation}
 X_b=(2\pi)^3q^{\Delta_{\chi,0}}\sum_{n,m\ge0}\biggl[\frac{(2E^b_{nm}-D^b_{n+1,m+1})E^b_{nm}}{d^b_{\mathcal E}(n,m)}
 -\frac{(V^b_{nm})^2}{d^b_0(n,m)}\biggr].
\end{equation}
The side denominator is negative at $n=0$ and must be kept. The total mixed term is $L_X(V_\epsilon)=X_c+2X_s$. A
function of the form $A(u)B(v)C(uv)$ is generated by diagonal convolutions,
\begin{equation}
 [u^nv^m]A(u)B(v)C(uv)=\sum_{j=0}^{\min(n,m)}A_{n-j}B_{m-j}C_j,
 \label{eq:cyl:diagconv}
\end{equation}
and the finite polynomials are applied by shifts. This reduces the mixed calculation to two-index arrays and permits
the much larger $N_{\rm mix}$.

%--------------------------------------------------------------------
\subsection{The leading face-tail coefficient in full}\label{app:cyl:facetail}
The coefficient $A_{\rm face}$ of Eq.~\eqref{eq:cyl:tailpower} is fixed independently of any fitted cutoff sequence. It
is derived from exact face limits of the blocks but evaluated numerically through convergent two-index sums
(\file{code/c2_cylinder_d2/code/hamiltonian/exact_leading_tail.py}; ``exact'' in the file name describes the source of
the asymptotic formula, not an interval enclosure of its floating-point evaluation). With the regulated $H$ of
Eq.~\eqref{eq:cyl:Hseries} define
\begin{equation}
 \Phi_B(u,v)=H(u)(1-v)^{-1/8}H(uv),\qquad \Phi_A(u,v)=H(u)H(v)(1-uv)^{-1/8},
\end{equation}
with Taylor coefficients $B_{ij}$ and $A_{ij}$, and set
\begin{equation}
\begin{aligned}
 E^B_{ij}&=[u^iv^j]\frac{1+v}{2}\Phi_B, & D^B_{ij}&=[u^iv^j]\frac{(1+v)^3}{8}\Phi_B,\\
 E^A_{ij}&=[u^iv^j]\frac{1+uv}{2}\Phi_A, & D^A_{ij}&=[u^iv^j]\frac{(1+uv)^3}{8}\Phi_A.
\end{aligned}
\end{equation}
The two outer faces are organized by
\begin{equation}
 O=\sum_{i,j\ge0}\frac1{2i+2\Delta_{\chi,0}}\biggl[\frac{(2E^B_{ij}-D^B_{i,j+1})E^B_{ij}}{2j+1}-\frac{B_{ij}^2}{2j+2}\biggr],
\end{equation}
the adjacent middle face by
\begin{equation}
 M_{\rm adj}=\sum_{i,j\ge0}\biggl[\frac{(2E^A_{ij}-D^A_{i+1,j+1})E^A_{ij}}{(2i+1)(2j+1)}-\frac{A_{ij}^2}{(2i+2)(2j+2)}\biggr],
\end{equation}
and the nested middle face by the collapsed norms of Eq.~\eqref{eq:cyl:collapsed},
\begin{equation}
 \mathcal N_X=2\mathcal N_{\mathcal E}-\mathcal N_d-\mathcal N_0,\qquad
 M_{\rm nest}=\frac{H(1)^2\,\mathcal N_X}{(2\pi)^2q^{\Delta_{\chi,0}}}.
\end{equation}
The leading coefficient is
\begin{equation}
 \boxed{A_{\rm face}=-(2\pi)^4q^{2\Delta_{\chi,0}}\,\frac{2}{7\Gamma(\frac18)^2}\bigl(4O+2M_{\rm adj}+M_{\rm nest}\bigr).}
 \label{eq:cyl:Aface}
\end{equation}
The base study evaluates it with two-index cutoff 2048 and records $O$, $M_{\rm adj}$, $M_{\rm nest}$, their
combination and the sign convention; refinement at smaller and larger auxiliary cutoffs and a comparison with
independently fitted leading tails are kept as diagnostics. The continuum analysis of Sec.~\ref{sec:cyl} sets this
two-index cutoff to a common value for all points and propagates its shift under a change (the ``cutoff of the
leading-tail calculation'' in its error budget). Once Eq.~\eqref{eq:cyl:Aface} is fixed, the residual spectral models of Appendix~\ref{app:cyl:tails} account for
the subleading powers and exponentials.

%--------------------------------------------------------------------
\subsection{Minimal end-to-end algorithm}\label{app:cyl:algorithm}
For an implementation in another language the logical order is:
\begin{enumerate}[leftmargin=2em]
\item Fix the unit-field normalization, $\Delta_{\chi,0}=\frac{15}8$, $a_1=-\pi^2/4$ and the common lower-order $b_0$.
 Choose $r$ and $\lambda>0$. No value of $c_2$ enters the solver.
\item Generate the six Ising block coefficient arrays of Eqs.~\eqref{eq:cyl:blocksge} and~\eqref{eq:cyl:blocksd} through
 one index beyond the desired cube; check branches, reality and selected low coefficients against explicit form
 factors.
\item Construct $H_n$ from Eq.~\eqref{eq:cyl:Hseries}, apply the three Wick-pairing convolutions, evaluate
 Eq.~\eqref{eq:cyl:chain} for each pairing and channel, and return every cubical partial sum.
\item Compute the collapsed energies, norms and spin-two mixed terms of Eq.~\eqref{eq:cyl:collapsed} at a much larger
 cutoff, including the normalization subtraction.
\item Generate the two-index thermal kernels and polynomials of Appendix~\ref{app:cyl:mixed} and extrapolate their square
 cutoff separately.
\item Evaluate the compact counterterm integrals of Appendix~\ref{app:cyl:quadratures}, form $t_r(\lambda)$ and
 $-\kappa_r$, and assemble Eq.~\eqref{eq:cyl:assembly}, including the exact thermal-squared term.
\item Compute $A_{\rm face}$ from Appendix~\ref{app:cyl:facetail} and extrapolate the corrected chain sequence,
 Eq.~\eqref{eq:cyl:spectralfit}, over the predeclared power sets, exponential terms and windows.
\item Repeat at several $\lambda$ for each regulator order, take the regulator limit only after the spectral limits,
 and keep model sensitivity and fits with residual degrees of freedom.
\item Convert $h_0$ to $c_2=a_1/b_0+h_0/b_0^2$, and compare with the closed expression only at this final stage.
\end{enumerate}
A faithful reproduction also tests the second-order endpoint anomaly (Appendix~\ref{app:cyl:anomaly}), the magnetic
normalization (Appendix~\ref{app:cyl:magnetic}), the finite-box equivalence of the two implementations
(Appendix~\ref{app:cyl:checks}) and the sensitivity to the coefficient grid; matching a final decimal number alone is not an
adequate test.

%====================================================================
\section{\texorpdfstring{$c_1$}{c1} in three dimensions: derivation and conditional control of the omitted spectrum}\label{app:cone}
%====================================================================
This appendix gives the derivation and the method of the $c_1$ package (written by OpenAI GPT-6 astra, in its review
revision of 8 October 2026; Sec.~\ref{sec:sources}), edited to the notation of this note. Section~\ref{sec:d3} states
the ratio formula, the retained value and its comparison with BRRZ, the implementations, the hypotheses of the
estimated range, and the feasibility, box and gap tests; they are not repeated here. Throughout this appendix
\begin{equation}
 S\equiv\Ds^{\rm SR}=0.518148806(24),\qquad E\equiv\De^{\rm SR}=1.41262528(29),\qquad
 \delta=\frac{s_*-s}{2},
\end{equation}
as in Table~\ref{tab:conv}, and $\epsilon$ is the energy operator of the 3d Ising CFT. Sans-serif letters
($\mathsf S$, $\mathsf E$, $\mathsf M$, $\mathsf V$) denote block sums, never these two dimensions. All long decimal
strings below are numerical regression targets of the package, not physical significant digits.

The coefficient $c_1$ is the right derivative of $\De$ with respect to $\delta$ at $\delta=0$; the long-range side has $\delta>0$. With the retained value
of Sec.~\ref{sec:retained} the expansion reads
\begin{equation}
 \De\simeq1.41262528+0.27315242\,\delta=1.41262528+0.13657621\,(s_*-s),
\end{equation}
with the remainder of the expansion understood. It describes the local slope at $s_*$ only, not a controlled
extrapolation over a large range of $s$.

\subsection{Three meanings of precision}\label{app:cone:precision}
Three kinds of uncertainty are kept apart throughout.
\begin{itemize}
\item \emph{Numerical.} For a fixed finite list of primary operators with fixed dimensions and operator-product-expansion
 (OPE) coefficients, the integrals can be evaluated to many digits. Numerical stability then tests the descendant
 truncation, the quadrature and the subtraction. By itself it says nothing about omitted primaries or about errors of
 the input correlator.
\item \emph{Parametric.} Letting the retained dimensions and couplings vary gives a second uncertainty. The spectrum
 tables do not provide a calibrated joint probability distribution for their error estimates, so the calculation uses
 explicit ranges and crossing constraints, not a statistical error model.
\item \emph{Physical.} A physical interval must also control the missing conformal families and the possible
 incompleteness of the selected low-dimensional spectrum. The interval of Sec.~\ref{sec:range} is conditional on
 explicit hypotheses about that missing spectrum and on numerical allowances for nonlinear effects. Whenever an
 interval for $c_1$ is written, its assumptions are part of the statement.
\end{itemize}

\subsection{Reduction to short-range conformal perturbation theory}\label{app:cone:cpt}
\paragraph{The product theory and the perturbing operator.} Near the crossover the fixed point is described by the
product of the short-range Ising CFT and a generalized free scalar $\chi$~\cite{Behan:2017dwr,Behan:2017emf}. All
scalar two-point functions are normalized to unity,
\begin{equation}
 \langle\sigma(x)\sigma(0)\rangle=|x|^{-2S},\qquad
 \langle\epsilon(x)\epsilon(0)\rangle=|x|^{-2E},\qquad
 \langle\chi(x)\chi(0)\rangle=|x|^{-2\Delta_\chi},
\end{equation}
with
\begin{equation}
 \Delta_\chi=3-S-\delta,\qquad \mathcal O=\sigma\chi,\qquad \Delta_{\mathcal O}=3-\delta .
\end{equation}
The deformation is $g\mu^\delta\int\dd^3x\,\sigma\chi$, with $\mu$ the renormalization scale, and the relevant
thermal coupling is tuned to the critical trajectory. The generalized free field has Wick-factorized correlators; the
Ising factors do not. With a dimensionless renormalized coupling $g$,
\begin{equation}
 \beta(g)=-\delta g+b_0g^3+O(g^5),\qquad \gamma_{\mathcal E}(g)=a_0g^2+O(g^4).
\end{equation}
There is no $g^2$ term in $\beta$, because the undeformed product theory forbids the relevant three-point function
($\langle\mathcal O\mathcal O\mathcal O\rangle=\langle\sigma\sigma\sigma\rangle\langle\chi\chi\chi\rangle=0$). To
first order in $\delta$ the infrared fixed point has $g_*^2=\delta/b_0+O(\delta^2)$, so
$\De=E+\gamma_{\mathcal E}(g_*)=E+(a_0/b_0)\,\delta+O(\delta^2)$, which is $c_1=a_0/b_0$. Rescaling $g$ by a constant
rescales $a_0$ and $b_0$ by the same factor and leaves the ratio unchanged.

\paragraph{Why the spin channel supplies the denominator.} In the crossover description the long-range spin dimension
is protected~\cite{Behan:2017emf}: $\Ds=(3-s)/2=S+\delta$. If $\gamma_\sigma(g)=a_\sigma g^2+O(g^4)$, the fixed point
gives $\Ds=S+a_\sigma\delta/b_0+O(\delta^2)$, and matching to $S+\delta$ gives $a_\sigma=b_0$. Both coefficients are
therefore first-order anomalous-dimension integrals of short-range four-point functions: the energy external state
supplies $a_0$ and the spin external state supplies $b_0$.

\subsection{The integrated four-point functional}\label{app:cone:functional}
\paragraph{Normalization and scale extraction.} With $e$ a unit vector, recall from Sec.~\ref{sec:d3} the
four-point function with one external state at infinity, $H_O$, and write the integral $I_O$ as
\begin{equation}
\begin{gathered}
 H_O(y)=\lim_{R\to\infty}R^{2\Delta_O}\langle O(0)\sigma(y)\sigma(e)O(Re)\rangle,\qquad O\in\{\sigma,\epsilon\},\\
 F_O(y)=|e-y|^{-2(3-S)}\bigl[H_O(y)-|e-y|^{-2S}\bigr],\qquad I_O=\AC\!\int_{\mathbb R^3}\dd^3y\,F_O(y).
\end{gathered}
 \label{eq:cone:FO}
\end{equation}
The vacuum term removes the identity contraction of the two inserted spins. The symbol $\AC$ denotes analytic
continuation through the power-divergent short-distance terms, with the relevant deformation tuned critically; it is not
a prescription to integrate an unregulated divergent expression numerically. The normalization
$a_0=-2\pi I_\epsilon$, $b_0=-2\pi I_\sigma$ follows from a second-order amplitude with a Gaussian source. Let
\begin{equation}
 D_O(x_1,x_2)=\langle O(\infty)\sigma(x_1)\sigma(x_2)O(0)\rangle-|x_1-x_2|^{-2S},
\end{equation}
where the state at infinity has unit overlap with the unperturbed state at the origin, and define
\begin{equation}
 A_{O,2}(\delta)=\frac12\AC\!\int\dd^3x_1\dd^3x_2\,e^{-|x_1|^2-|x_2|^2}\,|x_1-x_2|^{-2(3-S-\delta)}D_O(x_1,x_2).
 \label{eq:cone:A2}
\end{equation}
Set $x_2=te$, $x_1=ty$ and integrate over the direction of $e$. The measure is $4\pi t^5\dd t\,\dd^3y$, the correlator
scales as $D_O(ty,te)=t^{-2S}D_O(y,e)$, and the kernel supplies $t^{-6+2S+2\delta}$. With the factor $\frac12$ the
remaining scale integral is
\begin{equation}
 2\pi\int_0^\infty\dd t\,t^{2\delta-1}e^{-t^2(1+|y|^2)}=\pi\Gamma(\delta)(1+|y|^2)^{-\delta},
\end{equation}
so that $A_{O,2}(\delta)=\pi\Gamma(\delta)\,I_O(\delta;1)$ with the source-dependent integral
\begin{equation}
 I_O(\delta;1)=\AC\!\int\dd^3y\,(1+|y|^2)^{-\delta}|e-y|^{-2(3-S-\delta)}\bigl[H_O(y)-|e-y|^{-2S}\bigr],
 \qquad I_O(0;1)=I_O .
 \label{eq:cone:Ireg}
\end{equation}
The second argument records the equal Gaussian weights. This amplitude probes the local insertion at the origin
against a fixed external state; the source suppresses the second short-distance region, at infinity, so its pole is
cancelled by a \emph{single} local operator-renormalization factor, not by twice that factor. With $O_R=Z_OO_{\rm bare}$
and
\begin{equation}
 Z_O=1+\frac{a_O}{2\delta}g^2+O(g^4),\qquad \gamma_O=-\beta(g)\,\partial_g\log Z_O=a_Og^2+O(g^4),
\end{equation}
finiteness of $Z_O[1+g^2A_{O,2}+\cdots]$ requires
\begin{equation}
 A_{O,2}(\delta)=-\frac{a_O}{2\delta}+O(1),\qquad a_\sigma=b_0,\quad a_\epsilon=a_0 .
\end{equation}
Since $\pi\Gamma(\delta)I_O(\delta;1)=\pi I_O/\delta+O(1)$, this gives $a_O=-2\pi I_O$, the normalization of
Sec.~\ref{sec:d3}. The common factor $-2\pi$ cancels from $c_1$, but keeping it makes the comparison with the
published normalization possible. The finite, source-dependent term is not needed at first order. All these identities
use the same critical analytic continuation; the Gaussian source does not replace the subtraction of the
short-distance power terms.

\paragraph{The folded integration region.} Rotational invariance reduces the vector integral to a radius $r=|y|$ and a
cosine $\eta=\widehat y\cdot e$. Write
\begin{equation}
 z=r\bigl(\eta+\mathrm i\sqrt{1-\eta^2}\bigr),\qquad v=|1-z|^2=1-2r\eta+r^2,\qquad \dd\mu=2\pi r^2\dd r\,\dd\eta .
\end{equation}
Here $z$ only encodes the rotational invariants of a three-dimensional configuration; the integral is not a
two-dimensional CFT integral. The conformal maps and the channel identity of Appendix~\ref{app:cone:folding} fold the full
integration into
\begin{equation}
 \mathcal R=\{|z|<1,\ |z|<|1-z|\}=\{|z|<1,\ \operatorname{Re}z<\tfrac12\},
 \label{eq:cone:region}
\end{equation}
that is, $0<r<1$ and $-1<\eta<\min\{1,(2r)^{-1}\}$. The upper $\eta$ limit changes at $r=\frac12$, which is a breakpoint
of the numerical integration.

\paragraph{Block sums.} Let $g_{\Delta,\ell}^{d_a,d_b}(z,\bar z)$ be the conformal block with leading OPE behaviour
$r^\Delta P_\ell(\eta)$ as $z\to0$, where $P_\ell$ is the Legendre polynomial. For the ordered four-point function the
external-dimension differences are $d_a=\Delta_1-\Delta_2$ and $d_b=\Delta_3-\Delta_4$; the mixed ordering
$(\epsilon,\sigma,\sigma,\epsilon)$ has $(d_a,d_b)=(E-S,S-E)$. In the $f$ normalization of the published spectrum
(Appendix~\ref{app:cone:inputs}), with the rising Pochhammer symbol $(a)_\ell=\Gamma(a+\ell)/\Gamma(a)$, define
\begin{equation}
\begin{aligned}
 w_\ell&=\frac{\ell!}{(1/2)_\ell},\qquad
 \mathsf S(z)=\sum_{A\ne\mathbf 1,\ \mathbb Z_2\text{-even}}w_{\ell_A}f_{\sigma\sigma A}^2\,g_{\Delta_A,\ell_A}^{0,0}(z),\\
 \mathsf E(z)&=\sum_{A\ne\mathbf 1,\ \mathbb Z_2\text{-even}}w_{\ell_A}f_{\sigma\sigma A}f_{\epsilon\epsilon A}\,
 g_{\Delta_A,\ell_A}^{0,0}(z),\qquad
 \mathsf M(z)=\sum_{B,\ \mathbb Z_2\text{-odd}}w_{\ell_B}f_{\sigma\epsilon B}^2\,g_{\Delta_B,\ell_B}^{E-S,S-E}(z).
\end{aligned}
 \label{eq:cone:sums}
\end{equation}
For example, the direct (odd-channel) representation of the mixed correlator is $H_\epsilon(z)=r^{-S-E}\mathsf M(z)$,
and its crossed (even-channel) representation is $|1-z|^{-2S}[1+\mathsf E(1-z)]$. The two coincide for an exact
crossing-symmetric correlator, not necessarily for independently truncated sums.

\subsection{Derivation of the folding and its assumptions}\label{app:cone:folding}
Three steps take the integral $I_O$ of Eq.~\eqref{eq:cone:FO} over $\mathbb R^3$ to the region $\mathcal R$ and account
for the coefficients and the extra term of the folded expression.

\paragraph{Inversion about the origin.} For $y=x/|x|^2$ the Jacobian is $|x|^{-6}$, and
\begin{equation}
 H_O\Bigl(\frac{x}{|x|^2}\Bigr)=|x|^{2S}H_O(x),\qquad \Bigl|e-\frac{x}{|x|^2}\Bigr|=\frac{|e-x|}{|x|}.
\end{equation}
The first identity uses the identical operators $O$ at zero and infinity, which the inversion exchanges. Since
$2(3-S)+2S=6$, the kernel and correlator factors cancel the Jacobian exactly, and
\begin{equation}
 I_O=2\AC\!\int_{|y|<1}\dd^3y\,F_O(y).
 \label{eq:cone:foldorigin}
\end{equation}
The cancellation uses $\Delta_\chi=3-S$ at $\delta=0$; it is not an identity between integrands at finite $\delta$.

\paragraph{The crossed representation on the rest of the ball.} Split $|y|<1$ into $\mathcal R$ and
$\mathcal C=\{|e-y|<|y|<1\}$, and set $w=e-y$ on $\mathcal C$. The exact correlator admits
\begin{equation}
 H_O(e-w)=|w|^{-2S}\bigl[1+X_O(w)\bigr],\qquad X_\sigma=\mathsf S,\quad X_\epsilon=\mathsf E .
\end{equation}
For $O=\epsilon$ this equality selects the even $\sigma\sigma$--$\epsilon\epsilon$ OPE channel; it is \emph{not} a
permutation symmetry that interchanges $\sigma$ and $\epsilon$. The vacuum term cancels, $F_O(e-w)=|w|^{-6}X_O(w)$, and
the new domain is $\mathcal D=\{|w|<|e-w|<1\}\subset\mathcal R$. Hence
\begin{equation}
 I_O=2\AC\!\int_{\mathcal R}\dd^3y\,F_O(y)+2\AC\!\int_{\mathcal D}\dd^3w\,|w|^{-6}X_O(w).
 \label{eq:cone:foldchannel}
\end{equation}

\paragraph{Inversion about the second spin insertion.} Let $v=|e-w|^2$ and
\begin{equation}
 T(w)=e+\frac{w-e}{|w-e|^2},\qquad |T(w)|^2=\frac{|w|^2}{v},\qquad |e-T(w)|^2=\frac1v,\qquad
 |\det\partial T|=v^{-3}.
\end{equation}
The map $T$ sends $\mathcal R$ to itself (it exchanges the two conditions $|w|<1$ and $|w|<|e-w|$) and exchanges its
parts with $v<1$ and $v>1$, namely $\mathcal D$ and its complement, up to boundaries of zero measure. In the complex
representative it is $w\mapsto w/(w-1)$ up to conjugation. Both $X_O$ contain pairs of equal external dimensions and
only even-spin exchanges, so the pair-exchange factor $(-1)^\ell$ is one and $X_O(Tw)=X_O(w)$. Moreover
$|T(w)|^{-6}|\det\partial T|=|w|^{-6}$. The two parts of the last integral of Eq.~\eqref{eq:cone:foldchannel} are
therefore equal, and
\begin{equation}
 \boxed{\,I_O=\AC\!\int_{\mathcal R}\dd^3y\,\bigl[2F_O(y)+|y|^{-6}X_O(y)\bigr].\,}
 \label{eq:cone:fold}
\end{equation}
The changes of variables are applied with consistent subtractions and analytic continuation, not to bare divergent
integrals. They establish Eq.~\eqref{eq:cone:fold} for the exact correlator. Replacing its channel sums by
independently truncated data introduces a correlator approximation that no change of variables can remove.

\paragraph{Folded integrands and per-family kernels.} With $H_\sigma=r^{-2S}(1+\mathsf S)$ in the direct channel and
$H_\epsilon=r^{-S-E}\mathsf M$, the implemented linear functionals are
\begin{align}
 I_\sigma&=\AC\!\int_{\mathcal R}\dd\mu\,\Bigl\{2v^{-(3-S)}\bigl[r^{-2S}(1+\mathsf S)-v^{-S}\bigr]+r^{-6}\mathsf S\Bigr\},
 \label{eq:cone:Is}\\
 I_\epsilon&=\AC\!\int_{\mathcal R}\dd\mu\,\Bigl\{2v^{-(3-S)}\bigl[r^{-S-E}\mathsf M-v^{-S}\bigr]+r^{-6}\mathsf E\Bigr\}.
 \label{eq:cone:Ie}
\end{align}
Both the identity subtraction and the signs of the mixed products must be kept; in particular the energy integrand is
not a sum of positive even-channel coefficients. The per-family kernels are
\begin{equation}
\begin{aligned}
 K_b(\Delta,\ell)&=-2\pi w_\ell\AC\!\int_{\mathcal R}\dd\mu\,\bigl[2v^{-(3-S)}r^{-2S}+r^{-6}\bigr]g_{\Delta,\ell}^{0,0},\qquad
 K_e(\Delta,\ell)=-2\pi w_\ell\AC\!\int_{\mathcal R}\dd\mu\,r^{-6}g_{\Delta,\ell}^{0,0},\\
 K_o(\Delta,\ell)&=-4\pi w_\ell\int_{\mathcal R}\dd\mu\,v^{-(3-S)}r^{-S-E}g_{\Delta,\ell}^{E-S,S-E},
\end{aligned}
 \label{eq:cone:kernels}
\end{equation}
so that, with $f_A=f_{\sigma\sigma A}$, $h_A=f_{\epsilon\epsilon A}$ and $k_B=f_{\sigma\epsilon B}$,
\begin{equation}
 b_0=b_{\mathbf 1}+\sum_AK_b(\Delta_A,\ell_A)f_A^2,\qquad
 a_0=a_{\mathbf 1}+\sum_AK_e(\Delta_A,\ell_A)f_Ah_A+\sum_BK_o(\Delta_B,\ell_B)k_B^2 .
 \label{eq:cone:decomp}
\end{equation}
The nonspectral constants come from the remaining terms of Eqs.~\eqref{eq:cone:Is} and~\eqref{eq:cone:Ie},
\begin{equation}
 b_{\mathbf 1}=-4\pi\!\int_{\mathcal R}\dd\mu\,v^{-(3-S)}\bigl[r^{-2S}-v^{-S}\bigr]\approx-26.0638858592,\qquad
 a_{\mathbf 1}=4\pi\!\int_{\mathcal R}\dd\mu\,v^{-3}\approx44.4132198049
\end{equation}
(both integrals converge: $v>\frac14$ on $\mathcal R$, and $r^{-2S}$ is integrable at the origin because $2S<3$). They are bookkeeping terms, not separate physical
observables.

\subsection{Conformal blocks and numerical integration}\label{app:cone:blocks}
\paragraph{Radial expansion.} With the conformal radial coordinate
\begin{equation}
 \rho(z)=\frac{z}{(1+\sqrt{1-z})^2},\qquad \varrho=|\rho|,\qquad \xi=\frac{\operatorname{Re}\rho}{|\rho|},
\end{equation}
where the branch of the square root is continued from $z=0$, the blocks are evaluated as
\begin{equation}
 g_{\Delta,\ell}^{d_a,d_b}(z,\bar z)=(4\varrho)^\Delta\sum_{n=0}^{N}\varrho^n\sum_jB_{n,j}P_j(\xi),\qquad
 B_{0,j}=\delta_{j\ell}.
 \label{eq:cone:rhoblock}
\end{equation}
The coefficients solve the quadratic Casimir equation; the matrix recurrence evaluated is given in
Appendix~\ref{app:cone:recurrences}. The organization in the radial coordinate follows Ref.~\cite{Hogervorst:2013}; it
converges fast in favourable OPE configurations. On the closure of $\mathcal R$ the largest radial coordinate is
$\max|\rho|=2-\sqrt3\simeq0.267949$, reached at the corner $z=e^{\mathrm i\pi/3}$. This small maximum is why the folded domain
is used: equivalent domains with a larger maximum radial coordinate magnify omitted-family and descendant effects, so
their truncated values should not be averaged as independent measurements.

A primary is the lowest state of a conformal family, and the terms with $n>0$ account for its descendants. Increasing
$N$ completes a retained family but does not add a primary. Descendant convergence and completeness of the primary
spectrum are therefore different error questions.

\paragraph{Analytic continuation at small radius.} The $r^{-6}$ terms of Eqs.~\eqref{eq:cone:Is}
and~\eqref{eq:cone:Ie} are power divergent at the origin. Choose a subtraction radius $r_0<\frac12$, normally
$r_0=0.2$, so that the small ball has the complete angular range. On it a separate small-$z$ series is used,
\begin{equation}
 g_{\Delta,\ell}^{0,0}(z)=\sum_{n,j}A_{n,j}r^{\Delta+n}P_j(\eta)
\end{equation}
(Appendix~\ref{app:cone:recurrences}). The angular integral keeps only $j=0$, and gives the explicit small-ball
contribution
\begin{equation}
 \AC\!\int_{r<r_0}\dd\mu\,r^{-6}g_{\Delta,\ell}^{0,0}=4\pi\sum_nA_{n,0}\,\frac{r_0^{\Delta+n-3}}{\Delta+n-3}.
 \label{eq:cone:AC}
\end{equation}
The sum uses 32 levels of the $z$-series. Terms whose scalar coefficient vanishes are not evaluated as artificial
$0/0$ expressions. A genuine resonance with a nonzero coefficient would need a separate logarithmic treatment; the
implementation rejects an active scalar-series term with $|\Delta+n-3|<10^{-10}$. This is a diagnostic threshold, not a
way to assign a finite part. It is inactive for the retained inputs, where the nearest active power is about $0.41$
from resonance. At the conserved stress tensor the null-state consistency checks enforce the required cancellation.

For the rest of the integral, Gauss--Legendre quadrature is used in $\eta$, and $r$ is split into $(0,r_0)$,
$(r_0,\frac12)$ and $(\frac12,1)$, with the map $r=r_0u^4$ on the first interval to resolve the integrable endpoint and
linear maps on the others. The principal run uses $N=24$ and $100\times100$ nodes in each radial and angular piece. The
$r^{-6}$ contribution is left out of the small-ball quadrature, because Eq.~\eqref{eq:cone:AC} already includes it. The
critical analytic continuation, rather than a cutoff-dependent finite part with unspecified thermal tuning, is part of
the definition of the observable. The relevant test is therefore agreement between equivalent analytically continued
representations, not agreement between a cutoff integral and the intended critical quantity.

\paragraph{Numerical checks and their scope.} The radial block is compared with the independent $z$-series at complex
points and, for scalar exchange with equal external dimensions, with the diagonal hypergeometric expression
\begin{equation}
 g_{\Delta,0}(z,z)=\Bigl(\frac{z^2}{1-z}\Bigr)^{\Delta/2}
 {}_3F_2\!\left(\begin{matrix}\Delta/2,\ \Delta/2,\ \Delta/2-1/2\\[1pt](\Delta+1)/2,\ \Delta-1/2\end{matrix};
 \frac{z^2}{4(z-1)}\right).
 \label{eq:cone:3F2}
\end{equation}
The largest discrepancies are $\sci{5.6}{-17}$ against the $z$-series and $\sci{2.3}{-16}$ against
Eq.~\eqref{eq:cone:3F2}, floating-point effects at fixed input data. The
package's two integrated forms (Sec.~\ref{sec:retained}) agree on the retained coefficients; this checks the algebra
and the bookkeeping on the \emph{same} folded domain, including the spin-channel identity that reduces one form to the
other, and is not an independent validation of the folding.

\subsection{Input data and tensor conventions}\label{app:cone:inputs}
\paragraph{The selected 60-primary snapshot.} The retained selection transcribes the stable extracted operators of
dimension below 24 in Tables~3 to~7 of Ref.~\cite{Simmons-Duffin:2016wlq}: 37 internally $\mathbb Z_2$-even and 23
internally $\mathbb Z_2$-odd primaries (Table~\ref{tab:cone:spectrum}). It is a selection of published stable entries,
not a statement that every exchanged primary below 24 is present. The leading data are replaced coherently by the
stress-tensor bootstrap inputs of Ref.~\cite{Chang:2024whx}:
\begin{center}\small
\begin{tabular}{@{}lllll@{}}
\toprule
$S$ & $E$ & $\lambda_{\sigma\sigma\epsilon}$ & $\lambda_{\epsilon\epsilon\epsilon}$ & $c=c_T/c_B$\\
\midrule
$0.518148806(24)$ & $1.41262528(29)$ & $1.05185373(11)$ & $1.53244304(58)$ & $0.946538675(42)$\\
\bottomrule
\end{tabular}
\end{center}
Here $c_B$ is the stress-tensor two-point normalization of one free real scalar, and $c=c_T/c_B$ is the dimensionless
normalization used in the Ward identities. Shared inputs stay shared under variations: the same
$\lambda_{\sigma\sigma\epsilon}$ enters the even $\epsilon$ channel and the odd $\sigma$ channel. The stress-tensor
couplings are derived from $S$, $E$ and $c$, not treated as independent measured numbers:
\begin{equation}
 f_{\sigma\sigma T}=S\sqrt{\frac{3}{8c}},\qquad f_{\epsilon\epsilon T}=E\sqrt{\frac{3}{8c}},
 \label{eq:cone:ward}
\end{equation}
so the Ward ratio is $f_{\epsilon\epsilon T}/f_{\sigma\sigma T}=E/S$. A simultaneous sign change of both stress
couplings is an operator convention and does not change the products that enter.

\paragraph{What the printed errors mean.} In Table~1 of Ref.~\cite{Chang:2024whx}, the rigorous entries among these
five inputs are $S$ and $E$, under that bootstrap setup; the printed errors of the two OPE coefficients and of $c$ are
not established there as rigorous bounds. The errors of the higher families in Ref.~\cite{Simmons-Duffin:2016wlq}
describe spreads of extracted data over 60 spectra at the fixed derivative cutoff $\Lambda=43$; they contain no
certified $\Lambda\to\infty$ uncertainty. For example, Ref.~\cite{Henriksson:2022} reports an extraction near $7.297$
at $\Lambda=30$ for the even scalar listed at $7.2535(51)$ in the retained table. This does not determine the true
dimension, but it shows why the printed spreads cannot be read as complete errors. The strong sensitivity of $c_1$ to
the $\epsilon'$ family and to the odd spin-three coupling near $\Delta=4.638$ makes this limitation relevant to the
final interval; the package supplies no calibrated independent cross-check of the quoted error of that coupling.

\paragraph{Normalization conversion and signs.} The $\lambda$ convention meant here is the tensor convention of
Kos, Poland and Simmons-Duffin~\cite{Kos:2014bka} and of Reehorst~\cite{Reehorst:2021}, related to the table convention by $\lambda=2^{\ell/2}f$; the symbol is
not a convention-independent OPE normalization. The block contraction factor in the $\lambda$ convention is
$\ell!/[2^\ell(1/2)_\ell]$; multiplied by the conversion $2^\ell$ of the two couplings, it gives $w_\ell=\ell!/(1/2)_\ell$
in Eq.~\eqref{eq:cone:sums}. For scalar exchange the conventions agree.

The signs of $f_Ah_A$ are physical relative signs. Nine rows, all in the selected $[\epsilon\epsilon]_0$ trajectory
(family \texttt{ee0}), have a negative product; replacing all couplings by absolute values would change the mixed
correlator and the answer. The overall sign of one exchanged primary cancels from its product and is immaterial.

\paragraph{Odd spins in the reflected mixed ordering.} The lightcone-normalized blocks of
Ref.~\cite{Simmons-Duffin:2016wlq} behave as $g^{\rm SD}_{\Delta,\ell}\sim(-1)^\ell w_\ell(4\varrho)^\Delta P_\ell(\xi)$,
whereas the blocks of Eq.~\eqref{eq:cone:rhoblock} have leading coefficient $+1$. In the ordering
$(\epsilon,\sigma,\sigma,\epsilon)$ the two OPE pairs are reflected, and reversing the ordered scalar pair contributes a
second factor $(-1)^\ell$ in the OPE product, which cancels the sign of the lightcone block. The coefficient of the
mixed block is therefore $+w_\ell f_{\sigma\epsilon B}^2$, as in $\mathsf M$, with no further $(-1)^\ell$. For
$0<z<1$ the reflected radial-state decomposition must have positive leading squared amplitudes; its all-descendant
positivity is formulated with the prefactor $Q$ of Eq.~\eqref{eq:cone:Q}, not asserted for every coefficient of the
raw mixed block. The generalized-free prefactor and positive-state tests of Appendix~\ref{app:cone:validation} check this
convention. As an independent calibration of the stress family, a free real scalar has $S=\frac12$, $c=1$, and
Eq.~\eqref{eq:cone:ward} gives $w_2f_{\sigma\sigma T}^2=\frac83\cdot\frac3{32}=\frac14$.

\subsection{Retained-spectrum evaluation}\label{app:cone:retained}
The package stores the per-family kernels and contributions, not only the final ratio, which exposes normalization
changes and large cancellations. For example, the even $\epsilon$ family contributes about $108.0654$ to $a_0$ and the
odd $\sigma$ family about $-147.3028$, before the other families and the nonspectral constant $a_{\mathbf 1}$ are
added; the comparatively small $a_0$ of Sec.~\ref{sec:retained} should not be reconstructed from rounded individual
terms.

\paragraph{Numerical convergence at fixed inputs.} Raising the descendant order from 24 to 28 and the quadrature from
$100^2$ to $128^2$ nodes changes the ratio by about $\sci{1.7}{-12}$; moving $r_0$ from $0.2$ to $0.15$ changes it by
about $\sci{3}{-14}$. Quadrature refinement is therefore not the route to a substantially smaller physical error.

\paragraph{Comparison with the published normalization.} The component-by-component agreement with BRRZ in
Sec.~\ref{sec:retained} uses the same normalization $a_0=-2\pi I_\epsilon$, $b_0=-2\pi I_\sigma$. These comparisons
are compatible at the published accuracy; they are not independent measurements to be combined into a smaller error.
The precision of the retained integral and the uncertainty of the full CFT remain separate.

\subsection{Retained-input variations}\label{app:cone:variations}
In this and the following subsections $a_0$, $b_0$ denote the retained reference coefficients of
Sec.~\ref{sec:retained}, $c_1^{(0)}=a_0/b_0$, and $\Delta a$, $\Delta b$ denote shifts away from them.

\paragraph{Linearized integrated objective.} The first variation is
\begin{equation}
 \dd c_1=\frac{\dd a_0-c_1^{(0)}\dd b_0}{b_0}.
 \label{eq:cone:variation}
\end{equation}
By Eq.~\eqref{eq:cone:decomp} the coupling derivatives are elementary: for an even primary the variations with $f_A$
and $h_A$ are proportional to $K_eh_A-2c_1^{(0)}K_bf_A$ and $K_ef_A$, and for an odd primary the derivative is
$2K_ok_B$. Dimension derivatives are finite differences of the family kernels. The five shared low inputs $S$, $E$,
$\lambda_{\sigma\sigma\epsilon}$, $\lambda_{\epsilon\epsilon\epsilon}$ and $c$ are varied coherently through
Eq.~\eqref{eq:cone:ward}: varying an external dimension changes the powers in the integrands and all affected blocks,
as well as the associated exchanged operator, so these inputs are varied globally and not counted again as independent
entries of the $\sigma$, $\epsilon$ or $T$ rows.

For every other input $u_k$ with quoted error $e_k$, the scaled displacement is $x_k=(u_k-u_k^{(0)})/(2e_k)$. The range
$|x_k|\le1$ is a twice-quoted-error hypothesis, not a two-standard-deviation statement. The coupling derivatives are
analytic. For an exchanged dimension, the coefficient of the integrated objective is the centred kernel difference at
$\Delta\pm e_k$ multiplied by the full scale $2e_k$. For a shared low input it is the centred difference of the full
ratio at $u_k^{(0)}\pm2e_k$. Nonlinear corrections to this affine objective are allowed separately
(Appendix~\ref{app:cone:optimization}).

\paragraph{Retained-input bound for the denominator.} The conversion of a numerator range into a ratio interval also
needs a bound on the denominator. Every retained independent input is placed within twice its quoted error, keeping
shared quantities coherent. For each nonleading family the dimension is sampled at $\Delta+ne_\Delta$, $n=-2,\dots,2$,
together with the coupling corners $f\pm2e_f$ and, for an even family, $h\pm2e_h$; if the interval of $f$ contained
zero, $f=0$ would be added to the candidates, because $b_0$ depends on $f^2$ (no retained row triggers this). The
smallest and largest sampled shifts of each family's contribution to $b_0$ are kept and summed over families; the
shared low inputs are evaluated at their coherent $\pm2e_k$ endpoints and added. This calculation uses descendant order
20 and $80\times80$ nodes. The sampled envelope is
\begin{equation}
 -0.00142738811843\lesssim\Delta b_{\rm in}\lesssim0.00142355554030,
 \label{eq:cone:denbox}
\end{equation}
rounded outward to $|\Delta b_{\rm in}|\le0.001428$. It does not impose crossing or the $h_{\epsilon'}$ restriction of
Appendix~\ref{app:cone:tails}, so it remains applicable after those constraints shrink the input region. It is a sampled
numerical allowance, not an interval-arithmetic extremum over every continuous parameter value, and its adequacy is one
of the assumptions of the estimated range.

\subsection{Conditional control of the omitted spectrum}\label{app:cone:tails}
\paragraph{The hypotheses and their limitations.} The main scenario assumes the gaps
\begin{equation}
 G_{\rm even}=8,\qquad G_{\rm odd}=7
 \label{eq:cone:gaps}
\end{equation}
of Sec.~\ref{sec:range}: every unlisted primary exchanged in the relevant scalar OPEs has dimension at least the
corresponding $G$. This allows arbitrary omitted spectral weight above the thresholds, subject to crossing and
positivity; it does not assert that the retained list is complete up to its largest dimension, 24. Neither threshold is
proved, and additional unlisted low-dimensional operators remain a possible cause of failure of the narrow interval.
The input ranges are hypotheses as well: the actual retained CFT data are assumed to lie inside the
twice-quoted-error box. Exact crossing and reflection positivity are imposed on the full theory, not on the finite
retained approximation alone.

\paragraph{Evidence motivating the thresholds, not proving them.} The gap hypotheses refer to spatial-parity-even
primaries in the indicated internal $\mathbb Z_2$ sectors. Spatial and internal parity are different labels: a
spatial-parity-odd state does not contribute to these two-scalar OPEs merely because its internal parity is
appropriate. The selected list has 11 even primaries below dimension 8. Published spectral extractions and fuzzy-sphere
studies~\cite{Simmons-Duffin:2016wlq,Henriksson:2022,Zhu:2022,Fardelli:2026} give the following nearby unlisted
families; the entries are estimates, not gap bounds.
\begin{center}\small
\begin{tabular}{@{}llp{.2\textwidth}p{.47\textwidth}@{}}
\toprule
internal parity & spin & estimated dimension & evidence and status\\
\midrule
odd & 3 & $7.53(6)$ & fuzzy sphere~\cite{Fardelli:2026}; an additional trajectory is also discussed in extremal spectra\\
odd & 6 & $7.73(5)$; about $7.3$ & Ref.~\cite{Fardelli:2026}; lower extremal-functional estimate in Ref.~\cite{Henriksson:2022}\\
even & 0 & about $8.55$ to $8.92$ & differing extremal extractions compared in Ref.~\cite{Henriksson:2022}; not a certified interval\\
even & 2 & about $8.6$ & higher trajectory in the spectral plots of Ref.~\cite{Henriksson:2022}; approximate reading\\
even & 6 & about $8.78$ & described as an unstable extraction in Ref.~\cite{Henriksson:2022}; tentative\\
\bottomrule
\end{tabular}
\end{center}
The fuzzy-sphere results of Ref.~\cite{Fardelli:2026} explain why an odd threshold of 8 cannot simply be kept: allowed
unlisted odd states are estimated below 8. A threshold of 7 leaves only a margin of about $0.3$ to $0.7$ to the nearby
spin-six estimates, and the associated OPE data do not justify assuming that missing states just above the threshold
have negligible weight. None of these approximate couplings is used as a precision input; the remainder bounds must
accommodate them. The comparisons certify neither gap. In particular, Figures~1 and~2 of
Ref.~\cite{Simmons-Duffin:2016wlq} contain unstable extracted entries, and the main scenario assumes that no such
unretained entry below its threshold is an additional genuine exchanged primary; calling an entry an artifact does not
establish that assumption. The evidence motivates testing $(8,7)$, and the feasibility and sensitivity tests
(Sec.~\ref{sec:howcond} and Appendix~\ref{app:cone:sensitivity}) expose its tension and its consequences.

\paragraph{The published constraint on $\lambda_{\epsilon\epsilon\epsilon'}$.} Let $\epsilon'$ be the next retained
internally even scalar after $\epsilon$, with reference dimension $3.82968$ and $h_{\epsilon'}=f_{\epsilon\epsilon\epsilon'}=1.5360$;
for a scalar the $f$ and $\lambda$ conventions coincide. Reehorst's conditional bootstrap analysis gives
$\lambda_{\epsilon\epsilon\epsilon'}=1.5362(12)$~\cite{Reehorst:2021}. Its assumptions include an isolated
nondegenerate $\epsilon'$, the next even scalar above dimension 6 and the next odd scalar above dimension 3, which are
compatible with the spectral assumptions used here. To allow for the rounding of the printed bound, the retained-input
box is intersected with
\begin{equation}
 1.5349\le h_{\epsilon'}\le1.5375 .
 \label{eq:cone:reehorst}
\end{equation}
This applies a published conditional constraint; it is not a new semidefinite-programming computation, nor a weighted
average of independent experiments. If its sparsity or nondegeneracy hypotheses are not assumed, the restriction must
be removed and the envelope recomputed.

\paragraph{The identical-channel positive remainder.} For an identical external scalar of dimension $D\in\{S,E\}$, write
its reduced correlator as $P+R$, where $P$ is the retained block sum including the identity and $R$ the missing sum. In
radial variables, reflection positivity gives nonnegative expansion coefficients on the collinear axis,
$R(r,1)=\sum_pc_pr^p$ with $c_p\ge0$; each descendant adds a nonnegative integer level to the primary dimension, so with
a missing-primary gap $G$ every $p\ge G$. For $0<r_1<r_2$ this gives
\begin{equation}
 0\le R(r_1,1)\le\Bigl(\frac{r_1}{r_2}\Bigr)^GR(r_2,1),
 \label{eq:cone:radsupp}
\end{equation}
and the same gap controls the full omitted family. Choose a real $0<x<\frac12$ and let $A=[x/(1-x)]^{2D}$,
$r_1=\rho(x)$, $r_2=\rho(1-x)$. Crossing gives $P(x)+R(x)=A[P(1-x)+R(1-x)]$, with $R(x)$ evaluated at radius $r_1$ and
$R(1-x)$ at $r_2$. Solving for the remainder at the larger radius and using Eq.~\eqref{eq:cone:radsupp},
\begin{equation}
 R(r_2,1)=A^{-1}R(r_1,1)+A^{-1}P(x)-P(1-x)\le A^{-1}\Bigl(\frac{r_1}{r_2}\Bigr)^GR(r_2,1)+A^{-1}P(x)-P(1-x),
\end{equation}
and therefore
\begin{equation}
 R(r_2,1)\le\frac{A^{-1}P(x)-P(1-x)}{1-A^{-1}(r_1/r_2)^G},
 \label{eq:cone:anchor}
\end{equation}
provided the denominator is positive. The numerator must be maximized over the allowed retained inputs and include the
nonlinear and numerical allowances. A negative upper bound is a consistency failure, not permission to replace a
sizeable negative number by zero. Different reference points give valid bounds at different radii: if $B_i$ bounds the
reference norm at $r_i$, then
\begin{equation}
 R(r,1)\le\min_{i:\,r\le r_i}B_i\Bigl(\frac{r}{r_i}\Bigr)^G .
 \label{eq:cone:anchormin}
\end{equation}
The minimum is justified because each reference bound is separately necessary; it is not a statistical combination.
The radial-expansion framework is that of Ref.~\cite{Hogervorst:2013}.

\paragraph{Even mixed and odd remainders.} The even mixed remainder is not positive, but the Cauchy--Schwarz inequality
bounds it by the two identical-channel norms,
\begin{equation}
 \bigl|R_{\sigma\sigma;\epsilon\epsilon}(r,\xi)\bigr|\le\sqrt{R_{\sigma^4}(r,1)\,R_{\epsilon^4}(r,1)} .
 \label{eq:cone:cauchy}
\end{equation}
This is why the analysis uses all three scalar correlators, including $\langle\epsilon\epsilon\epsilon\epsilon\rangle$,
although that correlator does not enter the slope functional directly. In the odd mixed channel positivity applies to
\begin{equation}
 Q(\rho,\bar\rho)\,R_{\rm odd},\qquad Q(\rho,\bar\rho)=\Bigl|\frac{1+\rho}{1-\rho}\Bigr|^{2(E-S)},
 \label{eq:cone:Q}
\end{equation}
not to the raw mixed-block coefficients. At real odd-channel reference points $x>\frac12$, mixed crossing and
Eq.~\eqref{eq:cone:cauchy} bound this normalized odd norm; radial suppression is applied to the normalized quantity, and
the bound at a complex integration point is divided by $Q$ again. Applying positivity directly to the raw series,
without this prefactor, would be invalid.

Let $U_\sigma(u)$ and $U_\epsilon(u)$ be the identical-channel envelopes of Eq.~\eqref{eq:cone:anchormin}, and
$d_m^{\min}(x)$ the smallest allowed retained mixed residual, with the direct-minus-crossed normalization of
Eq.~\eqref{eq:cone:residuals} and including its nonlinear allowance. At a real odd reference point $x>\frac12$ mixed
crossing reads
\begin{equation}
 0=d_m(x)+x^{-S-E}R_{\rm odd}(x)-(1-x)^{-2S}R_{\sigma\sigma;\epsilon\epsilon}(1-x).
\end{equation}
Before the common numerical padding (Appendix~\ref{app:cone:optimization}), the normalized odd reference norm at
$u_x=\rho(x)$ is therefore bounded by
\begin{equation}
 B_{\rm odd}(x)=x^{S+E}\,Q(u_x,u_x)\Bigl[-d_m^{\min}(x)+(1-x)^{-2S}\sqrt{U_\sigma(\rho(1-x))\,U_\epsilon(\rho(1-x))}\Bigr].
 \label{eq:cone:oddanchor}
\end{equation}
Each valid odd reference supplies a radial bound, and the envelope used at radius $u$ is
\begin{equation}
 U_o(u)=\min_{x:\,u\le\rho(x)}B_{\rm odd}(x)\Bigl(\frac{u}{\rho(x)}\Bigr)^{G_{\rm odd}} .
 \label{eq:cone:oddenv}
\end{equation}
The raw odd remainder at a complex point is bounded by $U_o(|\rho|)/Q(\rho,\bar\rho)$. As in the identical channel, a
materially negative upper reference norm is a consistency failure. The identical-channel inequality and the mixed
normalization were tested on exact generalized-free correlators; these tests check the algebra of the envelope and its
implementation, not the completeness assumptions about the Ising spectrum.

\subsection{Joint crossing and input optimization}\label{app:cone:optimization}
\paragraph{The crossing residuals.} The residuals are unnormalized differences of the direct and crossed correlators,
not divided by an average correlator magnitude. With the further sum
$\mathsf V(z)=\sum_{A\ne\mathbf 1,\ \mathbb Z_2\text{-even}}w_{\ell_A}h_A^2g_{\Delta_A,\ell_A}^{0,0}(z)$, they are, for
each retained input configuration,
\begin{equation}
\begin{aligned}
 d_\sigma(z)&=|z|^{-2S}\bigl[1+\mathsf S(z)\bigr]-|1-z|^{-2S}\bigl[1+\mathsf S(1-z)\bigr],\\
 d_\epsilon(z)&=|z|^{-2E}\bigl[1+\mathsf V(z)\bigr]-|1-z|^{-2E}\bigl[1+\mathsf V(1-z)\bigr],\\
 d_m(z)&=|z|^{-S-E}\mathsf M(z)-|1-z|^{-2S}\bigl[1+\mathsf E(1-z)\bigr].
\end{aligned}
 \label{eq:cone:residuals}
\end{equation}
The identity contributions are explicit, and the external dimensions and all occurrences of shared couplings vary
coherently. These definitions fix the signs and normalization of the reference bounds, the Jacobian and the affine
constraints below.

\paragraph{The affine necessary-constraint problem.} Crossing constrains which retained-input shifts and omitted
contributions can coexist. The retained inputs enter through the scaled variables $x_k$ of
Appendix~\ref{app:cone:variations}; except for the intersected $h_{\epsilon'}$ range, $-1\le x_k\le1$ is a displacement of
at most twice a quoted error. Shared low inputs are parametrized once, and there are 154 variables. For residual type
$q\in\{\sigma^4,\epsilon^4,\text{mixed}\}$ at point $z_j$,
\begin{equation}
 d_{qj}(x)=d^{(0)}_{qj}+\sum_kJ_{qjk}x_k+e_{qj}(x),\qquad |e_{qj}(x)|\le p_{qj}.
 \label{eq:cone:affine}
\end{equation}
The columns of $J$ are centred secants of the retained residual at $x_k=\pm1$, with the other inputs at their
reference values. The bound on $e_{qj}$ is a sampled numerical allowance, not a proved uniform bound. For each
nonleading family the affine remainder is evaluated at scaled dimension displacements $-1,-\frac12,0,\frac12,1$ and
scaled coupling displacements $-1,0,1$; the largest absolute remainder at each crossing point is multiplied by $1.5$ and
summed over families. For each shared low input, twice the absolute difference between the mean of its two endpoint
residuals and the reference residual is added. Finally the floor $\sci{2}{-10}(1+|d^{(0)}|)$ is added. The factors
$1.5$ and $2$ and the floor are chosen allowances, not consequences of a proved uniform Hessian bound. (The ``margins''
of Sec.~\ref{sec:range} are these paddings, the objective allowance $\sci{5}{-5}$ and the sampled denominator, reference
and quadrature allowances.) Adding separate
per-family and low-input samples also does not enclose cross terms between a shared low input and a family input
uniformly; the nonlinear tests of Appendix~\ref{app:cone:validation} sample such combined shifts but do not close this
logical gap.

The array has 92 point entries: 35 equally spaced real identical-channel references from $0.15$ to $0.49$, 15 equally
spaced odd references from $0.64$ to $0.85$, and the grid $z=0.2+0.1i_1+0.12i_2\,\mathrm i$ with $i_1=0,\dots,6$ and
$i_2=0,\dots,5$. Some configurations serve more than one role, so 92 counts entries, not independent constraints.
Blocks are evaluated at descendant level 40 for the principal optimization.

Exact crossing requires a nonzero omitted contribution $t_{qj}$ to cancel the retained residual. With a conditional
bound $t_{qj}\in[-B^-_{qj},B^+_{qj}]$, Eq.~\eqref{eq:cone:affine} becomes two affine inequalities once the nonlinear
padding is included. Real-axis positivity gives one-sided improvements; complex-point bounds use magnitudes. The
truncated correlator is never forced to satisfy crossing with its tail set to zero.

The allowed residual ranges follow from the reference envelopes. Put $u=|\rho(z)|$, $u'=|\rho(1-z)|$ and
$Q_z=Q(\rho(z),\overline{\rho(z)})$. For an identical external scalar of dimension $D\in\{S,E\}$, with $U_S=U_\sigma$
and $U_E=U_\epsilon$, define
\begin{equation}
 T_D^d=|z|^{-2D}U_D(u),\qquad T_D^c=|1-z|^{-2D}U_D(u').
\end{equation}
In general the retained residual lies in $[L_D,U^{\rm res}_D]=[-(T_D^d+T_D^c),\,T_D^d+T_D^c]$. On the real axis a one-sided improvement
needs a sign condition: for $x=\min\{z,1-z\}<\frac12$ put
\begin{equation}
 q_D(x;G)=\Bigl[\frac{\rho(x)}{\rho(1-x)}\Bigr]^G\Bigl(\frac{1-x}{x}\Bigr)^{2D}\le1 .
 \label{eq:cone:realsign}
\end{equation}
Then the direct positive tail on the side of the smaller radius is at most $q_D$ times the crossed one, and the exact
retained residual there, $T^c-T^d$ in terms of the actual tails, is nonnegative. At the nominal points the largest
values are about $0.6629$ for $D=S$ and $0.7121$ for $D=E$, excluding the symmetric point $z=\frac12$, where the residual
vanishes identically. The implementation checks Eq.~\eqref{eq:cone:realsign} and stops if it fails, so a different gap
or point selection cannot silently keep the sharpening. Under this condition the ranges are
\begin{equation}
 \bigl[L_D,U^{\rm res}_D\bigr]=\begin{cases}[0,\,T_D^c],&0<z<\frac12,\\[2pt] [-T_D^d,\,0],&\frac12<z<1.\end{cases}
\end{equation}
For the mixed residual put
\begin{equation}
 T_m^o=|z|^{-S-E}\frac{U_o(u)}{Q_z},\qquad T_m^e=|1-z|^{-2S}\sqrt{U_\sigma(u')\,U_\epsilon(u')} .
\end{equation}
The general range is $[-T_m^o-T_m^e,\,T_m^o+T_m^e]$; on the real axis its upper end is reduced to $T_m^e$, because the
direct odd remainder is nonnegative. These bounds constrain the retained residuals, not the missing residual of opposite
sign. Writing any allowed range as $[L_{qj},U_{qj}]$, the necessary inequalities on the affine part are
\begin{equation}
 J_{qj}\cdot x\le U_{qj}-d^{(0)}_{qj}+p_{qj},\qquad -J_{qj}\cdot x\le-L_{qj}+d^{(0)}_{qj}+p_{qj}.
 \label{eq:cone:lp}
\end{equation}
Each pair is divided by $\max\{|d^{(0)}_{qj}|,\max_k|J_{qjk}|,p_{qj},10^{-10}\}$ before it is passed to the solver,
which does not change the feasible region. The result is a linear program, solved with the HiGHS backend of SciPy. It is
not a full semidefinite bootstrap calculation, and feasibility of these finitely many necessary inequalities is not an
existence proof of a CFT.

\paragraph{Iterating input constraints and reference bounds.} Starting from the input box and
Eq.~\eqref{eq:cone:reehorst}, each reference numerator of Eq.~\eqref{eq:cone:anchor} is maximized over the current
feasible input region. Equations~\eqref{eq:cone:anchor} to~\eqref{eq:cone:oddenv} then give omitted-tail bounds, which
tighten the allowed crossing residuals and hence the input region; the references are recomputed in the smaller region
and the step is repeated. New inequalities are intersected with the existing ones, not averaged. The nominal calculation
makes ten constraint updates; the archived doubled-padding run makes eight (Appendix~\ref{app:cone:sensitivity} compares
matched ten-update runs). A fixed number of updates is not a convergence criterion. The intermediate intervals of the
iteration use symmetric magnitude envelopes for the tails; they are stages of the construction, not the reported result,
which uses the sign-aware odd integral of Appendix~\ref{app:cone:sign}.

The reference bounds include a chosen multiplicative factor $1.001$, intended to allow for tiny external-dimension
effects; this too is a numerical allowance, not an interval-certified treatment of all continuous dependence on the
dimensions. The final problem has 552 scaled inequalities in 154 variables, plus the variable bounds. For the
contribution of the retained inputs the objective is
\begin{equation}
 F_{\rm in}=\frac{\Delta a_{\rm retained}-c_1^{(0)}\Delta b_{\rm retained}}{b_0},
\end{equation}
and its final linear extrema are
\begin{equation}
 -0.001584618275\lesssim F_{\rm in,lin}\lesssim0.001256144319 .
 \label{eq:cone:inputlp}
\end{equation}
A further $\sci{5}{-5}$ is allowed on each side for the nonlinearity of the integrated objective; the separately
sampled sum of objective remainders is about $\sci{2.3952}{-5}$, below this allowance. No covariance matrix is
inferred by fitting crossing residuals, and no independent-Gaussian interpretation is given to the box.

\paragraph{The central input vector is excluded.} Inserting $x=0$ into the \emph{final} nominal constraints $Ax\le B$
violates the 68 rows of Sec.~\ref{sec:howcond} (those with $B_i<-10^{-9}$, all in the $\sigma\sigma\sigma\sigma$
channel); the smallest scaled right-hand side is about $-0.727632$. The central tabulated data are therefore not a
feasible correlator input for the nominal relaxation. This concerns the vector of input data; its integrated value
$c_1^{(0)}$ can, and does, lie inside the final interval. The displacement needed for this fixed problem is measured by
\begin{equation}
 t_*=\min_{x,t}t\quad\text{subject to}\quad Ax\le B,\quad -t\le x_k\le t,\quad t\ge0,
 \label{eq:cone:minnorm}
\end{equation}
with the $h_{\epsilon'}$ restriction kept. The saved nominal matrices give $t_*=0.5735097802$, that is
$2t_*=1.1470195604$ quoted errors. This is the smallest \emph{maximum} displacement over the inputs; it does not say
that every input moves by that amount and has no Gaussian interpretation. Equation~\eqref{eq:cone:minnorm} freezes the
nominal reference bounds, secants and padding, so it does not give the feasibility threshold of a rebuilt input box,
whose arrays and iteration change with the box (Sec.~\ref{sec:howcond}). The exclusion shows a tension among the central
extracted spectrum, the missing-spectrum hypotheses and the chosen numerical relaxation. Under a valid remainder
construction, underestimated tabulated errors or missing low-dimensional spectral weight are natural explanations; the
calculation does not decide between them, does not establish a new operator, and does not prove that these
possibilities exhaust the limitations of the numerical envelopes. Nor does it turn an empirical error spread into a
rigorous physical range by shifting central values until crossing is feasible.

\subsection{Sign-aware integration of the odd missing contribution}\label{app:cone:sign}
\paragraph{Why absolute values lose information.} Integrating an absolute pointwise tail bound gives a valid
conditional magnitude estimate, but it can be loose. In the normalized odd channel
\begin{equation}
 Q\,R_{\rm odd}(\varrho,\xi)=\sum_{p,j}c_{p,j}\varrho^pP_j(\xi),\qquad c_{p,j}\ge0,\quad p\ge G_{\rm odd},
 \label{eq:cone:states}
\end{equation}
where the radial-state energy $p$ is a primary dimension plus an integer descendant level, so it is not restricted to
integers. A reference norm bounds $\sum_{p,j}c_{p,j}r_{\rm ref}^p$. For one such state, the Legendre polynomial is
integrated against the actual odd contribution to $a_0$, whose kernel carries an overall minus sign. Separating a
positive integration weight $W$, which includes $4\pi\,\dd\mu\,v^{-(3-S)}r^{-S-E}/Q$, define
\begin{equation}
 M_j(p;r_{\rm ref})=-\int_{\mathcal R}W(z)\Bigl(\frac{|\rho(z)|}{r_{\rm ref}}\Bigr)^pP_j\Bigl(\frac{\operatorname{Re}\rho}{|\rho|}\Bigr).
\end{equation}
A reference norm $B$ then bounds the integrated omitted contribution by $B$ times the extrema of
$M_j(p;r_{\rm ref})$ over the allowed states. The angular integration is done before the worst case is taken, and the
positive and negative extrema are in general unequal.

\paragraph{Coverage of spins and continuous exponents.} The main calculation evaluates $0\le j\le64$, radial energies
from $G_{\rm odd}$ to $G_{\rm odd}+60$ in steps of $0.1$, and 100 quadrature nodes in each variable. Three analytic
bounds cover the rest. For higher spins, the Legendre bound implied by Eq.~18.14.7 of Ref.~\cite{NIST:DLMF},
\begin{equation}
 |P_j(\xi)|\le\min\Bigl\{1,\ \sqrt{\frac{2}{\pi(j+1/2)\sqrt{1-\xi^2}}}\Bigr\},
 \label{eq:cone:legendre}
\end{equation}
is used with $j+\frac12$ replaced by $J+\frac32$, a common upper bound for every $j\ge J+1$; the cap at 1 prevents an
artificial divergence of the angular estimate at the endpoints. Between sampled exponents with step $h$, the second
derivative with respect to $p$ brings down $\log^2(|\rho|/r_{\rm ref})$, and since $|\rho|/r_{\rm ref}<1$ a uniform
interpolation allowance is
\begin{equation}
 \frac{h^2}{8}\int_{\mathcal R}W(z)\Bigl(\frac{|\rho|}{r_{\rm ref}}\Bigr)^{G_{\rm odd}}\log^2\Bigl(\frac{|\rho|}{r_{\rm ref}}\Bigr).
 \label{eq:cone:interp}
\end{equation}
Above the last sampled exponent, the absolute kernel is bounded by its magnitude at the gap times
$[(2-\sqrt3)/r_{\rm ref}]^{60}$; only reference radii larger than $2-\sqrt3$ are used for this bound. These are
analytic envelopes evaluated by floating-point quadrature. They do not certify the input or nonlinear assumptions, but
they control the tails in the state indices instead of stopping a scan at a convenient value.

\paragraph{Numerical effect.} In the nominal scenario the symmetric magnitude estimate of the odd contribution to $a_0$
is $0.01652521405$, while the sign-aware upper end is only $0.004677262775$:
\begin{equation}
 \Delta a_{\rm odd}\in[-0.01652521405,\,+0.004677262775]
 \label{eq:cone:oddint}
\end{equation}
under the hypotheses, about $[-0.0013485,\,+0.0003817]$ after division by $b_0$. The best upper reference in this run is
$x=0.76$, with $r_{\rm ref}\simeq0.34237382$, and the sampled maximum occurs at $p=7$, $j=2$. Raising the quadrature to
160 nodes and the explicit spin cutoff to 96 and halving the exponent step changes the positive allowance to about
$0.004672754193$; the sampled maximum itself is stable to about $\sci{5}{-15}$, and most of the change is the
intentionally reduced interpolation allowance. The main interval uses the larger, original allowance. Because the odd
remainder enters with this definite sign, lowering the odd gap mainly lowers the lower end of the range
(Sec.~\ref{sec:howcond}).

\subsection{The conditional result and its error budget}\label{app:cone:budget}
\paragraph{Combining numerator and denominator corrections.} The nominal tail estimates are
\begin{equation}
 |\Delta b_{\rm tail}|\le\sci{7.02687010}{-5},\qquad |\Delta a_{\rm even}|\le0.004552089294,
\end{equation}
together with the asymmetric odd interval of Eq.~\eqref{eq:cone:oddint}. With Eq.~\eqref{eq:cone:inputlp}, the
allowances in units of $c_1$ (numerator contributions divided by $b_0$) are
\begin{center}\small
\begin{tabular}{@{}lrr@{}}
\toprule
contribution & lower & upper\\
\midrule
crossing-constrained linear input variation & $-0.00158462$ & $+0.00125614$\\
nonlinear integrated-input allowance & $-0.00005000$ & $+0.00005000$\\
missing even contribution to $a_0$ & $-0.00037145$ & $+0.00037145$\\
missing odd contribution to $a_0$ & $-0.00134846$ & $+0.00038166$\\
missing contribution to $b_0$, times $-c_1^{(0)}$ & $-0.00000157$ & $+0.00000157$\\
\bottomrule
\end{tabular}
\end{center}
These are allowed ranges, not independent standard deviations, and they are added linearly. The denominator is not
frozen in the conversion to the ratio; exactly,
\begin{equation}
 c_1-c_1^{(0)}=\frac{\Delta a-c_1^{(0)}\Delta b}{b_0+\Delta b}.
 \label{eq:cone:exactratio}
\end{equation}
The retained-input bound $|\Delta b_{\rm in}|\le0.001428$ of Eq.~\eqref{eq:cone:denbox} plus the tail bound keeps the
denominator away from zero. With the total scaled numerator range $[F_-,F_+]$ and $d_b=0.001428+|\Delta b_{\rm tail}|$,
the conversion is
\begin{equation}
 \Bigl[c_1^{(0)}+\frac{F_-}{1-d_b/b_0},\ c_1^{(0)}+\frac{F_+}{1-d_b/b_0}\Bigr],\qquad\text{provided}\quad
 F_-\le0\le F_+,\quad 0\le d_b<b_0 .
 \label{eq:cone:ratioguard}
\end{equation}
The nominal range is $[F_-,F_+]=[-0.003356090758,\,0.002060824936]$, so both conditions hold. Taking the smaller
allowed denominator for both ends is conservative only because the two ends have opposite signs; it is not a general
rule for converting ratios. The code checks both conditions with explicit exceptions and does not use a first-order
expansion of the denominator for the quoted range. The result is the nominal interval
$[0.269795922707,\,0.275213500746]$; the archived doubled-padding run with eight updates gives
$[0.269273823430,\,0.275410362633]$ (Sec.~\ref{sec:range} quotes both to seven digits).

\paragraph{Nominal and enlarged padding.} The outward-rounded estimated range $[0.2690,0.2760]$ of Sec.~\ref{sec:range}
covers these two runs. The retained value and the estimated range are reported separately. The range extends about
$-0.0042$ and $+0.0028$ from the rounded retained value $0.2732$, not symmetrically; its rounded midpoint, $0.2725$, is
not a new best fit or a feasible optimized correlator. The range has no claimed confidence level, and its rounding
covers the specified doubled-padding test, not arbitrary changes of the input box, the gaps or the numerical
allowances. All endpoints remain conditional on the omitted-primary gaps, on the validity of the retained-input ranges
and of the nondegenerate-$\epsilon'$ constraint, and on the adequacy of the sampled nonlinear, denominator and
quadrature allowances. The finite crossing relaxation supplies necessary conditions only; its feasible points and
extrema need not be realized by an exact CFT. No assumption-independent physical interval is established.

\subsection{Sensitivity scans: method and further rows}\label{app:cone:sensitivity}
The scans of Sec.~\ref{sec:howcond} rerun the same final method with ten constraint updates each; they are not
alternative estimators to average. In the box scan each quoted error $e_k$ is replaced by $me_k$, and the integrated
sensitivities, the sampled denominator box, the crossing secants and padding, and the reference-bound iteration are all
recomputed. The physical $h_{\epsilon'}$ window of Eq.~\eqref{eq:cone:reehorst} is held fixed, not scaled with the table
errors, so its limits in the scaled variable are recomputed from the changed parameter scale. In the gap scan the
central spectrum and the twice-quoted-error box are unchanged; the same level-40 basis is reused, while the radial
suppression, the reference bounds, the crossing constraints and the integrated tails are recomputed at each threshold.
A third scan multiplies every crossing padding by a constant with ten updates, which isolates the effect of the padding
from that of the iteration count. The rows not shown in Sec.~\ref{sec:howcond} are:
\begin{center}\small
\begin{tabular}{@{}lll@{}}
\toprule
scan & setting & ten-update interval\\
\midrule
box, $m$ (box radius $\pm2m$ quoted errors) & $0.5$ & infeasible\\
 & $0.7$ & $[0.27050311,\,0.27455152]$\\
 & $0.9$ & $[0.27003354,\,0.27506550]$\\
 & $1.5$ & $[0.26848108,\,0.27600079]$\\
 & $3$ & $[0.26380302,\,0.27869935]$\\
\midrule
gaps $(G_{\rm even},G_{\rm odd})$ & $(6.5,7)$ & $[0.26850350,\,0.27679517]$\\
 & $(7,6)$ & $[0.26573717,\,0.27697080]$\\
 & $(6,6)$ & $[0.26419915,\,0.27833854]$\\
 & $(8,5.5)$ & $[0.26156563,\,0.27688176]$\\
\midrule
crossing-padding multiplier & $1$ (nominal) & $[0.26979592,\,0.27521350]$\\
 & $2$ & $[0.26927464,\,0.27540968]$\\
 & $3$ & $[0.26887007,\,0.27557227]$\\
\bottomrule
\end{tabular}
\end{center}
Every feasible range listed contains the retained value $c_1^{(0)}$, but that does not remove the infeasibility of the
central input vector. The scans keep the chosen $\sci{5}{-5}$ allowance for the integrated objective while recomputing
the other quantities; enlarging the box does not certify this fixed allowance on the enlarged domain, so the wider-box
outputs are numerical stress diagnostics, not validated confidence intervals. The sampled nonlinear maxima, the
incomplete control of mixed input cross terms and the absence of interval arithmetic apply throughout.

The iterations are finite and not claimed to be converged. At the end of the original histories the nominal lower end
still changed by order $10^{-6}$ per update, and the eight-update $(6,6)$ history by order $\sci{1.5}{-4}$, so the
ten-update $(6,6)$ row above differs correspondingly from that archived history. Assuming valid outer allowances,
further intersections can tighten the feasible region; stopping earlier is no reason to claim a sharper result.
Improving the precision requires resolving the tension between the spectrum and the inputs and certifying the numerical
envelopes, not rounding the nominal endpoints more finely.

\subsection{Validation}\label{app:cone:validation}
\paragraph{Nonlinear crossing and solver checks.} The nonlinear test uses a fixed pseudorandom seed, 6087, and 34
configurations: sampled interior points, parameter corners and extrema of the linear optimization. It recomputes the
full nonlinear retained residual and compares the affine error with its allocated padding. The largest ratio is
$0.666645$: below one, but about two thirds of the allowance, so the chosen factor $1.5$ provides essentially the
sampled margin rather than a large safety reserve. This is evidence that the padding is adequate on the tested
configurations, not a proof that it covers every point of the continuous box. The test of the integrated objective
samples nine dimension values and the coupling corners family by family, with a separate treatment of the global low
inputs; the summed sampled remainder is $\sci{2.3952}{-5}$, against the allowance $\sci{5}{-5}$. Summing such samples
cannot by itself give a certified statement. The two final objective linear programs have primal--dual gaps of order
$10^{-18}$ and largest scaled constraint violations of about $\sci{8.87}{-10}$. These tests establish solver
consistency for the floating-point affine problem; they do not certify that the problem contains every physical
correlator consistent with the stated gaps.

\paragraph{Descendants, integration and positive radial states.} Raising the descendant order of the crossing basis
from 40 to 48 changes the largest central residual by $\sci{5.75}{-11}$, a Jacobian entry by $\sci{3.30}{-14}$ and a
padding entry by $\sci{7.10}{-14}$; the bases at orders 32, 40 and 48 are all kept. Changing the quadrature of the
integrated envelopes from 80 to 220 nodes changes the even contribution to $a_0$ by about $\sci{1.01}{-7}$, or
$\sci{8.3}{-9}$ after division by $b_0$; the spin and odd magnitudes are much more stable. Separate low-level checks of the
normalized odd radial expansion support the use of Eq.~\eqref{eq:cone:Q}; they supplement the positivity argument, they
do not replace it.

\paragraph{Executable guards and the remaining logical limits.} The final runner reads $c_1^{(0)}$ and $b_0$ from the
current retained-spectrum output, derives the outward-rounded denominator allowance from the input-box output, and
derives the scaled $h_{\epsilon'}$ limits from the actual parameter scale; the chosen objective allowance and the
physical $h_{\epsilon'}$ window are explicit settings in one configuration file. The numerical guards raise exceptions
that stay active when Python assertions are disabled, among them the sign condition of Eq.~\eqref{eq:cone:realsign} and
the conditions of Eq.~\eqref{eq:cone:ratioguard}, and the table of the estimated range is generated only after checking
that it contains the computed nominal and doubled-padding intervals. The verification script tests the level-40 and
level-48 arrays (residual, Jacobian and padding), the normalized radial-state output, the generalized-free remainder
tests, the integrated-tail quadrature and the angular and exponent envelopes, as well as the central values and the
optimality of the linear programs. These are numerical regression checks: no passing test establishes completeness of
the spectrum, a uniform bound on the nonlinearity in the inputs, or the existence of a CFT at an extremum of the linear
program.

One limitation of the even mixed-tail construction is its dependence, through Eq.~\eqref{eq:cone:cauchy}, on the
positive $\epsilon^4$ remainder. The retained $\langle\epsilon\epsilon\epsilon\epsilon\rangle$ correlator can have a
substantially larger crossing mismatch than $\langle\sigma\sigma\sigma\sigma\rangle$ or the mixed correlator. It does
not enter the slope functional directly, but it weakens the auxiliary norm bound, so a small mixed residual alone
cannot replace the full remainder analysis.

\subsection{The implemented conformal-block recurrences}\label{app:cone:recurrences}
\paragraph{The small-$z$ series.} For external differences $d_a$, $d_b$ write
\begin{equation}
 g_{\Delta,\ell}^{d_a,d_b}(z)=\sum_{n\ge0}\sum_{j\ge0}A_{n,j}r^{\Delta+n}P_j(\eta),\qquad A_{0,j}=\delta_{j\ell}.
\end{equation}
Let $C_{\Delta,\ell}=\Delta(\Delta-3)+\ell(\ell+1)$ and $D_{n,j}=(\Delta+n)(\Delta+n-3)+j(j+1)-C_{\Delta,\ell}$. The
coefficients obey
\begin{equation}
\begin{aligned}
 D_{n,j}A_{n,j}={}&\frac{j}{2j-1}(\Delta+n+j-2-d_a)(\Delta+n+j-2+d_b)A_{n-1,j-1}\\
 &+\frac{j+1}{2j+3}(\Delta+n-j-3-d_a)(\Delta+n-j-3+d_b)A_{n-1,j+1},
\end{aligned}
 \label{eq:cone:zrec}
\end{equation}
where the first term is absent for $j=0$ and out-of-range coefficients vanish. At a null-state denominator the
numerator must vanish consistently; otherwise the routine raises an error. Only the scalar angular coefficient
$A_{n,0}$ enters the analytic continuation of Eq.~\eqref{eq:cone:AC}; for block validation the series is also evaluated
at complex points with all angular components.

\paragraph{The radial matrix recurrence.} The coefficients $B_{n,j}$ of Eq.~\eqref{eq:cone:rhoblock} are computed by
matrices acting on vectors of Legendre coefficients (matrices carry hats). Let $\hat M$ represent multiplication by
$\xi$ and $\hat D$ differentiation with respect to $\xi$:
\begin{equation}
 \hat M_{j+1,j}=\frac{j+1}{2j+1},\qquad \hat M_{j-1,j}=\frac{j}{2j+1},\qquad
 \hat D_{k,j}=\begin{cases}2k+1,&k=j-1,\,j-3,\,\dots\ge0,\\0,&\text{otherwise}.\end{cases}
\end{equation}
With the Chebyshev matrix polynomials $T_m(\hat M)$ and $U_m(\hat M)$, built by the usual second-order recurrence,
define
\begin{equation}
 \hat A=2\hat M^2-\hat I,\qquad \hat B=2\hat M(\hat I-\hat M^2)\hat D,\qquad
 a=-\frac{d_a}2,\quad b=\frac{d_b}2,\quad c=a+b,\quad p=ab .
\end{equation}
(These $a$, $b$, $c$, $p$ and the $u$ below are local symbols: here $c$ is not $c_T/c_B$, and $p$ and $u$ are not
the radial energy and radius of the preceding subsections.) Let $\boldsymbol B_n$ be the vector of $B_{n,j}$;
$\boldsymbol B_0$ has only its entry $\ell$ equal to one. For each
$1\le m\le n$ put $u=\Delta+n-m$ and accumulate a right-hand side acting on $\boldsymbol B_{n-m}$:
\begin{itemize}
\item for even $m=2(k+1)$, add $-4\bigl[u\bigl(-\hat AU_k(\hat A)+U_{k-1}(\hat A)-\tfrac12\hat I\bigr)+U_k(\hat A)\hat B\bigr]\boldsymbol B_{n-m}$,
 with $U_{-1}=0$;
\item for odd $m$, add $8c\bigl[u\,T_m(\hat M)+(\hat M^2-\hat I)U_{m-1}(\hat M)\hat D\bigr]\boldsymbol B_{n-m}$;
\item for every $m$, add $16p\,(-1)^{m-1}m\,T_m(\hat M)\boldsymbol B_{n-m}$.
\end{itemize}
The $j$th component of the total right-hand side, divided by $D_{n,j}$, is $B_{n,j}$. The same quadratic Casimir
eigenvalue governs both series, but the recurrences and the evaluation coordinates differ. The implementation uses
matrices of size $\ell+N+8$ and returns angular spins up to $\ell+N$. Denominators smaller than $10^{-9}$ are treated as
null candidates, whose numerator residual must be below $10^{-8}$; these are floating-point consistency thresholds, not
regulators of a genuine nonzero null-state obstruction, and the code fails rather than keep an inconsistent
coefficient.

\subsection{The retained input spectrum}\label{app:cone:spectrum}
Table~\ref{tab:cone:spectrum} lists the central values actually evaluated, grouped by internal $\mathbb Z_2$ parity and
by the trajectory labels \texttt{ss0}, \texttt{ee0}, \texttt{ss1}, \texttt{se0} and \texttt{other}. These labels are
bookkeeping inherited from the published selection; they do not prove a classification of the CFT spectrum. All rows
come from the stable-spectrum tables of Ref.~\cite{Simmons-Duffin:2016wlq}, except that the leading scalars and the
stress tensor (marked $\dagger$) are replaced by the inputs of Ref.~\cite{Chang:2024whx} and by
Eq.~\eqref{eq:cone:ward}. For even rows the columns are $(f,h)=(f_{\sigma\sigma A},f_{\epsilon\epsilon A})$, and for odd
rows the single coefficient is $f_{\sigma\epsilon B}$, all in the $f$ normalization of Appendix~\ref{app:cone:inputs}. The
overall sign of a primary is conventional, but the signs of the even products are kept. The decimal places specify the
numerical inputs, not a guaranteed physical accuracy. The package's input files also hold every quoted dimension and
coupling error, the values before replacement and the source table of each row; the stress-tensor row has zero
independent errors because its variations are inherited from the shared low inputs, which is not a claim that the
stress couplings are exact.

{\small
\renewcommand{\arraystretch}{1.05}
\begin{longtable}{@{}lrrrr@{}}
\caption{The retained input spectrum of the $c_1$ package: 37 internally $\mathbb Z_2$-even and 23 internally
$\mathbb Z_2$-odd primaries, with central values. $\dagger$: replaced by Ref.~\cite{Chang:2024whx} ($\epsilon$, $\sigma$)
or by the Ward identity, Eq.~\eqref{eq:cone:ward} ($T$, at $\Delta=3$).}\label{tab:cone:spectrum}\\
\toprule
family & $\ell$ & $\Delta$ & $f$ & $h$\\
\midrule\endfirsthead
\toprule
family & $\ell$ & $\Delta$ & $f$ & $h$\\
\midrule\endhead
\bottomrule\endfoot
\multicolumn{5}{@{}l}{\emph{internally $\mathbb Z_2$-even}: $f=f_{\sigma\sigma A}$, $h=f_{\epsilon\epsilon A}$}\\
\texttt{ss0}$^\dagger$ & $2$ & $3$ & $0.3261376622$ & $0.8891467104$\\
\texttt{ss0} & $4$ & $5.022665$ & $0.069076$ & $0.24792$\\
\texttt{ss0} & $6$ & $7.028488$ & $0.0157416$ & $0.066136$\\
\texttt{ss0} & $8$ & $9.031023$ & $0.003685$ & $0.017318$\\
\texttt{ss0} & $10$ & $11.0324141$ & $0.00087562$ & $0.0044811$\\
\texttt{ss0} & $12$ & $13.033286$ & $0.00020992$ & $0.00115174$\\
\texttt{ss0} & $14$ & $15.033838$ & $\sci{5.065}{-5}$ & $0.00029484$\\
\texttt{ss0} & $16$ & $17.034258$ & $\sci{1.228}{-5}$ & $\sci{7.517}{-5}$\\
\texttt{ss0} & $18$ & $19.034564$ & $\sci{2.98935}{-6}$ & $\sci{1.91408}{-5}$\\
\texttt{ss0} & $20$ & $21.0347884$ & $\sci{7.2954}{-7}$ & $\sci{4.8632}{-6}$\\
\texttt{ss0} & $22$ & $23.034983$ & $\sci{1.78412}{-7}$ & $\sci{1.23201}{-6}$\\
\texttt{ee0} & $4$ & $6.42065$ & $0.0019552$ & $-0.110247$\\
\texttt{ee0} & $6$ & $8.4957$ & $0.000472$ & $-0.0431$\\
\texttt{ee0} & $8$ & $10.562$ & $0.0001084$ & $-0.0139$\\
\texttt{ee0} & $10$ & $12.5659$ & $\sci{2.598}{-5}$ & $-0.004437$\\
\texttt{ee0} & $12$ & $14.633$ & $\sci{6.1}{-6}$ & $-0.001224$\\
\texttt{ee0} & $14$ & $16.6174$ & $\sci{1.417}{-6}$ & $-0.0003791$\\
\texttt{ee0} & $16$ & $18.678$ & $\sci{3.547}{-7}$ & $\sci{-9.72}{-5}$\\
\texttt{ee0} & $18$ & $20.654$ & $\sci{7.99}{-8}$ & $\sci{-2.84}{-5}$\\
\texttt{ee0} & $20$ & $22.651$ & $\sci{1.83}{-8}$ & $\sci{-7.58}{-6}$\\
\texttt{ss1} & $0$ & $3.82968$ & $0.053012$ & $1.536$\\
\texttt{ss1} & $2$ & $5.50915$ & $0.0105745$ & $0.69023$\\
\texttt{ss1} & $4$ & $7.38568$ & $0.00237745$ & $0.22975$\\
\texttt{ss1} & $6$ & $9.32032$ & $0.00055657$ & $0.06949$\\
\texttt{ss1} & $8$ & $11.2751$ & $0.00013251$ & $0.0198$\\
\texttt{ss1} & $10$ & $13.241$ & $\sci{3.234}{-5}$ & $0.005459$\\
\texttt{ss1} & $12$ & $15.2301$ & $\sci{7.64}{-6}$ & $0.001538$\\
\texttt{ss1} & $14$ & $17.1944$ & $\sci{1.93}{-6}$ & $0.000386$\\
\texttt{ss1} & $16$ & $19.195$ & $\sci{4.568}{-7}$ & $0.0001107$\\
\texttt{ss1} & $18$ & $21.172$ & $\sci{1.153}{-7}$ & $\sci{2.798}{-5}$\\
\texttt{ss1} & $20$ & $23.167$ & $\sci{2.74}{-8}$ & $\sci{7.45}{-6}$\\
\texttt{other}$^\dagger$ & $0$ & $1.41262528$ & $1.05185373$ & $1.53244304$\\
\texttt{other} & $2$ & $7.0758$ & $0.0004773$ & $0.21882$\\
\texttt{other} & $4$ & $8.941$ & $0.0001173$ & $0.08635$\\
\texttt{other} & $6$ & $10.975$ & $\sci{2.437}{-5}$ & $0.02775$\\
\texttt{other} & $0$ & $6.8956$ & $0.0007338$ & $0.1279$\\
\texttt{other} & $0$ & $7.2535$ & $0.000162$ & $0.1874$\\
\midrule
\multicolumn{5}{@{}l}{\emph{internally $\mathbb Z_2$-odd}: $f=f_{\sigma\epsilon B}$}\\
\texttt{se0} & $2$ & $4.180305$ & $0.38915941$ & \\
\texttt{se0} & $3$ & $4.63804$ & $0.1385$ & \\
\texttt{se0} & $4$ & $6.112674$ & $0.1077052$ & \\
\texttt{se0} & $5$ & $6.709778$ & $0.04191549$ & \\
\texttt{se0} & $6$ & $8.08097$ & $0.0286902$ & \\
\texttt{se0} & $7$ & $8.747293$ & $0.01161255$ & \\
\texttt{se0} & $8$ & $10.0623$ & $0.00745$ & \\
\texttt{se0} & $9$ & $10.77075$ & $0.003115$ & \\
\texttt{se0} & $10$ & $12.0492$ & $0.00194$ & \\
\texttt{se0} & $11$ & $12.787668$ & $0.000823634$ & \\
\texttt{se0} & $12$ & $14.0383$ & $0.0004983$ & \\
\texttt{se0} & $13$ & $14.80006$ & $0.000215$ & \\
\texttt{se0} & $14$ & $16.0305$ & $0.0001291$ & \\
\texttt{se0} & $15$ & $16.81009$ & $\sci{5.587}{-5}$ & \\
\texttt{se0} & $16$ & $18.025$ & $\sci{3.13}{-5}$ & \\
\texttt{se0} & $17$ & $18.81794$ & $\sci{1.44219}{-5}$ & \\
\texttt{se0} & $18$ & $20.01947$ & $\sci{8.442}{-6}$ & \\
\texttt{se0} & $19$ & $20.8246$ & $\sci{3.69}{-6}$ & \\
\texttt{se0} & $20$ & $22.0152$ & $\sci{2.131}{-6}$ & \\
\texttt{se0} & $21$ & $22.83035$ & $\sci{9.512}{-7}$ & \\
\texttt{other}$^\dagger$ & $0$ & $0.518148806$ & $1.05185373$ & \\
\texttt{other} & $0$ & $5.2906$ & $0.057235$ & \\
\texttt{other} & $2$ & $6.9873$ & $0.017413$ & \\
\end{longtable}
}

%====================================================================
\section{Reproduction}\label{app:repro}
%====================================================================
Everything in this note is reproduced from the folder \file{SRI/} of the public repository; paths below are relative
to it unless they are said to be in a package.

\subsection{Layout}
\begin{itemize}
\item \file{code/<pkg>/}: the code and records of each package. Production code is included verbatim, at its git
 commit or as a dated snapshot with a sha256 list (\file{public/SNAPSHOT.sha256}); the tools added for the public replay
 are in \file{public/} and \file{public.mk}.
 \begin{itemize}
 \item \file{dual_d1}: the compendium of the dual of the chain ($\CSR_3$ to $\CSR_8$, the certificates, the record of
  $\CSR_8$), pruned to what the replay and the cited checks need;
 \item \file{indep_d1}: the independent re-derivation of the inputs of $\CSR_5$ to $\CSR_8$;
 \item \file{integ_d1}: the magnetization-scheme calculation of $\CSR_8$ (Appendix~\ref{app:integ}): code, frozen
  records and the engine outputs they are rebuilt from;
 \item \file{c2_closed_d2}: the final-method study of $c_2$ and the three Monte Carlo configurations of the review;
 \item \file{c2_cylinder_d2}: the cylinder-Hamiltonian solver, its cluster records and the continuum analysis;
 \item \file{c1_d3}: the $c_1$ package (review revision of 8 October 2026) with the review's two re-implementations in
  \file{review_indep/};
 \item \file{closure_d1}: the checks of the two arguments of Appendix~\ref{app:clos} (no quoted number): an
  independent check of the closure argument's finite algebra (matrix-word charges, the radial-power identity, the
  $\Og\Og$ product), and in \file{prescription/} the prescription argument's own checker (rerun by
  \texttt{make test}) with an independent re-derivation of the $\beta$ recurrence.
 \end{itemize}
\item \file{docs/<pkg>/}: the documentation of each package (README, REPRODUCE, PROVENANCE, DATA, REVIEWS,
 requirements, \file{results_map.md} with every quoted number, its record, rule and tolerance, and TESTS with the
 export tests on the symmetry cluster).
\item \file{results/quoted_numbers.json}: every quoted number of the paper with its macro and its text as printed, the
 TeX of the two closed-form equations of the dual, and the cells of the table of short-range coefficients;
 \file{results/sr_values.json}: machine-readable values for other folders (Sec.~\ref{app:values}).
\item \file{verify/}: \file{sri_lib.py} (the replay helper: PASS/FAIL/INFO/SKIP table, rounding rule) and
 \file{check_sr_values.py}; \file{note/}: this note and \file{tools/check_note.py}.
\end{itemize}

\subsection{Commands and cost}
From \file{SRI/}:
\begin{verbatim}
make test          # fast tests of every package              about 12 min
make replay        # every quoted number from shipped records about 17 min
make fetch-data    # the hosted inputs (Zenodo data record; see below)
make replay-full   # the heavier reruns                       about 35 min
make check-values  # results/sr_values.json against the tree
make note          # tools/check_note.py of this note
make -C code/<pkg> -f public.mk replay   # one package (also test, replay-full)
\end{verbatim}
Measured on the symmetry cluster (seconds, Intel Xeon Gold 6148 / AMD EPYC 7532 nodes, Python 3.12, unpacked in
node-local storage; \file{docs/<pkg>/TESTS_<pkg>.md}):
\begin{center}\small
\begin{tabular}{@{}llll@{}}
\toprule
package & \texttt{test} & \texttt{replay} & \texttt{replay-full}\\
\midrule
\file{dual_d1} & 23 / 26 & 65 / 70 & 544 / 556 (a)\\
\file{indep_d1} & 258 / 308 (b) & 229 / 264 & (c)\\
\file{integ_d1} & 249 / 249 (b) & 78 / 85 & (c)\\
\file{c2_closed_d2} & 4 / 4 & 39 / 41 & --\\
\file{c2_cylinder_d2} & 170 / 152 & 3 / 3 & 10 / 9\\
\file{c1_d3} & 16 / 14 & 469 / 484 (d) & 1474 / 1486\\
\bottomrule
\end{tabular}
\end{center}
(a) After \texttt{fetch-data}. (b) On 8 cores. (c) Prints the production rerun. (d) Plus 66\,s (Intel) for the
two re-implementations in \file{review_indep/}. Further targets: \file{indep_d1} \texttt{replay-raw RAW=<dir>} (349\,s; the assembly from the raw per-sector engine
outputs), \file{c2_closed_d2} \texttt{rerun-mc-smoke} (24\,s) and \texttt{mc-batch-check} (169\,s; reruns the first
batch of one Monte Carlo chunk with the production flags; bitwise identical on AMD, within $\sci{1.6}{-11}$ of the
largest entry on Intel).

\paragraph{Rule.} A replay recomputes each quoted number of its package from shipped records, rounds it half-up to the
precision at which the paper prints it, and compares it with the printed text in \file{results/quoted_numbers.json}
(for a value with its uncertainty in parentheses, the uncertainty is in units of the last digit). It prints a PASS/FAIL
table and exits non-zero on any FAIL. Numbers that the note prints but the paper does not, and literature values,
give INFO rows. The regenerated records are compared with the shipped ones within the tolerances of
Table~\ref{tab:repro}.

\subsection{Map of the quoted numbers}
Table~\ref{tab:repro} gives, for each result, the package, the record (path in the package), the target that checks it
and the tolerance beyond the rounding rule; \file{results/quoted_numbers.json} gives the paper's macro for each. In the
table and below, $G$ stands for \file{results/continuum_2026-10/generations/20261015T000632Z_99019_1e66b9}, the current
generation of the cylinder analysis (named by \file{results/continuum_2026-10/CURRENT}).

{\small
\renewcommand{\arraystretch}{1.2}
\begin{longtable}{@{}>{\raggedright\arraybackslash}p{.205\textwidth}>{\raggedright\arraybackslash}p{.165\textwidth}>{\raggedright\arraybackslash}p{.29\textwidth}>{\raggedright\arraybackslash}p{.26\textwidth}@{}}
\caption{Reproduction map. Target: \texttt{make -C code/<package> -f public.mk <target>}, or \texttt{make <target>}
from \texttt{SRI/}.}\label{tab:repro}\\
\toprule
result & package, target & record (in the package) & tolerance\\
\midrule\endfirsthead
\toprule
result & package, target & record (in the package) & tolerance\\
\midrule\endhead
\bottomrule\endfoot
$\CSR_1$, $\CSR_2$ & \file{dual_d1}, replay & Refs.~\cite{Benedetti:2024wgx,Benedetti:2025nzp}; registry \file{data/coefficients.json} & INFO (literature)\\
$\CSR_3$ to $\CSR_7$ & \file{dual_d1}, test, replay & Eq.~\eqref{eq:closed} with mpmath at 50 digits; \file{data/coefficients.json} & rounding\\
Eq.~\eqref{eq:closed} & \file{dual_d1}, test, replay & \file{data/coefficients.json}, the closed forms as the compendium's TeX states them (verbatim blocks in
\file{public/compendium_extract.json}), the candidates proved by the $\CSR_6$, $\CSR_7$ input maps & exact equality of every rational coefficient\\
certificates of $\CSR_5$ to $\CSR_7$ & \file{dual_d1}, replay; replay-full & $\CSR_5$ ($D_3$) and $\CSR_6$ certificates, $\CSR_6$ and $\CSR_7$ input maps, $\CSR_7$ certificate at weights 2 to 5 (replay); the complete certificate task, with the hosted weight-6 part (replay-full) & exact zero residual (exit 0); input-map reports equal the recorded ones\\
$\pi_0\CSR_8=88$ & \file{dual_d1}, test, replay & \file{data/archived/C8_rational_completion/result.json} \file{/C8_weight_zero} & exact: the record, the full-Jacobian script on the 35 rational inputs and two separate rational sums agree\\
$\CSR_8$ & \file{dual_d1}, test, replay & \file{data/archived/C8_zeta_candidate/result.json} \file{/C8_assembled_without_candidate} & rounding; recomputed from its 36 inputs within $10^{-95}$\\
``about 85 digits'' & \file{dual_d1}, test & the record and its inputs & printing floor about $\sci{2}{-85}$ (INFO)\\
Eq.~\eqref{eq:C8cand} & \file{dual_d1}, test, replay & \file{data/coefficients.json} ($n=8$), \file{result.json} & exact coefficients; residual $\sci{3.085}{-86}$ below the printing floor\\
fresh $P_6$ evaluations & \file{dual_d1}, replay-full & the 66\,404-term $P_6$ expression at 70 and 100 digits & $10^{-(\text{digits}-15)}$\\
largest input pull through $\CSR_7$ & \file{indep_d1}, replay & \file{RESULTS_independent.json} against \file{spec/comparison/archive_comparison.json}; \file{RESULTS_archive_comparison.json} & INFO (not printed by the paper; $0.915\to0.9$); regenerated records below\\
$\CSR_8$ independent and its pull & \file{indep_d1}, replay (replay-raw) & \file{RESULTS_C8_independent.json} \file{/step1/outputs/C8}; \file{RESULTS_C8_archive_comparison.json}; \file{public/pooled_sector_record.json} & INFO (not printed by the paper); central values $10^{-24}$ relative, linear and replicate errors $10^{-12}$, sector-pooled errors $10^{-6}$ of the quoted error\\
$\CSR_8$, magnetization scheme (Eq.~\eqref{eq:C8integ}) & \file{integ_d1}, replay & \file{results/C8_record_v2.json} \file{/C8}, \file{/C8_error}; every record from two bubbles to $\CSR_8$ rebuilt from \file{results/m*_runs/} by the verbatim gates & rounding; records byte-identical (sha256)\\
numbers of Appendix~\ref{app:integ} (inputs, gates, comparison) & \file{integ_d1}, replay & the records and the outputs of the verbatim tools & 25-digit inputs equal to the record; ``to $10^{-n}$'' means in $[10^{-n-1/2},10^{-n+1})$; counts equal\\
engine builds and production sectors & \file{integ_d1}, test & \file{results/builds.json}; sectors of four production chunks rerun & build ids equal; $|\text{rerun}-\text{shipped}|\le4\varepsilon B$ per coefficient ($B$ the engine's bound)\\
$c_2$ closed form (Eq.~\eqref{eq:c2closed}) & \file{c2_closed_d2}, test, replay & \file{outputs/summary.json} \file{/C2}; the study's eight-stage rerun & rounding; 50-digit evaluation within $10^{-45}$\\
$c_2$ Monte Carlo & \file{c2_closed_d2}, test, replay & \file{mc/run_{A,B,C}/out_*.txt} (142 chunks); \file{public/mc_combined.json} \file{/combined/C2}, \file{/combined/C2_se} & rounding; recombination $10^{-12}$ relative; \texttt{summ.py} output equal to the production logs\\
$c_2$ cylinder, its relative difference from the closed form and $N_{\max}$ (also the $c_2$ row of the paper's table of perturbative Hamiltonian truncation) & \file{c2_cylinder_d2}, replay & \file{G/C2_estimate.json} \file{/final/C2}, \file{/final/dC2}, from \file{G/continuum_points.json} & rounding ($N_{\max}$: INFO, not printed by the paper); fits recomputed twice (verbatim and independent) within $10^{-9}$ relative\\
same, from the raw records & \file{c2_cylinder_d2}, replay-full & 16 production, 25 extraction-check and 17 archived records, auxiliary inputs & discrete choices equal; chain $10^{-7}$ absolute; fit diagnostics $10^{-6}+10^{-3}|x|$\\
$c_1$ ($d=3$) & \file{c1_d3}, replay & \file{outputs/tables/current_summary.json} \file{/retained/C1}; the 20-stage workflow rerun from scratch & rounding; rerun within $\sci{5}{-10}$\\
$c_1$ estimated range & \file{c1_d3}, replay & \file{refinement/final_8_7_pad1_summary.json} and \file{refinement/final_8_7_pad2_summary.json} (\file{/interval}) of the rerun; \file{current_summary.json} \file{/recommended_working_summary/interval} & rounding; envelopes within $\sci{2}{-8}$\\
two independent implementations of $c_1$ (Table~\ref{tab:main}) & \file{c1_d3}, replay & \file{review_indep/} (two scripts) on \file{data/inputs/retained_spectrum.json} & each within $10^{-12}$ of the package; their agreement, a bound, INFO (not printed by the paper)\\
feasibility and sensitivity (Sec.~\ref{sec:d3}) & \file{c1_d3}, replay-full & the review-revision diagnostics and the 11 gap and margin and 8 box scenarios & status and endpoints within $\sci{2}{-8}$\\
$c_1$ BRRZ & \file{c1_d3}, replay & Ref.~\cite{Behan:2017emf} & INFO (literature)\\
numbers of orders (eight computed, seven in closed form, two published) & \file{dual_d1}, test & \file{data/coefficients.json} & equal; INFO for those the paper does not print\\
\end{longtable}
}

\subsection{Values for other folders}\label{app:values}
\file{results/sr_values.json} publishes each value with all the digits of its record, its status, and the file and
JSON pointer where it lives in the tree; \texttt{make check-values} resolves every entry. The stable pointers are
\begin{center}\small
\begin{tabular}{@{}ll>{\raggedright\arraybackslash}p{.46\textwidth}@{}}
\toprule
pointer & quantity & file (in \file{SRI/code/})\\
\midrule
\texttt{/d1/C/1} to \texttt{/d1/C/7} & $\CSR_1$ to $\CSR_7$ & \file{dual_d1/data/coefficients.json}\\
\texttt{/d1/C/8} & $\CSR_8$ (record) & \file{dual_d1/data/archived/C8_zeta_candidate/result.json}\\
\texttt{/d1/magnetization/C8} & $\CSR_8$, magnetization scheme & \file{integ_d1/results/C8_record_v2.json}\\
\texttt{/d2/c2/closed} & $c_2$, closed form & \file{c2_closed_d2/outputs/summary.json}\\
\texttt{/d2/c2/mc} & $c_2$, Monte Carlo & \file{c2_closed_d2/public/mc_combined.json}\\
\texttt{/d2/c2/cylinder} & $c_2$, cylinder & \file{c2_cylinder_d2/G/C2_estimate.json}\\
\texttt{/d3/c1/value} & $c_1$ & \file{c1_d3/outputs/tables/current_summary.json}\\
\texttt{/d3/c1/range} & estimated range of $c_1$ & \file{c1_d3/outputs/tables/current_summary.json}\\
\texttt{/d3/S}, \texttt{/d3/E\_SR} & $\Ds^{\rm SR}$, $\De^{\rm SR}$ ($d=3$) & \file{c1_d3/data/inputs/low_inputs.json}\\
\bottomrule
\end{tabular}
\end{center}
The resummation folder reads $\CSR_1$ to $\CSR_8$, $c_2$ (closed form), $c_1$ with its estimated range, and the two $d=3$
inputs at these pointers.

\subsection{What cannot be reproduced}\label{app:cannot}
\begin{itemize}
\item \emph{The compendium's raw integrals cannot be regenerated}, from this folder or from the compendium's original
 archive: 21 ordered-integration stage programs are missing from the archive, so the documented raw regenerations stop
 at the first missing stage. What is offered instead is the replay of every quoted number from the shipped records,
 the exact certificate checks of $\CSR_5$ to $\CSR_7$, and the independent re-derivation of every input of $\CSR_5$
 to $\CSR_8$ (\file{indep_d1}).
\item Six inputs of $\CSR_8$ ($W,S,L_3,M_{21},M_{12},N_2$) are single numerical records with no exact expression; no
 $\CSR_8$ certificate exists.
\item Hosted inputs: the weight-6 $\CSR_7$ certificate (57\,MB, needed only by \file{dual_d1} \texttt{replay-full}) and
 the raw per-sector outputs of \file{indep_d1} (131\,MB, as a 31\,MB archive; needed only by \texttt{replay-raw}) are in
 the project's Zenodo data record, DOI~\href{https://doi.org/10.5281/zenodo.23466766}{10.5281/zenodo.23466766};
 \texttt{make fetch-data} downloads and checks them (the certificate can also be taken from the compendium's original
 archive).
\item \file{dual_d1} needs a case-sensitive work directory: the compendium has two input files whose names differ
 only in case.
\item The base study's checks of the closed expression of $c_2$ (the integer-interval evaluation to 51 decimal places,
 its tail inequalities, elliptic integration and local checks; Appendix~\ref{app:cyl:closedchecks}) are described but
 their code is not in the release; the value of the closed expression is reproduced by \file{c2_closed_d2} (50 digits,
 not as an interval enclosure). The closure argument's own script of finite checks was not
 supplied and is not part of the release (Appendix~\ref{app:clos:checks}).
\item The release has no git history, so it cannot show the order in which the magnetization-scheme calculation
 committed its designs and its data (Appendix~\ref{app:integ:disclosures}): \file{integ_d1} checks that the files its
 tools ask git about equal the committed blobs, and the code inspections checked the order on the history.
\item Production reruns (not needed for any replay): \file{indep_d1} about $\sci{2.2}{4}$ core-hours on the symmetry
 cluster; the engine runs of \file{integ_d1} about $\sci{1.6}{4}$ core-hours on the symmetry cluster; the Monte Carlo of \file{c2_closed_d2} about 135 core-hours; the cylinder records of \file{c2_cylinder_d2}
 about 1.5 node-hours on the symmetry cluster, about 470 service units on SDSC Expanse, and, again on the symmetry
 cluster, the two rebuilt $N=1024$ grids (2.6 to 2.7 hours each) and twelve extraction-check jobs (0.3 to 1.6 hours
 each). A Monte Carlo rerun is statistically equivalent, not identical, unless the same binary, weights and seeds are
 used.
\end{itemize}

\subsection{Numbers of this note that the paper does not quote}
The decimal strings of Table~\ref{tab:main} and Sec.~\ref{sec:closed2} are evaluations of the closed forms
(mpmath, 60 digits; the same digits as \file{results/sr_values.json}). The inputs of Table~\ref{tab:C8inputs}, the
pulls, the estimators and the checks of Sec.~\ref{sec:indep} are in the result records of \file{code/indep_d1/} (and
its \file{RESULTS.md}, \file{RESULTS_C8.md}); the eighth-order sub-results of Sec.~\ref{sec:C8} in those of
\file{code/dual_d1/}; the constants of Sec.~\ref{sec:closed2} in \file{code/c2_closed_d2/outputs/summary.json}; the
per-order table and the error budget of Sec.~\ref{sec:cyl}, and the superseded first analysis, in
\file{G/C2_estimate.json}, \file{G/REPORT.md} and \file{results/continuum_2026-10/evidence/rev4/} of
\file{code/c2_cylinder_d2/}; and the intervals and sensitivity tables of Sec.~\ref{sec:d3} in
\file{code/c1_d3/outputs/tables/} (the box scan of Sec.~\ref{sec:d3} at $m=2$ is the package's; the review's own scan
gave $[0.26698,0.27693]$). The NRG residuals of Sec.~\ref{sec:d1caveats} are checked in the NRG folder of the
repository~\cite{NRGnote}. The magnetization-scheme values of Sec.~\ref{sec:integ} and Appendix~\ref{app:integ}
other than the one the paper quotes are in the records of \file{code/integ_d1/results/}
(Appendix~\ref{app:integ:records}), and the replay of \file{integ_d1} checks each of them.

%====================================================================
\section{The eighth order from the integrable point: the magnetization scheme}\label{app:integ}
%====================================================================
This appendix gives the third calculation of $\CSR_8$ summarized in Sec.~\ref{sec:integ}: its derivation, method,
checks and records. It starts from the wall gas of Appendix~\ref{app:dualA:gas}, with the wall weight $W_n$ of
Eq.~\eqref{eq:dualA:moments}, but it uses nothing of the matched two-coupling prescription of Sec.~\ref{sec:dual} (no
hard core, no finite interval, no endpoint source) and no input of the compendium or of the independent
re-derivation. Its code and records are the package \file{code/integ_d1/} (Appendix~\ref{app:repro}). Its result is
numerical in the sense of Sec.~\ref{sec:words}.

\paragraph{Notation.} As elsewhere $\db=1-s$; $v$ is the charge of the wall weight $W_n$ (at $\db=0$ a pair of
opposite walls at distance $r$ has weight $r^{-2v}$), and $w=1-v$. In this appendix $h_{\rm m}$ is a uniform magnetic
field (not the current coupling $h$ of the dual), $\bar m$ is the magnetization, the mean of $m(X)$, $\xi$ is the
expansion variable of the magnetization, $\hat u$, $\hat z$ and $\hat F$ are the couplings and the beta function of
the magnetization scheme (Appendix~\ref{app:integ:scheme}), and $f_{ijk}$ is the coefficient of $\hat u^i\hat
z^j\db^k$ in $\hat F$. A ``bubble'' is a pair of consecutive opposite walls, which bounds one reversed domain.

%--------------------------------------------------------------------
\subsection{The integrable point}\label{app:integ:point}
At $s=1$ the wall gas, with the weights $\prod_{i<j}(X_j-X_i)^{2vq_iq_j}$, is the Coulomb gas of Anderson, Yuval and
Hamann for the anisotropic Kondo model~\cite{Anderson:1969xx,Anderson:1970kyb}, and equivalently that of the ohmic
spin--boson model with dissipation strength $\alpha=v$~\cite{Leggett:1987zz}. These models are solvable by Bethe
ansatz~\cite{Andrei:1983,Tsvelick:1983}. For $\db>0$ the interaction $f_{\db}$ of Eq.~\eqref{eq:dualA:fdb} breaks the
integrability, and the long-range fixed point lies at $\db>0$. The integrable point can therefore supply the $\db=0$
part of the flow exactly, to all orders in the wall fugacity, while the dependence on $\db$ needs perturbation theory
in kink clusters.

The observable is the magnetization of the gas on the infinite line in a uniform field $h_{\rm m}$ coupled to $m(X)$.
Relative to the all-up configuration, a configuration of walls carries the extra weight $e^{-2h_{\rm m}\Xi}$, where
$\Xi$ is the total length of the reversed domains (as in $\Xi_{2m}$ of Eq.~\eqref{eq:dualA:moments}). No hard core is
imposed: the short-distance divergences near $v=1$ are regulated by analytic continuation in $v$
(Appendix~\ref{app:integ:clusters}).

%--------------------------------------------------------------------
\subsection{The magnetization at \texorpdfstring{$s=1$}{s=1} from the Bethe ansatz}\label{app:integ:BA}
\paragraph{Bethe-ansatz equations.} The spin sector of the anisotropic Kondo model is an inhomogeneous six-vertex
(XXZ) chain in its gapless regime, with anisotropy $\gamma$, $N$ bulk sites at rapidity $0$ and the impurity at
rapidity $\Theta$; $\Theta\to\infty$ is the weak-coupling scaling limit~\cite{Andrei:1983,Tsvelick:1983}. With the
kernels
\begin{equation}
 a_n(\theta)=\frac1{2\pi}\,\frac{\gamma\sin(n\gamma)}{\cosh(\gamma\theta)-\cos(n\gamma)},\qquad
 \hat a_n(\omega)=\int a_n(\theta)\,e^{\mathrm i\omega\theta}\dd\theta
 =\frac{\sinh\bigl((\pi-n\gamma)\omega/\gamma\bigr)}{\sinh(\pi\omega/\gamma)},
\end{equation}
the roots fill $\theta<\theta_F$ in a field, and the impurity root density $\sigma_{\rm imp}$ solves
\begin{equation}
 \sigma_{\rm imp}(\theta)+\int_{-\infty}^{\theta_F}a_2(\theta-\theta')\,\sigma_{\rm imp}(\theta')\dd\theta'
 =a_1(\theta-\Theta)\quad(\theta<\theta_F),\qquad
 \bar m=1-2\int_{-\infty}^{\theta_F}\sigma_{\rm imp}(\theta)\dd\theta .
 \label{eq:integ:BAE}
\end{equation}
The impurity magnetization is $\frac12-\int\sigma_{\rm imp}$, and $\bar m\in[0,1]$ is twice it. The zero-field bulk
density is $1/(4\cosh(\pi\theta/2))$, so spinon energies scale as $e^{-\pi\theta/2}$, and $h_{\rm m}/T_K\propto
e^{-\pi L/2}$ with $L=\theta_F-\Theta$ and $T_K$ the Kondo scale.

\paragraph{Wiener--Hopf solution.} With $x=\theta_F-\theta$ and $\varsigma(x)=\sigma_{\rm imp}(\theta_F-x)$,
Eq.~\eqref{eq:integ:BAE} becomes $\varsigma(x)+\int_0^\infty a_2(x-y)\varsigma(y)\dd y=a_1(L-x)$ for $x>0$. With
$p=\pi/\gamma$ the kernel factorizes,
\begin{equation}
\begin{gathered}
 1+\hat a_2(\omega)=\frac{2\sinh((p-1)\omega)\cosh\omega}{\sinh(p\omega)}=\Phi_+(\omega)\Phi_-(\omega),\\
 \Phi_-(\omega)=\sqrt{\frac{2\pi(p-1)}{p}}\,
 \frac{\Gamma(1+\mathrm ip\omega/\pi)\,e^{-\mathrm i\omega c/\pi}}{\Gamma(1+\mathrm i(p-1)\omega/\pi)\,
 \Gamma(\frac12+\mathrm i\omega/\pi)},
\end{gathered}
\end{equation}
with $c=p\log p-(p-1)\log(p-1)$ and $\Phi_+(\omega)=\Phi_-(-\omega)$, which tends to 1 at infinity. The standard
argument gives $\int_0^\infty\varsigma=[e^{\mathrm i\omega L}\hat a_1/\Phi_-]_+(0)/\Phi_+(0)$, where $[\cdot]_+$ is
the part analytic in the upper half plane. With $\omega=\mathrm i\pi t/p$ the Gamma functions combine into one
Mellin--Barnes integral,
\begin{equation}
 \bar m=1+\frac1{\sqrt\pi}\,\frac1{2\pi\mathrm i}\int_{c_0-\mathrm i\infty}^{c_0+\mathrm i\infty}
 \frac{\Gamma(t)\,\Gamma\bigl(\frac12-(1-v)t\bigr)}{\Gamma(1+vt)}\,\xi^{-t}\dd t,\qquad -1<c_0<0,
 \label{eq:integ:MB}
\end{equation}
with $v=1-\gamma/\pi$ and $\xi=e^{\gamma(L+c/\pi)}\propto(h_{\rm m}/T_K)^{-2(1-v)}$. The identification
$v=1-\gamma/\pi$ is forced by the exponent: perturbation theory in the wall fugacity $y$ proceeds in powers of
$y^2h_{\rm m}^{-2(1-v)}$, the Bethe ansatz in powers of $h_{\rm m}^{-2\gamma/\pi}$. The free-fermion point
$\gamma=\pi/2$ is then the Toulouse point $v=\frac12$, and the isotropic limit $\gamma\to0$ is $v\to1$.

\paragraph{The two expansions.} Closing the contour to the left picks up the poles of $\Gamma(t)$ at $t=-n$ and gives
the large-field series
\begin{equation}
 \bar m=1-\frac1{\sqrt\pi}\sum_{n\ge1}\frac{(-1)^{n+1}\,\Gamma\bigl(\frac12+n(1-v)\bigr)}{n!\,\Gamma(1-nv)}\,\xi^n .
 \label{eq:integ:largefield}
\end{equation}
Closing to the right picks up the pole of $\Gamma(t)$ at $t=0$, which cancels the 1, and the poles at
$t=s_m=(m+\frac12)/(1-v)$, and gives a series odd in $h_{\rm m}$,
\begin{equation}
 \bar m=\frac1{\sqrt\pi}\sum_{m\ge0}\frac{(-1)^m\,\Gamma(1+s_m)}{m!\,(m+\frac12)\,\Gamma(1+vs_m)}\,\xi^{-s_m}.
\end{equation}
At $v=0$ Eq.~\eqref{eq:integ:largefield} sums to $\bar m=(1+\xi)^{-1/2}$, the free spin. At $v=\frac12$ only odd
powers survive, since $1/\Gamma(1-\frac n2)=0$ for even $n$, and $\bar m=\frac2\pi\arctan(h_{\rm m}/\Gamma_0)$, the
resonant level of width $\Gamma_0$.

\paragraph{Checks.}
\begin{itemize}
\item \emph{The Wiener--Hopf algebra.} The integral equation~\eqref{eq:integ:BAE} was solved numerically by Nystr\"om
 discretization and compared with Eq.~\eqref{eq:integ:MB}, including the shift $c$. At $\gamma=0.4$, $0.7$ and $1.2$
 ($v=0.873$, $0.777$, $0.618$) and $L=-6$, $-2$, $0$, $2$ and $6$, where $\bar m$ runs from $0.9996$ to $\sci{4}{-5}$,
 the differences in $\bar m$ are at most $\sci{8}{-13}$.
\item \emph{Against the wall gas, with no Bethe ansatz.} The ratio $R=c_2/c_1^2$ of the coefficients of $\xi^2$ and
 $\xi$ in Eq.~\eqref{eq:integ:largefield} does not depend on the normalization of $\xi$:
 \begin{equation}
  R_{\rm BA}=\frac{\sqrt\pi\,\Gamma(\frac52-2v)\,\Gamma(1-v)^2}{2\,\Gamma(1-2v)\,\Gamma(\frac32-v)^2}.
 \end{equation}
 On the gas side $R$ is the connected two-bubble cluster with the field weight,
 \begin{equation}
 \begin{gathered}
  R_{\rm gas}=\frac{(4v-3)\,A_2}{2\,\Gamma(2-2v)^2},\qquad
  A_2=\Gamma(3-4v)\int_0^1\bigl(t(1-t)\bigr)^{-2v}\Bigl[I(t)-\frac12\Bigr]\dd t,\\
  I(t)=\int_0^\infty\Bigl[\frac{(t+g)^{2v}(g+1-t)^{2v}}{g^{2v}(1+g)^{2v}}-1\Bigr]\dd g,
 \end{gathered}
 \end{equation}
 which converges for $v<\frac12$ (where $R_{\rm BA}$ is the analytic continuation of the formula). The two agree
 exactly at $v=0$ ($R=\frac32$) and to $\sci{2}{-13}$, $\sci{3}{-13}$, $\sci{4}{-12}$, $\sci{4}{-10}$ and $\sci{1}{-9}$
 at $v=0.1$, $0.2$, $0.3$, $0.4$ and $0.45$, each within the quadrature error of the gas integral.
\item \emph{Higher clusters.} The later stages reproduce the exact two-, three- and four-bubble clusters at $\db=0$
 (Appendix~\ref{app:integ:gates}).
\end{itemize}
The derivation covers $\frac12<v<1$ directly (real roots, $0<\gamma<\frac\pi2$); for $v<\frac12$ the formula is the
analytic continuation in $v$, which the second check tests. The Bethe-ansatz equations and the identification of the
gas with the anisotropic Kondo model are taken from the literature; everything from the impurity integral equation
on is derived here.

%--------------------------------------------------------------------
\subsection{What the integrable point can and cannot test}\label{app:integ:gauge}
At $\db=0$ the change $u\to H/4$ brings the flow of Sec.~\ref{sec:invariants} to coordinates in which $H=4u$
exactly. The freedom that remains is the shifts $z\to z+u\,\chi(u,z)$, with $u$ readjusted so that $H=4u$ still
holds. Write $F_{ij}$ for the coefficient of $u^iz^j$ of $F$ in these coordinates; $F(0,z)=z$. Through total degree
four in $(u,z)$, the order that $\CSR_8$ needs, $F$ has ten coefficients with $i\ge1$. The linearized action of the
shifts on them has rank eight (in exact arithmetic, at random flows), so exactly two combinations are invariant:
\begin{equation}
 F_{10}=1,\qquad I_2=F_{20}-\tfrac12F_{11}-F_{12}.
\end{equation}
In the inputs of the hard-core scheme (Sec.~\ref{sec:inputs1}), $I_2=b+\frac f4-\frac a2-\frac c8-A$. With the closed
forms of $a$, $b$, $c$, $f$ (Sec.~\ref{sec:inputs1}) and $A=-8-4\zeta(2)-2\zeta(3)+16\log2$
(Appendix~\ref{app:dualB:D3}) it vanishes identically: the rational, $\zeta(2)$, $\zeta(3)$ and $\log2$ parts cancel
separately, and $\log3$, $\log^22$ and $\Li_2(-\frac12)$ cancel between $b$ and $f/4$. The Bethe-ansatz flow of
Appendix~\ref{app:integ:scheme}, whose $\db=0$ part is already in these coordinates, gives $F_{10}=f_{100}=1$ and
$I_2=f_{200}-\frac12f_{110}-f_{120}=0$ (to $10^{-54}$ in the first fit and $10^{-63}$ in the frozen matching). The agreement is a check of the compendium's $s=1$ sector by
integrability.

All other $s=1$ coefficients through degree four can be changed by the shifts: among them $J_9$, and with it
$A_9[\varepsilon^0\ell^1]$, and the current-type inputs $L_3$, $N_2$ and $K_v$ of $\CSR_8$ (the direction of $J_9$ lies
in the span of the linearized action). They are coordinate choices of the hard-core scheme, and only $\CSR_8$ as a
whole is physical. The $s=1$ solution alone therefore cannot test $A_9$ or $J_9$; the calculation below instead
computes $\CSR_8$ in its own scheme.

%--------------------------------------------------------------------
\subsection{The magnetization scheme}\label{app:integ:scheme}
\paragraph{Couplings.} For the gas of bare charge $v$ in the field $h_{\rm m}$, at $0<\db<1$, define
\begin{equation}
 K_h=v\,(2h_{\rm m})^{-\db},\qquad 1+\hat z=K_h\,\bar m^{\,2-\db},\qquad
 \hat u=\frac{2-\db}{4}\,\frac{\dd\log\bar m}{\dd\log h_{\rm m}} .
 \label{eq:integ:couplings}
\end{equation}
Then, identically by these definitions,
\begin{equation}
 h_{\rm m}\frac{\dd\hat z}{\dd h_{\rm m}}=(1+\hat z)\,\hat H,\qquad \hat H=4\hat u-\db,
 \label{eq:integ:Hhat}
\end{equation}
and the only function to be computed is $\hat F$ in
\begin{equation}
 h_{\rm m}\frac{\dd\hat u}{\dd h_{\rm m}}=2\hat u\,\hat F(\hat u,\hat z,\db),\qquad
 \hat F=\hat z+\sum_{i+j+k\ge1}f_{ijk}\,\hat u^i\hat z^j\db^k,\qquad f_{0j0}=0,\quad f_{100}=1 .
 \label{eq:integ:Fhat}
\end{equation}
The flow is in $\log h_{\rm m}$, and Eqs.~\eqref{eq:integ:Hhat} and~\eqref{eq:integ:Fhat} have the form of
Eq.~\eqref{eq:dualA:essential} with $(\hat u,\hat z,\hat F,\hat H)$ in place of $(u,z,F,H)$.

\paragraph{Fixed point and invariants.} $\hat H=0$ gives $\hat u_*=\db/4$ exactly, and $\hat F(\db/4,\hat z_*,\db)=0$
fixes $\hat z_*$. Since $\hat H_{\hat z}=0$ and $\hat H_{\hat u}=4$, the trace and determinant of
Eq.~\eqref{eq:dualA:stability} per unit $\log h_{\rm m}$ are $T_h=\frac\db2\hat F_{\hat u}$ and
$D_h=-2\db(1+\hat z_*)\hat F_{\hat z}$ at the fixed point. The field has dimension $1-\Ds=1-\db/2$, so per unit
logarithm of length
\begin{equation}
 T=\Bigl(1-\frac\db2\Bigr)\frac\db2\,\hat F_{\hat u},\qquad
 D=-\Bigl(1-\frac\db2\Bigr)^{\!2}\,2\db\,(1+\hat z_*)\,\hat F_{\hat z},
 \label{eq:integ:TD}
\end{equation}
and Eqs.~\eqref{eq:eigen} and~\eqref{eq:conv} give the $\CSR_n$. The leading $\hat F=\hat z+\hat u$ gives
$T=\db/2+O(\db^2)$ and $D=-2\db+O(\db^2)$, hence $\CSR_1=-2$, $\CSR_2=\frac12$. The next orders are, for example,
\begin{equation}
 \CSR_3=\frac{39}{16}+2f_{001}-2f_{011}-\frac12f_{110},\qquad
 \CSR_4=f_{101}+\frac12f_{200}-f_{001}f_{110}-\frac14f_{110}-\frac12,
\end{equation}
and $\CSR_8$ is affine in its two new coefficients, $\partial\CSR_8/\partial f_{202}=2$ and $\partial\CSR_8/\partial
f_{301}=\frac34$; the $\db=0$ coefficient $f_{400}$ enters as $f_{400}/4$, the counterpart of $J_9/4$ in
Eq.~\eqref{eq:C8assembly}.

\paragraph{Which coefficients, from how many kinks.} $\hat u\propto y^2$ counts bubbles, so $f_{ijk}$ comes from
connected clusters of $i+1$ bubbles, that is $2i+2$ kinks; $k$ is the number of logarithm insertions from expanding
$f_{\db}$, and the dependence on $\hat z$ comes from the expansion in $w$ about $v=1$. $\CSR_8$ needs the $f_{ijk}$ with
$i+j+k\le4$, 29 of them:
\begin{center}\small
\begin{tabular}{@{}>{\raggedright\arraybackslash}p{.27\textwidth}>{\raggedright\arraybackslash}p{.25\textwidth}>{\raggedright\arraybackslash}p{.24\textwidth}>{\raggedright\arraybackslash}p{.16\textwidth}@{}}
\toprule
coefficients & content & kinks (integration dimension) & first needed for\\
\midrule
$f_{ij0}$ & $\db=0$, from Eq.~\eqref{eq:integ:largefield} & none & $\CSR_3$\\
$f_{0jk}$, $k\ge1$ & one bubble, all orders in $\db$ & 2 (closed form in derivatives of $\Gamma$) & $\CSR_3$\\
$f_{101}$; $f_{111}$; $f_{121},f_{102},f_{112},f_{103}$ & two bubbles, one to three logarithms & 4 (2D) & $\CSR_4$; $\CSR_5$; $\CSR_6$ to $\CSR_8$\\
$f_{201}$, $f_{211}$ & three bubbles, one logarithm & 6 (4D) & $\CSR_6$, $\CSR_7$\\
$f_{202}$ & three bubbles, two logarithms & 6 (4D) & $\CSR_8$\\
$f_{301}$ & four bubbles, one logarithm & 8 (6D) & $\CSR_8$\\
\bottomrule
\end{tabular}
\end{center}
The $\db=0$ coefficients follow from Eq.~\eqref{eq:integ:largefield}: with $K_h=v$ at $\db=0$,
Eq.~\eqref{eq:integ:couplings} expresses $\hat u$ and $\hat z$ through $\xi$ at $v=1-w$, and eliminating $\xi$ gives
$\hat F(\hat u,\hat z,0)$ to all orders; through degree four the overdetermined matching has residual $10^{-63}$. They
include $f_{400}$: the $s=1$ part of $\CSR_8$ is exact in this scheme, and no nine-kink integral is needed. No
current-type input is needed either, since $\hat H$ is exact.

%--------------------------------------------------------------------
\subsection{Kink clusters and their continuation to \texorpdfstring{$v=1$}{v=1}}\label{app:integ:clusters}
\paragraph{Clusters.} The $\db$-dependent coefficients come from the connected Mayer clusters of $n$ bubbles per unit
length (field scaled to $2h_{\rm m}=1$),
\begin{equation}
 J_n=\frac1{n!}\int\prod_{a=1}^n\mathcal W_a\sum_{\text{connected graphs}}\ \prod_{(ab)}\varphi_{ab},
\end{equation}
where $\mathcal W_a$ is the weight of bubble $a$ (its pair factor and field factor), and the bond $\varphi_{ab}$ is
$C_{ab}-1$ for disjoint bubbles, with $C_{ab}$ the cross-pair factor of $W_n$ between them, and $-1$ for overlapping
ones, which alternation excludes. The expansion of $f_{\db}$ to order $k$ gives $k$ logarithm insertions. A symbolic
map, derived once per order, turns the Laurent coefficients of the $J_n$ at $v=1$ into the $f_{ijk}$.

\paragraph{Radial factor and sectors.} Every monomial of the integrand is homogeneous of degree $-2nv$ in the gaps (the
cross-pair factors are scale invariant), so after the overall size is scaled out the radial integral is exact:
$J_n^{(k)}=\sum_j\Gamma^{(j)}(\alpha_n)A_j^{(k)}$ with $\alpha_n=2n-1-2nv$, the logarithm insertions giving the
derivatives $\Gamma^{(j)}$, and $A_j^{(k)}$ angular integrals over a simplex of dimension $2n-2$. The angular domain is
split into primary Cheng--Wu sectors (the largest gap set to 1), refined by iterated Binoth--Heinrich splitting until
every sum of gaps is a monomial times a positive function: 187 sectors for three bubbles and 4777 for four, of which
3445 are computed, using a mirror symmetry of the bubble patterns.

\paragraph{Exponents.} In a sector every term is $\prod_kt_k^{a_k+b_kw}$ times an analytic function. A set $G$ of
gaps that shrinks together at rate $\rho$, possibly as several disjoint contacts, contributes
$\rho^{a_G+b_Gw}$ with $a_G=(|G|-1)+2S_G$ at $v=1$, where $S_G$ is the sum of $q_iq_k$ over all pairs whose separation
shrinks with $\rho$. A single kink--antikink contact gives $\rho^{-2+2w}$; two simultaneous contacts give
$\rho^{-3+4w}$ and three $\rho^{-4+6w}$. Every Taylor term whose integrated exponent is at most $-1$ at $w=0$ is
continued analytically, which produces poles in $w$. No term has $a_k\le-1$ with $b_k=0$ (a divergence independent of
$v$) for $n\le4$.

\paragraph{Continuation.} Each sector variable has $[0,1]=[0,\tau]\cup[\tau,1]$. On $[\tau,1]$ the integral is done by
Gauss--Legendre quadrature in $\log t$. On $[0,\tau]$ the Taylor coefficients of the analytic part are obtained by the
trapezoid rule on the circle $|t|=\tau$ and integrated analytically against $t^{c+n}\log^pt$, which continues every
non-integrable power. The radius $\tau$ is set by the analyticity radius of the regular part: sampled for three bubbles,
and for four bubbles a rigorous radius from Rouch\'e's theorem with a joint polydisc check. All quantities are
truncated Laurent series in $w$, so the Laurent coefficients of the clusters at $v=1$ come out directly. The engine is
written in C++, in double-double arithmetic (about 32 digits) for three bubbles; for four bubbles it runs in double
precision with double-double accumulators, the largest-bound sectors entirely in double-double. It gives a first-order
roundoff bound for every sector.

\paragraph{What the poles test.} The raw Laurent coefficients carry information in their pole parts as well as their
finite parts, and all of them are computed. The checks are that the $f_{ijk}$ have no negative powers of $w$ (the
analyticity relations, listed symbolically for each order), the exact $\db=0$ clusters from
Eq.~\eqref{eq:integ:largefield}, and the closed forms of $\CSR_3$ to $\CSR_7$.

%--------------------------------------------------------------------
\subsection{Gates and results}\label{app:integ:gates}
The plan fixed the stages, their checks and their stopping rules before any multi-bubble cluster with logarithm
insertions was computed; each stage gated the
next.
\begin{itemize}
\item \emph{One bubble.} The poles cancel in $\hat F(0,\hat z,\db)$ through $\db^3$, and $f_{001}=-\gamma_E$.
\item \emph{$\CSR_3$.} With $f_{001}=-\gamma_E$, $f_{011}=-\gamma_E-2\zeta(2)$ and $f_{110}=4\zeta(2)$,
 $\CSR_3=\frac{39}{16}+2\zeta(2)$, to $\sci{2}{-38}$ (absolute) numerically.
\item \emph{Two bubbles ($\CSR_4$, $\CSR_5$).} The cluster at $\db=0$, continued to $v=1$, matches the Bethe ansatz to
 $\sci{4}{-32}$, and all 36 analyticity relations hold to $10^{-28}$. The value of $f_{101}$ implied by $\CSR_4$ was
 printed in advance as a blind check; the computation gave $f_{101}=-36.49853744243762\ldots$. $\CSR_4$ agrees to
 $\sci{3}{-28}$ and $\CSR_5$ to $10^{-26}$ (absolute).
\item \emph{Three bubbles, one logarithm ($\CSR_6$, $\CSR_7$).} The engine matches the exact three-bubble cluster to
 $10^{-23}$ at real $v$ and $10^{-28}$ at complex $v$, and the analyticity relations hold to $10^{-27}$. With the results
 hashed before the comparison, $\CSR_6$ agrees to $\sci{4}{-24}$ and $\CSR_7$ to $\sci{2}{-24}$ (absolute).
\item \emph{Three bubbles, two logarithms ($f_{202}$).} The $\db^0$ part matches the Bethe ansatz to $\sci{4}{-23}$, the
 $\db^1$ part the frozen one-logarithm record to $\sci{9}{-24}$, and all twelve $\db^2$ analyticity relations hold to
 $\sci{4}{-24}$. $f_{202}=172415.83018563\pm\sci{1.7}{-7}$, the difference between two quadrature grids.
\item \emph{Four bubbles, one logarithm ($f_{301}$).} Three complete runs: a fine and a coarse grid at contour-radius
 ratio $0.2$, and a radius-check run at $0.3$; in each, the sectors with the largest roundoff bounds were recomputed
 entirely in double-double arithmetic (8 at ratio $0.2$, 14 at $0.3$). The $\db^0$ part matches the exact four-bubble
 cluster to $\sci{4}{-11}$, and all twelve $\db^1$ analyticity relations hold to $\sci{7}{-12}$. A confirmatory check of
 the $\db^0$ part at two complex values of $v$ gives $\sci{1.7}{-14}$ and $\sci{1.6}{-12}$ (the second under amendment
 A of Appendix~\ref{app:integ:disclosures}).
\item \emph{Roundoff, measured.} 56 further sectors were recomputed in double-double arithmetic at the fine grid: the
 16 with the next-largest bounds, 24 drawn by bound decile, and 16 from the two top deciles as a held-out sample fixed
 in advance. The difference between double and double-double is at most $\sci{2.5}{-3}$ of the first-order bound in the
 top 16 sectors and the top decile, which together carry about 98\% of the bound, at most $\sci{1.3}{-2}$ in the next
 decile, and reaches $0.41$ only in sectors whose bound is negligible.
\item \emph{Result.} $f_{301}=-514275.75\pm0.17$; the uncertainty, $0.174$, is the linear sum of the fine--coarse difference $0.016$, the
 fine--radius difference $0.022$ and the validated roundoff estimate $0.137$.
\end{itemize}
The inputs were hashed and frozen before $\CSR_8$ was assembled. The same assembly reproduces the closed forms of
Eq.~\eqref{eq:closed}:
\begin{center}\small
\begin{tabular}{@{}lllr@{}}
\toprule
& magnetization scheme & closed form & relative difference\\
\midrule
$\CSR_3$ & $5.72736813369645287294483$ & $5.72736813369645287294483$ & $\sci{3}{-39}$\\
$\CSR_4$ & $-12.19619149266966002908757$ & $-12.19619149266966002908757$ & $\sci{2}{-29}$\\
$\CSR_5$ & $29.23781158099003956882694$ & $29.23781158099003956882694$ & $\sci{3}{-28}$\\
$\CSR_6$ & $-58.66569747422530894971715$ & $-58.66569747422530894971715$ & $\sci{7}{-26}$\\
$\CSR_7$ & $51.06898861292810883118393$ & $51.06898861292810883118394$ & $\sci{4}{-26}$\\
\bottomrule
\end{tabular}
\end{center}
and gives
\begin{equation}
 \CSR_8=247.381179745789\pm0.131,
\end{equation}
quoted as $247.38118\pm0.13$ in Eq.~\eqref{eq:C8integ}. The uncertainty is the linear sum of the propagated input
uncertainties: $\frac34\times0.174=0.13$ from $f_{301}$, $\sci{3.3}{-7}$ from $f_{202}$ and $\sci{3.6}{-15}$ from all
other inputs. Table~\ref{tab:integ:inputs} lists the 29 inputs, from which Eq.~\eqref{eq:integ:TD} reassembles
$\CSR_3$ to $\CSR_8$: from the printed 25-digit values $\CSR_8$ equals the record to all its printed digits, and
$\CSR_3$ to $\CSR_7$ come out to $10^{-25}$ to $\sci{2}{-22}$ relative (the table above used the unrounded inputs). The computation used about $\sci{2.7}{4}$ core-hours on the symmetry cluster, including pilots and failed jobs; the
engine runs behind the records, which a rerun repeats, took about $\sci{1.6}{4}$.

{\small
\renewcommand{\arraystretch}{1.1}
\begin{longtable}{@{}lllll@{}}
\caption{The 29 inputs of $\CSR_8$ in the magnetization scheme, as frozen in the record of $\CSR_8$
(Appendix~\ref{app:integ:records}). Index $ijk$: $\hat u^i\hat z^j\db^k$, with $i+1$ bubbles and $k$ logarithms.
The $\db^0$ entries come from the exact Bethe-ansatz formula (matching residual $10^{-63}$); the one-bubble entries are exact or
high-precision numerics; the two-bubble entries are good to about $10^{-26}$ relative and the three-bubble $\db^1$
entries to about $10^{-22}$. (a) $f_{202}$ and $f_{301}$ are printed in full only for exact reassembly; they are
accurate only to $\pm\sci{1.7}{-7}$ and $\pm0.174$.}\label{tab:integ:inputs}\\
\toprule
& monomial & value & source & closed form\\
\midrule\endfirsthead
\toprule
& monomial & value & source & closed form\\
\midrule\endhead
\bottomrule\endfoot
$f_{100}$ & $\hat u$ & $1$ & Bethe ansatz (exact) & $1$\\
$f_{110}$ & $\hat u\hat z$ & $6.579736267392905745889661$ & Bethe ansatz (exact) & $4\zeta(2)$\\
$f_{120}$ & $\hat u\hat z^2$ & $42.00883821061607434137304$ & Bethe ansatz (exact) & \\
$f_{130}$ & $\hat u\hat z^3$ & $174.5896405924634503829158$ & Bethe ansatz (exact) & \\
$f_{200}$ & $\hat u^2$ & $45.29870634431252721431787$ & Bethe ansatz (exact) & \\
$f_{210}$ & $\hat u^2\hat z$ & $777.0075210550932975552262$ & Bethe ansatz (exact) & \\
$f_{220}$ & $\hat u^2\hat z^2$ & $8099.313348527777075847784$ & Bethe ansatz (exact) & \\
$f_{300}$ & $\hat u^3$ & $3848.445467772377268238629$ & Bethe ansatz (exact) & \\
$f_{310}$ & $\hat u^3\hat z$ & $115632.5319783609240553502$ & Bethe ansatz (exact) & \\
$f_{400}$ & $\hat u^4$ & $545193.0244029872263588049$ & Bethe ansatz (exact) & \\
$f_{001}$ & $\db$ & $-0.5772156649015328606065121$ & one bubble & $-\gamma_E$\\
$f_{011}$ & $\hat z\db$ & $-3.867083798597985733551342$ & one bubble & $-\gamma_E-2\zeta(2)$\\
$f_{021}$ & $\hat z^2\db$ & $-8.098095746334830014543783$ & one bubble & \\
$f_{031}$ & $\hat z^3\db$ & $-13.46681348232748267372698$ & one bubble & \\
$f_{002}$ & $\db^2$ & $5.292133223877124175286$ & one bubble & \\
$f_{012}$ & $\hat z\db^2$ & $43.26016741947440932389841$ & one bubble & \\
$f_{022}$ & $\hat z^2\db^2$ & $194.2323046570014639460987$ & one bubble & \\
$f_{003}$ & $\db^3$ & $-61.14566451935350750983253$ & one bubble & \\
$f_{013}$ & $\hat z\db^3$ & $-767.6867345496617072942533$ & one bubble & \\
$f_{101}$ & $\hat u\db$ & $-36.49853744243762329795634$ & two bubbles & \\
$f_{111}$ & $\hat u\hat z\db$ & $-404.1590010240865956167455$ & two bubbles & \\
$f_{121}$ & $\hat u\hat z^2\db$ & $-2982.066378636498938272215$ & two bubbles & \\
$f_{102}$ & $\hat u\db^2$ & $810.9918099349166764628352$ & two bubbles & \\
$f_{112}$ & $\hat u\hat z\db^2$ & $14885.33541859945105628608$ & two bubbles & \\
$f_{103}$ & $\hat u\db^3$ & $-23498.57164495448058791384$ & two bubbles & \\
$f_{201}$ & $\hat u^2\db$ & $-3157.008605154717628460989$ & three bubbles & \\
$f_{211}$ & $\hat u^2\hat z\db$ & $-75994.36349692199322461628$ & three bubbles & \\
$f_{202}$ & $\hat u^2\db^2$ & $172415.8301856328384692596$ & three bubbles & (a)\\
$f_{301}$ & $\hat u^3\db$ & $-514275.7548978331544915535$ & four bubbles & (a)\\
\end{longtable}
}

%--------------------------------------------------------------------
\subsection{Comparison, disclosures and review}\label{app:integ:disclosures}
\paragraph{Comparison.} The comparison stage was run after the last code inspection had approved, and it refused to
run unless the record of $\CSR_8$ was the current frozen one, bound to the current records of its inputs and to a
passing confirmatory check. With $\Delta$ the magnetization-scheme value minus the reference:
\begin{center}\small
\begin{tabular}{@{}llrrr@{}}
\toprule
reference & value & $\Delta$ & $\Delta/0.131$ & $\Delta$/own error\\
\midrule
record, Eq.~\eqref{eq:C8num} & $247.3812048049935\ldots$ & $\sci{-2.506}{-5}$ & $-0.0002$ & --\\
independent, Sec.~\ref{sec:indep} & $242.825\pm2.371$ & $+4.556$ & $+34.8$ & $+1.92$\\
\bottomrule
\end{tabular}
\end{center}
(The comparison stage of the calculation used the rounded $242.83\pm2.37$, which gives $\Delta=+4.551$.)

\paragraph{Disclosures.} The review of the code approved quoting the value on condition that the following are stated.
\begin{itemize}
\item The central value was frozen, by hash, before any comparison. A first version, $247.38118\pm0.077$ with the same
 central value, is superseded because its roundoff estimate had not been validated.
\item The session coordinating the calculation saw the comparison when that first version was frozen. The session
 running the computation was exposed to it later: after the central values and the design of the additional runs for
 the error bar were frozen, but before the first record of the confirmatory check was computed (from finished runs,
 against its fixed criterion, which it failed) and before the error bar was evaluated.
\item Two stopping rules fixed in advance fired during that work. The author decided both amendments after seeing the
 data; the failing first records stay frozen.
 \begin{itemize}
 \item[A.] \emph{Confirmatory check.} The second complex point missed its criterion $10^{-12}$ by roundoff
  ($\sci{1.6}{-12}$, against that run's own first-order roundoff bound $\sci{1.7}{-9}$). A double-double rerun (about
  $\sci{1.7}{4}$ core-hours) was not made; the criterion was amended to at most $10^{-11}$ and within the run's own
  roundoff bound.
 \item[B.] \emph{Roundoff validation.} The global rule, a stop if in any sampled sector the difference between the
  double and double-double results exceeds $0.1$ of the first-order bound, fired on sectors of negligible bound; applied as written it would
  have made the uncertainty of $\CSR_8$ about 16. The rule was amended to apply the same factor-ten margin decile by
  decile, with a stop only for a ratio above $0.1$ among the sectors ranked 9 to 24 by bound (the top 8, recomputed entirely in
  double-double, are checked separately) or in the two top deciles, and it
  was tested on the fresh held-out sample, fixed before its data, which stayed within the margin.
 \end{itemize}
\end{itemize}

\paragraph{Review.} OpenAI Codex reviewed the plan in three rounds, approving it before any multi-bubble cluster with logarithm insertions
was computed, and
inspected the code six times; the sixth inspection approved quoting the value with the disclosures above. After each
of the last two rounds of fixes, every current record was rebuilt byte for byte with the revised checks.

\paragraph{What the calculation establishes and what it does not.}
\begin{itemize}
\item It gives $\CSR_8=247.38118\pm0.13$ in a scheme that shares no input with the compendium or with the independent
 re-derivation, and that never uses $A_9$. It agrees with the record to $\sci{2.5}{-5}$. Since $\CSR_8-\frac14J_9$ of
 the independent re-derivation agrees with the compendium (Sec.~\ref{sec:indep}), the low independent value lies in its
 $J_9$, and within $J_9$ in $A_9[\varepsilon^0\ell^1]$, the one piece that deviates; whether that deviation is a
 fluctuation with an underestimated error or a systematic error is not known.
\item With a different regulator (analytic continuation in $v$ instead of a hard core) and different responses (the
 magnetization on the infinite line instead of finite-interval responses with an endpoint source), it reproduces the
 closed forms of $\CSR_3$ to $\CSR_7$ to $\sci{7}{-26}$ and the record of $\CSR_8$ within its uncertainty. This is
 numerical support for the prescription of Sec.~\ref{sec:dual} and for the compendium's integrations, not a proof of
 either.
\item The uncertainty is an empirical numerical estimate from three runs and measured roundoff: not a standard
 deviation and not a bound, and it cannot exclude an error common to all three runs. At $\sci{5}{-4}$ relative it can
 neither identify nor certify the rational coefficients of Eq.~\eqref{eq:C8cand}.
\item It assumes, in its own coordinates, the same two-coupling description of the gas as the prescription (the
 analyticity relations test it order by order), the identification $\De=\Delta_-$ of Eq.~\eqref{eq:eigen}, the
 Bethe-ansatz equations of the anisotropic Kondo model and its identification with the gas at $s=1$. It does not test
 whether the gas, or the dual, describes the critical chain.
\end{itemize}

%--------------------------------------------------------------------
\subsection{Records}\label{app:integ:records}
The records of the calculation are frozen and identified by their sha256 hashes; they are in
\file{code/integ_d1/results/}, \texttt{make test} checks the hashes, and \texttt{make replay} rebuilds each current
record byte for byte from the shipped engine outputs. The
record of $\CSR_8$ contains Table~\ref{tab:integ:inputs}, the values of $\CSR_3$ to $\CSR_7$ and the amendment texts,
and it is bound to the records of its inputs:
\begin{center}\small
\begin{tabular}{@{}>{\raggedright\arraybackslash}p{.37\textwidth}l@{}}
\toprule
record & sha256\\
\midrule
$\CSR_8$ (\file{C8_record_v2.json}) & \scriptsize\texttt{0d463b29d9ec474cacf2ab96afff62ae56b0ec262d766684ebbf0b7b8b20ec0c}\\
two bubbles (\file{M2_flow_coefficients.json}) & \scriptsize\texttt{33cc2e1e9a6e9a57510b0b3c6a7ef419258317be892774c3054d04852161e833}\\
three bubbles, one logarithm (\file{M3_flow_coefficients.json}) & \scriptsize\texttt{4328cab931ea96e953e9e6ba3789dfbfc0859c94f26fc8d38c1b1161eb271bce}\\
$f_{202}$ (\file{M4A_flow_coefficients.json}) & \scriptsize\texttt{584562ca72485e0f57c2f1e8ac6a05dcc5cc07b2acc73659a39587bd61641048}\\
$f_{301}$ (\file{M4B_flow_coefficients_v2.json}) & \scriptsize\texttt{90542a62db08d875995f2944e84ec7d6ff8987acc1b50892cd459d7e632d9899}\\
confirmatory check (\file{M4B_cx_check_v2.json}) & \scriptsize\texttt{66287b2af2f4c5352a1c15dc07edf6b381c2215f54b5b979fe726f4febf0f1fb}\\
\bottomrule
\end{tabular}
\end{center}

%====================================================================
\phantomsection
\addcontentsline{toc}{section}{References}
\begin{thebibliography}{99}
\bibitem{Paper}
\emph{The long-range Ising CFT: perturbative and nonperturbative approaches}, the paper that this
release accompanies (2026); source in \file{paper/} of the public repository.

\bibitem{NRGnote}
\emph{Thermal dimension of the one-dimensional long-range Ising model from bosonic NRG of the sub-ohmic spin--boson
model}, note in the folder \file{NRG/} of the public repository (\file{NRG/note/lri_nrg_1d_note.pdf}), 2026.

\bibitem{Benedetti:2024wgx}
D.~Benedetti, E.~Lauria, D.~Maz\'a\v c and P.~van Vliet,
\emph{One-dimensional Ising model with $1/r^{1.99}$ interaction},
Phys.\ Rev.\ Lett.\ \textbf{134}, 201602 (2025),
\href{https://arxiv.org/abs/2412.12243}{arXiv:2412.12243}.

\bibitem{Benedetti:2025nzp}
D.~Benedetti, E.~Lauria, D.~Maz\'a\v c and P.~van Vliet,
\emph{A strong--weak duality for the 1d long-range Ising model},
SciPost Phys.\ \textbf{20}, 029 (2026),
\href{https://arxiv.org/abs/2509.05250}{arXiv:2509.05250}.

\bibitem{Anderson:1969xx}
P.~W.~Anderson and G.~Yuval,
\emph{Exact results in the Kondo problem: equivalence to a classical one-dimensional Coulomb gas},
Phys.\ Rev.\ Lett.\ \textbf{23}, 89 (1969).

\bibitem{Anderson:1970kyb}
P.~W.~Anderson, G.~Yuval and D.~R.~Hamann,
\emph{Exact results in the Kondo problem. II. Scaling theory, qualitatively correct solution, and some new results on
one-dimensional classical statistical models}, Phys.\ Rev.\ B \textbf{1}, 4464 (1970).

\bibitem{Leggett:1987zz}
A.~J.~Leggett, S.~Chakravarty, A.~T.~Dorsey, M.~P.~A.~Fisher, A.~Garg and W.~Zwerger,
\emph{Dynamics of the dissipative two-state system},
Rev.\ Mod.\ Phys.\ \textbf{59}, 1 (1987).

\bibitem{Andrei:1983}
N.~Andrei, K.~Furuya and J.~H.~Lowenstein,
\emph{Solution of the Kondo problem},
Rev.\ Mod.\ Phys.\ \textbf{55}, 331 (1983).

\bibitem{Tsvelick:1983}
A.~M.~Tsvelick and P.~B.~Wiegmann,
\emph{Exact results in the theory of magnetic alloys},
Adv.\ Phys.\ \textbf{32}, 453 (1983).

\bibitem{Kosterlitz:1976zz}
J.~M.~Kosterlitz,
\emph{Phase transitions in long-range ferromagnetic chains},
Phys.\ Rev.\ Lett.\ \textbf{37}, 1577 (1976).

\bibitem{Behan:2017dwr}
C.~Behan, L.~Rastelli, S.~Rychkov and B.~Zan,
\emph{Long-range critical exponents near the short-range crossover},
Phys.\ Rev.\ Lett.\ \textbf{118}, 241601 (2017),
\href{https://arxiv.org/abs/1703.03430}{arXiv:1703.03430}.

\bibitem{Behan:2017emf}
C.~Behan, L.~Rastelli, S.~Rychkov and B.~Zan,
\emph{A scaling theory for the long-range to short-range crossover and an infrared duality},
J.\ Phys.\ A \textbf{50}, 354002 (2017),
\href{https://arxiv.org/abs/1703.05325}{arXiv:1703.05325}.

\bibitem{Ardonne:2010zz}
E.~Ardonne and G.~Sierra,
\emph{Chiral correlators of the Ising conformal field theory},
J.\ Phys.\ A \textbf{43}, 505402 (2010),
\href{https://arxiv.org/abs/1008.2863}{arXiv:1008.2863}.

\bibitem{Simmons-Duffin:2016wlq}
D.~Simmons-Duffin,
\emph{The lightcone bootstrap and the spectrum of the 3d Ising CFT},
JHEP \textbf{03}, 086 (2017),
\href{https://arxiv.org/abs/1612.08471}{arXiv:1612.08471}.

\bibitem{Chang:2024whx}
C.-H.~Chang, V.~Dommes, R.~S.~Erramilli, A.~Homrich, P.~Kravchuk, A.~Liu, M.~S.~Mitchell, D.~Poland and
D.~Simmons-Duffin, \emph{Bootstrapping the 3d Ising stress tensor},
JHEP \textbf{03}, 136 (2025),
\href{https://arxiv.org/abs/2411.15300}{arXiv:2411.15300}.

\bibitem{Kos:2014bka}
F.~Kos, D.~Poland and D.~Simmons-Duffin,
\emph{Bootstrapping mixed correlators in the 3D Ising model},
JHEP \textbf{11}, 109 (2014),
\href{https://arxiv.org/abs/1406.4858}{arXiv:1406.4858}.

\bibitem{Reehorst:2021}
M.~Reehorst,
\emph{Rigorous bounds on irrelevant operators in the 3d Ising model CFT},
JHEP \textbf{09}, 177 (2022),
\href{https://arxiv.org/abs/2111.12093}{arXiv:2111.12093}.

\bibitem{Henriksson:2022}
J.~Henriksson, S.~R.~Kousvos and M.~Reehorst,
\emph{Spectrum continuity and level repulsion: the Ising CFT from infinitesimal to finite $\varepsilon$},
JHEP \textbf{02}, 218 (2023),
\href{https://arxiv.org/abs/2207.10118}{arXiv:2207.10118}.

\bibitem{Zhu:2022}
W.~Zhu, C.~Han, E.~Huffman, J.~S.~Hofmann and Y.-C.~He,
\emph{Uncovering conformal symmetry in the 3D Ising transition: state-operator correspondence from a quantum fuzzy
sphere regularization}, Phys.\ Rev.\ X \textbf{13}, 021009 (2023),
\href{https://arxiv.org/abs/2210.13482}{arXiv:2210.13482}.

\bibitem{Fardelli:2026}
G.~Fardelli, A.~L.~Fitzpatrick and E.~Katz,
\emph{Improving 3d Ising OPE coefficients with fuzzy sphere conformal generators},
\href{https://arxiv.org/abs/2602.04958}{arXiv:2602.04958}.

\bibitem{Behan:2018hfx}
C.~Behan,
\emph{Bootstrapping the long-range Ising model in three dimensions},
J.\ Phys.\ A \textbf{52}, 075401 (2019),
\href{https://arxiv.org/abs/1810.07199}{arXiv:1810.07199}.

\bibitem{Behan:2023ile}
C.~Behan, E.~Lauria, M.~Nocchi and P.~van Vliet,
\emph{Analytic and numerical bootstrap for the long-range Ising model},
JHEP \textbf{03}, 136 (2024),
\href{https://arxiv.org/abs/2311.02742}{arXiv:2311.02742}.
\bibitem{Benedetti:2020rrq}
D.~Benedetti, R.~Gurau, S.~Harribey and K.~Suzuki,
\emph{Long-range multi-scalar models at three loops},
J.\ Phys.\ A \textbf{53}, 445008 (2020),
\href{https://arxiv.org/abs/2007.04603}{arXiv:2007.04603}.

\bibitem{Shiraishi:2023}
D.~Shiraishi,
\emph{Spence--Kummer's trilogarithm functional equation and its underlying geometry},
\href{https://arxiv.org/abs/2307.09414}{arXiv:2307.09414}.

\bibitem{Brown:2011uzk}
F.~Brown,
\emph{Mixed Tate motives over $\mathbb Z$},
Ann.\ Math.\ \textbf{175}, 949 (2012),
\href{https://arxiv.org/abs/1102.1312}{arXiv:1102.1312}.

\bibitem{Brown:2009qja}
F.~Brown,
\emph{Multiple zeta values and periods of moduli spaces $\mathfrak M_{0,n}$},
Ann.\ Sci.\ \'Ec.\ Norm.\ Sup\'er.\ \textbf{42}, 371 (2009),
\href{https://arxiv.org/abs/math/0606419}{arXiv:math/0606419}.

\bibitem{Brown:2013depth}
F.~Brown,
\emph{Depth-graded motivic multiple zeta values},
\href{https://arxiv.org/abs/1301.3053}{arXiv:1301.3053}.

\bibitem{Benedetti:2024mqx}
D.~Benedetti, R.~Gurau and S.~Harribey,
\emph{Addendum: Long-range multi-scalar models at three loops},
J.\ Phys.\ A \textbf{58}, 129401 (2025),
\href{https://arxiv.org/abs/2411.00805}{arXiv:2411.00805}.

\bibitem{Brunetti:2000}
R.~Brunetti and K.~Fredenhagen,
\emph{Microlocal analysis and interacting quantum field theories: renormalization on physical backgrounds},
Commun.\ Math.\ Phys.\ \textbf{208}, 623 (2000),
\href{https://arxiv.org/abs/math-ph/9903028}{arXiv:math-ph/9903028}.

\bibitem{Brunetti:2009}
R.~Brunetti, M.~D\"utsch and K.~Fredenhagen,
\emph{Perturbative algebraic quantum field theory and the renormalization groups},
Adv.\ Theor.\ Math.\ Phys.\ \textbf{13}, 1541 (2009),
\href{https://arxiv.org/abs/0901.2038}{arXiv:0901.2038}.

\bibitem{Keller:2009}
K.~J.~Keller,
\emph{Euclidean Epstein--Glaser renormalization},
J.\ Math.\ Phys.\ \textbf{50}, 103503 (2009),
\href{https://arxiv.org/abs/0902.4789}{arXiv:0902.4789}.

\bibitem{Dutsch:2014}
M.~D\"utsch, K.~Fredenhagen, K.~J.~Keller and K.~Rejzner,
\emph{Dimensional regularization in position space, and a forest formula for Epstein--Glaser renormalization},
J.\ Math.\ Phys.\ \textbf{55}, 122303 (2014),
\href{https://arxiv.org/abs/1311.5424}{arXiv:1311.5424}.

\bibitem{Duhr:2023}
C.~Duhr and F.~Porkert,
\emph{Feynman integrals in two dimensions and single-valued hypergeometric functions},
\href{https://arxiv.org/abs/2309.12772}{arXiv:2309.12772}.

\bibitem{NIST:DLMF}
\emph{NIST Digital Library of Mathematical Functions}, \url{https://dlmf.nist.gov}: Chapters 16 (generalized hypergeometric functions), 18 (\S18.14, Eq.~(18.14.7)) and 19 (elliptic integrals).

\bibitem{EliasMiro:2021}
J.~Elias Mir\'o and J.~Ingoldby,
\emph{Hamiltonian truncation with larger dimensions},
\href{https://arxiv.org/abs/2112.09049}{arXiv:2112.09049}.

\bibitem{EliasMiro:2022}
J.~Elias Mir\'o and J.~Ingoldby,
\emph{Effective Hamiltonians and counterterms for Hamiltonian truncation},
JHEP \textbf{07}, 052 (2023),
\href{https://arxiv.org/abs/2212.07266}{arXiv:2212.07266}.

\bibitem{Liu:2026}
Y.~Liu,
\emph{Spin-operator form factors of the critical Ising chain and their finite volume scaling limits},
\href{https://arxiv.org/abs/2601.00751}{arXiv:2601.00751}.

\bibitem{Hogervorst:2013}
M.~Hogervorst and S.~Rychkov,
\emph{Radial coordinates for conformal blocks},
Phys.\ Rev.\ D \textbf{87}, 106004 (2013),
\href{https://arxiv.org/abs/1303.1111}{arXiv:1303.1111}.
\end{thebibliography}
\end{document}
